\documentclass[10pt, letterpaper]{article}

% Inhaltsverzeichnis für Pakettypen (nur für Übersicht im Header, wird nicht im Dokument angezeigt)
% 1. Seitenlayout und Ränder
% 2. Sprache und Zeichensatz
% 3. Mathematik und Theorem-Umgebungen
% 4. Eigene Makros
% 5. Diagramme und Grafiken
% 6. Tabellen und Aufzählungen
% 7. Inhaltsverzeichnis
% 8. Abschnittsüberschriften
% 9. Abstrakt-Umgebung
% 10. Todos/Notizen
% 11. Rahmen/Box-Umgebungen
% 12. Python-Integration
% 13. Literaturverwaltung
% 14. Hyperlinks
% 15. Absatzeinstellungen
% 16. Umgebungen
% 17  Graphik
% 18  Extra
% 00. Titel und Autor

\usepackage{slashed}

% --- 1. Seitenlayout und Ränder ---
\usepackage[margin=3cm]{geometry}

% --- 2. Sprache und Zeichensatz ---
\usepackage[english]{babel}
\usepackage[T1]{fontenc}
\usepackage[utf8]{inputenc}

% --- 3. Mathematik und Theorem-Umgebungen ---
\usepackage{amsmath, amssymb, amsthm}
\usepackage{mathrsfs}
\DeclareMathOperator{\WF}{WF}

% --- 4. Eigene Makros ---
\usepackage{xcolor}
\newcommand{\SKP}{\langle\cdot,\cdot\rangle}
\newcommand{\id}{\operatorname{id}}
\newcommand{\sgn}{\operatorname{sgn}}
\newcommand{\Ad}{\operatorname{Ad}}
\newcommand{\ad}{\operatorname{ad}}
\newcommand{\R}{\mathbb{R}}
\newcommand{\N}{\mathbb{N}}
\newcommand{\Q}{\mathbb{Q}}
\newcommand{\Z}{\mathbb{Z}}
\newcommand{\QU}{\mathbb{H}}
\newcommand{\Aut}{\mathrm{Aut}}
\newcommand{\End}{\mathrm{End}}
\newcommand{\QUom}{\mathrm{Hom}}
\newcommand{\C}{\mathbb{C}}
\newcommand{\Cl}{\mathrm{Cl}}
\newcommand{\entwurf}[1]{\textcolor{red}{#1}}
\newcommand{\GL}{\mathrm{GL}}
\newcommand{\Mat}{\mathrm{Mat}}
\newcommand{\SL}{\mathrm{SL}}
\newcommand{\SO}{\mathrm{SO}}
\newcommand{\SU}{\mathrm{SU}}
\newcommand{\Sp}{\mathrm{Sp}}
\newcommand{\Spin}{\mathrm{Spin}}
\newcommand{\Pin}{\mathrm{Pin}}
\newcommand{\Lor}{\mathrm{Lor}}
\newcommand{\Ogroup}{\mathrm{O}}
\newcommand{\U}{\mathrm{U}}

% --- 5. Diagramme und Grafiken ---
\usepackage{graphicx}
\graphicspath{{../../Library/Images/}}
\usepackage[export]{adjustbox}
\usepackage{tikz}
\usetikzlibrary{decorations.pathreplacing, arrows.meta, positioning}
\usepackage{tikz-cd}

% --- 6. Tabellen und Aufzählungen ---
\usepackage{enumitem}
\setlist[itemize]{left=0.5cm}

\newenvironment{romanenum}[1][]
  {%
    \ifx&#1&
    \else
      \textbf{#1}\quad
    \fi
    \begin{enumerate}[label=\roman*)]
  }
  {%
    \end{enumerate}%
  }

% --- 7. Inhaltsverzeichnis ---
\usepackage{tocloft}
\renewcommand{\cftsecfont}{\footnotesize}
\renewcommand{\cftsubsecfont}{\footnotesize}
\renewcommand{\cftsubsubsecfont}{\footnotesize}
\renewcommand{\cftsecpagefont}{\footnotesize}
\renewcommand{\cftsubsecpagefont}{\footnotesize}
\renewcommand{\cftsubsubsecpagefont}{\footnotesize}
\usepackage{etoc}

% --- 8. Abschnittsüberschriften ---
\usepackage{titlesec}
\titleformat{\section}{\normalfont\large\bfseries}{\thesection}{1em}{}
\titleformat{\subsection}{\normalfont\normalsize\bfseries}{\thesubsection}{0.5em}{}
\titleformat{\subsubsection}{\normalfont\normalsize\bfseries}{\thesubsubsection}{0.5em}{}
\setcounter{secnumdepth}{4}

% --- 9. Abstrakt-Umgebung ---
\usepackage{changepage}
\renewenvironment{abstract}
  {
    \begin{adjustwidth}{1.5cm}{1.5cm}
    \small
    \textsc{Abstract. –}%
  }
  {
    \end{adjustwidth}
  }

% --- 10. Todos/Notizen ---
\usepackage{todonotes}
% Hellblau definieren
\definecolor{hellblau}{rgb}{0.3,0.6,1}
\newenvironment{note}{\par\begingroup\color{hellblau}\itshape}{\par\endgroup}

% --- 11. Rahmen/Box-Umgebungen ---
\usepackage{mdframed}
\usepackage{tcolorbox}
\colorlet{shadecolor}{gray!25}

\newenvironment{customTheorem}
  {\vspace{10pt}%
   \begin{mdframed}[
     backgroundcolor=gray!20,
     linewidth=0pt,
     innertopmargin=10pt,
     innerbottommargin=10pt,
     skipabove=\dimexpr\topsep+\ht\strutbox\relax,
     skipbelow=\topsep,
   ]}
  {\end{mdframed}
   \vspace{10pt}%
  }

% --- 12. Python-Integration ---
% (Deaktiviert in dieser Version, aktiviere bei Bedarf)
% \usepackage{pythontex}
% \usepackage[makestderr]{pythontex}

% --- 13. Literaturverwaltung ---
\usepackage{csquotes}
\usepackage[backend=biber, style=alphabetic, citestyle=alphabetic]{biblatex}
\addbibresource{bibliography.bib}

% --- 14. Hyperlinks ---
\usepackage{hyperref}
\hypersetup{
  colorlinks   = true,
  urlcolor     = blue,
  linkcolor    = blue,
  citecolor    = blue,
  frenchlinks  = true
}

% --- 15. Absatzeinstellungen ---
\usepackage[parfill]{parskip}
\sloppy

% --- 16. Umgebungen ---
\usepackage{thmtools}

\newcommand{\CustomHeading}[3]{%
  \par\medskip\noindent%
  \textbf{#1 #2} \textnormal{(#3)}.\enskip%
}

\newenvironment{DEF}[2]{\begin{unitbox}\CustomHeading{Definition}{#1}{#2}}{\end{unitbox}}
\newenvironment{PROP}[2]{\begin{unitbox}\CustomHeading{Proposition}{#1}{#2}}{\end{unitbox}}
\newenvironment{THEO}[2]{\begin{unitbox}\CustomHeading{Theorem}{#1}{#2}}{\end{unitbox}}
\newenvironment{LEM}[2]{\begin{unitbox}\CustomHeading{Lemma}{#1}{#2}}{\end{unitbox}}
\newenvironment{KORO}[2]{\begin{unitbox}\CustomHeading{Corollar}{#1}{#2}}{\end{unitbox}}
\newenvironment{REM}[2]{\begin{unitbox}\CustomHeading{Remark}{#1}{#2}}{\end{unitbox}}
\newenvironment{EXA}[2]{\begin{unitbox}\CustomHeading{Example}{#1}{#2}}{\end{unitbox}}
\newenvironment{STUD}[2]{\begin{unitbox}\CustomHeading{Study}{#1}{#2}}{\end{unitbox}}
\newenvironment{CONC}[2]{\begin{unitbox}\CustomHeading{Concept}{#1}{#2}}{\end{unitbox}}
\newenvironment{OTH}[2]{\begin{unitbox}\CustomHeading{Other}{#1}{#2}}{\end{unitbox}}
\newenvironment{EXE}[2]{\begin{unitbox}\CustomHeading{Exercise}{#1}{#2}}{\end{unitbox}}
\newenvironment{MOT}[2]{\begin{unitbox}\CustomHeading{Motivation}{#1}{#2}}{\end{unitbox}}
\newenvironment{PROOF}[2]{\begin{unitbox}\CustomHeading{Proof}{#1}{#2}}{\end{unitbox}}

% --- Unit Umgebung für Source-Inhalte ---
\usepackage{mdframed}
\newmdenv[
  linewidth=1pt,
  topline=false,
  bottomline=false,
  rightline=false,
  leftmargin=0cm,
  rightmargin=0cm,
  skipabove=10pt,
  skipbelow=10pt,
  innertopmargin=0.5\baselineskip,
  innerbottommargin=0.5\baselineskip,
  backgroundcolor=gray!10,
  linecolor=gray
]{unitbox}

\newenvironment{unit}[1]
  {\begin{unitbox}\textbf{Unit #1}\par\smallskip}
  {\end{unitbox}}

% --- 17. ... ---

% --- 18. Extras ---
\usepackage{stmaryrd}
\usepackage{bbold}  % falls du \mathbb{1} nutzen willst

% --- 00. Titel und Autor ---
\title{Mein Titel}
\author{Tim Jaschik}
\date{\today}

\begin{document}

\maketitle
\rule{\textwidth}{0.5pt}
\begin{abstract}
Kurze Beschreibung …
\end{abstract}
\rule{\textwidth}{0.5pt}
\vspace{0.5cm}

\tableofcontents

\pagebreak

\section{Lorentzkovarianz der Dirac Gleichung}

% ---------------------------------------------------
\subsection{Clifford-Algebraic Foundations}
% ---------------------------------------------------

\begin{DEF}{C1}{Clifford Algebra}
Let $(V,\eta)$ be a real vector space with a non-degenerate symmetric bilinear form $\eta:V\times V\to\R$.  
The \emph{Clifford algebra} associated with $(V,\eta)$ is
\[
\Cl(V,\eta) := T(V) \big/ \langle v\otimes v - \eta(v,v) 1 \mid v\in V\rangle.
\]
We embed $V\subset\Cl(V,\eta)$ via $v\mapsto [v]$, so that for all $v,w\in V$ one has the fundamental relation
\[
vw + wv = 2\eta(v,w)1.
\]
\end{DEF}

\begin{DEF}{C2}{Grading and Invertibility}
$\Cl(V,\eta)$ is naturally $\mathbb Z_2$-graded by even and odd elements:
\[
\Cl(V,\eta) = \Cl^0(V,\eta)\oplus\Cl^1(V,\eta),
\]
where $\Cl^0$ is generated by products of an even number of vectors.  
Elements $a\in\Cl(V,\eta)$ are invertible if there exists $a^{-1}$ with $aa^{-1}=a^{-1}a=1$.
\end{DEF}

\begin{DEF}{C3}{Pin and Spin Groups}
Inside $\Cl(V,\eta)$, define
\[
\Pin(V,\eta) := \langle u_1\cdots u_k \mid u_i\in V,\ \eta(u_i,u_i)=\pm 1\rangle,
\qquad
\Spin(V,\eta):=\Pin(V,\eta)\cap \Cl^0(V,\eta).
\]
\end{DEF}

\begin{PROP}{C4}{Adjoint Action and Canonical Homomorphism}
For each $s\in \Pin(V,\eta)$ define the adjoint action on $V$ by
\[
\rho_s(v):= s v s^{-1}.
\]
Then $\rho_s\in \Ogroup(V,\eta)$, and the map
\[
\pi:\Pin(V,\eta)\longrightarrow \Ogroup(V,\eta), \qquad \pi(s)=\rho_s,
\]
is a group homomorphism with kernel $\{\pm1\}$; its restriction
\[
\pi:\Spin(V,\eta)\longrightarrow \SO(V,\eta)
\]
is a double covering of the special orthogonal group.
\end{PROP}

\begin{PROOF}{P1}
Every $\Lambda\in \Ogroup(V,\eta)$ is a product of reflections $\rho_{u_i}$ with unit vectors $u_i$ (Cartan--Dieudonné theorem).  
By definition of $\rho_s(v)=s v s^{-1}$ this gives $\pi$ surjective, and $\pi(s)=\pi(-s)$ yields $\ker\pi=\{\pm1\}$.  
The even subgroup acts orientation-preserving.
\end{PROOF}

% ---------------------------------------------------
\subsection{Lie Algebras and Spinor Representations}
% ---------------------------------------------------

\begin{DEF}{L1}{Lie Algebra of the Spin Group}
The Lie algebra of $\Spin(V,\eta)$ is the bivector space
\[
\mathfrak{spin}(V,\eta):=\mathrm{span}\{\, v w \mid v,w\in V\}\cap \Cl^0(V,\eta)
\simeq \Lambda^2 V,
\]
with commutator bracket $[a,b]=ab-ba$.  
For an orthonormal basis $(e_\mu)$ of $V$ the elements
\[
E_{\mu\nu}:=\tfrac12\,e_\mu e_\nu,\quad \mu<\nu,
\]
form a basis of $\mathfrak{spin}(V,\eta)$.
\end{DEF}

\begin{PROP}{L2}{Canonical Isomorphism}
The differential of $\pi$ at the identity,
\[
d\pi:\mathfrak{spin}(V,\eta)\longrightarrow \mathfrak{so}(V,\eta),
\qquad
d\pi(E_{\mu\nu}) = M_{\mu\nu},
\]
is a Lie algebra isomorphism, where $M_{\mu\nu}$ act on $V$ by
\[
M_{\mu\nu}(v) = \eta(e_\mu,v)e_\nu - \eta(e_\nu,v)e_\mu.
\]
\end{PROP}

\begin{PROOF}{P2}
For $B=\frac12\omega_{\mu\nu}E^{\mu\nu}\in\mathfrak{spin}$ and $v\in V$,
\[
[B,v]=\frac12\omega_{\mu\nu}(e^\mu e^\nu v - v e^\mu e^\nu)
= \omega^\mu{}_\nu\,v^\nu e_\mu,
\]
which reproduces exactly the infinitesimal $\mathfrak{so}(1,3)$-action. 
Hence $d\pi$ is a Lie algebra isomorphism.
\end{PROOF}

\begin{DEF}{L3}{Dirac Representation}
Let $\gamma_\mu$ be a faithful complex representation of $\Cl(1,3)$ on $\C^4$ satisfying
\[
\{\gamma_\mu,\gamma_\nu\} = 2\eta_{\mu\nu}I.
\]
Define
\[
\Sigma^{\mu\nu} := \tfrac14[\gamma^\mu,\gamma^\nu].
\]
Then $\Sigma^{\mu\nu}$ obey the $\mathfrak{so}(1,3)$ commutation relations, giving a representation
\[
dS:\mathfrak{so}(1,3)\to \End(\C^4),\qquad 
dS(M_{\mu\nu}) = \tfrac12\Sigma_{\mu\nu}.
\]
Exponentiation yields the \emph{spinor representation}
\[
S:\Spin^+(1,3)\to \GL(\C^4),\qquad S(\exp \omega)=\exp\!\big(\tfrac12\,\omega_{\mu\nu}\Sigma^{\mu\nu}\big).
\]
\end{DEF}

% ---------------------------------------------------
\section{Spin-Lift and Lorentz Covariance}
% ---------------------------------------------------

\begin{LEM}{S1}{Intertwining Property}
For all $\widetilde\Lambda\in \Spin^+(1,3)$ and $\Lambda=\pi(\widetilde\Lambda)\in SO^+(1,3)$ holds
\[
S(\widetilde\Lambda)\,\gamma^\mu\,S(\widetilde\Lambda)^{-1}
= \Lambda^\mu{}_\nu\,\gamma^\nu.
\]
\end{LEM}

\begin{PROOF}{P3}
Infinitesimally:
\[
[dS(\omega),\gamma^\mu]
= \tfrac12\omega_{\rho\sigma}[\Sigma^{\rho\sigma},\gamma^\mu]
= \omega^\mu{}_\nu\,\gamma^\nu.
\]
Integrating this relation (via BCH formula) proves the finite version.
\end{PROOF}

\begin{PROP}{S2}{Lorentz Covariance of the Dirac Equation}
Let $\psi:\R^{1,3}\to\C^4$ solve the Dirac equation
\[
(i\gamma^\mu\partial_\mu - m)\psi = 0.
\]
For $(a,\Lambda)\in \R^{1,3}\rtimes SO^+(1,3)$ choose a lift $\widetilde\Lambda\in\Spin^+(1,3)$ with $\pi(\widetilde\Lambda)=\Lambda$.  
Define
\[
\psi'(x') = S(\widetilde\Lambda)\psi(x),\qquad x'=\Lambda x + a.
\]
Then $\psi'$ satisfies the same equation:
\[
(i\gamma^\mu\partial'_\mu - m)\psi'(x')=0.
\]
\end{PROP}

\begin{PROOF}{P4}
Using $\partial'_\mu=(\Lambda^{-1})^\nu{}_\mu\partial_\nu$ and the intertwining property,
\[
\begin{aligned}
(i\gamma^\mu\partial'_\mu - m)\psi'(x')
&= i\gamma^\mu(\Lambda^{-1})^\nu{}_\mu S(\widetilde\Lambda)\partial_\nu\psi(x) - mS(\widetilde\Lambda)\psi(x)\\
&= S(\widetilde\Lambda)\big(iS^{-1}\gamma^\mu S\partial_\mu - m\big)\psi(x)\\
&= S(\widetilde\Lambda)\big(i\gamma^\mu\partial_\mu - m\big)\psi(x)=0.
\end{aligned}
\]
\end{PROOF}

% ---------------------------------------------------
\section{Physical Commentary}
% ---------------------------------------------------

\begin{REM}{P5}{Conceptual Interpretation}
\begin{itemize}
\item The Klein--Gordon field transforms trivially under Lorentz transformations, hence its equation is directly invariant under the Poincaré group.
\item The Dirac field carries a non-trivial spinor representation.  Spinors are \emph{not} tensors of $SO^+(1,3)$ but linear representations of its double cover $\Spin^+(1,3)\cong SL(2,\C)$.
\item The Lorentz covariance of the Dirac operator requires the intertwining relation of Lemma~S1, achievable only via the spin-lift $\pi^{-1}(\Lambda)$.
\item Physically, the necessity of $\Spin$ reflects the intrinsic double-valuedness of spin-$\tfrac12$ states under $2\pi$ rotations.
\end{itemize}
\end{REM}

\begin{REM}{P6}{Geometric and Field-Theoretic Relevance}
On general curved spacetimes $(M,g)$, the above construction generalizes:
a spin structure is a principal $\Spin(1,3)$-bundle $P_{\Spin}$ projecting onto the orthonormal frame bundle $P_{SO}$,
and the Dirac operator $\slashed{D}$ acts on sections of the associated spinor bundle.
This geometric lift is essential for defining fermionic fields in general relativity and quantum field theory on curved manifolds.
\end{REM}


% =====================================================
%  SECTION 5: Hierarchical Definitions of Lorentz Covariance
% =====================================================

\subsection{Hierarchical Definitions of Lorentz Covariance}

\begin{MOT}{H1}{Top--Down Motivation}
Physical theories in relativistic spacetime are formulated under the requirement that their equations of motion have the same form in all inertial reference frames.  
Mathematically, this requirement is expressed through the invariance or covariance of the equations under the action of the Lorentz group or, more generally, the Poincaré group.  
Since physical fields take values in certain linear spaces on which the symmetry group acts, the concept of \emph{representation} is the formal tool linking geometry (group action on spacetime) and algebra (action on field components).
\end{MOT}

% -----------------------------------------------------
\subsection{Abstract Level: Symmetry and Invariance}
% -----------------------------------------------------

\begin{DEF}{H2}{Lorentz Invariance of a Theory}
Let $\mathcal{P} = \R^{1,3}\rtimes SO^+(1,3)$ denote the Poincaré group acting on Minkowski space via
\[
(a,\Lambda)\cdot x = \Lambda x + a.
\]
A field equation
\[
E(\partial,\Phi)=0,\qquad \Phi:\R^{1,3}\to V,
\]
is called \emph{Lorentz-invariant} if there exists a group action
\[
U:\mathcal{P}\to \Aut(\mathcal{F}),\qquad (U(a,\Lambda)\Phi)(x) = \rho(\Lambda)\Phi(\Lambda^{-1}(x-a)),
\]
such that for all $(a,\Lambda)\in\mathcal{P}$ one has
\[
E(\partial,\Phi)=0\quad \Longleftrightarrow\quad E'(\partial',U(a,\Lambda)\Phi)=0,
\]
where $\partial'_\mu = (\Lambda^{-1})^\nu{}_\mu\,\partial_\nu$.  
Here $\rho$ is a (possibly projective) representation of the Lorentz group on the internal space $V$.
\end{DEF}

\begin{REM}{H3}{On the Role of Representations}
The group action on spacetime coordinates $x\mapsto \Lambda x + a$ alone is not sufficient to describe how the \emph{components} of a field transform.  
For instance, while the coordinates of a vector change linearly under $\Lambda$, the components of a spinor require a lifted action through the Spin group.  
Hence, to express physical covariance, one must specify a \emph{representation}
\[
\rho:G\to \GL(V),\qquad G\in\{SO^+(1,3),\Spin^+(1,3)\},
\]
which tells how the internal indices of the field transform when the observer changes frame.
\end{REM}

\begin{DEF}{H4}{Representation}
Let $G$ be a Lie group and $V$ a vector space over $\R$ or $\C$.  
A \emph{representation} of $G$ on $V$ is a homomorphism
\[
\rho:G\longrightarrow \GL(V),\qquad \rho(g_1 g_2)=\rho(g_1)\rho(g_2).
\]
The pair $(V,\rho)$ defines a linear action of $G$ on $V$, given by $g\cdot v:=\rho(g)v$.  
If $\rho$ is not single-valued but satisfies $\rho(g_1)\rho(g_2)=\pm\rho(g_1g_2)$, it is called a \emph{projective representation}.  
Spinor representations of $SO^+(1,3)$ are of this projective type, becoming honest representations on $\Spin^+(1,3)$.
\end{DEF}

\begin{REM}{H5}{Why Representations Capture Physical Covariance}
The Lorentz transformations act on the coordinates of spacetime, while physical quantities (fields, currents, tensors) are sections of bundles whose fibres transform under specific finite-dimensional representations of the symmetry group.  
The requirement that the field equations remain form-invariant under these actions is exactly what the representation formalism encodes.  
Hence, Lorentz-covariance of an equation is nothing else than equivariance of the corresponding differential operator with respect to the group representation acting on its domain.
\end{REM}

% -----------------------------------------------------
\subsection{Tensorial Level: Direct Covariance under \texorpdfstring{$SO^+(1,3)$}{SO+(1,3)}}
% -----------------------------------------------------

\begin{DEF}{H6}{Tensor Covariance}
Let $\rho:SO^+(1,3)\to \GL(V)$ be a linear representation (e.g.\ the vector or tensor representation).  
A $V$-valued field $\Phi:\R^{1,3}\to V$ transforms as
\[
\Phi'(x') = \rho(\Lambda)\Phi(x),\qquad x'=\Lambda x.
\]
The field equation $E(\Phi)=0$ is \emph{tensorially Lorentz-covariant} if
\[
E(\Phi)=0 \implies E'(\Phi')=0,
\]
with $E'$ obtained from $E$ by replacing all indices according to the Lorentz transformation laws.
\end{DEF}

\begin{EXA}{H7}{Klein--Gordon Equation}
For a scalar field $\phi:\R^{1,3}\to\R$ with trivial representation $\rho(\Lambda)=1$, one has
\[
(\Box+m^2)\phi = (\eta^{\mu\nu}\partial_\mu\partial_\nu+m^2)\phi = 0.
\]
Since $\eta^{\mu\nu}$ and $\Box$ are Lorentz scalars under $SO^+(1,3)$, the equation satisfies
\[
(\Box'+m^2)\phi'(x')=(\Box+m^2)\phi(x)=0,
\]
and is therefore fully Lorentz-invariant under the \emph{tensorial} definition.
\end{EXA}

% -----------------------------------------------------
\subsection{Spinorial Level: Covariance under the Spin-Lift}
% -----------------------------------------------------

\begin{DEF}{H8}{Spinor Covariance}
Let $\pi:\Spin^+(1,3)\twoheadrightarrow SO^+(1,3)$ be the canonical covering map.  
A spinor field $\psi:\R^{1,3}\to V$ with a representation $S:\Spin^+(1,3)\to\GL(V)$ transforms as
\[
\psi'(x') = S(\widetilde\Lambda)\psi(x),\qquad \pi(\widetilde\Lambda)=\Lambda.
\]
A differential operator $D$ acting on such fields is called \emph{Spin-covariant} if
\[
D'(U(\widetilde\Lambda)\psi)=U(\widetilde\Lambda)\,D\psi,
\]
where $U(\widetilde\Lambda)$ denotes the induced field transformation.
\end{DEF}

\begin{EXA}{H9}{Dirac Equation}
The Dirac operator $D=i\gamma^\mu\partial_\mu-m$ satisfies
\[
S(\widetilde\Lambda)\gamma^\mu S(\widetilde\Lambda)^{-1} = \Lambda^\mu{}_\nu\gamma^\nu.
\]
Hence, for $\psi'(x')=S(\widetilde\Lambda)\psi(x)$ and $\partial'_\mu=(\Lambda^{-1})^\nu{}_\mu\partial_\nu$, one obtains
\[
(i\gamma^\mu\partial'_\mu-m)\psi'(x')
=S(\widetilde\Lambda)(i\gamma^\mu\partial_\mu-m)\psi(x)=0.
\]
The Dirac equation is therefore Lorentz-covariant under the lifted Spin-action, but not under $SO^+(1,3)$ itself.
\end{EXA}

% -----------------------------------------------------
\subsection{Comparative Hierarchy}
% -----------------------------------------------------

\begin{center}
\begin{tabular}{|c|c|c|c|c|}
\hline
\textbf{Level} & \textbf{Group} & \textbf{Type of Field} & \textbf{Equation} & \textbf{Covariance}\\
\hline
I & $\mathcal P$ & abstract field space $\mathcal F$ & $E(\Phi)=0$ & general invariance\\
\hline
II & $SO^+(1,3)$ & tensors, scalars, vectors & $(\Box+m^2)\phi=0$ & direct\\
\hline
III & $\Spin^+(1,3)$ & spinors & $(i\gamma^\mu\partial_\mu-m)\psi=0$ & via lift\\
\hline
\end{tabular}
\end{center}

\begin{REM}{H10}{Why the Dirac Equation Cannot be Level II Covariant}
Spinors are not tensors under $SO^+(1,3)$:  
there exists no representation $\rho:SO^+(1,3)\to GL(\C^4)$ satisfying $\rho(\Lambda)\gamma^\mu\rho(\Lambda)^{-1}=\Lambda^\mu{}_\nu\gamma^\nu$.  
Hence, the Dirac operator cannot be expressed in a tensorial form on $SO^+(1,3)$; its covariance requires the lifted Spin-action.  
Physically, this does not contradict relativity, because all measurable bilinear combinations (currents, stress tensors) are tensors on $SO^+(1,3)$, ensuring full Lorentz-invariance of observables.
\end{REM}

% -----------------------------------------------------
\subsection{Physical Criterion and Adequacy of the Representation Approach}
% -----------------------------------------------------

\begin{REM}{H11}{Why One Does Not Use Direct Left-Actions on Coordinates}
If one defined Lorentz-invariance purely as invariance under the \emph{left-action} $x\mapsto \Lambda x$ on coordinates, one would only describe the geometry of spacetime itself, not how fields residing on it respond to that transformation.  
Different physical quantities (scalars, vectors, spinors) transform differently, and these transformation laws are not captured by coordinate actions alone.  
The representation approach internalizes this dependence and hence fully characterizes the physical symmetry.
\end{REM}

\begin{REM}{H12}{Sufficiency of the Representation Framework}
Mathematically, the representation formalism provides a sufficient and rigorous encoding of the physical concept of Lorentz covariance because:
\begin{enumerate}[label=\roman*)]
\item it specifies how every component of a field transforms under changes of inertial frames;
\item it ensures that the differential operators built from $\eta_{\mu\nu}$ and $\partial_\mu$ are \emph{equivariant} with respect to the group action;
\item it reproduces the transformation properties of all observables as true tensors under $SO^+(1,3)$.
\end{enumerate}
Thus, requiring covariance with respect to an appropriate representation is the precise mathematical formulation of the physical postulate of relativity.
\end{REM}

% -----------------------------------------------------
\subsection{Conclusion of the Hierarchy}
% -----------------------------------------------------

\begin{CONC}{H13}{Summary}
\begin{itemize}
\item The notion of Lorentz-invariance has several mathematical refinements:
  \begin{enumerate}[label=\roman*)]
  \item abstract invariance under the Poincaré group (level I);
  \item tensorial covariance under $SO^+(1,3)$ (level II);
  \item spinorial covariance under $\Spin^+(1,3)$ (level III).
  \end{enumerate}
\item The Klein--Gordon equation realizes (ii) directly; the Dirac equation realizes (iii) through the Spin-lift.
\item Both formulations satisfy the physical requirement of Lorentz invariance; their difference lies solely in the internal representation space in which their fields live.
\end{itemize}
\end{CONC}

% =====================================================


% =====================================================
%  SECTION 6: A Functorial / Bundle-Theoretic Notion of Physical Covariance
% =====================================================


\pagebreak

\subsection{A Functorial / Bundle-Theoretic Notion of Physical Covariance}

\begin{MOT}{N0}{Top--Down Aim}
We formulate \emph{physical covariance} as a functorial (natural) property of field equations with respect to the automorphism group of the spacetime structure.  
This framework treats scalar, tensor and spinor fields on the same footing: fields are sections of vector bundles \emph{associated} to a principal bundle encoding the spacetime structure group.  
Equivariance is imposed at the level of bundle automorphisms, not ad hoc at the level of components.  
In Minkowski space this reproduces both Klein--Gordon and Dirac as \emph{natural differential operators} associated to the orthonormal (resp.\ spin) frame bundle.
\end{MOT}

% -----------------------------------------------------
\subsection{Spacetime, Structure Group, and Principal Bundles}
% -----------------------------------------------------

\begin{DEF}{N1}{Relativistic Spacetime with Structure}
Let $(M,\eta)$ be a time-oriented, oriented Lorentzian manifold (Minkowski: $M=\R^{1,3}$, $\eta=\mathrm{diag}(+,-,-,-)$).  
Let $O(M,\eta)\to M$ denote the orthonormal frame bundle with structure group $SO^+(1,3)$.  
A \emph{spin structure} is a principal $\Spin^+(1,3)$-bundle $\pi:P_{\Spin}\to M$ together with a $2\!:\!1$ bundle morphism $\zeta:P_{\Spin}\to O(M,\eta)$ intertwining the right-actions via the covering $\Spin^+(1,3)\xrightarrow{\ \pi\ }SO^+(1,3)$.
\end{DEF}

\begin{DEF}{N2}{Associated Vector Bundles and Fields}
Let $H$ be either $SO^+(1,3)$ or $\Spin^+(1,3)$, and let $P\to M$ be a principal $H$-bundle (orthonormal frames or spin frames).  
Given a finite-dimensional representation $\rho:H\to \GL(V)$, the \emph{associated vector bundle} is
\[
E_\rho := P\times_\rho V \longrightarrow M.
\]
A (classical) \emph{field of type $\rho$} is a section $\Phi\in\Gamma(M,E_\rho)$.
\end{DEF}

\begin{REM}{N3}{Unification of Scalars, Tensors, Spinors}
Scalars and tensors correspond to representations of $SO^+(1,3)$ (via the standard action on tensor powers of $\R^{1,3}$), while Dirac spinors correspond to representations of $\Spin^+(1,3)$ that do \emph{not} descend to $SO^+(1,3)$.  
The bundle-theoretic definition treats both cases identically: pick the appropriate structure group $H$ and representation $\rho$.
\end{REM}

% -----------------------------------------------------
\subsection{Automorphisms and Natural Actions on Fields}
% -----------------------------------------------------

\begin{DEF}{N4}{Spacetime Automorphisms and Bundle Lifts}
Let $\mathrm{Iso}(M,\eta)$ be the group of isometries of $(M,\eta)$ (Poincaré group in the flat case).  
A \emph{lift} of $g\in \mathrm{Iso}(M,\eta)$ to $P$ is a principal bundle automorphism
\[
\widehat g:P\to P,\qquad \widehat g(p\cdot h)=\widehat g(p)\cdot h,\quad \text{covering } g:M\to M.
\]
For $P=O(M,\eta)$ such lifts exist canonically.  
For $P=P_{\Spin}$ lifts exist iff $(M,\eta)$ is spin and $g$ preserves the spin structure (true for $\mathrm{Iso}(M,\eta)$ in Minkowski space).
\end{DEF}

\begin{DEF}{N5}{Induced Action on Associated Fields}
Given $\widehat g\in\Aut(P)$ covering $g\in\mathrm{Iso}(M,\eta)$, there is a naturally induced vector bundle automorphism
\[
\widehat g_\rho: E_\rho \to E_\rho,\qquad [p,v] \mapsto [\widehat g(p),v],
\]
yielding a linear action
\[
U_\rho(g):\Gamma(M,E_\rho)\to\Gamma(M,E_\rho),\qquad (U_\rho(g)\Phi)(x):=
\widehat g_\rho\!\left(\Phi\big(g^{-1}x\big)\right).
\]
\end{DEF}

\begin{REM}{N6}{Why Not ``Left Actions on Coordinates'' Alone}
A left action on $M$ encodes how points move, but does not prescribe how \emph{field components} transform.  
The bundle lift $\widehat g$ together with $\rho$ precisely determines how internal indices are transported.  
This is the correct mathematical capture of the physical change of observer for any field type (scalar, tensor, spinor).
\end{REM}

% -----------------------------------------------------
\subsection{Natural Differential Operators (Equivariance)}
% -----------------------------------------------------

\begin{DEF}{N7}{Natural (Equivariant) Differential Operators}
Let $D:\Gamma(M,E_\rho)\to\Gamma(M,F_\sigma)$ be a linear differential operator between associated bundles built from the same principal bundle $P$ (and possibly a connection).  
We say $D$ is \emph{natural (equivariant)} if for all $g\in\mathrm{Iso}(M,\eta)$ with lift $\widehat g$ one has
\[
U_\sigma(g)\circ D \;=\; D\circ U_\rho(g).
\]
\end{DEF}

\begin{REM}{N8}{Physical Covariance as Naturality}
The above identity is the bundle-theoretic (coordinate-free) expression of ``same equation in all inertial frames'': the diagram
\[
\begin{tikzcd}
\Gamma(M,E_\rho) \arrow[r,"D"] \arrow[d,"U_\rho(g)"']
& \Gamma(M,F_\sigma) \arrow[d,"U_\sigma(g)"]\\
\Gamma(M,E_\rho) \arrow[r,"D"']
& \Gamma(M,F_\sigma)
\end{tikzcd}
\]
commutes.  
No component-level ad hoc arguments are required; covariance is functorial.
\end{REM}

% -----------------------------------------------------
\subsection{KG and Dirac as Natural Operators}
% -----------------------------------------------------

\begin{EXA}{N9}{Klein--Gordon}
Take $P=O(M,\eta)$, $\rho$ the trivial representation on $V=\R$ (scalar fields).  
Let $D:=\Box + m^2:\Gamma(M,\R)\to\Gamma(M,\R)$, where $\Box=\nabla^\mu\nabla_\mu$ is built from the Levi--Civita connection of $\eta$.  
Then $D$ is natural with respect to $\mathrm{Iso}(M,\eta)$ (it is an invariant scalar operator).  
Hence solutions are carried into solutions by $U_\rho(g)$ for all isometries $g$.
\end{EXA}

\begin{EXA}{N10}{Dirac}
Assume a spin structure $P_{\Spin}\to M$ and the complex Dirac representation $\rho: \Spin^+(1,3)\to \GL(\C^4)$.  
Form $S:=P_{\Spin}\times_\rho \C^4$ (spinor bundle).  
Let $\nabla^S$ be the spin connection and define the Dirac operator
\[
\slashed D := \gamma^\mu \nabla^S_\mu : \Gamma(M,S)\to\Gamma(M,S),\qquad (\text{flat case: } \slashed D = \gamma^\mu\partial_\mu).
\]
Then $\slashed D+m$ is a natural operator with respect to all $g\in\mathrm{Iso}(M,\eta)$ that lift to $P_{\Spin}$ (in Minkowski: all Poincaré transformations).  
Equivariance follows from the fact that the $\Spin$-lift intertwines the Clifford action:
\[
U_\rho(g)\circ(\slashed D+m) = (\slashed D+m)\circ U_\rho(g).
\]
\end{EXA}

\begin{PROP}{N11}{Unified Covariance}
Both Klein--Gordon and Dirac equations are special cases of naturality of $D$ with respect to lifted isometries:  
\[
D\Phi=0 \ \Longrightarrow\ D(U_\rho(g)\Phi)=0\qquad (\forall g\in\mathrm{Iso}(M,\eta)).
\]
The only difference lies in the chosen structure group $H$ and representation $\rho$:  
for KG, $(H,\rho)=(SO^+, \text{trivial})$; for Dirac, $(H,\rho)=(\Spin^+, \text{Dirac})$.
\end{PROP}

\begin{PROOF}{N11p}
By Def.~N7 and the construction of $U_\rho(g)$ in Def.~N5, the commuting square implies solution spaces are mapped into each other by symmetry.  
The KG and Dirac constructions fit the hypotheses (Examples N9, N10).
\end{PROOF}

% -----------------------------------------------------
\subsection{Consequences and Conceptual Clarity}
% -----------------------------------------------------

\begin{REM}{N12}{``Spinors are not tensors'' reinterpreted}
In this framework, the statement means: spinor fields are sections of bundles associated to the principal \emph{spin} bundle, not to the orthonormal frame bundle.  
They therefore transform under lifts $\widehat g$ in $\Aut(P_{\Spin})$, not directly under $\Aut(O(M,\eta))$.  
Because $\zeta:P_{\Spin}\to O(M,\eta)$ covers the same spacetime isometries, no physical symmetry is changed—only the correct carrier of the internal action is used.
\end{REM}

\begin{REM}{N13}{Why this is not ad hoc}
The choice of $H$ (orthonormal vs.\ spin) is dictated by the \emph{existence of a representation} for the field’s internal indices.  
Projective actions of $SO^+$ correspond to honest actions of its double cover $\Spin^+$.  
The functorial/natural definition imposes a single covariance requirement: equivariance under \emph{all} lifted spacetime isometries.  
Thus the same axiom uniformly captures scalars, tensors, and spinors.
\end{REM}

\begin{REM}{N14}{Wigner’s View (Quantum)}
In relativistic quantum theory, physical one-particle states carry unitary \emph{projective} representations of the Poincaré group; by a standard theorem these are unitary representations of its universal cover.  
Fermions (spin-$\tfrac12$) thus necessarily involve the spin cover.  
This is the Hilbert-space analogue of the classical bundle picture above.
\end{REM}

\begin{CONC}{N15}{Takeaway}
\begin{enumerate}[label=\roman*)]
\item \textbf{Physical covariance} is the \emph{naturality} of field equations under spacetime isometries, expressed as equivariance with respect to induced actions on associated bundles.
\item \textbf{KG} and \textbf{Dirac} are both natural operators; the difference is solely the structure group/representation pair $(H,\rho)$ used to build the associated bundle.
\item Spinors not being tensors is \emph{expected}: they require the spin principal bundle; nonetheless, the underlying spacetime symmetries remain exactly the Lorentz/Poincaré ones.
\end{enumerate}
\end{CONC}


\pagebreak

% =====================================================
%  SECTION 7: Why Representations? — From QM Axioms to SRT × QM and Spin
% =====================================================

\subsection{Representations as the Correct Mathematical Carrier of Symmetry}

\begin{MOT}{R0}{Aim}
We explain, in a formal and top--down manner, why \emph{representations} of symmetry groups are fundamental in quantum theory (and a fortiori in SRT~$\times$~QM).  
The key points are: (i) what a symmetry means on a Hilbert space of states, (ii) why such symmetries act (at least) projectively and hence lift to genuine representations of a covering group, and (iii) how this forces the use of the spin cover for relativistic symmetries. We then give a systematic historical pathway and end with a critical historical remark.
\end{MOT}

% -----------------------------------------------------
\subsection{Quantum Symmetries on Hilbert Space: Rays and Projective Actions}
% -----------------------------------------------------

\begin{DEF}{R1}{Hilbert Space of States and Rays}
Let $\mathcal H$ be a complex Hilbert space. Physical pure states are identified with \emph{rays}, i.e.\ one-dimensional subspaces $[\psi]=\{\lambda\psi:\lambda\in\C^\times\}$ for $\psi\in\mathcal H\setminus\{0\}$. The projective space $\mathbb P(\mathcal H)$ is the set of rays.
\end{DEF}

\begin{DEF}{R2}{Quantum Symmetry}
A bijection $\mathsf S:\mathbb P(\mathcal H)\to\mathbb P(\mathcal H)$ is a \emph{(projective) quantum symmetry} if it preserves transition probabilities:
\[
\big|\langle \psi,\phi\rangle\big|^2 = \big|\langle \mathsf S\psi,\mathsf S\phi\rangle\big|^2
\quad\text{(well defined on rays).}
\]
\end{DEF}

\begin{THEO}{R3}{Wigner’s Theorem}
Every projective quantum symmetry $\mathsf S$ is implemented by either a unitary or an antiunitary operator $U:\mathcal H\to\mathcal H$ such that $\mathsf S[\psi]=[U\psi]$. Two implementers differ by a phase: if $U$ and $U'$ implement the same ray map, then $U'=e^{i\theta}U$.
\end{THEO}

\begin{REM}{R4}{Consequences}
If a topological (Lie) group $G$ acts by quantum symmetries on rays, then there exists a map $U:G\to\mathcal U(\mathcal H)$ such that
\[
U(g_1)U(g_2)=\omega(g_1,g_2)\,U(g_1g_2),
\qquad \omega:G\times G\to \U(1),
\]
i.e.\ a \emph{projective unitary representation}. The multiplier $\omega$ encodes the obstruction to strict multiplicativity.
\end{REM}

\begin{THEO}{R5}{Bargmann’s Lift to a Cover}
Let $G$ be a connected Lie group. Any continuous projective unitary representation $U$ of $G$ lifts to a genuine unitary representation of the universal covering group $\widetilde G$:
\[
\exists\ \widetilde U:\widetilde G\to\mathcal U(\mathcal H)\quad\text{unitary,}\qquad
U = \widetilde U\circ \mathrm{pr}\quad\text{mod phases}.
\]
\end{THEO}

\begin{CONC}{R6}{Why Representations are Fundamental}
States as rays force symmetries to act projectively; Wigner~$\Rightarrow$ every symmetry on rays is unitary/antiunitary; Bargmann~$\Rightarrow$ projective actions are honest actions of a cover.  
\emph{Therefore:} the mathematically correct carrier of symmetry in QM is a (unitary) \emph{representation} of a suitable \emph{covering} of the classical symmetry group. This is not optional; it follows from the axioms (rays, probability preservation, continuity).
\end{CONC}

% -----------------------------------------------------
\subsection{Relativistic Symmetry: Poincaré Group and Its Spin Cover}
% -----------------------------------------------------

\begin{DEF}{R7}{Relativistic Symmetry Groups}
In Minkowski spacetime, the isometry group is the Poincaré group
\[
\mathcal P=\R^{1,3}\rtimes SO^+(1,3).
\]
The connected Lorentz group satisfies $\pi_1\big(SO^+(1,3)\big)\cong \Z_2$, hence its universal cover is the double cover
\[
\Spin^+(1,3)\cong SL(2,\C),\qquad \pi:\Spin^+\twoheadrightarrow SO^+.
\]
\end{DEF}

\begin{PROP}{R8}{Spin as Forced by SRT $\times$ QM}
Any continuous projective unitary representation of $SO^+(1,3)$ (or $\mathcal P$) lifts to a genuine unitary representation of its universal cover. Because $\pi_1\cong\Z_2$, there is exactly one nontrivial connected cover: the spin cover. Therefore halfinteger spin appears \emph{inevitably} once we combine SRT with the quantum postulates for symmetries.
\end{PROP}

\begin{REM}{R9}{Classical vs.\ Quantum Carriers}
Classical (tensorial) fields transform under $SO^+(1,3)$. Quantum states need unitary (projective) representations; these live naturally on the cover $\Spin^+(1,3)$. Both encode the same spacetime isometries (via $\pi$), but the quantum carrier detects the nontrivial topology ($\Z_2$) and thus allows Spin-$\tfrac12$.
\end{REM}

% -----------------------------------------------------
\subsection{Field Equations as Natural/Equivariant Operators}
% -----------------------------------------------------

\begin{DEF}{R10}{Equivariant (Natural) Differential Operators}
Let $P\to M$ be a principal $H$-bundle ($H=SO^+$ or $\Spin^+$), $\rho:H\to\GL(V)$ a finite-dimensional representation, and $E_\rho=P\times_\rho V$. A differential operator $D:\Gamma(M,E_\rho)\to\Gamma(M,E_\rho)$ is \emph{$H$-natural} if for every spacetime isometry $g$ with lift $\widehat g\in\Aut(P)$,
\[
U_\rho(g)\circ D = D\circ U_\rho(g),
\qquad (U_\rho(g)\Phi)(x) := \widehat g_\rho\big(\Phi(g^{-1}x)\big).
\]
\end{DEF}

\begin{EXA}{R11}{KG vs.\ Dirac as Two Instances}
\begin{itemize}
\item \emph{Klein--Gordon:} $H=SO^+$, $\rho$ trivial on $\R$, $D=\Box+m^2$ is $SO^+$-natural.
\item \emph{Dirac:} $H=\Spin^+$, $\rho$ Dirac representation on $\C^4$, $D=\slashed D+m$ is $\Spin^+$-natural and, via the intertwining $S(\widetilde\Lambda)\gamma^\mu S^{-1}=\Lambda^\mu{}_\nu\gamma^\nu$, yields covariance under projected Lorentz transformations.
\end{itemize}
Thus both equations are \emph{physically covariant}; they differ only in the \emph{structure group/representation} pair $(H,\rho)$ appropriate to their carrier fields.
\end{EXA}

% -----------------------------------------------------
\subsection{Systematic Historical Trajectory (Formal Sketch)}
% -----------------------------------------------------

\begin{STUD}{R12}{Formal-Historical Milestones}
\begin{enumerate}[label=\roman*)]
\item \textbf{Pre-1930s (QM Foundations).} States as rays; symmetries must preserve transition probabilities $\Rightarrow$ Wigner’s theorem (unitary/antiunitary implementers).
\item \textbf{1928 (Dirac).} Search for a \emph{linear}, relativistically covariant wave equation with relativistic spectrum $\Rightarrow$ Clifford relations $\{\gamma^\mu,\gamma^\nu\}=2\eta^{\mu\nu}$ and the Dirac operator $i\gamma^\mu\partial_\mu-m$. Algebraically, this uncovers the spinorial carrier.
\item \textbf{1931--1954 (Wigner--Bargmann).} Projective representations of symmetry groups lift to the universal cover; classification of unitary irreducible reps of the Poincaré group (mass, spin). Spin-$\tfrac12$ arises \emph{necessarily} from the nontrivial topology of $SO^+$.
\item \textbf{Geometric consolidation.} Spin structures, spinor bundles, and natural (equivariant) operators provide a coordinate-free setting where KG and Dirac are on equal footing as natural operators for different $(H,\rho)$.
\end{enumerate}
\end{STUD}

% -----------------------------------------------------
\subsection{Critical Historical Comment}
% -----------------------------------------------------

\begin{REM}{R13}{Critical Assessment}
It is historically true that Dirac’s equation emerged from a constructive (``ansatz'') demand of linearity and relativistic covariance. However, representation theory clarifies ex post that this was not an ad hoc embellishment: once SRT is combined with the quantum symmetry postulate (states as rays; symmetries preserve transition probabilities; continuity), the correct mathematical language is (projective) \emph{representations}, which—by Bargmann’s lift—forces the universal cover. Because $SO^+(1,3)$ has $\pi_1\cong\Z_2$, the \emph{unique} connected cover is the spin group. Thus spinorial carriers are not optional; they are the inevitable manifestation of relativistic quantum symmetry. Dirac’s equation is then the simplest natural first-order operator realising the spin-$\tfrac12$ representation.
\end{REM}

\begin{CONC}{R14}{Takeaway}
Representations are fundamental in QM because states are rays and symmetries act projectively; Wigner and Bargmann force a lift to a covering group. In SRT~$\times$~QM this uniquely selects $\Spin^+(1,3)$ as the correct carrier for halfinteger spin. Field equations are \emph{natural} operators with respect to the lifted spacetime isometries; KG and Dirac differ only by the representation type, not by the underlying physical symmetry.
\end{CONC}



\pagebreak

% =====================================================
%  SECTION 8: Lorentz vs. Spin — Physical vs. Representational Symmetry
% =====================================================

\section{Lorentz vs. Spin: Physical vs. Representational Symmetry}

\begin{MOT}{S0}{Aim}
We distinguish between the \emph{physical} symmetry of spacetime (the Lorentz group) and the \emph{representational} symmetry acting on quantum states (the Spin group).  
While the Spin group is mathematically indispensable for a faithful realization of relativistic quantum theory, the physically observable transformations correspond to the Lorentz group itself.  
We give a formal comparison, clarify the geometric relation, and close with a critical reflection.
\end{MOT}

% -----------------------------------------------------
\subsection{Two Notions of Symmetry}
% -----------------------------------------------------

\begin{DEF}{S1}{Physical vs. Representational Symmetry}
\begin{itemize}
\item The \emph{physical symmetry group} of special relativity is
\[
G_{\mathrm{phys}} := SO^+(1,3),
\]
the connected component of the Lorentz group preserving orientation and time-orientation.  
It acts on spacetime points by isometries $x\mapsto \Lambda x$, preserving $\eta_{\mu\nu}$.
\item The \emph{representational symmetry group} is the universal covering group
\[
G_{\mathrm{rep}} := \Spin^+(1,3)\cong SL(2,\C),
\]
together with the covering map $\pi:G_{\mathrm{rep}}\twoheadrightarrow G_{\mathrm{phys}}$, $\ker \pi=\{\pm 1\}$.
It acts linearly on carrier spaces (Hilbert spaces or spinor bundles) via representations $S:G_{\mathrm{rep}}\to \GL(V)$.
\end{itemize}
\end{DEF}

\begin{REM}{S2}{Conceptual Distinction}
$G_{\mathrm{phys}}$ encodes the \emph{observable spacetime transformations}.  
$G_{\mathrm{rep}}$ encodes how states or fields must transform in order to realise those same transformations consistently at the quantum level.  
The two actions are related by the projection $\pi$:
\[
\begin{tikzcd}
G_{\mathrm{rep}}=\Spin^+(1,3) \arrow[r, "\pi", two heads] \arrow[dr, bend right=10, "S(\cdot)"']
& G_{\mathrm{phys}}=SO^+(1,3) \arrow[d, "\Lambda\cdot x"] \\
& \text{Spacetime or field space}
\end{tikzcd}
\]
Thus $\Spin^+(1,3)$ carries no additional spacetime symmetry—it refines the same symmetry at the representational level.
\end{REM}

% -----------------------------------------------------
\subsection{Observable Symmetry: Lorentz Group}
% -----------------------------------------------------

\begin{DEF}{S3}{Lorentz Group as Observable Symmetry}
The Lorentz group is defined by
\[
SO^+(1,3) = \{\,\Lambda\in GL(4,\R) \mid \Lambda^\top \eta \Lambda = \eta,\ \det\Lambda=1,\ \Lambda^0{}_0>0\,\}.
\]
All measurable phenomena of relativity (time dilation, aberration, Doppler shift) depend solely on this condition.  
Hence $SO^+(1,3)$ remains the \emph{physically observable} symmetry of spacetime.
\end{DEF}

\begin{REM}{S4}{Geometric Completeness of $SO^+$}
The Lorentz group already exhausts all possible linear isometries of Minkowski space.  
No extension of the symmetry group can change physical predictions as long as transformations of spacetime coordinates are concerned.
\end{REM}

% -----------------------------------------------------
\subsection{Spin Group as Universal Representation Carrier}
% -----------------------------------------------------

\begin{DEF}{S5}{Spin Group as Universal Cover}
Because $\pi_1(SO^+(1,3))\cong\Z_2$, the universal covering group is a unique double cover
\[
\pi:\Spin^+(1,3)\to SO^+(1,3).
\]
The Spin group acts by conjugation on the Clifford algebra $\Cl(1,3)$:
\[
s\,v\,s^{-1} = \Lambda(v),\quad \Lambda=\pi(s).
\]
This makes $\Spin^+(1,3)$ the minimal group realising the Lorentz action internally on spinorial objects.
\end{DEF}

\begin{PROP}{S6}{Equivalence of Physical Transformations}
For any $\Lambda\in SO^+(1,3)$ and any lift $s\in \pi^{-1}(\Lambda)$,
\[
s\text{ and }-s\text{ induce the same transformation of spacetime coordinates.}
\]
Hence $\Spin^+$ adds no new physical spacetime operation—it distinguishes only internal algebraic phases on carriers of halfinteger spin.
\end{PROP}

\begin{REM}{S7}{Observable vs. Unobservable Action}
The two elements $\pm s$ act identically on all tensors and observable quantities built from $\psi$ (e.g.\ $\bar\psi\gamma^\mu\psi$), but differ by a sign on the spinor field itself:
\[
\psi \mapsto s\psi,\quad \psi\mapsto -s\psi.
\]
Since wavefunctions are not directly observable, the doubled structure is a representational refinement, not an extra symmetry of nature.
\end{REM}

% -----------------------------------------------------
\subsection{Why the Physical Symmetry is Still Lorentz}
% -----------------------------------------------------

\begin{PROP}{S8}{Reasons for Retaining $SO^+$ as Physical Group}
\begin{enumerate}[label=\roman*)]
\item The Lorentz group acts on measurable quantities (positions, momenta, electromagnetic fields).  
\item The Spin group has no new spacetime action: $\pi:\Spin^+\to SO^+$ is surjective, and $\ker\pi$ acts trivially on spacetime.
\item The difference between $s$ and $-s$ manifests only in global phases of fermionic states, which are unobservable in single-particle experiments.
\end{enumerate}
Thus the empirically accessible symmetry group of spacetime remains $SO^+(1,3)$.
\end{PROP}

\begin{REM}{S9}{Physical vs. Representational Realism}
One may adopt two complementary viewpoints:
\begin{itemize}
\item \textbf{Physical Realism:} Nature’s spacetime symmetry is $SO^+(1,3)$; the Spin group merely realises it in the appropriate representation spaces.
\item \textbf{Representational Realism:} Since all fundamental fields (including fermions) require a spin structure, $\Spin^+(1,3)$ expresses the “true” symmetry acting in the microphysical domain; $SO^+$ is the classical projection.
\end{itemize}
Both are equivalent empirically, differing only in the level at which symmetry is considered primary (geometry vs.\ algebraic realisation).
\end{REM}

% -----------------------------------------------------
\subsection{Philosophical and Historical Comment}
% -----------------------------------------------------

\begin{REM}{S10}{Historical Perspective}
Einstein’s postulates of SRT determined $SO^+(1,3)$ as the group of physically measurable spacetime transformations.  
Dirac’s discovery (1928) and Wigner’s representation analysis (1939) revealed that, at the quantum level, this group must be represented projectively, thereby forcing the use of the universal cover $\Spin^+(1,3)$.  
The Spin group is therefore not an additional symmetry but the \emph{unique algebraic refinement} enabling a faithful quantum realisation of the same physical invariances.
\end{REM}

\begin{CONC}{S11}{Summary}
\begin{itemize}
\item The Lorentz group $SO^+(1,3)$ is the physical symmetry of spacetime—the group of observable inertial transformations.
\item The Spin group $\Spin^+(1,3)$ is its universal cover, necessary for implementing those same transformations linearly on quantum states.
\item The difference is unobservable at the classical level but indispensable for quantum consistency (existence of spin-$\tfrac12$).  
\item Hence, the Lorentz symmetry remains the empirical principle; the Spin structure is its mathematical completion within quantum theory.
\end{itemize}
\end{CONC}


\pagebreak

% =====================================================
%  SECTION 9: Two Roads to the Spin Group — Dirac and Wigner–Bargmann
% =====================================================

\section{Two Roads to the Spin Group: Dirac and Wigner–Bargmann}

\begin{MOT}{T0}{Aim}
We reconstruct the two independent conceptual paths that lead to the appearance of the Spin group as the correct symmetry carrier in relativistic quantum theory:  
(i) Dirac’s constructive approach via the search for a linear relativistic wave equation, and  
(ii) the Wigner–Bargmann representation-theoretic classification of relativistic quantum states.  
Both roads converge to the same mathematical structure $\Spin^+(1,3)\cong SL(2,\C)$.
\end{MOT}

% -----------------------------------------------------
\subsection{Dirac’s Constructive Path (1928)}
% -----------------------------------------------------

\begin{STUD}{T1}{From Relativistic Linearity to Clifford Algebra}
\begin{enumerate}[label=\roman*)]
\item \textbf{Problem:} Find a first-order differential equation for $\psi(x)$ such that its square reproduces the relativistic energy relation
\[
E^2 - \mathbf{p}^2 = m^2 \quad \Rightarrow\quad (\Box + m^2)\psi = 0
\]
but which is linear in $\partial_t$.

\item \textbf{Dirac’s Ansatz:}
Seek matrices $\alpha^i,\beta$ such that
\[
H = \boldsymbol{\alpha}\cdot\mathbf{p} + \beta m,\qquad H^2 = \mathbf{p}^2 + m^2.
\]
This requires
\[
\{\alpha^i,\alpha^j\}=2\delta^{ij},\quad \{\alpha^i,\beta\}=0,\quad \beta^2=1.
\]

\item \textbf{Clifford Algebra:}
Defining $\gamma^0=\beta$, $\gamma^i=\beta\alpha^i$ yields
\[
\{\gamma^\mu,\gamma^\nu\}=2\eta^{\mu\nu},
\]
i.e.\ the Clifford algebra $\Cl(1,3)$, whose minimal complex representation has dimension~$4$.

\item \textbf{Covariance Requirement:}
Lorentz covariance demands
\[
S(\Lambda)\gamma^\mu S(\Lambda)^{-1} = \Lambda^\mu{}_\nu \gamma^\nu.
\]
This is precisely the defining property of the $\Spin^+(1,3)$ representation:
\[
S:\Spin^+(1,3)\to \GL(4,\C),\quad \pi(S)=\Lambda.
\]

\item \textbf{Outcome:}
The Dirac equation
\[
(i\gamma^\mu\partial_\mu - m)\psi=0
\]
is the simplest \emph{first-order natural operator} for the spinor representation of $\Spin^+(1,3)$.
\end{enumerate}
\end{STUD}

% -----------------------------------------------------
\subsection{Wigner–Bargmann Path (1939–1954)}
% -----------------------------------------------------

\begin{STUD}{T2}{From Quantum Symmetries to Spin}
\begin{enumerate}[label=\roman*)]
\item \textbf{Axioms:} In quantum theory, physical symmetries $G$ act as bijections of the projective Hilbert space preserving transition probabilities.

\item \textbf{Wigner’s Theorem:} Such symmetries are implemented by unitary (or antiunitary) operators $U(g)$, defined projectively:
\[
U(g_1)U(g_2) = e^{i\theta(g_1,g_2)}U(g_1g_2).
\]

\item \textbf{Bargmann’s Lift:} Every continuous projective unitary representation of $G$ is a genuine representation of its universal covering group $\widetilde G$.

\item \textbf{Application to SRT:} For $G=SO^+(1,3)$ one has $\pi_1(G)\cong\Z_2$, hence $\widetilde G=\Spin^+(1,3)$.

\item \textbf{Classification (Wigner 1939):}  
Irreducible unitary representations of the Poincaré group $\mathcal P=\R^{1,3}\rtimes SO^+(1,3)$ are labelled by $(m^2,s)$,  
where $s$ arises from the representation of $\Spin^+(1,3)$ acting on the stabiliser of the momentum.

\item \textbf{Outcome:}  
Halfinteger spins correspond exactly to representations that do not descend to $SO^+(1,3)$ but live on its spin cover.
\end{enumerate}
\end{STUD}

% -----------------------------------------------------
\subsection{Synthesis and Interpretation}
% -----------------------------------------------------

\begin{CONC}{T3}{Convergence of Both Paths}
\begin{itemize}
\item \textbf{Dirac’s route:} algebraic necessity of a linear relativistic operator $\Rightarrow$ Clifford algebra $\Rightarrow$ Spin covariance.
\item \textbf{Wigner–Bargmann route:} topological necessity of unitary projective representations $\Rightarrow$ universal cover $\Spin^+(1,3)$.
\item \textbf{Unified picture:} The Spin group is simultaneously the algebraic and topological completion of the Lorentz group required by SRT~$\times$~QM.
\end{itemize}
\end{CONC}

\begin{REM}{T4}{Historical Comment}
Dirac discovered the spin structure empirically through the attempt to linearize the Klein–Gordon operator;  
Wigner and Bargmann later demonstrated that the same structure follows inevitably from the general axioms of quantum mechanics and relativity, independent of any ansatz.  
Hence, the emergence of $\Spin^+(1,3)$ is not an accident of Dirac’s construction but a structural necessity of relativistic quantum theory.
\end{REM}


\pagebreak


\section{Where the Lorentz Group enters Dirac’s Construction}
% =====================================================
%  SECTION 9.1: Where the Lorentz Group enters Dirac’s Construction
% =====================================================

\begin{REM}{T5}{Where the Lorentz Group enters Dirac’s Construction}
Dirac’s derivation of his relativistic wave equation was guided from the outset by the demand of \emph{Lorentz covariance}—that is, by the requirement that the equation should have the same form in all inertial frames related by Lorentz transformations.  
However, the role of the Lorentz group entered his reasoning in an implicit and somewhat indirect way, becoming explicit only when he analysed how the matrices $\gamma^\mu$ must transform.

\medskip

\textbf{(i) Implicit Lorentz structure in the starting point.}  
The very motivation for Dirac’s equation was to obtain a linear differential equation whose square reproduces the relativistic energy--momentum relation
\[
E^2-\mathbf{p}^2 = m^2,
\]
which is the invariant condition
\[
p_\mu p^\mu = m^2
\]
under the Lorentz group $SO^+(1,3)$.  
Thus the background symmetry of the entire construction was already that of Minkowski space, although Dirac did not at first write it in group-theoretic form.

\medskip

\textbf{(ii) The linearisation step.}  
Dirac introduced operators $\gamma^\mu$ such that
\[
(\gamma^\mu p_\mu - m)(\gamma^\nu p_\nu + m) = p_\mu p^\mu - m^2,
\]
which leads to the Clifford relations
\[
\{\gamma^\mu,\gamma^\nu\} = 2\eta^{\mu\nu}.
\]
This defines an action of the Lorentz metric on the internal index space of the $\gamma^\mu$ matrices.  
Up to this point, Lorentz invariance manifests only through the use of the invariant metric $\eta_{\mu\nu}$.

\medskip

\textbf{(iii) Explicit appearance of the Lorentz transformations.}  
In order for the equation
\[
(i\gamma^\mu\partial_\mu - m)\psi = 0
\]
to hold in all inertial frames, Dirac required that under a Lorentz transformation $x'=\Lambda x$, the field $\psi$ should transform linearly as
\[
\psi'(x') = S(\Lambda)\psi(x),
\]
and that the operator $i\gamma^\mu\partial_\mu$ should transform covariantly, i.e.
\[
S(\Lambda)\gamma^\mu S(\Lambda)^{-1} = \Lambda^\mu{}_\nu\gamma^\nu.
\]
This condition expresses the demand of Lorentz covariance on the ``representation side'': the matrices $S(\Lambda)$ act on the values of the field, while $\Lambda$ acts on spacetime coordinates.

\medskip

\textbf{(iv) The algebraic consequence.}  
Dirac initially treated $S(\Lambda)$ as if it were a representation of the Lorentz group itself, expecting
\[
S(\Lambda_1)S(\Lambda_2)=S(\Lambda_1\Lambda_2).
\]
When he examined the consistency of … for successive transformations, it turned out that this relation holds only \emph{up to sign}:
\[
S(\Lambda_1)S(\Lambda_2)=\pm\,S(\Lambda_1\Lambda_2).
\]
That is, the matrices satisfying the Clifford covariance condition do not form an honest representation of $SO^+(1,3)$, but rather of a group that double--covers it.

\medskip

\textbf{(v) Emergence of the Spin group.}  
The set of all such matrices $S(\Lambda)$ forms a Lie group whose adjoint action on the Clifford algebra reproduces the Lorentz action on vectors:
\[
s\,v\,s^{-1} = \Lambda(v),
\qquad \Lambda = \pi(s),
\]
where $\pi:\Spin^+(1,3)\to SO^+(1,3)$ is a $2\!:\!1$ surjective homomorphism.  
Hence $S(\Lambda)$ realises a representation of $\Spin^+(1,3)$, not of $SO^+(1,3)$.

\medskip

\textbf{(vi) Conceptual meaning.}  
Dirac thus worked \emph{in the representation space} of what he presumed to be the Lorentz group;  
the requirement of covariance … revealed that this presumption cannot be realised within $SO^+(1,3)$ itself.  
The matrices implementing Lorentz covariance necessarily form a representation of its universal covering group $\Spin^+(1,3)\cong SL(2,\C)$.

\medskip

\textbf{(vii) Retrospective interpretation.}  
From a modern viewpoint one can say:
\begin{quote}
Dirac discovered the Spin group by attempting to prove that his equation was Lorentz covariant.  
He expected $S(\Lambda)$ to represent the Lorentz group on the field values, but the algebraic consistency of the Clifford covariance forced him to pass to its double cover.  
The Spin group therefore emerged not as an additional hypothesis, but as the minimal algebraic realisation making Lorentz covariance possible for linear first-order equations.
\end{quote}

\end{REM}

\pagebreak






\section{Clifford Algebras and Spinor Groups}
In this chapter, we generalise the quaternions by studying the real Clifford algebras, and our account of these is heavily influenced by the classic paper of Atiyah, Bott \& Shapiro [3]; Porteous [23, 24] also provides an accessible description, as does Curtis [7] but there are some errors and omissions in that account. Lawson \& Michelsohn [19] provides a more sophisticated introduction which shows how central Clifford algebras have become to modern geometry and topology; they also appear in Quantum Theory in connection with the Dirac operator. There is also a theory of Clifford Analysis in which the field of complex numbers is replaced by a Clifford algebra and a suitable class of Clifford analytic functions generalising complex analytic functions is studied; motivation for this is provided by the above applications. The groups of units in Clifford algebras contain the spinor groups which we define and also show how they provide double coverings of the special orthogonal groups.

We restrict attention to real Clifford algebras associated to positive definite inner products on $\mathbb{R}^{n}$. There are also complex and indefinite Clifford algebras discussed in [3, 19, 23, 24].

\subsection{Real Clifford Algebras}


The sequence of real division algebras $\mathbb{R}, \mathbb{C}, \mathbb{H}$ can be extended by introducing the (real) Clifford algebras $\mathrm{Cl}_{n}$ for $n \geqslant 0$, where the first few satisfy
$$
\mathrm{Cl}_{0}=\mathbb{R}, \quad \mathrm{Cl}_{1} \cong \mathbb{C}, \quad \mathrm{Cl}_{2} \cong \mathbb{H},
$$
as $\mathbb{R}$-algebras.

We begin by describing $\mathrm{Cl}_{n}$ as an $\mathbb{R}$-vector space and define the product in terms of a particular basis. There are elements $e_{1}, e_{2}, \ldots, e_{n} \in \mathrm{Cl}_{n}$ for which
$$
\left\{\begin{aligned}
e_{s} e_{r} & =-e_{r} e_{s} \quad \text { if } s \neq r, \\
e_{r}^{2} & =-1 .
\end{aligned}\right.
$$
Moreover, $\mathrm{Cl}_{n}$ has an $\mathbb{R}$-basis consisting of the elements $e_{i_{1}} e_{i_{2}} \cdots e_{i_{r}}$ corresponding to increasing sequences $1 \leqslant i_{1}<i_{2}<\cdots<i_{r} \leqslant n$ with $0 \leqslant r \leqslant n$. Thus
$$
\operatorname{dim}_{\mathbb{R}} \mathrm{Cl}_{n}=2^{n} .
$$
When $r=0$, the element $e_{i_{1}} e_{i_{2}} \cdots e_{i_{r}}$ is taken to be 1 .




\begin{DEF}{SG-B02-01-01}{Reelle Clifford Algebra \(\Cl_n\)}
Fix \(n\in\mathbb{N}_0\) and let \(\R^n\) carry the negative Euclidean form
\(\beta(x,y)=\sum_{r=1}^n x_r y_r\) with the Clifford sign convention
\(v\cdot v=-\beta(v,v)\,1\).
The \emph{real Clifford algebra} \(\Cl_n:=\Cl(\R^n,\beta)\) is the unital
\(\R\)-algebra generated by symbols \(e_1,\dots,e_n\) subject to the relations
\[
e_r e_s+e_s e_r=-2\delta_{rs}\,1
\quad\Longleftrightarrow\quad
\begin{cases}
e_r^2=-1,\\
e_r e_s=-e_s e_r\quad(r\neq s).
\end{cases}
\]
\end{DEF}



\begin{CONC}{SG-B02-01-02}{PBW-basis and dimension of \(\Cl_n\)}
For each increasing multi-index \(I=\{i_1<\dots<i_k\}\subset\{1,\dots,n\}\) set
\(e_I:=e_{i_1}\cdots e_{i_k}\) and \(e_\varnothing:=1\). Then:
\begin{enumerate}
\item The family \(\{e_I:\ I\subset\{1,\dots,n\}\}\) is an \(\R\)-basis of \(\Cl_n\).
\item In particular, \(\dim_\R \Cl_n=2^n\).
\end{enumerate}
\end{CONC}



\begin{LEM}{SG-B02-01-03}{Product of basis elements}
If \(I,J\) are increasing multi-indices, then
\[
e_I\,e_J
= (-1)^{N(I,J)}\,(-1)^{|I\cap J|}\,e_{\,I\,\triangle\,J},
\]
where \(I\triangle J=(I\setminus J)\cup(J\setminus I)\) is the symmetric
difference and \(N(I,J)=\#\{(i,j)\in I\times J\mid i>j\}\) counts the swaps
needed to order the concatenation. The factor \((-1)^{|I\cap J|}\) comes from
replacing each \(e_r e_r\) by \(-1\).
\end{LEM}





Doing calculations with these basis elements is easy as the following example illustrates.



Example 5.1


Let $i, j, k, \ell$ be distinct numbers in the range 1 to $n$. Then
$$
e_{i} e_{j} e_{k} e_{\ell}=e_{k} e_{\ell} e_{i} e_{j} .
$$


Repeatedly using the relations of (5.1) we obtain
$$
e_{i} e_{j} e_{k} e_{\ell}=-e_{i} e_{k} e_{j} e_{\ell}=e_{k} e_{i} e_{j} e_{\ell}=-e_{k} e_{i} e_{\ell} e_{j}=e_{k} e_{\ell} e_{i} e_{j} .
$$



\begin{LEM}{SG-B02-01-04}{Kommutator-Beziehung zwischen Basisvektoren der Clifford-Algebra}
Let \(i,j,k,\ell\) be pairwise distinct indices in \(\{1,\dots,n\}\) in the real
Clifford algebra \(\Cl_n\) with relations \(e_r e_s=-\,e_s e_r\) for \(r\neq s\).
Then
\[
e_i e_j e_k e_\ell \;=\; e_k e_\ell e_i e_j .
\]
\end{LEM}



\begin{PROOF}{SG-B02-01-05}{P: Kommutator-Beziehung zwischen Basisvektoren der Clifford-Algebra}
Using \(e_a e_b=-e_b e_a\) whenever \(a\neq b\), we commute adjacent unequal
generators step by step:
\[
e_i e_j e_k e_\ell
= -\,e_i e_k e_j e_\ell
= \;e_k e_i e_j e_\ell
= -\,e_k e_i e_\ell e_j
= \;e_k e_\ell e_i e_j,
\]
as claimed.
\end{PROOF}



Proposition 5.2

There are isomorphisms of $\mathbb{R}$-algebras
$$
\mathrm{Cl}_{1} \cong \mathbb{C}, \quad \mathrm{Cl}_{2} \cong \mathbb{H} .
$$


Proof

For $\mathrm{Cl}_{1}$, the function
$$
\mathrm{Cl}_{1} \longrightarrow \mathbb{C} ; \quad x+y e_{1} \longmapsto x+y i \quad(x, y \in \mathbb{R}),
$$
is an $\mathbb{R}$-linear ring isomorphism.

Similarly, for $\mathrm{Cl}_{2}$, the function
$$
\mathrm{Cl}_{2} \longrightarrow \mathbb{H} ; \quad t 1+x e_{1}+y e_{2}+z e_{1} e_{2} \longmapsto t 1+x i+y j+z k \quad(t, x, y, z \in \mathbb{R}),
$$
is an $\mathbb{R}$-linear ring isomorphism.


\begin{PROP}{SG-B02-01-06}{$\Cl_1 \cong \mathbb{C}, \qquad \Cl_2 \cong \mathbb{H}$}
There are isomorphisms of $\mathbb{R}$-algebras
\[
\Cl_1 \cong \mathbb{C}, \qquad \Cl_2 \cong \mathbb{H}.
\]
\end{PROP}






\begin{PROOF}{SG-B02-01-07}{P: $\Cl_1 \cong \mathbb{C}, \qquad \Cl_2 \cong \mathbb{H}$}
\textbf{Case 1}: \(\Cl_1 \cong \C\)

By definition, \(\Cl_1\) is generated (as an \(\R\)-algebra) by \(1,e_1\) with
\(e_1^2=-1\). Every element has a unique form \(x + y e_1\) with \(x,y\in\R\).
Define
\[
\Phi_1:\Cl_1\to\C,\qquad \Phi_1(x+y e_1)=x+yi.
\]
Then \(\Phi_1\) is \(\R\)-linear, bijective (by the basis \(\{1,e_1\}\)),
and multiplicative since \(\Phi_1(e_1)^2=i^2=-1=\Phi_1(e_1^2)\) and the
universal property of \(\Cl_1\) gives a unique algebra homomorphism sending
\(e_1\mapsto i\). Hence \(\Phi_1\) is an \(\R\)-algebra isomorphism.

\medskip
\textbf{Case 2:} \(\Cl_2 \cong \QU\)

\(\Cl_2\) is generated by \(1,e_1,e_2\) with relations
\(e_1^2=e_2^2=-1\) and \(e_1e_2=-e_2e_1\).
Every element has a unique form \(t 1 + x e_1 + y e_2 + z e_1e_2\),
\(t,x,y,z\in\R\). Define
\[
\Phi_2:\Cl_2\to\QU,\qquad
\Phi_2\bigl(t 1 + x e_1 + y e_2 + z e_1e_2\bigr)=t\cdot 1 + x i + y j + z k.
\]
This is \(\R\)-linear and bijective (by the basis \(\{1,e_1,e_2,e_1e_2\}\)).
It is multiplicative because the defining relations are preserved:
\[
\Phi_2(e_1)^2=i^2=-1=\Phi_2(e_1^2),\qquad
\Phi_2(e_2)^2=j^2=-1=\Phi_2(e_2^2),
\]
and
\[
\Phi_2(e_1)\Phi_2(e_2)=ij=k=-ji=\Phi_2(-e_2e_1).
\]
By the universal property of the Clifford algebra, there is a unique
algebra homomorphism sending \(e_1\mapsto i\), \(e_2\mapsto j\);
this homomorphism is \(\Phi_2\) and is bijective. Thus \(\Cl_2\cong\QU\).
\end{PROOF}











We can order the basis monomials in the $e_{r}$ by declaring $e_{i_{1}} e_{i_{2}} \cdots e_{i_{r}}$ to be numbered
$$
1+2^{i_{1}-1}+2^{i_{2}-1}+\cdots+2^{i_{r}-1},
$$
which should be interpreted as 1 when $r=0$. Every integer $k$ for which $1 \leqslant k \leqslant 2^{n}$ has a unique binary expansion
$$
k=k_{0}+2 k_{1}+\cdots+2^{j} k_{j}+\cdots+2^{n} k_{n},
$$
where each $k_{j}=0,1$. This provides a one to one correspondence between such numbers $k$ and the basis monomials of $\mathrm{Cl}_{n}$. Here are the orderings of the bases for the first few Clifford algebras.
\begin{center}
\begin{tabular}{|c|c|}
\hline
$\mathrm{Cl}_{1}$ & $1, e_{1}$ \\
$\mathrm{Cl}_{2}$ & $1, e_{1}, e_{2}, e_{1} e_{2}$ \\
$\mathrm{Cl}_{3}$ & $1, e_{1}, e_{2}, e_{1} e_{2}, e_{3}, e_{1} e_{3}, e_{2} e_{3}, e_{1} e_{2} e_{3}$ \\
\hline
\end{tabular}
\end{center}
Using the left regular representation over $\mathbb{R}$ associated with this basis of $\mathrm{Cl}_{n}$, we can realise $\mathrm{Cl}_{n}$ as a subalgebra of $\mathrm{M}_{2^{n}}(\mathbb{R})$.




\begin{DEF}{SG-B02-01-08}{Ordnungsstruktur der Basis in reellen Clifford-Algebren}
Sei $\mathrm{Cl}_n$ die reelle Clifford-Algebra über $\mathbb{R}^n$ mit Erzeugern
$e_1,\ldots,e_n$ und Relationen
\[
e_i e_j + e_j e_i = 2\delta_{ij}.
\]
\begin{enumerate}
\item
Die \textbf{Basis-Monomiale} der Form
\[
e_{i_1} e_{i_2} \cdots e_{i_r}, \quad 1 \leq i_1 < i_2 < \cdots < i_r \leq n,
\]
werden durch eine binäre Ordnung nummeriert mittels
\[
e_{i_1} e_{i_2} \cdots e_{i_r}
\;\longmapsto\;
1 + 2^{i_1 - 1} + 2^{i_2 - 1} + \cdots + 2^{i_r - 1},
\]
wobei für $r = 0$ der leere Ausdruck als $1$ interpretiert wird.

\item
Jede ganze Zahl $k$ mit $1 \le k \le 2^n$ besitzt eine eindeutige Binärdarstellung
\[
k = k_0 + 2k_1 + 2^2 k_2 + \cdots + 2^n k_n, \quad k_j \in \{0,1\}.
\]
Dies liefert eine \emph{bijektive Korrespondenz} zwischen den Indizes $k$ und den Basis-Monomialen von $\mathrm{Cl}_n$:
\[
k \longleftrightarrow e_{i_1} e_{i_2} \cdots e_{i_r},
\]
wobei $i_j$ diejenigen Indizes sind, für die $k_{i_j} = 1$ gilt.
\end{enumerate}
\end{DEF}



\begin{DEF}{SG-B02-01-09}{Beispiele für die Basisordnung kleiner Clifford-Algebren}
Für die ersten Clifford-Algebren ergibt sich die folgende Reihenfolge der Basisvektoren:
\begin{center}
\begin{tabular}{|c|l|}
\hline
$\mathrm{Cl}_1$ & $1,\ e_1$ \\[2pt]
$\mathrm{Cl}_2$ & $1,\ e_1,\ e_2,\ e_1 e_2$ \\[2pt]
$\mathrm{Cl}_3$ & $1,\ e_1,\ e_2,\ e_1 e_2,\ e_3,\ e_1 e_3,\ e_2 e_3,\ e_1 e_2 e_3$ \\
\hline
\end{tabular}
\end{center}
\end{DEF}



\begin{DEF}{SG-B02-01-10}{Realisierung der Clifford-Algebra als Matrixalgebra}
Unter Verwendung der \textbf{linken regulären Darstellung} von $\mathrm{Cl}_n$ über $\mathbb{R}$, die durch die obige Basis definiert ist, kann man
\[
\mathrm{Cl}_n \hookrightarrow M_{2^n}(\mathbb{R})
\]
als Unteralgebra der reellen Matrixalgebra realisieren. Dabei wirkt jedes Element $x \in \mathrm{Cl}_n$ durch \emph{linke Multiplikation} auf der Basis
\[
L_x : \mathrm{Cl}_n \longrightarrow \mathrm{Cl}_n,\quad y \mapsto x y,
\]
woraus eine Darstellung $L : \mathrm{Cl}_n \to \mathrm{End}_{\mathbb{R}}(\mathrm{Cl}_n) \cong M_{2^n}(\mathbb{R})$ entsteht.
\end{DEF}



Example 5.3


For $\mathrm{Cl}_{1}$ we have the basis $\left\{1, e_{1}\right\}$ and find that
$$
\rho(1)=I_{2}, \quad \rho\left(e_{1}\right)=\left[\begin{array}{rr}
0 & -1 \\
1 & 0
\end{array}\right],
$$
while the general formula is
$$
\rho\left(x+y e_{1}\right)=\left[\begin{array}{rr}
x & -y \\
y & x
\end{array}\right] \quad(x, y \in \mathbb{R}) .
$$
Applying $\rho$ to the basis $\left\{1, e_{1}, e_{2}, e_{1} e_{2}\right\}$ we obtain the following elements in $\mathrm{M}_{4}(\mathbb{R})$ :
$$
\begin{array}{lrl}
\rho(1)=I_{4}, & & \rho\left(e_{1}\right)=\left[\begin{array}{rrrr}
0 & -1 & 0 & 0 \\
1 & 0 & 0 & 0 \\
0 & 0 & 0 & -1 \\
0 & 0 & 1 & 0
\end{array}\right], \\
\rho\left(e_{2}\right)=\left[\begin{array}{rrrr}
0 & 0 & -1 & 0 \\
0 & 0 & 0 & 1 \\
1 & 0 & 0 & 0 \\
0 & -1 & 0 & 0
\end{array}\right], & \rho\left(e_{1} e_{2}\right)=\left[\begin{array}{rrrr}
0 & 0 & 0 & -1 \\
0 & 0 & -1 & 0 \\
0 & 1 & 0 & 0 \\
1 & 0 & 0 & 0
\end{array}\right] .
\end{array}
$$
In all cases the matrices $\rho\left(e_{i_{1}} e_{i_{2}} \cdots e_{i_{r}}\right)$ are generalised permutation matrices all of whose entries are $0, \pm 1$ and exactly one non-zero entry in each row and column. These are always orthogonal matrices of determinant 1 .



\begin{EXA}{SG-B02-01-11}{Explizite linke reguläre Darstellung für $\mathrm{Cl}_1$ und $\mathrm{Cl}_2$}
Wir betrachten die reellen Clifford-Algebren $\mathrm{Cl}_1$ und $\mathrm{Cl}_2$ in der durch die geordnete Basis beschriebenen Darstellung.

\paragraph{Fall $\mathrm{Cl}_1$:}
Die Basis sei $\{1, e_1\}$.  
Die linke reguläre Darstellung $\rho : \mathrm{Cl}_1 \to M_2(\mathbb{R})$ ist gegeben durch
\[
\rho(1) = I_2, 
\qquad 
\rho(e_1) = 
\begin{bmatrix}
0 & -1 \\[2pt]
1 & 0
\end{bmatrix}.
\]
Für ein allgemeines Element $x + y e_1 \in \mathrm{Cl}_1$ gilt
\[
\rho(x + y e_1) = 
\begin{bmatrix}
x & -y \\[2pt]
y & x
\end{bmatrix},
\qquad x,y \in \mathbb{R}.
\]
Diese Darstellung identifiziert $\mathrm{Cl}_1$ mit $\mathbb{C}$ über der Basis $\{1, e_1\}$ durch die Zuordnung
$1 \mapsto 1$ und $e_1 \mapsto i$.

\paragraph{Fall $\mathrm{Cl}_2$:}
Für die Basis $\{1, e_1, e_2, e_1 e_2\}$ erhält man durch dieselbe linke Multiplikation die folgenden Matrizen in $M_4(\mathbb{R})$:
\[
\rho(1) = I_4, \qquad
\rho(e_1) =
\begin{bmatrix}
0 & -1 & 0 & 0 \\
1 & 0 & 0 & 0 \\
0 & 0 & 0 & -1 \\
0 & 0 & 1 & 0
\end{bmatrix},
\]
\[
\rho(e_2) =
\begin{bmatrix}
0 & 0 & -1 & 0 \\
0 & 0 & 0 & 1 \\
1 & 0 & 0 & 0 \\
0 & -1 & 0 & 0
\end{bmatrix},
\qquad
\rho(e_1 e_2) =
\begin{bmatrix}
0 & 0 & 0 & -1 \\
0 & 0 & -1 & 0 \\
0 & 1 & 0 & 0 \\
1 & 0 & 0 & 0
\end{bmatrix}.
\]

\emph{Beobachtung.}
In allen Fällen sind die Matrizen $\rho(e_{i_1} e_{i_2} \cdots e_{i_r})$ \textbf{verallgemeinerte Permutationsmatrizen}:
\begin{itemize}
  \item Sie besitzen genau einen nichtverschwindenden Eintrag in jeder Zeile und Spalte,
  \item alle Einträge gehören zu $\{0, \pm 1\}$,
  \item und sie sind \emph{orthogonal} mit $\det = 1$.
\end{itemize}
Damit liefert die linke reguläre Darstellung eine \textbf{orthogonale Einbettung} von $\mathrm{Cl}_n$ in $M_{2^n}(\mathbb{R})$.
\end{EXA}



The Clifford algebras $\mathrm{Cl}_{n}$ are characterised by an important universal property. 



Theorem 5.4 (Universal property of a Clifford algebra)


Let $A$ be an $\mathbb{R}$-algebra. If $f: \mathbb{R}^{n} \longrightarrow A$ is an $\mathbb{R}$-linear transformation for which
$$
f(\mathbf{x})^{2}=-|\mathbf{x}|^{2} 1 \quad\left(\mathbf{x} \in \mathbb{R}^{n}\right),
$$
then there is a unique homomorphism of $\mathbb{R}$-algebras $F: \mathrm{Cl}_{n} \longrightarrow A$ for which $F \circ j_{n}=f$, i.e., for all $\mathbf{x} \in \mathbb{R}^{n}, F\left(j_{n}(\mathbf{x})\right)=f(\mathbf{x})$.

This can be indicated in the following commutative diagram
\[
\begin{tikzcd}[row sep=large, column sep=large]
& \mathbb{R}^n \arrow[dl, "j_n"'] \arrow[dr, "f"] & \\
\mathrm{Cl}_n \arrow[rr, dotted, "{\exists!\,F}"'] & & A
\end{tikzcd}
\]
in which ' $\exists$ ! $F$ ' stands for 'there exists a unique $F$ ' and the dotted arrow indicates a function $F$ which solves the following equation amongst functions, $F \circ j_{n}=f$.


Proof

First notice that there is an $\mathbb{R}$-linear transformation
$$
j_{n}: \mathbb{R}^{n} \longrightarrow \mathrm{Cl}_{n} ; \quad j_{n}\left(\sum_{r=1}^{n} x_{r} \mathbf{e}_{r}\right)=\sum_{r=1}^{n} x_{r} e_{r},
$$
for which
$$
j_{n}\left(\sum_{r=1}^{n} x_{r} \mathbf{e}_{r}\right)^{2}=-\sum_{r=1}^{n} x_{r}^{2}=-\left|\sum_{r=1}^{n} x_{r} \mathbf{e}_{r}\right|^{2} .
$$


The homomorphism $F: \mathrm{Cl}_{n} \longrightarrow A$ is defined by setting $F\left(e_{r}\right)=f\left(\mathbf{e}_{r}\right)$ and showing that it extends it to a ring homomorphism given on basis monomials by
$$
F\left(e_{i_{1}} \cdots e_{i_{r}}\right)=f\left(\mathbf{e}_{i_{1}}\right) \cdots f\left(\mathbf{e}_{i_{r}}\right) .
$$
The details are left as an exercise.



\begin{THEO}{SG-B02-01-12}{Universelle Eigenschaft der Clifford-Algebra}
Sei $A$ eine $\mathbb{R}$\nobreakdash-Algebra und $f : \mathbb{R}^n \to A$ eine $\mathbb{R}$\nobreakdash-lineare Abbildung, die für alle $\mathbf{x} \in \mathbb{R}^n$ die Bedingung
\[
f(\mathbf{x})^2 = -\lvert \mathbf{x}\rvert^2 \, 1_A
\]
erfüllt.
Dann existiert genau ein $\mathbb{R}$\nobreakdash-Algebrenhomomorphismus 
\[
F : \mathrm{Cl}_n \longrightarrow A
\]
mit der Eigenschaft
\[
F \circ j_n = f,
\]
d.\,h. für alle $\mathbf{x} \in \mathbb{R}^n$ gilt $F(j_n(\mathbf{x})) = f(\mathbf{x})$.

Dies lässt sich durch das folgende kommutative Diagramm veranschaulichen:
\[
\begin{tikzcd}[row sep=large, column sep=large]
& \mathbb{R}^n \arrow[dl, "j_n"'] \arrow[dr, "f"] & \\
\mathrm{Cl}_n \arrow[rr, dotted, "{\exists!\,F}"'] & & A
\end{tikzcd}
\]
\end{THEO}


\begin{PROOF}{SG-B02-01-13}{P: Universelle Eigenschaft der Clifford-Algebra}
\begin{enumerate}
\item
Zunächst sei 
\[
j_n : \mathbb{R}^n \to \mathrm{Cl}_n, \qquad 
j_n\!\left(\sum_{r=1}^n x_r \mathbf{e}_r\right)
= \sum_{r=1}^n x_r e_r,
\]
die kanonische $\mathbb{R}$\nobreakdash-lineare Einbettung der Erzeuger in die Clifford-Algebra.  
Dann gilt unmittelbar
\[
j_n(\mathbf{x})^2 = -\lvert \mathbf{x}\rvert^2 \, 1
\quad\text{für alle }\mathbf{x}\in\mathbb{R}^n.
\]
\item
Sei $F : \mathrm{Cl}_n \to A$ die Abbildung, die durch
\[
F(e_r) = f(\mathbf{e}_r), \quad r=1,\ldots,n
\]
auf den Erzeugern definiert wird und auf den Basis-Monomialen
\[
F(e_{i_1} e_{i_2} \cdots e_{i_r})
:= f(\mathbf{e}_{i_1}) f(\mathbf{e}_{i_2}) \cdots f(\mathbf{e}_{i_r})
\]
fortgesetzt wird.
Da $f$ dieselben Relationen erfüllt wie die Erzeuger $e_r$ von $\mathrm{Cl}_n$, ist $F$ wohldefiniert und ein Algebrenhomomorphismus.
\item
Die Eindeutigkeit folgt aus der universellen Eigenschaft der Clifford-Algebra:  
Jeder Homomorphismus $F$ wird durch die Bilder der Erzeuger eindeutig bestimmt.
\end{enumerate}
\end{PROOF}



\begin{CONC}{SG-B02-01-14}{Konsequenz der Universelle Eigenschaft der Clifford-Algebra}
Das Theorem charakterisiert $\mathrm{Cl}_n$ bis auf eindeutigen Isomorphismus durch folgende universelle Eigenschaft:
\[
\mathrm{Cl}_n 
\simeq
\text{freieste $\mathbb{R}$\nobreakdash-Algebra, in der $v^2 = -|v|^2$ für alle $v \in \mathbb{R}^n$ gilt.}
\]
Diese Eigenschaft ermöglicht es, jedes lineare System mit der Relation $f(\mathbf{x})^2=-|\mathbf{x}|^2$ eindeutig in die Clifford-Algebra einzubetten.
\end{CONC}



Corollary 5.5

Let $U$ be an $\mathbb{R}$-algebra and $j: \mathbb{R}^{n} \longrightarrow U$ be an $\mathbb{R}$-linear transformation for which
$$
j(\mathbf{x})^{2}=-|\mathbf{x}|^{2} 1 \quad\left(\mathbf{x} \in \mathbb{R}^{n}\right) .
$$
Suppose that $U$ and $j$ have the universal property enjoyed by $\mathrm{Cl}_{n}$. If $A$ is an $\mathbb{R}$-algebra and $h: \mathbb{R}^{n} \longrightarrow A$ is an $\mathbb{R}$-linear transformation for which
$$
h(\mathbf{x})^{2}=-|\mathbf{x}|^{2} 1 \quad\left(\mathbf{x} \in \mathbb{R}^{n}\right),
$$
then there is a unique homomorphism of $\mathbb{R}$-algebras $H: U \longrightarrow A$ satisfying $H \circ j=f$, i.e., for all $\mathbf{x} \in \mathbb{R}^{n}, H(j(\mathbf{x}))=h(\mathbf{x})$.

Then there is a unique $\mathbb{R}$-algebra isomorphism $\Phi: \mathrm{Cl}_{n} \longrightarrow U$ which satisfies $\Phi \circ j_{n}=j$.



\begin{KORO}{SG-B02-01-15}{Eindeutigkeit bis Isomorphie der Clifford-Algebra}
Sei $U$ eine $\mathbb{R}$\nobreakdash-Algebra und $j:\mathbb{R}^n\to U$ eine
$\mathbb{R}$\nobreakdash-lineare Abbildung mit
\[
j(\mathbf{x})^2=-\lvert\mathbf{x}\rvert^2\,1_U \qquad (\mathbf{x}\in\mathbb{R}^n),
\]
die (zusammen mit $U$) dieselbe universelle Eigenschaft wie $\mathrm{Cl}_n$ besitzt: 
Für jede $\mathbb{R}$\nobreakdash-Algebra $A$ und jede $\mathbb{R}$\nobreakdash-lineare
$h:\mathbb{R}^n\to A$ mit $h(\mathbf{x})^2=-\lvert\mathbf{x}\rvert^2 1_A$ existiert genau ein
$\mathbb{R}$\nobreakdash-Algebrenhomomorphismus $H:U\to A$ mit $H\circ j=h$.

Dann existiert genau ein $\mathbb{R}$\nobreakdash-Algebrenisomorphismus
\[
\Phi:\mathrm{Cl}_n\longrightarrow U
\]
mit
\[
\Phi\circ j_n = j,
\]
wobei $j_n:\mathbb{R}^n\to \mathrm{Cl}_n$ die kanonische Einbettung ist.
\end{KORO}


Proof

\begin{PROOF}{SG-B02-01-16}{P: Eindeutigkeit bis Isomorphie der Clifford-Algebra}
The algebra homomorphism $\Phi: \mathrm{Cl}_{n} \longrightarrow U$ is constructed using the universal property $\mathrm{Cl}_{n}$ applied to the linear transformation $j: \mathbb{R}^{n} \longrightarrow U$. On the other hand, applying the universal property of $U$ to $j_{n}: \mathbb{R}^{n} \longrightarrow \mathrm{Cl}_{n}$, we obtain an algebra homomorphism $\Theta: U \longrightarrow \mathrm{Cl}_{n}$. These homomorphisms satisfy the equations
$$
\Phi \circ j_{n}=j, \quad \Theta \circ j=j_{n} .
$$
Notice that the following equations also hold:
$$
\Theta \circ \Phi \circ j_{n}=\Theta \circ j=j_{n}, \quad \Phi \circ \Theta \circ j=\Phi \circ j_{n}=j .
$$
By the uniqueness part of the universality of $\mathrm{Cl}_{n}$ applied to the linear transformation $j_{n}: \mathbb{R}^{n} \longrightarrow \mathrm{Cl}_{n}$ together with the identity $\operatorname{Id}_{\mathrm{Cl}_{n}} \circ j_{n}=j_{n}$, we have $\Theta \circ \Phi=\mathrm{Id}_{\mathrm{Cl}_{n}}$. Similarly, combining the uniqueness part of the universality of $U$ applied to the linear transformation $j: \mathbb{R}^{n} \longrightarrow U$ with the identity $\operatorname{Id}_{U} \circ j=j$, we obtain $\Phi \circ \Theta=\operatorname{Id}_{U}$.
\[
\begin{tikzcd}[row sep=large, column sep=huge]
{} & \mathbb{R}^n & {} \\%
\mathrm{Cl}_n & U & \mathrm{Cl}_n %
% --- Pfeile ---
\arrow[from=1-2, to=2-1, "j_n"']           % R^n -> Cl_n (links)
\arrow[from=1-2, to=2-2, "j"]              % R^n -> U (vertikal)
\arrow[from=1-2, to=2-3, "j_n"]            % R^n -> Cl_n (rechts)
\arrow[from=2-1, to=2-2, dotted, "{\exists!\,\Phi}"]   % Cl_n -> U
\arrow[from=2-2, to=2-3, dotted, "{\exists!\,\Theta}"] % U -> Cl_n
\arrow[from=2-1, to=2-3, bend right=25, "\mathrm{Id}_{\mathrm{Cl}_n}"'] % gebogener Pfeil unten
\end{tikzcd}
\]
\[
\begin{tikzcd}[row sep=large, column sep=huge]
{} & \mathbb{R}^n & {} \\%
U  & \mathrm{Cl}_n & U %
% --- Pfeile (keine Leerzeile dazwischen!) ---
\arrow[from=1-2, to=2-1, "j"']            % R^n -> U (links)
\arrow[from=1-2, to=2-2, "j_n"]           % R^n -> Cl_n (senkrecht)
\arrow[from=1-2, to=2-3, "j"]             % R^n -> U (rechts)
\arrow[from=2-1, to=2-2, dotted, "{\exists!\,\Theta}"] % U -> Cl_n
\arrow[from=2-2, to=2-3, dotted, "{\exists!\,\Phi}"]   % Cl_n -> U
\arrow[from=2-1, to=2-3, bend right=25, "\mathrm{Id}_U"'] % gebogener Pfeil unten
\end{tikzcd}
\]
Hence $\Phi$ is an algebra isomorphism with inverse $\Phi^{-1}=\Theta$.
\end{PROOF}

This uniqueness up to isomorphism is an important characteristic of objects satisfying such universal properties. Another example we discuss is that of the quotient topology, see Proposition 8.3 and Corollary 8.4.



Example 5.6


There is an $\mathbb{R}$-linear transformation
$$
\alpha_{0}: \mathbb{R}^{n} \longrightarrow \mathrm{Cl}_{n} ; \quad \alpha_{0}(\mathbf{x})=-j_{n}(\mathbf{x})=j_{n}(-\mathbf{x}) .
$$
Then in $\mathrm{Cl}_{n}$,
$$
\alpha_{0}(\mathbf{x})^{2}=j_{n}(-\mathbf{x})^{2}=-|\mathbf{x}|^{2},
$$
so by Theorem 5.4 there is a unique algebra homomorphism $\alpha: \mathrm{Cl}_{n} \longrightarrow \mathrm{Cl}_{n}$ for which
$$
\alpha\left(j_{n}(\mathbf{x})\right)=\alpha_{0}(\mathbf{x}) .
$$
Since $j_{n}\left(\mathbf{e}_{r}\right)=e_{r}$, this implies that
$$
\alpha\left(e_{r}\right)=-e_{r} .
$$
Notice that for $1 \leqslant i_{1}<i_{2}<\cdots<i_{k} \leqslant n$,
$$
\alpha\left(e_{i_{1}} e_{i_{2}} \cdots e_{i_{k}}\right)=(-1)^{k} e_{i_{1}} e_{i_{2}} \cdots e_{i_{k}}=\left\{\begin{aligned}
e_{i_{1}} e_{i_{2}} \cdots e_{i_{k}} & \text { if } k \text { is even, } \\
-e_{i_{1}} e_{i_{2}} \cdots e_{i_{k}} & \text { if } k \text { is odd. }
\end{aligned}\right.
$$

It is easy to see that $\alpha$ is an isomorphism and hence an automorphism, often called the canonical automorphism of $\mathrm{Cl}_{n}$.




\begin{DEF}{SG-B02-01-17}{Kanonische (Grad-)Automorphism \(\alpha\) der \(\Cl_n\)}
Sei \(j_n:\R^n\to\Cl_n\) die kanonische Einbettung. Definiere die
\(\R\)-lineare Abbildung
\[
\alpha_0:\R^n\longrightarrow \Cl_n,\qquad \alpha_0(\mathbf x):=j_n(-\mathbf x)=-\,j_n(\mathbf x).
\]
Da \(\alpha_0(\mathbf x)^2=j_n(-\mathbf x)^2=-\lVert\mathbf x\rVert^2\cdot 1\) gilt, liefert die
Universal­eigenschaft der Clifford-Algebra einen \emph{eindeutigen}
\(\R\)-Algebrenhomomorphismus
\[
\alpha:\Cl_n\longrightarrow \Cl_n\quad\text{mit}\quad \alpha\circ j_n=\alpha_0.
\]
Man nennt \(\alpha\) den \emph{kanonischen Automorphismus} (auch: \emph{Grad-Involution}).
\end{DEF}




\begin{PROP}{SG-B02-01-18}{Eigenschaften der kanonischen Involution}
Für \(\alpha\) aus Grad-Automorphismus gelten:
\begin{enumerate}
  \item \(\alpha(e_r)=-e_r\) für alle Erzeuger \(e_r=j_n(\mathbf e_r)\).
  \item Für jedes Monom \(e_{i_1}\cdots e_{i_k}\) (mit strikt wachsenden Indizes)
  \[
  \alpha\!\left(e_{i_1}\cdots e_{i_k}\right)=(-1)^k\,e_{i_1}\cdots e_{i_k}.
  \]
  \item \(\alpha\) ist ein Algebrenautomorphismus und sogar eine Involution:
  \(\;\alpha(xy)=\alpha(x)\alpha(y)\) und \(\alpha^2=\id_{\Cl_n}\).
  \item Die Zerlegung nach Parität wird respektiert:
  \[
  \Cl_n=\Cl_n^{\bar 0}\oplus \Cl_n^{\bar 1},\qquad
  \alpha|_{\Cl_n^{\bar 0}}=\id,\quad \alpha|_{\Cl_n^{\bar 1}}=-\id,
  \]
  insbesondere ist die gerade Unteralgebra \(\Cl_n^{\bar 0}\) der Fixpunktbereich von \(\alpha\).
\end{enumerate}
\end{PROP}




\begin{PROOF}{SG-B02-01-19}{P: Eigenschaften der kanonischen Involution}
(1) folgt unmittelbar aus \(\alpha\circ j_n=\alpha_0\).

(2) ergibt sich aus (1) und der Homomorphieeigenschaft von \(\alpha\):
\(\alpha(e_{i_1}\cdots e_{i_k})=\prod_{r=1}^k\alpha(e_{i_r})=(-1)^k e_{i_1}\cdots e_{i_k}\).

(3) \(\alpha\) ist per Konstruktion ein Algebrenhomomorphismus. Ferner gilt
\((\alpha\circ\alpha)\circ j_n=\alpha\circ\alpha_0=j_n\), also \(\alpha^2=\id\) wegen Eindeutigkeit.

(4) ist eine unmittelbare Umschreibung von (2).
\end{PROOF}



We record the explicit forms of the next few Clifford algebras. Consider the $\mathbb{R}$-algebra $\mathrm{M}_{2}(\mathbb{H})$ of dimension 16. There is an $\mathbb{R}$-linear transformation
$$
\begin{aligned}
& \theta_{4}: \mathbb{R}^{4} \longrightarrow \mathrm{M}_{2}(\mathbb{H}) ; \\
& \quad \theta_{4}\left(x_{1} \mathbf{e}_{1}+x_{2} \mathbf{e}_{2}+x_{3} \mathbf{e}_{3}+x_{4} \mathbf{e}_{4}\right)=\left[\begin{array}{cc}
x_{1} i+x_{2} j+x_{3} k & x_{4} k \\
x_{4} k & x_{1} i+x_{2} j-x_{3} k
\end{array}\right]
\end{aligned}
$$
Direct calculation shows that $\theta_{4}$ satisfies the condition of Theorem 5.4, hence there is a unique $\mathbb{R}$-algebra homomorphism $\Theta_{4}: \mathrm{Cl}_{4} \rightarrow \mathrm{M}_{2}(\mathbb{H})$ for which $\Theta_{4} \circ j_{4}=\theta_{4}$. It can be shown that this is an isomorphism of $\mathbb{R}$-algebras, so
$$
\mathrm{Cl}_{4} \cong \mathrm{M}_{2}(\mathbb{H}) .
$$
Since $\mathbb{R} \subseteq \mathbb{R}^{2} \subseteq \mathbb{R}^{3} \subseteq \mathbb{R}^{4}$, we obtain compatible injective homomorphisms
$$
\Theta_{1}: \mathrm{Cl}_{1} \longrightarrow \mathrm{M}_{2}(\mathbb{H}), \quad \Theta_{2}: \mathrm{Cl}_{2} \longrightarrow \mathrm{M}_{2}(\mathbb{H}), \quad \Theta_{3}: \mathrm{Cl}_{3} \longrightarrow \mathrm{M}_{2}(\mathbb{H}),
$$
which turn out to have images
$$
\begin{aligned}
& \operatorname{im} \Theta_{1}=\left\{z I_{2}: z \in \mathbb{C}\right\}, \\
& \operatorname{im} \Theta_{2}=\left\{q I_{2}: q \in \mathbb{H}\right\}, \\
& \operatorname{im} \Theta_{3}=\left\{\left[\begin{array}{cc}
q_{1} & 0 \\
0 & q_{2}
\end{array}\right]: q_{1}, q_{2} \in \mathbb{H}\right\} .
\end{aligned}
$$
This shows that there is an isomorphism of $\mathbb{R}$-algebras
$$
\mathrm{Cl}_{3} \cong \mathbb{H} \times \mathbb{H},
$$
where the latter is the direct product of Definition 4.8. By direct calculation, we also obtain
$$
\mathrm{Cl}_{5} \cong \mathrm{M}_{4}(\mathbb{C}), \quad \mathrm{Cl}_{6} \cong \mathrm{M}_{8}(\mathbb{R}), \quad \mathrm{Cl}_{7} \cong \mathrm{M}_{8}(\mathbb{R}) \times \mathrm{M}_{8}(\mathbb{R}) .
$$



These results are summarised in Table 5.1, whose last column gives their dimensions.

\begin{table}[h]
\begin{center}
\begin{tabular}{|c|c|c|}
\hline
$\mathrm{Cl}_{0}$ & $\mathbb{R}$ & 1 \\
$\mathrm{Cl}_{1}$ & $\mathbb{C}$ & 2 \\
$\mathrm{Cl}_{2}$ & $\mathbb{H}$ & 4 \\
$\mathrm{Cl}_{3}$ & $\mathbb{H} \times \mathbb{H}$ & 8 \\
$\mathrm{Cl}_{4}$ & $\mathrm{M}_{2}(\mathbb{H})$ & 16 \\
$\mathrm{Cl}_{5}$ & $\mathrm{M}_{4}(\mathbb{C})$ & 32 \\
$\mathrm{Cl}_{6}$ & $\mathrm{M}_{8}(\mathbb{R})$ & 64 \\
$\mathrm{Cl}_{7}$ & $\mathrm{M}_{8}(\mathbb{R}) \times \mathrm{M}_{8}(\mathbb{R})$ & 128 \\
$\mathrm{Cl}_{8}$ & $\mathrm{M}_{16}(\mathbb{R})$ & 256 \\
\hline
\end{tabular}
\end{center}
\end{table}




\begin{EXA}{SG-B02-01-20}{Explizite Realisierung von \(\Cl_4\)}
Sei \(\Mat_2(\QU)\) die \(\R\)-Algebra der \(2\times2\)-Quaternionenmatrizen (Dimension \(16\) über \(\R\)).
Definiere die \(\R\)-lineare Abbildung
\[
\theta_4:\R^4\longrightarrow \Mat_2(\QU),\qquad
\theta_4\!\left(\sum_{r=1}^4 x_r\,\mathbf e_r\right)
=
\begin{bmatrix}
x_1 i + x_2 j + x_3 k & x_4 k\\
x_4 k & x_1 i + x_2 j - x_3 k
\end{bmatrix}.
\]
Dann gilt \(\theta_4(\mathbf x)^2=-\|\mathbf x\|^2\cdot I\) für alle \(\mathbf x\in\R^4\).
Folglich existiert aufgrund der Universal­eigenschaft von \(\Cl_4\) ein eindeutiger
\(\R\)-Algebrenhomomorphismus \(\Theta_4:\Cl_4\to \Mat_2(\QU)\) mit \(\Theta_4\circ j_4=\theta_4\).
Dieser Homomorphismus ist ein Isomorphismus, also
\[
\Cl_4 \;\cong\; \Mat_2(\QU).
\]
\end{EXA}



\begin{EXA}{SG-B02-01-21}{Kompatible Einbettungen \(\Cl_1,\Cl_2,\Cl_3\hookrightarrow \Mat_2(\QU)\)}
Über den Turm \(\R\subset\R^2\subset\R^3\subset\R^4\) induziert … kompatible injektive Homomorphismen
\[
\Theta_1:\Cl_1\hookrightarrow \Mat_2(\QU),\qquad
\Theta_2:\Cl_2\hookrightarrow \Mat_2(\QU),\qquad
\Theta_3:\Cl_3\hookrightarrow \Mat_2(\QU),
\]
deren Bilder explizit lauten
\[
\mathrm{im}\,\Theta_1=\{\,z\,I_2: z\in\C\,\},\qquad
\mathrm{im}\,\Theta_2=\{\,q\,I_2: q\in\QU\,\},\qquad
\mathrm{im}\,\Theta_3=\Bigl\{\begin{bmatrix}q_1&0\\[2pt]0&q_2\end{bmatrix}: q_1,q_2\in\QU\Bigr\}.
\]
Insbesondere folgt der \(\R\)-Algebrenisomorphismus
\[
\Cl_3 \;\cong\; \QU\times \QU.
\]
\end{EXA}



\begin{EXA}{SG-B02-01-23}{Explizite isomorphe Realisierungen für kleine $n$}
Durch direkte Rechnung (z.\,B. via explizite Darstellungen oder rekursive Struktur­resultate) erhält man:
\[
\Cl_5 \;\cong\; \Mat_4(\C),\qquad
\Cl_6 \;\cong\; \Mat_8(\R),\qquad
\Cl_7 \;\cong\; \Mat_8(\R)\times\Mat_8(\R),\qquad
\Cl_8 \;\cong\; \Mat_{16}(\R).
\]


\[
\begin{array}{|c|c|c|}
\hline
\Cl_0 & \R & 1\\
\Cl_1 & \C & 2\\
\Cl_2 & \QU & 4\\
\Cl_3 & \QU\times\QU & 8\\
\Cl_4 & \Mat_2(\QU) & 16\\
\Cl_5 & \Mat_4(\C) & 32\\
\Cl_6 & \Mat_8(\R) & 64\\
\Cl_7 & \Mat_8(\R)\times \Mat_8(\R) & 128\\
\Cl_8 & \Mat_{16}(\R) & 256\\
\hline
\end{array}
\]
\end{EXA}


\begin{REM}{bott}{Periodizität (Hinweis)}
Die obigen Isomorphismen reflektieren die reelle Bott-Periodizität:
\(\Cl_{n+8}\cong \Cl_n\otimes \Mat_{16}(\R)\).
Insbesondere wiederholt sich die Morita-Äquivalenzklasse der \(\R\)-Algebren modulo \(8\).
\end{REM}



To go beyond this we use the following periodicity result, in which $\mathrm{M}_{m}\left(\mathrm{Cl}_{n}\right)$ denotes the ring of $m \times m$ matrices with entries in $\mathrm{Cl}_{n}$.



Theorem 5.7


For $n \geqslant 0$,
$$
\mathrm{Cl}_{n+8} \cong \mathrm{M}_{16}\left(\mathrm{Cl}_{n}\right) .
$$
We will not give the proof (which can be found in [3,23,24]) but note that it involves the useful observation that for a unital ring $R$, there is an isomorphism of rings
$$
\mathrm{M}_{m}\left(\mathrm{M}_{n}(R)\right) \cong \mathrm{M}_{m n}(R),
$$
obtained by expressing each $m \times m$ matrix of $n \times n$ matrices as an $m n \times m n$ matrix. Thus we obtain for example
$$
\begin{aligned}
& \mathrm{Cl}_{10} \cong \mathrm{M}_{16}(\mathbb{H}), \\
& \mathrm{Cl}_{12} \cong \mathrm{M}_{16}\left(\mathrm{M}_{2}(\mathbb{H})\right) \cong \mathrm{M}_{32}(\mathbb{H}), \\
& \mathrm{Cl}_{14} \cong \mathrm{M}_{16}\left(\mathrm{M}_{8}(\mathbb{R})\right) \cong \mathrm{M}_{128}(\mathbb{R}) .
\end{aligned}
$$


\begin{THEO}{SG-B02-01-24}{Bott-Periodizität der reellen Clifford-Algebren}
Für alle \(n\ge 0\) gilt der \emph{periodische Isomorphismus}
\[
\Cl_{n+8}\;\cong\;\Mat_{16}(\Cl_n),
\]
d.\,h. die Clifford-Algebra wächst bei Erhöhung des Index um \(8\)
um einen Faktor von \(16\times16\)-Matrizen über \(\Cl_n\).

\textbf{Bemerkung.}  
Der Beweis (vgl. \cite{ABS64,LawsonMichelsohn89,Karoubi68})
verwendet die Beobachtung, dass für jede unitalen Algebra \(R\)
ein kanonischer Ringisomorphismus existiert:
\[
\Mat_m(\Mat_n(R))\;\cong\;\Mat_{mn}(R),
\]
wobei jede \(m\times m\)-Matrix aus \(n\times n\)-Blöcken als eine
einzige \(mn\times mn\)-Matrix über \(R\) aufgefasst wird.

\textbf{Beispiele.}  Durch sukzessives Anwenden ergibt sich:
\[
\begin{aligned}
\mathrm{Cl}_{10} &\cong \mathrm{Mat}_{16}(\mathbb{H}),\\
\mathrm{Cl}_{12} &\cong \mathrm{Mat}_{16}(\mathrm{Mat}_{2}(\mathbb{H})) \cong \mathrm{Mat}_{32}(\mathbb{H}),\\
\mathrm{Cl}_{14} &\cong \mathrm{Mat}_{16}(\mathrm{Mat}_{8}(\mathbb{R})) \cong \mathrm{Mat}_{128}(\mathbb{R}).
\end{aligned}
\]

\textbf{Folgerung.}
Die Morita-Äquivalenzklasse von \(\Cl_n\) ist \(8\)-periodisch:
\[
\Cl_{n+8}\;\simeq_{\text{Morita}}\;\Cl_n.
\]
\end{THEO}



In the next section we will make use of some more structure in $\mathrm{Cl}_{n}$. 


First there is a conjugation $\overline{()}: \mathrm{Cl}_{n} \longrightarrow \mathrm{Cl}_{n}$ defined by
$$
\overline{e_{i_{1}} e_{i_{2}} \cdots e_{i_{k}}}=(-1)^{k} e_{i_{k}} e_{i_{k-1}} \cdots e_{i_{1}}
$$
whenever $1 \leqslant i_{1}<i_{2}<\cdots<i_{k} \leqslant n$, and satisfying
$$
\begin{aligned}
\overline{x+y} & =\bar{x}+\bar{y}, \\
\overline{t x} & =t \bar{x},
\end{aligned}
$$
for $x, y \in \mathrm{Cl}_{n}$ and $t \in \mathbb{R}$. Notice that if $n>1, \overline{()}$ is not a ring homomorphism since whenever $r<s$,
$$
\overline{e_{r} e_{s}}=e_{s} e_{r}=-e_{r} e_{s}=-\overline{e_{r}} \overline{e_{s}} \neq \overline{e_{r}} \overline{e_{s}}
$$
However, it is a ring anti-homomorphism in the sense that for all $x, y \in \mathrm{Cl}_{n}$,
$$
\overline{x y}=\bar{y} \bar{x} \quad\left(x, y \in \mathrm{Cl}_{n}\right) .
$$
When $n=1$ or $2, \overline{()}$ agrees with the conjugations already defined in $\mathbb{C}$ and $\mathbb{H}$.




\begin{DEF}{SG-B02-01-25}{Konjugation in der Clifford-Algebra}
Es sei \( \Cl_n \) eine reelle Clifford-Algebra mit Erzeugern \( e_1,\dots,e_n \).
Wir definieren die \emph{Konjugation}
\[
\overline{(\,)}:\Cl_n \longrightarrow \Cl_n
\]
durch ihre Wirkung auf Basis­monome:
\[
\overline{e_{i_1} e_{i_2} \cdots e_{i_k}}
:= (-1)^k\, e_{i_k} e_{i_{k-1}} \cdots e_{i_1},
\qquad
1 \le i_1 < i_2 < \cdots < i_k \le n,
\]
und setzen für beliebige \(x,y\in\Cl_n\) und \(t\in\R\)
\[
\overline{x+y}=\bar x+\bar y,
\qquad
\overline{t x}=t\,\bar x.
\]

\textbf{Beobachtung.}
Für \(n>1\) ist \(\overline{(\,)}\) \emph{kein} Ringhomomorphismus, da etwa für \(r<s\)
\[
\overline{e_r e_s}
= e_s e_r
= -e_r e_s
= -\overline{e_r}\,\overline{e_s}
\ne \overline{e_r}\,\overline{e_s}.
\]
Stattdessen gilt die \textbf{Anti\-Homomorphie\-Eigenschaft}
\[
\overline{xy} = \bar y\,\bar x
\qquad(x,y\in\Cl_n).
\]

\textbf{Bemerkung.}
Für \(n=1\) bzw. \(n=2\) stimmt diese Konjugation mit den bekannten Konjugationen
in den Körpern \(\C\) bzw. \(\QU\) überein:
\[
\overline{x+yi}=x-yi, \qquad
\overline{a+bi+cj+dk}=a-bi-cj-dk.
\]
\end{DEF}



Second there is the canonical automorphism $\alpha: \mathrm{Cl}_{n} \longrightarrow \mathrm{Cl}_{n}$ defined in Example 5.6. We can use $\alpha$ to define a $\pm$-grading on $\mathrm{Cl}_{n}$ :
$$
\mathrm{Cl}_{n}^{+}=\left\{u \in \mathrm{Cl}_{n}: \alpha(u)=u\right\}, \quad \mathrm{Cl}_{n}^{-}=\left\{u \in \mathrm{Cl}_{n}: \alpha(u)=-u\right\} .
$$


\begin{DEF}{SG-B02-01-26}{Kanonische Automorphismus und $\pm$-Graduierung der Clifford-Algebra}
Sei 
\[
\alpha:\Cl_n\longrightarrow\Cl_n
\]
der \emph{kanonische Automorphismus} der Clifford-Algebra (vgl. Example 5.6), der auf den Erzeugern durch
\[
\alpha(e_r)=-e_r, \qquad r=1,\dots,n
\]
definiert ist und somit für ein Monom \(e_{i_1}\cdots e_{i_k}\)
\[
\alpha(e_{i_1}\cdots e_{i_k})=(-1)^k\,e_{i_1}\cdots e_{i_k}
\]
liefert.

Damit induziert $\alpha$ eine $\mathbb{Z}_2$-Graduierung der Clifford-Algebra in gerade und ungerade Teile:
\[
\Cl_n^+ := \{\,u\in\Cl_n \mid \alpha(u)=u\,\},
\qquad
\Cl_n^- := \{\,u\in\Cl_n \mid \alpha(u)=-u\,\}.
\]
\textbf{Interpretation:}
\begin{itemize}
  \item $\Cl_n^+$ besteht aus allen Linearkombinationen von Monomen gerader Länge;
  \item $\Cl_n^-$ besteht aus allen Linearkombinationen von Monomen ungerader Länge.
\end{itemize}
Diese Zerlegung erfüllt
\[
\Cl_n = \Cl_n^+ \oplus \Cl_n^-,
\qquad
\Cl_n^+\Cl_n^+\subseteq\Cl_n^+,\quad
\Cl_n^+\Cl_n^-\subseteq\Cl_n^-,\quad
\Cl_n^-\Cl_n^-\subseteq\Cl_n^+.
\]
\end{DEF}



Proposition 5.8

i) Every element $v \in \mathrm{Cl}_{n}$ can be uniquely expressed as $v=v^{+}+v^{-}$with $v^{+} \in \mathrm{Cl}_{n}^{+}$and $v^{-} \in \mathrm{Cl}_{n}^{-}$. Hence $\mathrm{Cl}_{n}=\mathrm{Cl}_{n}^{+} \oplus \mathrm{Cl}_{n}^{-}$as an $\mathbb{R}$-vector space.

ii) This decomposition is multiplicative in the sense that
$$
\begin{cases}u v \in \mathrm{Cl}_{n}^{+} & \text {if } u, v \in \mathrm{Cl}_{n}^{+} \text {or } u, v \in \mathrm{Cl}_{n}^{-} \\ u v, v u \in \mathrm{Cl}_{n}^{+} & \text {if } u \in \mathrm{Cl}_{n}^{+} \text {and } v \in \mathrm{Cl}_{n}^{-}\end{cases}
$$



\begin{DEF}{SG-B02-01-27}{$\pm$-Zerlegung und Multiplikationsregeln in der Clifford-Algebra}
\leavevmode

\textbf{(i) Additive Zerlegung.}
Jedes Element \(v\in\Cl_n\) lässt sich eindeutig schreiben als
\[
v = v^+ + v^-,
\qquad
v^+ \in \Cl_n^+, \quad v^- \in \Cl_n^-,
\]
wobei
\[
v^{\pm} = \tfrac{1}{2}\bigl(v \pm \alpha(v)\bigr)
\]
die Projektionsanteile bezüglich des kanonischen Automorphismus
\(\alpha(e_r)=-e_r\) sind.  
Damit gilt die direkte Summenzerlegung
\[
\Cl_n = \Cl_n^+ \oplus \Cl_n^-
\]
als $\R$-Vektorraum.

\medskip
\textbf{(ii) Multiplikationsregeln.}
Die $\pm$-Graduierung ist \emph{multiplikativ} im Sinne von:
\[
\begin{cases}
u v \in \Cl_n^+ & \text{falls } u,v\in\Cl_n^+ \text{ oder } u,v\in\Cl_n^-,\\[4pt]
u v,\; v u \in \Cl_n^- & \text{falls } u\in\Cl_n^+,\, v\in\Cl_n^-.
\end{cases}
\]
Insbesondere ist $\Cl_n^+$ eine Unteralgebra von $\Cl_n$, während
$\Cl_n^-$ ein $\Cl_n^+$\nobreakdash-Modul ist.


(i) The elements
$$
v^{+}=\frac{1}{2}(v+\alpha(v)), \quad v^{-}=\frac{1}{2}(v-\alpha(v))
$$
satisfy $\alpha\left(v^{+}\right)=v^{+}, \alpha\left(v^{-}\right)=-v^{-}$and $v=v^{+}+v^{-}$, hence $v^{+} \in \mathrm{Cl}_{n}^{+}$and $v^{-} \in \mathrm{Cl}_{n}^{-}$. On the other hand, if $v=v^{\prime}+v^{\prime \prime}$ with $v^{\prime} \in \mathrm{Cl}_{n}^{+}$and $v^{\prime \prime} \in \mathrm{Cl}_{n}^{-}$, then $\alpha(v)=v^{\prime}-v^{\prime \prime}$ and so
$$
\frac{1}{2}(v+\alpha(v))=v^{\prime}, \quad \frac{1}{2}(v-\alpha(v))=v^{\prime \prime} .
$$
So this expression is the only one with such properties and therefore defines the stated vector space direct sum decomposition.

(ii) This is easily checked using the fact that $\alpha$ is a ring homomorphism.
\end{DEF}



Notice that for bases of $\mathrm{Cl}_{n}^{ \pm}$we have the monomials
$$
\left\{\begin{array}{cl}
e_{j_{1}} \cdots e_{j_{2 m}} \in \mathrm{Cl}_{n}^{+} & \left(1 \leqslant j_{1}<\cdots<j_{2 m} \leqslant n\right), \\
e_{j_{1}} \cdots e_{j_{2 m+1}} \in \mathrm{Cl}_{n}^{-} & \left(1 \leqslant j_{1}<\cdots<j_{2 m+1} \leqslant n\right) .
\end{array}\right.
$$



\begin{DEF}{SG-B02-01-28}{Basen der geraden und ungeraden Komponenten}
Die $\mathbb{Z}_2$-Graduierung durch den kanonischen Automorphismus
\(\alpha(e_r)=-e_r\) induziert eine natürliche Basisaufteilung der Clifford-Algebra.
Für die Unterräume \(\Cl_n^\pm\) gilt:

\[
\begin{cases}
e_{j_1}\cdots e_{j_{2m}} \in \Cl_n^+,
& (1\le j_1<\cdots<j_{2m}\le n),\\[4pt]
e_{j_1}\cdots e_{j_{2m+1}} \in \Cl_n^-,
& (1\le j_1<\cdots<j_{2m+1}\le n).
\end{cases}
\]

Das heißt:
\begin{itemize}
  \item \(\Cl_n^+\) wird von allen Monomen \emph{gerader Länge} aufgespannt,
  \item \(\Cl_n^-\) wird von allen Monomen \emph{ungerader Länge} aufgespannt.
\end{itemize}
Insbesondere bildet die Menge aller geordneten Produkte
\(e_{j_1}\cdots e_{j_k}\) mit \(1\le j_1<\dots<j_k\le n\)
eine Orthonormalbasis von \(\Cl_n\),
und die Parität \(k\bmod 2\) bestimmt die Zugehörigkeit zu
\(\Cl_n^+\) bzw. \(\Cl_n^-\).
\end{DEF}


For later use, it is useful to record the following identity which holds for all $u \in \mathrm{Cl}_{n}$ :
$$
\alpha(\bar{u})=\overline{\alpha(u)}
$$
This is verified by checking that it holds on the monomial basis in the $e_{i}$.


\begin{LEM}{SG-B02-01-29}{Kommutationsrelation zwischen kanonischer Automorphismus und Konjugation}
Für alle $u \in \Cl_n$ gilt die Identität
\[
\alpha(\bar{u}) = \overline{\alpha(u)}.
\]
\end{LEM}


\begin{PROOF}{SG-B02-01-30}{P: Kommutationsrelation zwischen kanonischer Automorphismus und Konjugation}
Da sowohl $\alpha$ als auch die Konjugation $\overline{(\cdot)}$ durch ihre Wirkung auf die Basis-Monomiale $e_{i_1}\cdots e_{i_k}$ bestimmt sind, genügt es, die Gleichung für diese zu überprüfen.  
Für $u = e_{i_1}\cdots e_{i_k}$ gilt:
\[
\alpha(u) = (-1)^k e_{i_1}\cdots e_{i_k},
\qquad
\bar{u} = (-1)^k e_{i_k}\cdots e_{i_1}.
\]
Daraus folgt unmittelbar
\[
\alpha(\bar{u}) = (-1)^k \bar{u} = (-1)^k (-1)^k e_{i_k}\cdots e_{i_1}
= \overline{(-1)^k e_{i_1}\cdots e_{i_k}}
= \overline{\alpha(u)}.
\]
Somit kommutieren $\alpha$ und $\overline{(\cdot)}$ miteinander.
\end{PROOF}



Remark 5.9


It is worth noting that the composite of $\alpha$ and $\overline{()}$ (in either order) is another anti-homomorphism of $\mathrm{Cl}_{n}$; then Id, $\alpha, \overline{()}$ and $\bar{\alpha}=\alpha \circ \overline{()}$ form a finite group of order 4 which is not cyclic since the square of every element is the identity.


\begin{DEF}{SG-B02-01-31}{Gruppe der Involutionen in der Clifford-Algebra}
Die beiden Abbildungen, der kanonische Automorphismus \(\alpha\) und die Konjugation \(\overline{(\cdot)}\),
erzeugen zusammen mit der Identität weitere Struktursymmetrien der Clifford-Algebra.

\textbf{Beobachtung.}  
Die Kompositionen \(\alpha\circ\overline{(\cdot)}\) und \(\overline{(\cdot)}\circ\alpha\)
fallen zusammen (vgl. Lemma~\ref{SG-B02-01-29}) und definieren eine weitere
\emph{Anti\-Homomorphie}
\[
\bar{\alpha} := \alpha \circ \overline{(\cdot)} = \overline{(\cdot)} \circ \alpha : \Cl_n \to \Cl_n.
\]
Somit bilden die vier Abbildungen
\[
\{\mathrm{Id},\, \alpha,\, \overline{(\cdot)},\, \bar{\alpha}\}
\]
eine endliche Gruppe der Ordnung \(4\) unter Komposition.

\textbf{Struktureigenschaft.}  
Diese Gruppe ist nicht zyklisch, da jedes ihrer Elemente eine \emph{Involution} ist, d.\,h.
\[
\alpha^2 = \overline{(\cdot)}^2 = \bar{\alpha}^2 = \mathrm{Id}.
\]
Damit ist die Struktur isomorph zur \emph{Kleinschen Vierergruppe} \(V_4 \cong \Z_2 \times \Z_2\).
\end{DEF}


Finally, we introduce an inner product $\cdot$ and a norm $\left|\mid\right.$ on $\mathrm{Cl}_{n}$ by defining the distinct monomials $e_{i_{1}} e_{i_{2}} \cdots e_{i_{k}}$ with $1 \leqslant i_{1}<i_{2}<\cdots<i_{k} \leqslant n$ to be an orthonormal basis, i.e.,
$$
\left(e_{i_{1}} e_{i_{2}} \cdots e_{i_{k}}\right) \cdot\left(e_{j_{1}} e_{j_{2}} \cdots e_{j_{\ell}}\right)= \begin{cases}1 & \text { if } \ell=k \text { and } i_{r}=j_{r} \text { for all } r \\ 0 & \text { otherwise }\end{cases}
$$
and
$$
|x|=\sqrt{x \cdot x}
$$
We leave it to the reader to check that || is a norm. Perhaps a more illuminating way to define the inner product $\cdot$ is by using the formula
$$
u \cdot v=\frac{1}{2} \operatorname{Re}(\bar{u} v+\bar{v} u)
$$
where for $w \in \mathrm{Cl}_{n}$ we define its real part $\operatorname{Re} w$ to be the coefficient of 1 when $w$ is expanded as an $\mathbb{R}$-linear combination of the basis monomials $e_{i_{1}} \cdots e_{i_{r}}$. It is easily verified that for any $u, v \in \mathrm{Cl}_{n}$ and $w \in j_{n} \mathbb{R}^{n}$,
$$
(w u) \cdot(w v)=|w|^{2}(u \cdot v)
$$

In particular, when $|w|=1$ left multiplication by $w$ defines an $\mathbb{R}$-linear transformation on $\mathrm{Cl}_{n}$ which is an isometry. 




\begin{DEF}{SG-B02-01-32}{Skalarprodukt und Norm auf der Clifford-Algebra}
Auf der Clifford-Algebra $\Cl_n$ wird ein Skalarprodukt
\(\cdot : \Cl_n \times \Cl_n \to \R\)
durch die Forderung definiert, dass die Basis-Monomiale
\[
e_{i_1} e_{i_2} \cdots e_{i_k}, 
\qquad 1 \le i_1 < i_2 < \dots < i_k \le n,
\]
\textbf{orthonormal} sind, d.\,h.
\[
(e_{i_1}\cdots e_{i_k}) \cdot (e_{j_1}\cdots e_{j_\ell})
=
\begin{cases}
1, & \text{falls } k=\ell \text{ und } i_r=j_r \text{ für alle } r,\\[3pt]
0, & \text{sonst.}
\end{cases}
\]
Die zugehörige Norm ist gegeben durch
\[
\lVert x\rVert := \sqrt{x \cdot x}.
\]
Dies definiert eine reelle Norm auf $\Cl_n$.

Eine äquivalente und konzeptionell aufschlussreichere Formulierung des Skalarprodukts lautet
\[
u \cdot v = \frac{1}{2}\,\Re(\bar{u}v + \bar{v}u),
\]
wobei für $w \in \Cl_n$ der \emph{Realteil} $\Re(w)$ der Koeffizient von $1$
in der Basisentwicklung von $w$ ist, d.\,h.
\[
w = a_0 \cdot 1 + \sum_I a_I e_I 
\quad\Longrightarrow\quad 
\Re(w) = a_0.
\]
Diese Definition stimmt mit der orthonormalen Basisdefinition überein und
macht die Invarianz des Skalarprodukts gegenüber der Clifford-Struktur
sichtbar.
\end{DEF}




\begin{PROP}{SG-B02-01-33}{Isometrieeigenschaft der Linksmultiplikation}
Für alle $u,v\in \Cl_n$ und jedes $w\in j_n(\R^n)$ gilt
\[
(wu)\cdot(wv)=\lvert w\rvert^2\,(u\cdot v).
\]
Insbesondere ist, falls $\lvert w\rvert=1$, die Linksmultiplikation
\[
L_w:\Cl_n\to\Cl_n,\qquad x\mapsto wx,
\]
eine \emph{$\R$-lineare Isometrie} bezüglich des Clifford-Skalarprodukts.
\end{PROP}




The norm | | gives rise to a metric on $\mathrm{Cl}_{n}$ which makes the group of units $\mathrm{Cl}_{n}^{\times}$into a topological group and the above embeddings of $\mathrm{Cl}_{n}$ into matrix rings are then continuous, so $\mathrm{Cl}_{n}^{\times}$is a matrix group. 

Unfortunately, these embeddings are not norm preserving in general. For example, $2+e_{1} e_{2} e_{3} \in \mathrm{Cl}_{3}$ has $\left|2+e_{1} e_{2} e_{3}\right|=\sqrt{5}$, but the corresponding matrix in $\mathrm{M}_{8}(\mathbb{R})$ has norm $\sqrt{3}$. 

However, by defining for each $w \in \mathrm{Cl}_{n}$
$$
\|w\|=\sup \left\{|w x|: x \in \mathrm{Cl}_{n},|x|=1\right\}=\max \left\{|w x|: x \in \mathrm{Cl}_{n},|x|=1\right\}
$$
we obtain another equivalent norm on $\mathrm{Cl}_{n}$ for which the above embedding $\mathrm{Cl}_{n} \longrightarrow \mathrm{M}_{2^{n}}(\mathbb{R})$ does preserve norms. For $w \in j_{n} \mathbb{R}^{n}$ we do have $\|w\|=|w|$ and more generally, for $w_{1}, \ldots, w_{k} \in j_{n} \mathbb{R}^{n}$,
$$
\left\|w_{1} \cdots w_{k}\right\|=\left|w_{1} \cdots w_{k}\right|=\left|w_{1}\right| \cdots\left|w_{k}\right| .
$$
For $x, y \in \mathrm{Cl}_{n}$,
$$
\|x y\| \leqslant\|x\|\|y\|
$$
without equality in general.




\begin{DEF}{SG-B02-01-34}{Topologische Struktur und Operatornorm auf der Clifford-Algebra}
Die durch das Clifford-Skalarprodukt induzierte Norm $|\cdot|$ definiert eine Metrik auf $\Cl_n$, durch die $\Cl_n$ ein metrischer Raum wird.  
Die Menge der Einheiten
\[
\Cl_n^\times := \{u \in \Cl_n : \exists v \in \Cl_n, \, uv = vu = 1\}
\]
ist damit eine \emph{topologische Gruppe}, und die natürlichen Einbettungen
\[
\Cl_n \hookrightarrow \Mat_{2^n}(\R)
\]
sind stetige Algebra-Homomorphismen. Somit ist $\Cl_n^\times$ eine \emph{Matrixgruppe}.
\end{DEF}




\begin{EXA}{SG-B02-01-35}{Nicht-Isometrie der Standardeinbettung $\Cl_n \hookrightarrow \Mat_{2^n}(\R)$}
Die kanonische Einbettung $\Cl_n \hookrightarrow \Mat_{2^n}(\R)$ ist im Allgemeinen nicht normerhaltend.

\textbf{Beispiel:}  
In $\Cl_3$ gilt für das Element $2 + e_1 e_2 e_3$
\[
\lvert 2 + e_1 e_2 e_3\rvert = \sqrt{5},
\]
während die zugehörige Matrixdarstellung in $\Mat_8(\R)$ die Norm $\sqrt{3}$ besitzt.
\end{EXA}




\begin{DEF}{SG-B02-01-36}{Operatornorm auf der Clifford-Algebra (Normerhaltende kan. Einbettung)}
Zur Herstellung der Normerhaltung definiert man für jedes $w \in \Cl_n$ die \emph{Operatornorm}
\[
\lVert w\rVert := \sup\{\,|wx| : x \in \Cl_n,\ |x| = 1\,\} 
= \max\{\,|wx| : x \in \Cl_n,\ |x| = 1\,\}.
\]
Diese Norm ist zur Standardnorm $|\cdot|$ äquivalent, und die kanonische Einbettung
\[
\Cl_n \longrightarrow \Mat_{2^n}(\R)
\]
ist bezüglich $\lVert\cdot\rVert$ normerhaltend.
\end{DEF}




\begin{PROP}{SG-B02-01-37}{Eigenschaften der Operatornorm (Normerhaltend)}
Für alle $x,y \in \Cl_n$ gilt:
\begin{enumerate}[label=(\roman*)]
\item $\|xy\| \le \|x\|\,\|y\|$ (Submultiplikativität);
\item Für $w \in j_n(\R^n)$ gilt $\|w\| = |w|$;
\item Für $w_1,\ldots,w_k \in j_n(\R^n)$ gilt
\[
\|w_1 \cdots w_k\| = |w_1 \cdots w_k| = |w_1|\cdots|w_k|.
\]
\end{enumerate}
\end{PROP}





\pagebreak

\subsection{Clifford Groups}


Using the injective linear transformation $j_{n}: \mathbb{R}^{n} \longrightarrow \mathrm{Cl}_{n}$ we can identify $\mathbb{R}^{n}$ with a subspace of $\mathrm{Cl}_{n}$, i.e.,

$$
\sum_{r=1}^{n} x_{r} \mathbf{e}_{r} \longleftrightarrow j_{n}\left(\sum_{r=1}^{n} x_{r} \mathbf{e}_{r}\right)=\sum_{r=1}^{n} x_{r} e_{r} .
$$

From now on we do this without further comment, writing elements as $x \in \mathbb{R}^{n}$ rather than $\mathbf{x}$.



\begin{DEF}{SG-B02-02-01}{Einbettung des euklidischen Raumes in die Clifford-Algebra}
Die kanonische injektive lineare Abbildung
\[
j_n : \R^n \longrightarrow \Cl_n, 
\qquad 
j_n\!\left(\sum_{r=1}^n x_r \mathbf{e}_r\right) = \sum_{r=1}^n x_r e_r,
\]
ermöglicht es, den euklidischen Raum $\R^n$ als \emph{linearen Unterraum} der Clifford-Algebra $\Cl_n$ aufzufassen.

Somit können wir jedes $x \in \R^n$ mit seinem Bild in $\Cl_n$ identifizieren, d.\,h.
\[
x = \sum_{r=1}^n x_r \mathbf{e}_r
\quad\longleftrightarrow\quad
j_n(x) = \sum_{r=1}^n x_r e_r,
\]
und wir schreiben im Folgenden der Einfachheit halber direkt $x \in \R^n \subset \Cl_n$.
\end{DEF}



Notice that $\mathbb{R}^{n} \subseteq \mathrm{Cl}_{n}^{-}$, so for $x \in \mathbb{R}^{n}, u \in \mathrm{Cl}_{n}^{+}$and $v \in \mathrm{Cl}_{n}^{-}$,
$$
x u, u x \in \mathrm{Cl}_{n}^{-}, \quad x v, v x \in \mathrm{Cl}_{n}^{+} .
$$




\begin{PROP}{SG-B02-02-02}{\(\R^n\subset \Cl_n^{-}\) und Paritätsregeln}
Mit der Grad-Involution \(\alpha:\Cl_n\to\Cl_n\) (gegeben durch \(\alpha(e_i)=-e_i\) und multiplicativ fortgesetzt) gilt:
\begin{enumerate}
\item \(\R^n\subset \Cl_n^{-}\), d.\,h. für jeden Vektor \(x\in \R^n\) ist \(\alpha(x)=-x\).
\item Für \(u\in \Cl_n^{+}\) (gerader Teil) und \(v\in \Cl_n^{-}\) (ungerader Teil) gilt
\[
xu,\,ux\in \Cl_n^{-}
\quad\text{und}\quad
xv,\,vx\in \Cl_n^{+}
\qquad(x\in \R^n).
\]
\end{enumerate}
\end{PROP}



\begin{PROOF}{SG-B02-02-03}{P: \(\R^n\subset \Cl_n^{-}\) und Paritätsregeln}
Da \(\alpha(e_i)=-e_i\) für alle Erzeuger und \(\R^n=\mathrm{span}_\R\{e_1,\dots,e_n\}\), folgt für jedes
\(x\in \R^n\) die Relation \(\alpha(x)=-x\). Somit \(x\in \Cl_n^{-}\) und (1) ist gezeigt.

Für (2) benutze \(\alpha\) als Algebra-Automorphismus und die Charakterisierung
\(\Cl_n^{\pm}=\{a\in\Cl_n\mid \alpha(a)=\pm a\}\).
Ist \(u\in \Cl_n^{+}\), so \(\alpha(u)=u\), und für \(x\in\R^n\subset\Cl_n^{-}\) gilt
\[
\alpha(xu)=\alpha(x)\alpha(u)=(-x)u=-(xu),
\]
also \(xu\in \Cl_n^{-}\); analog \(ux\in \Cl_n^{-}\).
Ist \(v\in \Cl_n^{-}\), so \(\alpha(v)=-v\), und damit
\[
\alpha(xv)=\alpha(x)\alpha(v)=(-x)(-v)=xv,
\]
also \(xv\in \Cl_n^{+}\); analog \(vx\in \Cl_n^{+}\).
\end{PROOF}


Definition 5.10


For $n \geqslant 1$, the Clifford group $\Gamma_{n}$ is the subgroup
$$
\Gamma_{n}=\left\{u \in \mathrm{Cl}_{n}^{\times}: \forall x \in \mathbb{R}^{n}, \alpha(u) x u^{-1} \in \mathbb{R}^{n}\right\} \leqslant \mathrm{Cl}_{n}^{\times} .
$$


\begin{DEF}{SG-B02-02-04}{Clifford-Gruppe}
Für \(n \geq 1\) sei \(\Cl_n^{\times}\) die Gruppe der invertierbaren Elemente der Clifford-Algebra \(\Cl_n\).  
Dann heißt
\[
\Gamma_n
:= 
\left\{
u \in \Cl_n^{\times} \;\middle|\;
\forall\, x \in \R^n:\;
\alpha(u)\, x\, u^{-1} \in \R^n
\right\}
\subseteq \Cl_n^{\times}
\]
die \emph{Clifford-Gruppe} (oder \emph{Clifford group}).

\medskip
\noindent
Dabei bezeichnet \(\alpha:\Cl_n \to \Cl_n\) die kanonische Involution
\(\alpha(e_i)=-e_i\), welche die \(\mathbb{Z}_2\)-Graduierung der Clifford-Algebra bestimmt.
\end{DEF}



Notice that $\mathbb{R}^{\times} \leqslant \Gamma_{n}$ is a normal subgroup and is the centre of $\Gamma_{n}$. Also for any non-zero $u \in \mathbb{R}^{n}$ and $x \in \mathbb{R}^{n}$ it is easily seen that $\alpha(u) x u^{-1} \in \mathbb{R}^{n}$, hence $u \in \Gamma_{n}$.



\begin{PROP}{SG-B02-02-05}{Zentrum und Untergruppen der Clifford-Gruppe}
Sei \(\Gamma_n\) die Clifford-Gruppe gemäß Definition \ref{SG-B02-01-42}. Dann gilt:
\begin{enumerate}
\item Die multiplikative Gruppe der reellen Zahlen \(\R^{\times}\) ist eine normale Untergruppe von \(\Gamma_n\)
und bildet zugleich das Zentrum \(Z(\Gamma_n)\) von \(\Gamma_n\):
\[
\R^{\times} \;\trianglelefteq\; \Gamma_n,
\qquad
Z(\Gamma_n)=\R^{\times}.
\]
\item Für jedes von Null verschiedene \(u\in\R^n\) und jedes \(x\in\R^n\) gilt
\[
\alpha(u)\,x\,u^{-1}\in\R^n,
\]
also insbesondere \(u\in\Gamma_n\).
\end{enumerate}
\end{PROP}



\begin{PROOF}{SG-B02-02-06}{P: Zentrum und Untergruppen der Clifford-Gruppe}
\emph{Zu (1):}  
Für jedes \(r\in\R^{\times}\) und \(x\in\R^n\) gilt \(\alpha(r)=r\), also
\[
\alpha(r)x r^{-1}=r x r^{-1}=x\in\R^n.
\]
Damit ist \(\R^{\times}\subseteq\Gamma_n\).
Zudem kommutiert jedes \(r\in\R^{\times}\) mit allen Elementen von \(\Cl_n\),
insbesondere mit \(\Gamma_n\), also ist \(\R^{\times}\) zentral und normal.

\smallskip
\emph{Zu (2):}  
Für \(u\in\R^n\setminus\{0\}\) gilt \(u^2=-|u|^2\in\R^{\times}\),
also \(u^{-1}=-\frac{u}{|u|^2}\).
Da \(\alpha(u)=-u\), folgt für \(x\in\R^n\)
\[
\alpha(u)x u^{-1}=(-u)x\!\left(-\frac{u}{|u|^2}\right)=\frac{u x u}{|u|^2}.
\]
Das bekannte Identitätssystem der Clifford-Multiplikation impliziert,
dass \(u x u \in \R^n\), also \(\alpha(u)x u^{-1}\in\R^n\).  
Somit \(u\in\Gamma_n\).
\end{PROOF}



Proposition 5.11


\begin{PROP}{SG-B02-02-07}{Clifford Gruppe ist abgeschlossene Untergruppe von $\Cl_n^\times$}
$\Gamma_{n}$ is a closed subgroup of $\mathrm{Cl}_{n}^{\times}$.
\end{PROP}

Proof

\begin{PROOF}{SG-B02-02-08}{P: Clifford Gruppe ist abgeschlossene Untergruppe von $\Cl_n^\times$}
There is a continuous action
$$
\mathrm{Cl}_{n}^{\times} \times \mathrm{Cl}_{n} \longrightarrow \mathrm{Cl}_{n} ; \quad(u, v)=\alpha(u) v u^{-1} .
$$
Notice that for each $u \in \mathrm{Cl}_{n}^{\times}$, the function
$$
\mathrm{Cl}_{n} \longrightarrow \mathrm{Cl}_{n} ; \quad v \longmapsto \alpha(u) v u^{-1},
$$
is a linear isomorphism. Since $\mathbb{R}^{n} \subseteq \mathrm{Cl}_{n}$ is an $\mathbb{R}$-subspace, it is closed and so by Section 1.9, its stabiliser
$$
\operatorname{Stab}_{\mathrm{Cl}_{n}^{\times}}\left(\mathbb{R}^{n}\right) \leqslant \mathrm{Cl}_{n}^{\times}
$$
is a closed subgroup of $\mathrm{Cl}_{n}^{\times}$. But this stabiliser is clearly $\Gamma_{n}$.
\end{PROOF}





Corollary 5.12

\begin{KORO}{SG-B02-02-09}{Clifford Gruppe ist Matrixgruppe für $n\geq 1$}
For $n \geqslant 1, \Gamma_{n}$ is a matrix group.
\end{KORO}


\begin{PROOF}{SG-B02-02-10}{P: Clifford Gruppe ist Matrixgruppe für $n\geq 1$}
Notice that for each $u \in \mathrm{Cl}_{n}^{\times}$, the function
$$
\mathrm{Cl}_{n} \longrightarrow \mathrm{Cl}_{n} ; \quad v \longmapsto \alpha(u) v u^{-1},
$$
is a linear isomorphism. For each $u \in \mathrm{Cl}_{n}^{\times}$, we can restrict the function $\mathrm{Cl}_{n} \longrightarrow \mathrm{Cl}_{n}$ to
$$
\rho_{u}: \mathbb{R}^{n} \longrightarrow \mathbb{R}^{n} ; \quad \rho_{u}(v)=\alpha(u) v u^{-1}
$$
Since we already know that $\mathrm{Cl}_{n}^{\times}$ is a matrix group we have the corollary.
\end{PROOF}


Proposition 5.13

For $u \in \Gamma_{n}, \alpha(u)$ and $\bar{u}$ are also in $\Gamma_{n}$.



\begin{PROP}{SG-B02-02-11}{Stabilität der Clifford-Gruppe unter Involutionen}
Für jedes \(u \in \Gamma_n\) gilt:
\[
\alpha(u) \in \Gamma_n
\quad \text{und} \quad
\bar{u} \in \Gamma_n .
\]
\end{PROP}


Proof

\begin{PROOF}{SG-B02-02-12}{P: Stabilität der Clifford-Gruppe unter Involutionen}
For $u \in \Gamma_{n}$ and $x \in \mathbb{R}^{n}, \alpha(x)=-x=\bar{x} \in \mathbb{R}^{n}$. By Equations (5.8) and (5.5) together with the fact that $\alpha \circ \alpha=\mathrm{Id}$, we have
$$
\begin{aligned}
\alpha(\alpha(u)) x \alpha(u)^{-1} & =\alpha(\alpha(u)) \alpha(-x) \alpha(u)^{-1} \\
& =\alpha\left(\alpha(u)(-x) u^{-1}\right) \in \mathbb{R}^{n},
\end{aligned}
$$
and
$$
\begin{aligned}
\alpha(\bar{u}) x \bar{u}^{-1} & =\alpha(\bar{u}) \overline{(-x)} \overline{u^{-1}} \\
& =\overline{\alpha(u)} \overline{(-x)} \overline{u^{-1}} \\
& =\overline{u^{-1}(-x) \alpha(u)} \\
& =\overline{\alpha\left(\alpha\left(u^{-1}\right) \alpha(-x) u\right)} \in \mathbb{R}^{n}
\end{aligned}
$$
since $\mathbb{R}^{n}$ is closed under $\alpha, \overline{()}$ and $\rho_{\alpha\left(u^{-1}\right)}$.
\end{PROOF}



Proposition 5.14


For each $u \in \Gamma_{n}$, the function $\rho_{u}: \mathbb{R}^{n} \longrightarrow \mathbb{R}^{n}$ is an $\mathbb{R}$-linear isometry.

\begin{PROP}{SG-B02-02-13}{Adjungierte Wirkung der Clifford-Gruppe als Isometrie}
Für jedes \(u \in \Gamma_n\) definiert die Abbildung
\[
\rho_u : \R^n \longrightarrow \R^n,
\qquad
\rho_u(v) := \alpha(u)\, v\, u^{-1},
\]
eine \(\R\)-lineare Isometrie des euklidischen Raumes \((\R^n,|\cdot|)\).
\end{PROP}

Proof

\begin{PROOF}{SG-B02-02-14}{P: Adjungierte Wirkung der Clifford-Gruppe als Isometrie}
For any $x \in \mathbb{R}^{n}$,
$$
\begin{aligned}
\left|\rho_{u}(x)\right|^{2} & =-\rho_{u}(x)^{2} \\
& =-\left((-u) x u^{-1}\right)^{2} \\
& =-(-u) x u^{-1}(-u) x u^{-1} \\
& =-u x^{2} u^{-1} \\
& =u|x|^{2} u^{-1} \\
& =|x|^{2}
\end{aligned}
$$
Hence $\left|\rho_{u}(x)\right|=|x|$ for all $x \in \mathbb{R}^{n}$.
\end{PROOF}

This means that for $u \in \Gamma_{n}$, if we express $\rho_{u}$ in terms of the standard basis $\left\{e_{1}, \ldots, e_{n}\right\}$, then we have $\rho_{u} \in \mathrm{O}(n)$. Hence there is a group homomorphism
$$
\rho: \Gamma_{n} \longrightarrow \mathrm{O}(n) ; \quad \rho(u)=\rho_{u} .
$$
In fact, $\rho$ is also continuous. 


\begin{PROP}{SG-B02-02-15}{Darstellung der Clifford-Gruppe in der orthogonalen Gruppe}
Für jedes \(u \in \Gamma_n\) sei
\[
\rho_u : \R^n \longrightarrow \R^n, \qquad \rho_u(v) = \alpha(u)\, v\, u^{-1}.
\]
Dann gilt:

\begin{enumerate}[label=(\roman*)]
\item \(\rho_u \in \mathrm{O}(n)\), d.\,h. \(\rho_u\) ist eine orthogonale lineare Abbildung.
\item Die Zuordnung
\[
\rho : \Gamma_n \longrightarrow \mathrm{O}(n), \qquad u \longmapsto \rho_u,
\]
ist ein Gruppenhomomorphismus.
\item Diese Abbildung \(\rho\) ist stetig und damit ein stetiger Homomorphismus von Matrixgruppen.
\end{enumerate}
\end{PROP}



\begin{PROOF}{SG-B02-02-16}{P: Darstellung der Clifford-Gruppe in der orthogonalen Gruppe}
\textbf{(i)}  
Aus Proposition~\ref{SG-B02-01-47} folgt, dass jedes \(\rho_u\) eine \(\R\)-lineare Isometrie von \(\R^n\) ist, also \(\rho_u \in \mathrm{O}(n)\).

\textbf{(ii)}  
Für \(u,v \in \Gamma_n\) und \(x \in \R^n\) gilt:
\[
\rho_{uv}(x)
= \alpha(uv)\, x\, (uv)^{-1}
= \alpha(u)\,\alpha(v)\, x\, v^{-1}\, u^{-1}
= \alpha(u)\, \big(\alpha(v)\, x\, v^{-1}\big)\, u^{-1}
= \rho_u\big(\rho_v(x)\big).
\]
Daraus folgt \(\rho_{uv} = \rho_u \circ \rho_v\), also \(\rho\) ist ein Gruppenhomomorphismus.

\textbf{(iii)}  
Da die Multiplikation und Inversion in \(\Cl_n^\times\) stetig sind und die Abbildung \(u \mapsto \alpha(u)\) stetig ist, folgt aus
\[
\rho_u(v) = \alpha(u)\, v\, u^{-1}
\]
die Stetigkeit von \(\rho: \Gamma_n \to \mathrm{O}(n)\).
Somit ist \(\rho\) ein stetiger Homomorphismus von Matrixgruppen.
\end{PROOF}




We also need to identify the kernel of $\rho$.




Proposition 5.15


$$
\operatorname{ker} \rho=\mathbb{R}^{\times}=\{t 1: t \in \mathbb{R}, t \neq 0\}
$$



Proof

Suppose that $u \in \operatorname{ker} \rho$. Then for every $x \in \mathbb{R}^{n}, \alpha(u) x u^{-1}=x$, i.e., $\alpha(u) x=x u$. Writing $u=u^{+}+u^{-}$with $u^{ \pm} \in \mathrm{Cl}_{n}^{ \pm}$, we have
$$
\begin{aligned}
u^{+} x & =x u^{+} \\
-u^{-} x & =x u^{-}
\end{aligned}
$$
For each $r=1, \ldots, n$, we can write
$$
u^{+}=a_{r}^{+}+e_{r} b_{r}^{-}, \quad u^{-}=a_{r}^{-}+e_{r} b_{r}^{+}
$$
where $a_{r}^{ \pm}, b_{r}^{ \pm} \in \mathrm{Cl}_{n}^{ \pm}$do not involve $e_{r}$ in their expansions in terms of the monomial bases of Equation (5.7).

On taking $x=e_{r}$, by Equation (5.12a) together with the identities
$$
e_{r} a_{r}^{+}=a_{r}^{+} e_{r}, \quad e_{r} b_{r}^{-}=-b_{r}^{-} e_{r}
$$
we obtain
$$
a_{r}^{+} e_{r}+b_{r}^{-}=a_{r}^{+} e_{r}-b_{r}^{-}
$$
Now by comparing the parts involving or not involving $e_{r}$ we find that $b_{r}^{-}=0$. Similarly, Equation (5.12b) gives
$$
-a_{r}^{-} e_{r}+b_{r}^{+}=-a_{r}^{-} e_{r}-b_{r}^{+}
$$
from which we obtain $b_{r}^{+}=0$. So we have $u^{+}=a_{r}^{+}$and $u^{-}=a_{r}^{-}$, showing that neither of these involves $e_{r}$. But as this is true for all values of $r$, we must have $u^{-}=0$ and $u^{+}=t 1$ for some $t \in \mathbb{R}$, hence $u=t 1$.




\begin{PROP}{SG-B02-02-17}{Kern der orthogonalen Darstellung der Clifford-Gruppe}
Für den stetigen Gruppenhomomorphismus
\[
\rho : \Gamma_n \longrightarrow \mathrm{O}(n), \qquad 
\rho(u)(x) = \alpha(u)\, x\, u^{-1},
\]
gilt
\[
\ker(\rho) = \R^{\times} = \{\, t\,1 : t \in \R,\ t \neq 0 \,\}.
\]
\end{PROP}



\begin{PROOF}{SG-B02-02-18}{P: Kern der orthogonalen Darstellung der Clifford-Gruppe}
Sei \(u \in \ker(\rho)\). Dann gilt für alle \(x \in \R^n\):
\[
\rho_u(x) = \alpha(u)\,x\,u^{-1} = x.
\]
Multiplizieren wir diese Gleichung rechts mit \(u\), so folgt:
\[
\alpha(u)\,x = x\,u \qquad (\forall\,x \in \R^n).
\]

\textbf{(1) Fall: \(u \in \R^\times\).}  
Dann ist \(\alpha(u)=u\) und die obige Gleichung wird zur Identität
\(u x = x u\), also \(\rho_u = \mathrm{id}\).  
Damit gilt \(\R^\times \subseteq \ker(\rho)\).

\textbf{(2) Umkehrung.}  
Angenommen, \(u\) erfülle \(\alpha(u)\,x = x\,u\) für alle \(x \in \R^n\).
Schreibe \(u = u^+ + u^-\) gemäß der Zerlegung
\(\Cl_n = \Cl_n^+ \oplus \Cl_n^-\).
Da \(\alpha(u^+) = u^+\) und \(\alpha(u^-) = -u^-\), erhalten wir:
\[
(u^+ - u^-)\,x = x\,(u^+ + u^-),
\]
also
\[
u^+x - u^-x = x u^+ + x u^-.
\]
Da \(x \in \R^n \subset \Cl_n^-\) gilt \(x u^+ \in \Cl_n^-\) und \(x u^- \in \Cl_n^+\).
Durch Vergleich der \(\Cl_n^+\)- und \(\Cl_n^-\)-Komponenten folgt
\[
u^- = 0,
\]
und \(u^+\) kommutiert mit allen \(x \in \R^n\).
Nach der bekannten Struktur der Clifford-Algebra ist das Zentrum
\(Z(\Cl_n) = \R\), also \(u^+ = t\,1\) für ein \(t \in \R\).
Daher \(u = t\,1 \in \R^\times\).

\textbf{(3) Schluss.}  
Damit ist gezeigt:
\[
\ker(\rho) = \R^\times.
\]
\end{PROOF}



Proposition 5.16


For $u \in \Gamma_{n}, u \bar{u} \in \mathbb{R}^{\times}$and $\bar{u} u=u \bar{u}$. If $v \in \Gamma_{n}$ as well, then
$$
u v \overline{u v}=u \bar{u} v \bar{v}
$$


Proof

For $u \in \Gamma_{n}, u \bar{u} \in \Gamma_{n}$ by Proposition 5.13. We will show that $u \bar{u} \in \operatorname{ker} \rho$.

Let $x \in \mathbb{R}^{n}$. Then using the notation $\bar{\alpha}=\alpha \circ \overline{()}=\overline{()} \circ \alpha$ of Remark 5.9 we have
$$
\begin{aligned}
\rho_{u \bar{u}}(x) & =\alpha(u \bar{u}) x(u \bar{u})^{-1} \\
& =\alpha(u) \alpha(\bar{u}) x \overline{u^{-1}} u^{-1} \\
& =\alpha(u)\left(\alpha(\bar{u}) x \overline{u^{-1}}\right) u^{-1} \\
& =\alpha(u) \alpha\left(\bar{u} \bar{x} \alpha\left(\overline{u^{-1}}\right)\right) u^{-1}
\end{aligned}
$$
since $\alpha(\bar{x})=x$. Hence
$$
\begin{aligned}
\rho_{u \bar{u}}(x)= & \alpha(u) \alpha\left(\overline{\alpha\left(u^{-1}\right) x u}\right) u^{-1} \\
= & \alpha(u) \alpha\left(-\alpha\left(u^{-1}\right) x u\right) u^{-1} \\
& \quad\left[\text { since } \bar{y}=-y \text { for } y \in \mathbb{R}^{n}\right] \\
= & \alpha(u)\left(\alpha\left(u^{-1}\right) x u\right) u^{-1} \\
& \quad\left[\text { since } \alpha(y)=-y \text { for } y \in \mathbb{R}^{n}\right] \\
= & \alpha\left(u u^{-1}\right) x\left(u u^{-1}\right)=x
\end{aligned}
$$
Thus $\rho_{u \bar{u}}(x)=x$. We also have
$$
\bar{u} u=u^{-1} u \bar{u} u=u \bar{u} u^{-1} u
$$
since $u \bar{u} \in \mathbb{R}^{\times}$, hence $\bar{u} u=u \bar{u}$.\\
If $v \in \Gamma_{n}$, then
$$
u v \overline{u v}=u v \bar{v} \bar{u}=u \bar{u} v \bar{v},
$$
since $v \bar{v} \in \mathbb{R}^{\times}$and so commutes with every element of $\mathrm{Cl}_{n}$.



\begin{PROP}{SG-B02-02-19}{Norm-Element und Multiplikativität in Clifford Gruppe}
Für \(u \in \Gamma_n\) gilt \(u\bar u \in \R^{\times}\) und \(\bar u\,u = u\,\bar u\).
Sind außerdem \(v \in \Gamma_n\), so
\[
uv\,\overline{uv} \;=\; u\bar u \, v\bar v.
\]
\end{PROP}



\begin{PROOF}{SG-B02-02-20}{P: Norm-Element und Multiplikativität in Clifford Gruppe}
Zunächst liegt nach Proposition~\ref{SG-B02-02-11} auch \(\bar u \in \Gamma_n\). Somit ist \(u\bar u \in \Gamma_n\).
Wir zeigen, dass \(u\bar u \in \ker \rho\). Für \(x \in \R^n\) gilt mit \(\bar x = -x\) und \(\alpha(x)=-x\):
\begin{align*}
\rho_{u\bar u}(x)
&= \alpha(u\bar u)\, x\, (u\bar u)^{-1}
 = \alpha(u)\,\alpha(\bar u)\, x \, \overline{u^{-1}}\, u^{-1} \\
&= \alpha(u)\,\big(\alpha(\bar u)\, x \, \overline{u^{-1}}\big)\,u^{-1}
 = \alpha(u)\,\alpha\!\big(\bar u\,\bar x\, \alpha(\overline{u^{-1}})\big)\,u^{-1} \\
&= \alpha(u)\,\alpha\!\big(\bar u\,(-x)\,\overline{\alpha(u^{-1})}\big)\,u^{-1}
 = \alpha(u)\,\big(\alpha(u^{-1})\, x \, u\big)\,u^{-1}
 = x.
\end{align*}
Also \(\rho_{u\bar u}=\mathrm{id}\) und damit \(u\bar u \in \ker\rho\). Mit Proposition~\ref{SG-B02-02-17} folgt
\(u\bar u \in \R^{\times}\). Da reelle Skalare zentral sind, ist \(\bar u\,u = u\,\bar u\).

Für \(v \in \Gamma_n\) ergibt sich schließlich
\[
uv\,\overline{uv}
 = uv\,\bar v\,\bar u
 = u\,(v\bar v)\,\bar u
 = (v\bar v)\,u\bar u
 = u\bar u\, v\bar v,
\]
weil \(v\bar v \in \R^{\times}\) ebenso zentral ist wie \(u\bar u\).
\end{PROOF}



\begin{DEF}{SG-B02-02-21}{Normabbildung der Clifford Gruppe}
Die \emph{Normabbildung} der Clifford-Gruppe ist definiert durch
\[
\nu : \Gamma_n \longrightarrow \R^{\times}, 
\qquad 
\nu(u) := u\,\bar u.
\]
mit: Für jedes \(u \in \Gamma_n\) gilt \(\nu(u) = \bar u\,u \in \R^{\times}\).
\end{DEF}



Proposition 5.17


\begin{PROP}{SG-B02-02-22}{Normabbildung der Clifford Gruppe ist stetiger Gruppenhomomorphismus}
The function
$$
\nu: \Gamma_{n} \longrightarrow \mathbb{R}^{\times} ; \quad \nu(u)=u \bar{u}
$$
is a continuous group homomorphism.
\end{PROP}




Proof

\begin{PROOF}{SG-B02-02-23}{P: Normabbildung der Clifford Gruppe ist stetiger Gruppenhomomorphismus}
Notice that if $u \in \Gamma_{n}$,
$$
|u|^{2}=\operatorname{Re}(\bar{u} u)=\operatorname{Re}(u \bar{u})=u \bar{u}=\nu(u),
$$
hence $\nu$ is continuous and always takes positive values. Proposition 5.16 implies that it is a group homomorphism.
\end{PROOF} 




We will study the kernel of $\nu$ more thoroughly in the next section.


\pagebreak

\subsection{Pinor and Spinor Groups}


In this section we will describe for each $n \geqslant 1$, the compact connected spinor group $\operatorname{Spin}(n)$ which is a group of units in the Clifford algebra $\mathrm{Cl}_{n}$. Moreover, there is a surjective Lie homomorphism $\operatorname{Spin}(n) \longrightarrow \operatorname{SO}(n)$ whose kernel has two elements. We will first introduce the pinor group $\operatorname{Pin}(n)$ which has two connected components, one of them being $\operatorname{Spin}(n)$.

Definition 5.18

For $n \geqslant 1$, the pinor group $\operatorname{Pin}(n)$ is the kernel of $\nu: \Gamma_{n} \longrightarrow \mathbb{R}^{\times}$.

Notice that $\operatorname{Pin}(n)$ is a closed subgroup of $\Gamma_{n} \subseteq \mathrm{Cl}_{n}$ and is a bounded subset with respect to the metric induced from the norm ||, hence it is compact.


\begin{DEF}{SG-B02-03-01}{Pinor Gruppe}
Für \(n \ge 1\) sei die \emph{Pinor-Gruppe} definiert als der Kern der Normabbildung
\[
\operatorname{Pin}(n) := \ker\!\left(\nu : \Gamma_n \longrightarrow \R^{\times}\right)
= \{\,u \in \Gamma_n \mid \nu(u) = 1\,\}.
\]
\textbf{Bemerkung:}
\begin{itemize}
\item \(\operatorname{Pin}(n)\) ist eine abgeschlossene Untergruppe von \(\Gamma_n \subseteq \Cl_n\).
\item Mit der von der Norm \(\|\cdot\|\) induzierten Metrik ist \(\operatorname{Pin}(n)\) beschränkt
und somit kompakt.
\end{itemize}
\end{DEF}




Definition 5.19

For $n \geqslant 1$, the spinor group $\operatorname{Spin}(n)$ is
$$
\operatorname{Spin}(n)=\operatorname{Pin}(n) \cap \mathrm{Cl}_{n}^{+} \leqslant \operatorname{Pin}(n) .
$$
The restriction of $\alpha$ to a function $\alpha: \operatorname{Pin}(n) \longrightarrow \operatorname{Pin}(n)$ is a continuous group homomorphism for which
$$
\{u \in \operatorname{Pin}(n): \alpha(u)=u\}=\operatorname{Pin}(n) \cap \mathrm{Cl}_{n}^{+}
$$
hence $\operatorname{Spin}(n) \leqslant \operatorname{Pin}(n)$ is a closed subgroup. Later we will see that it is also a normal subgroup.


\begin{DEF}{SG-B02-03-02}{Spin Gruppe}
Für \(n \ge 1\) ist die \emph{Spinor-Gruppe} definiert durch
\[
\Spin(n) := \Pin(n) \cap \Cl_n^{+} \le \Pin(n).
\]
\textbf{Eigenschaften:}
\begin{enumerate}[label=(\roman*)]
\item Die Einschränkung der kanonischen Automorphismusabbildung
\[
\alpha : \Pin(n) \longrightarrow \Pin(n)
\]
ist ein stetiger Gruppenhomomorphismus.

\item Es gilt
\[
\{\,u \in \Pin(n) \mid \alpha(u) = u\,\} 
= \Pin(n) \cap \Cl_n^{+} = \Spin(n).
\]
\item Somit ist \(\Spin(n)\) eine abgeschlossene Untergruppe von \(\Pin(n)\).
Später wird gezeigt, dass \(\Spin(n)\) sogar eine \emph{normale Untergruppe} ist.
\end{enumerate}
\end{DEF}


Our principal goal in this section is to show that the restricted homomorphism $\rho: \operatorname{Pin}(n) \longrightarrow \mathrm{O}(n)$ is surjective and $\operatorname{Spin}(n)=\rho^{-1} \mathrm{SO}(n)$. We will do this by showing that $\operatorname{Pin}(n)$ is generated by certain elements $u \in \mathbb{R}^{n}$ for which $\rho_{u}$ is a reflection.

Within $\mathbb{R}^{n} \subseteq \mathrm{Cl}_{n}$ is the unit sphere

$$
\mathbb{S}^{n-1}=\left\{\mathbf{x} \in \mathbb{R}^{n}:|x|=1\right\}=\left\{\sum_{r=1}^{n} x_{r} e_{r}: \sum_{r=1}^{n} x_{r}^{2}=1\right\} .
$$



\begin{DEF}{SG-B02-03-03}{Einheitssphäre in der Clifford-Algebra}
In der Einbettung \(\R^n \subseteq \Cl_n\) ist die \emph{Einheitssphäre} definiert als
\[
\mathbb{S}^{n-1}
:= \{\,x \in \R^n : |x| = 1\,\}
= \left\{\,\sum_{r=1}^{n} x_r e_r \;\middle|\; \sum_{r=1}^{n} x_r^2 = 1\,\right\}.
\]
\textbf{Bemerkung:}
\(\mathbb{S}^{n-1}\) besteht also aus den Vektoren der Clifford-Algebra vom Betrag \(1\)
und bildet die Menge der \emph{Einheitsvektoren} in \(\R^n\).
\end{DEF}






Lemma 5.20


Let $u \in \mathbb{S}^{n-1} \subseteq \mathrm{Cl}_{n}$. Then $u$ is a unit in $\mathrm{Cl}_{n}$, i.e., $u \in \mathrm{Cl}_{n}^{\times}$, and $u^{-1} \in \mathbb{S}^{n-1}$.



\begin{LEM}{SG-B02-03-04}{Einheitsvektoren sind Einheiten}
Sei \(u \in \mathbb{S}^{n-1} \subseteq \Cl_n\). Dann gilt:
\begin{enumerate}[label=(\roman*)]
\item \(u\) ist eine \emph{Einheit} in der Clifford-Algebra, d.\,h.
      \(u \in \Cl_n^{\times}\).
\item Das Inverse von \(u\) ist gegeben durch
      \[
      u^{-1} = -u,
      \]
      und insbesondere gilt \(u^{-1} \in \mathbb{S}^{n-1}\).
\end{enumerate}
\end{LEM}


\begin{PROOF}{SG-B02-03-05}{P: Einheitsvektoren sind Einheiten}
Für \(u \in \R^n \subseteq \Cl_n\) gilt \(u^2 = -|u|^2 = -1\),
da \(|u| = 1\). Damit folgt
\[
u^{-1} = -u,
\]
und durch \(|u^{-1}| = |u| = 1\) erhalten wir \(u^{-1} \in \mathbb{S}^{n-1}\).
\end{PROOF}



Proof


Since $u \in \mathbb{R}^{n}$,
$$
(-u) u=u(-u)=-u^{2}=-\left(-|u|^{2}\right)=1,
$$
so $-u$ is the inverse of $u$. Of course $-u \in \mathbb{S}^{n-1}$. \(\square\)



More generally, for $u_{1}, \ldots, u_{k} \in \mathbb{S}^{n-1}$ we have
$$
\left(u_{1} \cdots u_{k}\right)^{-1}=(-1)^{k} u_{k} \cdots u_{1}=\overline{u_{1} \cdots u_{k}} .
$$


\begin{LEM}{SG-B02-03-06}{Inverse von Produkten von Einheitsvektoren}
Seien \(u_{1}, \ldots, u_{k} \in \mathbb{S}^{n-1} \subseteq \Cl_n\).
Dann gilt für ihr Produkt:
\[
\left(u_{1} \cdots u_{k}\right)^{-1}
= (-1)^{k}\, u_{k} \cdots u_{1}
= \overline{u_{1} \cdots u_{k}} .
\]
\end{LEM}



\begin{PROOF}{SG-B02-03-07}{P: Inverse von Produkten von Einheitsvektoren}
Da für jeden Einheitsvektor \(u_i\) gilt \(u_i^2 = -1\), folgt zunächst
\[
u_i^{-1} = -u_i.
\]
Für das Produkt \(u_1 \cdots u_k\) ergibt sich dann
\[
\left(u_{1} \cdots u_{k}\right)\left((-1)^{k} u_{k} \cdots u_{1}\right)
= (-1)^{k} u_{1} \cdots u_{k} u_{k} \cdots u_{1}
= (-1)^{k} (-1)^{k} = 1,
\]
woraus \(\left(u_{1} \cdots u_{k}\right)^{-1} = (-1)^{k} u_{k} \cdots u_{1}\) folgt.
Nach Definition der Clifford-Konjugation gilt zudem
\[
\overline{u_{1} \cdots u_{k}} = (-1)^{k} u_{k} \cdots u_{1},
\]
also stimmen beide Darstellungen überein.
\end{PROOF}




Recall that in a group $G$, the subgroup generated by the subset $S \subseteq G$ is the smallest subgroup of $G$ containing $S$, often denoted $\langle S\rangle \leqslant G$. A typical element of $\langle S\rangle$ is a product of elements of the form $s$ or $s^{-1}$ where $s \in S$, i.e.,
$$
s_{1}^{\varepsilon_{1}} \cdots s_{k}^{\varepsilon_{k}} \quad\left(\varepsilon_{1}, \ldots, \varepsilon_{k}= \pm 1\right) 
$$
Furthermore, if $H \leqslant G$ and $S \subseteq H$ then $\langle S\rangle \leqslant H$.



\begin{DEF}{SG-B02-03-08}{Erzeugte Untergruppe}
Sei \(G\) eine Gruppe und \(S \subseteq G\) eine Teilmenge.
Die von \(S\) \emph{erzeugte Untergruppe} ist die kleinste Untergruppe von \(G\),
die \(S\) enthält, und wird mit \(\langle S\rangle \le G\) bezeichnet.
Explizit gilt
\[
\langle S\rangle
=\Big\{\, s_1^{\varepsilon_1}\cdots s_k^{\varepsilon_k}
\;\Big|\; k\in\N,\; s_i\in S,\; \varepsilon_i\in\{\pm1\}\,\Big\}.
\]
Ist \(H \le G\) eine Untergruppe mit \(S\subseteq H\), so folgt \(\langle S\rangle \le H\).
\end{DEF}



\begin{CONC}{SG-B02-03-09}{Typische Elemente in erzeugten Untergruppen und Monotonie}
Ein typisches Element aus \(\langle S\rangle\) ist ein endliches Produkt
\(
s_1^{\varepsilon_1}\cdots s_k^{\varepsilon_k}
\) mit \(s_i\in S\) und \(\varepsilon_i=\pm1\).
Ferner gilt die Monotonie: aus \(S\subseteq H\le G\) folgt
\(\langle S\rangle\subseteq H\), da \(\langle S\rangle\) per Definition die
kleinste Untergruppe mit dieser Eigenschaft ist.
\end{CONC}



We will show that $\mathbb{S}^{n-1} \subseteq \operatorname{Pin}(n)$ and therefore $\left\langle\mathbb{S}^{n-1}\right\rangle \leqslant \operatorname{Pin}(n)$. We begin with the following useful result.



Lemma 5.21


Let $u, v \in \mathbb{R}^{n} \subseteq \mathrm{Cl}_{n}$. If $u \cdot v=0$, then $v u=-u v$.


\begin{LEM}{SG-B02-03-10}{Orthogonale Vektoren in $\Cl_n$ antikommutieren}
Seien \(u,v\in \R^n\subset \Cl_n\). Gilt \(u\cdot v=0\), so folgt
\[
vu=-uv.
\]
\end{LEM}


Proof

\begin{PROOF}{SG-B02-03-11}{P: Orthogonale Vektoren in $\Cl_n$ antikommutieren}
Writing $u=\sum_{r=1}^{n} x_{r} e_{r}$ and $v=\sum_{s=1}^{n} y_{s} e_{s}$ with $x_{r}, y_{s} \in \mathbb{R}$, we obtain
$$
\begin{aligned}
v u=\sum_{s=1}^{n} \sum_{r=1}^{n} y_{s} x_{r} e_{s} e_{r} & =\sum_{r=1}^{n} y_{r} x_{r} e_{r}^{2}+\sum_{r<s}\left(x_{s} y_{r}-x_{r} y_{s}\right) e_{r} e_{s} \\
& =-\sum_{r=1}^{n} y_{r} x_{r}-\sum_{r<s}\left(x_{r} y_{s}-x_{s} y_{r}\right) e_{r} e_{s} \\
& =-u . v-\sum_{r<s}\left(x_{r} y_{s}-x_{s} y_{r}\right) e_{r} e_{s} \\
& =-\sum_{r<s}\left(x_{r} y_{s}-x_{s} y_{r}\right) e_{r} e_{s} \\
& =v . u-\sum_{r<s}\left(x_{r} y_{s}-x_{s} y_{r}\right) e_{r} e_{s} \\
& =-\sum_{r=1}^{n} \sum_{s=1}^{n} x_{r} y_{s} e_{r} e_{s} \\
& =-u v
\end{aligned}
$$
\end{PROOF}


Notice that for $u \in \mathbb{S}^{n-1}$ and $x \in \mathbb{R}^{n}$, by Equation (5.13),
$$
\alpha(u) x u^{-1}=(-u) x(-u)=u x u .
$$
If $u \cdot x=0$, then Lemma 5.21 implies that
$$
\alpha(u) x u^{-1}=-u^{2} x=-(-1) x=x,
$$
since $u^{2}=-|u|^{2}=-1$. On the other hand, if $x=t u$ for some $t \in \mathbb{R}$, then
$$
\alpha(u) x u^{-1}=t u^{3}=-t u .
$$
So in particular $\alpha(u) x^{-1} \in \mathbb{R}^{n}$ which yields the following result.




\begin{LEM}{SG-B02-03-12}{Adjungierte Wirkung eines Einheitsvektors = Spiegelung}
Sei \(u\in \mathbb S^{n-1}\subset \R^n\subset \Cl_n\). Dann gilt für alle \(x\in\R^n\)
\[
\rho_u(x):=\alpha(u)\,x\,u^{-1}=u\,x\,u,
\]
und \(\rho_u\) ist die Spiegelung an der zu \(u\) orthogonalen Hyperbene:
\[
\rho_u(x)=
\begin{cases}
x,& \text{falls } u\cdot x=0,\\[2pt]
-x,& \text{falls } x=t\,u\text{ für ein }t\in\R.
\end{cases}
\]
Insbesondere \(\rho_u\in \mathrm O(n)\).
\end{LEM}



\begin{PROOF}{SG-B02-03-13}{P: Adjungierte Wirkung eines Einheitsvektors = Spiegelung}
Da \(\alpha(u)=-u\) und \(u^{-1}=-u\) (weil \(u^2=-|u|^2=-1\)) ist,
\[
\rho_u(x)=(-u)\,x\,(-u)=u\,x\,u.
\]
Ist \(u\cdot x=0\), so liefert Lemma~\ref{SG-B02-02-29} \(ux=-xu\) und damit
\[
\rho_u(x)=u\,x\,u=-(u^2)\,x=-(-1)\,x=x.
\]
Für \(x=t\,u\) folgt
\[
\rho_u(x)=t\,u^3 = t\,u(u^2)=t\,u(-1)=-t\,u=-x.
\]
Dies ist genau die Charakterisierung einer Spiegelung an \(u^\perp\). Damit ist \(\rho_u\) orthogonal.
\end{PROOF}



\begin{LEM}{SG-B02-03-14}{Einheitsvektoren liegen in \(\Pin(n)\)}
Für \(u\in \mathbb S^{n-1}\) gilt \(u\bar u=1\) und \(\alpha(u)\,x\,u^{-1}\in \R^n\) für alle \(x\in\R^n\).
Also \(\mathbb S^{n-1}\subset \ker\nu=\Pin(n)\).
\end{LEM}



\begin{PROOF}{SG-B02-03-15}{P: Einheitsvektoren liegen in \(\Pin(n)\)}
Aus der Proposition folgt \(\alpha(u)\,x\,u^{-1}\in\R^n\) für alle \(x\).
Zudem ist \(\bar u=-u\), also \(u\bar u=-u^2=|u|^2=1\). Damit \(u\in \Gamma_n\) und \(\nu(u)=1\), also \(u\in \Pin(n)\).
\end{PROOF}





Proposition 5.22


The subgroup of $\mathrm{Cl}_{n}^{\times}$generated by $\mathbb{S}^{n-1}$ is contained in $\operatorname{Pin}(n)$, i.e., $\left\langle\mathbb{S}^{n-1}\right\rangle \leqslant \operatorname{Pin}(n)$.


\begin{PROP}{SG-B02-03-16}{Von der Einheitssphäre erzeugte Untergruppe liegt in \(\Pin(n)\)}
Es gilt
\[
\langle \mathbb S^{n-1}\rangle \;\leqslant\; \Pin(n).
\]
\end{PROP}



\begin{PROOF}{SG-B02-03-17}{P: Von der Einheitssphäre erzeugte Untergruppe liegt in \(\Pin(n)\)}
Sei \(u\in \mathbb S^{n-1}\). Dann ist \(u\in\Gamma_n\) und \(\rho_u(x)=\alpha(u)xu^{-1}=u x u\in\R^n\) für alle \(x\in\R^n\), also wirkt \(u\) als Spiegelung (orthogonale Abbildung). Ferner gilt
\[
u\bar u = 1,
\]
da \(\bar u=-u\) und \(u^2=-|u|^2=-1\).

Für ein endliches Produkt \(w=u_1\cdots u_k\) mit \(u_i\in\mathbb S^{n-1}\) folgt aus der Multiplikativität des Norm-Elements (Prop.~\ref{SG-B02-02-19})
\[
w\,\overline{w}=u_1\bar u_1\cdots u_k\bar u_k=1,
\]
also \(\nu(w)=1\). Zudem ist \(\rho_w=\rho_{u_1}\circ\cdots\circ \rho_{u_k}\) eine Komposition orthogonaler Abbildungen, also \(\rho_w(\R^n)\subset \R^n\). Damit \(w\in\Gamma_n\) und \(\nu(w)=1\), also \(w\in\Pin(n)\).

Da ein typisches Element von \(\langle \mathbb S^{n-1}\rangle\) genau ein solches Produkt (mit Exponenten \(\pm1\)) ist und \(\mathbb S^{n-1}\) invers abgeschlossen ist (\(u^{-1}=-u\in\mathbb S^{n-1}\)), folgt insgesamt
\(\langle \mathbb S^{n-1}\rangle \subseteq \Pin(n)\).
\end{PROOF}



Remark 5.23

By Lemma 5.20, every element of $\left\langle\mathbb{S}^{n-1}\right\rangle$ is a product of elements of $\mathbb{S}^{n-1}$.


\begin{LEM}{SG-B02-03-18}{Erzeugung von \(\langle \mathbb S^{n-1}\rangle\)}
Jedes Element von \(\langle \mathbb S^{n-1}\rangle\) ist ein Produkt von endlich vielen Elementen aus \(\mathbb S^{n-1}\).
\end{LEM}



\begin{PROOF}{SG-B02-03-19}{P: Erzeugung von \(\langle \mathbb S^{n-1}\rangle\)}
Nach Lemma~\ref{SG-B02-02-29} gilt für \(u\in\mathbb S^{n-1}\), dass \(u^{-1}=-u\in \mathbb S^{n-1}\), d.\,h. die Sphäre ist invers abgeschlossen.  
Damit kann jedes Element der von \(\mathbb S^{n-1}\) erzeugten Untergruppe in der Form
\[
u_1 u_2 \cdots u_k, \qquad u_i\in\mathbb S^{n-1}, \ k\in\mathbb N,
\]
geschrieben werden.  
Dies folgt unmittelbar aus der allgemeinen Definition der von einer Teilmenge \(S\) einer Gruppe \(G\) erzeugten Untergruppe:
\[
\langle S\rangle = \{\, s_1^{\varepsilon_1}\cdots s_k^{\varepsilon_k} : s_i\in S,\ \varepsilon_i=\pm1,\ k\in\mathbb N\,\}.
\]
Da \(\mathbb S^{n-1}\) invers abgeschlossen ist, können alle Exponenten \(\varepsilon_i=-1\) durch Multiplikation weiterer Elemente aus \(\mathbb S^{n-1}\) ersetzt werden. Somit sind alle Elemente von \(\langle\mathbb S^{n-1}\rangle\) tatsächlich Produkte von Einheitsvektoren.
\end{PROOF}



Next we will investigate the restriction of the homomorphism $\rho: \operatorname{Pin}(n) \longrightarrow \mathrm{O}(n)$ to $\rho:\left\langle\mathbb{S}^{n-1}\right\rangle \longrightarrow \mathrm{O}(n)$. For each $u \in \mathbb{S}^{n-1}$, we have an $\mathbb{R}$-linear isometry
$$
\rho_{u}: \mathbb{R}^{n} \longrightarrow \mathbb{R}^{n} ; \quad \rho_{u}(x)=\alpha(u) x u^{-1}=u x u
$$
More generally, for $u \in\left\langle\mathbb{S}^{n-1}\right\rangle$, suppose that $u=u_{1} \cdots u_{r}$ for $u_{1}, \ldots, u_{r} \in \mathbb{S}^{n-1}$. Then
$$
\begin{aligned}
\rho_{u}(x) & =\alpha\left(u_{1} \cdots u_{r}\right) x\left(u_{1} \cdots u_{r}\right)^{-1} \\
& =\left((-1)^{r} u_{1} \cdots u_{r}\right) x\left((-1)^{r} u_{r} \cdots u_{1}\right) \\
& =u_{1} \cdots u_{r} x u_{r} \cdots u_{1}
\end{aligned}
$$



\begin{PROP}{SG-B02-03-20}{Beschränkung von \(\rho\) auf \(\langle\mathbb S^{n-1}\rangle\) und Produktformel und $\rho_u\in \Ogroup(n)$}
Für \(u\in\langle\mathbb S^{n-1}\rangle\) mit Darstellung \(u=u_1\cdots u_r\) (\(u_i\in\mathbb S^{n-1}\)) gilt für alle \(x\in\R^n\):
\[
\rho_u(x)=\alpha(u)\,x\,u^{-1}
= u_1\cdots u_r\,x\,u_r\cdots u_1.
\]
Insbesondere ist \(\rho_u\in \mathrm{O}(n)\).
\end{PROP}



\begin{PROOF}{SG-B02-03-21}{P: Beschränkung von \(\rho\) auf \(\langle\mathbb S^{n-1}\rangle\) und Produktformel und $\rho_u\in \Ogroup(n)$}
Da \(\alpha\) die Grad-Involution ist und \(\alpha(u_i)=-u_i\) für \(u_i\in\R^n\subset\Cl_n\), folgt
\[
\alpha(u)=\alpha(u_1\cdots u_r)=\alpha(u_1)\cdots\alpha(u_r)=(-1)^r u_1\cdots u_r.
\]
Ferner liefert Lemma~\ref{SG-B02-02-29} (Einheitsvektoren sind invertierbar mit \(u_i^{-1}=-u_i\)) die Inversenformel
\[
u^{-1}=(u_1\cdots u_r)^{-1}=(-1)^r\,u_r\cdots u_1.
\]
Zusammen ergibt dies für \(x\in\R^n\)
\[
\rho_u(x)=\alpha(u)\,x\,u^{-1}
= \big((-1)^r u_1\cdots u_r\big)\,x\,\big((-1)^r u_r\cdots u_1\big)
= u_1\cdots u_r\,x\,u_r\cdots u_1.
\]
Dass \(\rho_u\) orthogonal ist, folgt bereits aus der Isometrieaussage für \(\rho\) (vgl. Proposition~\ref{SG-B02-02-13}); daher \(\rho_u\in\mathrm{O}(n)\).
\end{PROOF}




Proposition 5.24

For $u \in \mathbb{S}^{n-1}, \rho_{u}: \mathbb{R}^{n} \longrightarrow \mathbb{R}^{n}$ is reflection in the hyperplane orthogonal to $u$.



\begin{PROP}{SG-B02-03-22}{Reflexion durch Einheitsvektor in der Clifford-Darstellung}
Für \(u \in \mathbb S^{n-1} \subset \R^n \subset \Cl_n\) ist die Abbildung
\[
\rho_u : \R^n \longrightarrow \R^n, 
\qquad
\rho_u(x) = \alpha(u)\,x\,u^{-1} = u\,x\,u,
\]
die Spiegelung an der Hyperebene orthogonal zu \(u\).
\end{PROP}



\begin{PROOF}{SG-B02-03-23}{P: Reflexion durch Einheitsvektor in der Clifford-Darstellung}
Da \(u\in\mathbb S^{n-1}\) gilt \(u^2=-1\) und \(u^{-1}=-u\).  
Für \(x\in\R^n\) können wir \(x\) in seine Komponenten parallel und orthogonal zu \(u\) zerlegen:
\[
x = (x\cdot u)u + x_\perp, \qquad x_\perp \cdot u = 0.
\]
Dann gilt nach Lemma~\ref{SG-B02-02-33} \(u\,x_\perp = -x_\perp\,u\). Somit:
\begin{align*}
\rho_u(x)
&= u\,x\,u
 = u\big((x\cdot u)u + x_\perp\big)u
 = (x\cdot u)\,u\,u\,u + u\,x_\perp\,u \\
&= (x\cdot u)\,u(-1) + (-x_\perp)u\,u
 = -(x\cdot u)u + x_\perp.
\end{align*}
Dies entspricht genau der Spiegelung an der Hyperebene \(u^\perp\):
\[
\rho_u(x) = x - 2(x\cdot u)\,u.
\]
Tatsächlich sind beide Formen äquivalent, da \(u\) Einheitslänge hat. Somit ist \(\rho_u\) die orthogonale Reflexion an der Ebene senkrecht zu \(u\).
\end{PROOF}



Proof


Recalling the defining equation (1.4) of a hyperplane reflection, Equations (5.15) show that
$$
\rho_{u}(x)=\left\{\begin{aligned}
x & \text { if } u \cdot x=0 \\
-x & \text { if } x=t u \text { for some } t \in \mathbb{R}
\end{aligned}\right.
$$
So $\rho_{u}: \mathbb{R}^{n} \longrightarrow \mathbb{R}^{n}$ is indeed reflection in the hyperplane orthogonal to $u$.




Proposition 5.25


i) $\rho:\left\langle\mathbb{S}^{n-1}\right\rangle \longrightarrow \mathrm{O}(n)$ is surjective and $\operatorname{ker} \rho=\{1,-1\}$.\\
ii) $\rho: \operatorname{Pin}(n) \longrightarrow \mathrm{O}(n)$ is surjective with $\operatorname{ker} \rho=\{1,-1\}$.\\
iii) $\left\langle\mathbb{S}^{n-1}\right\rangle=\operatorname{Pin}(n)$.



\begin{PROP}{SG-B02-03-24}{Hauptaussage zu \(\rho\) und \(\Pin(n)\)}
\begin{enumerate}[label=\textnormal{(\roman*)}]
\item \(\rho:\langle \mathbb S^{n-1}\rangle \to O(n)\) ist surjektiv und \(\ker\rho=\{1,-1\}\).
\item \(\rho:\Pin(n)\to O(n)\) ist surjektiv mit \(\ker\rho=\{1,-1\}\).
\item \(\langle \mathbb S^{n-1}\rangle=\Pin(n)\).
\end{enumerate}
\end{PROP}


Proof

(i) The observation in the proof of Proposition 5.24 shows that reflection in the hyperplane orthogonal to $u \in \mathbb{S}^{n-1}$ has the form $\rho_{u}$. Surjectivity is a consequence of Proposition 1.41.\\

(ii) If $t \in \operatorname{ker} \rho=\operatorname{Pin}(n) \cap \mathbb{R}^{\times}$, then $|t|=1$ then $|t|^{2}=t \bar{t}=t^{2}$, so $t= \pm 1$. On the other hand, $\pm 1 \in \operatorname{ker} \rho$.\\

(iii) Suppose that $v \in \operatorname{Pin}(n)$. Then $\rho_{v} \in \mathrm{O}(n)$ can be expressed as a product of hyperplane reflections, each of which has the form $\rho_{w}$ for some vector $w \in \mathbb{S}^{n-1}$. Hence $\rho_{v}=\rho_{w_{1} \cdots w_{k}}$ for suitable $w_{1}, \ldots, w_{k} \in \mathbb{S}^{n-1}$. But then
$$
(-1)^{k} v w_{k} \cdots w_{1} \in \operatorname{ker} \rho=\{ \pm 1\} .
$$
Thus $v= \pm w_{1} \cdots w_{k} \in\left\langle\mathbb{S}^{n-1}\right\rangle$.



\begin{PROOF}{SG-B02-03-25}{P: Hauptaussage zu \(\rho\) und \(\Pin(n)\)}
\emph{Vorbereitung.} Für \(u\in\mathbb S^{n-1}\subset\R^n\subset\Cl_n\) definiert
\[
\rho_u(x)=\alpha(u)\,x\,u^{-1}=u\,x\,u \qquad (x\in\R^n)
\]
eine orthogonale Abbildung; nach Prop.\,5.24 ist \(\rho_u\) die Spiegelung an \(u^\perp\).
Für ein Produkt \(g=u_1\cdots u_r\) mit \(u_i\in\mathbb S^{n-1}\) gilt
\[
\rho_g=\rho_{u_1}\circ\cdots\circ\rho_{u_r}.
\]
Damit ist \(\rho:\langle\mathbb S^{n-1}\rangle\to O(n)\) ein Gruppenhomomorphismus.

\smallskip
\emph{(i) Surjektivität und Kern auf \(\langle\mathbb S^{n-1}\rangle\).}
Nach dem Satz von Cartan–Dieudonné wird jede orthogonale Abbildung als Produkt
endlich vieler Spiegelungen dargestellt. Da \(\rho_{u}\) genau die Spiegelung ist,
ist \(\rho\) surjektiv.

Sei nun \(g\in\langle\mathbb S^{n-1}\rangle\) mit \(\rho_g=\id_{\R^n}\).
Dann gilt \(g\,x\,g^{-1}=x\) für alle \(x\in\R^n\), also \(gx=xg\) für alle Vektoren.
Daraus folgt, dass \(g\) im Zentrum von \(\Cl_n\) liegt. Das Zentrum von \(\Cl_n\)
ist (über \(\R\)) durch Skalare \(\R\cdot 1\) (und ggf. das Pseudoskalar-Element in
ungeraden Dimensionen) beschrieben; die einzigen zentralen Elemente, die zudem in
\(\langle\mathbb S^{n-1}\rangle\subset \Cl_n^\times\) liegen und die Normbedingung
\(\nu(g)=g\bar g=1\) erfüllen, sind \(g=\pm 1\).
Also \(\ker\rho=\{1,-1\}\).

The observation in the proof of Proposition 5.24 shows that reflection in the hyperplane orthogonal to $u \in \mathbb{S}^{n-1}$ has the form $\rho_{u}$. Surjectivity is a consequence of Proposition 1.41.

\smallskip
\emph{(ii) Surjektivität und Kern auf \(\Pin(n)\).}
Wir haben bereits \(\langle\mathbb S^{n-1}\rangle\le \Pin(n)\) gezeigt, also ist die
Beschränkung \(\rho:\Pin(n)\to O(n)\) surjektiv, weil (i) surjektiv ist.
Ist \(g\in\Pin(n)\) mit \(\rho_g=\id\), so vertauscht \(g\) mit allen
\(x\in\R^n\), also liegt \(g\) wiederum im Zentrum von \(\Cl_n\).
Aus der Normbedingung \(\nu(g)=1\) und \(g\in\Pin(n)\) folgt wie oben \(g=\pm 1\).
Also \(\ker\rho=\{1,-1\}\).

If $t \in \operatorname{ker} \rho=\operatorname{Pin}(n) \cap \mathbb{R}^{\times}$, then $|t|=1$ then $|t|^{2}=t \bar{t}=t^{2}$, so $t= \pm 1$. On the other hand, $\pm 1 \in \operatorname{ker} \rho$.

\smallskip
\emph{(iii) Gleichheit \(\langle\mathbb S^{n-1}\rangle=\Pin(n)\).}
Wir haben die Inklusion \(\langle\mathbb S^{n-1}\rangle\le \Pin(n)\).
Beide Gruppen werden durch \(\rho\) auf \(O(n)\) surjektiv abgebildet und besitzen
denselben Kern \(\{1,-1\}\). Somit induzieren beide dieselbe zweifache Überlagerung
von \(O(n)\). Wegen der Inklusion muss daher \(\langle\mathbb S^{n-1}\rangle=\Pin(n)\) gelten.

Suppose that $v \in \operatorname{Pin}(n)$. Then $\rho_{v} \in \mathrm{O}(n)$ can be expressed as a product of hyperplane reflections, each of which has the form $\rho_{w}$ for some vector $w \in \mathbb{S}^{n-1}$. Hence $\rho_{v}=\rho_{w_{1} \cdots w_{k}}$ for suitable $w_{1}, \ldots, w_{k} \in \mathbb{S}^{n-1}$. But then
$$
(-1)^{k} v w_{k} \cdots w_{1} \in \operatorname{ker} \rho=\{ \pm 1\} .
$$
Thus $v= \pm w_{1} \cdots w_{k} \in\left\langle\mathbb{S}^{n-1}\right\rangle$.
\end{PROOF}



At this point we have a great deal of information about elements of $\operatorname{Pin}(n)$ which we collect into a theorem.



Theorem 5.26

i) $\operatorname{Pin}(n)$ is the disjoint union of open subsets
$$
\operatorname{Pin}(n)=\operatorname{Pin}(n) \cap \mathrm{Cl}_{n}^{+} \cup \operatorname{Pin}(n) \cap \mathrm{Cl}_{n}^{-}=\operatorname{Spin}(n) \cup \operatorname{Pin}(n) \cap \mathrm{Cl}_{n}^{-} .
$$

ii) $\operatorname{Spin}(n)$ is a normal subgroup of $\operatorname{Pin}(n)$ and every element $v \in \operatorname{Spin}(n)$ can be expressed as an even-length product $v=v_{1} \cdots u_{2 \ell}$ where $v_{1}, \ldots, v_{2 \ell} \in \mathbb{S}^{n-1}$. iii) For any $w \in \mathbb{S}^{n-1}$,
$$
\operatorname{Pin}(n) \cap \mathrm{Cl}_{n}^{-}=w \operatorname{Spin}(n),
$$
and every element $v \in \operatorname{Pin}(n) \cap \mathrm{Cl}_{n}^{-}$can be expressed as an odd-length product $v=v_{1} \cdots u_{2 \ell+1}$ where $v_{1}, \ldots, v_{2 \ell+1} \in \mathbb{S}^{n-1}$.



\begin{THEO}{SG-B02-03-26}{Struktur von \(\Pin(n)\) und \(\Spin(n)\)}
\begin{enumerate}[label=\textnormal{(\roman*)}]
\item \(\displaystyle \Pin(n)=\Pin(n)\cap\Cl_n^{+}\;\dot\cup\;\Pin(n)\cap\Cl_n^{-}
=\Spin(n)\;\dot\cup\;\bigl(\Pin(n)\cap\Cl_n^{-}\bigr)\), und beide Summanden sind offen in \(\Pin(n)\).
\item \(\Spin(n)\) ist ein Normalteiler von \(\Pin(n)\). Jedes \(v\in\Spin(n)\) lässt sich als Produkt gerader Länge
\(v=v_1\cdots v_{2\ell}\) mit \(v_i\in\mathbb S^{n-1}\) schreiben.
\item Für jedes \(w\in\mathbb S^{n-1}\) gilt
\[
\Pin(n)\cap\Cl_n^{-}\;=\;w\,\Spin(n),
\]
und jedes \(u\in\Pin(n)\cap\Cl_n^{-}\) besitzt eine Darstellung ungerader Länge
\(u=v_1\cdots v_{2\ell+1}\) mit \(v_i\in\mathbb S^{n-1}\).
\end{enumerate}
\end{THEO}


Proof

\begin{PROOF}{SG-B02-03-27}{P: Struktur von \(\Pin(n)\) und \(\Spin(n)\)}
(i) Since $\operatorname{Pin}(n)=\left\langle\mathbb{S}^{n-1}\right\rangle$, every element $u \in \operatorname{Pin}(n)$ can be expressed as $u=u_{1} \cdots u_{k}$ for some $u_{1}, \ldots, u_{k} \in \mathbb{S}^{n-1}$. Since
$$
\alpha\left(u_{1} \cdots u_{k}\right)=(-1)^{k} u_{1} \cdots u_{k},
$$
we find that $u \in \operatorname{Pin}(n) \cap \mathrm{Cl}_{n}^{+}$if $k$ is even, while $u \in \operatorname{Pin}(n) \cap \mathrm{Cl}_{n}^{-}$if $k$ is odd. The function
$$
\widetilde{\alpha}: \operatorname{Pin}(n) \longrightarrow\{+1,-1\} ; \quad \widetilde{\alpha}(u)=\alpha(u) u^{-1},
$$
is continuous and
$$
\widetilde{\alpha}^{-1}\{1\}=\operatorname{Pin}(n) \cap \mathrm{Cl}_{n}^{+}, \quad \widetilde{\alpha}^{-1}\{-1\}=\operatorname{Pin}(n) \cap \mathrm{Cl}_{n}^{-},
$$
showing that these are disjoint open subsets.


(ii) If $u \in \operatorname{Spin}(n)=\operatorname{Pin}(n) \cap \mathrm{Cl}_{n}^{+}$and $v \in \operatorname{Pin}(n)$, then writing $u=u_{1} \cdots u_{2 k}$ and $v=v_{1} \cdots v_{2 \ell+1}$ with $u_{i}, v_{j} \in \mathbb{S}^{n-1}$ we find
$$
v u v^{-1}=-v_{1} \cdots u_{2 \ell+1} u_{1} \cdots u_{2 k} v_{2 \ell+1} \cdots v_{1} \in \operatorname{Pin}(n) \cap \mathrm{Cl}_{n}^{+} .
$$

(iii) If $u \in \operatorname{Pin}(n) \cap \mathrm{Cl}_{n}^{-}$then for each $w \in \mathbb{S}^{n-1}$ we have
$$
w^{-1} u=(-w) u \in \operatorname{Pin}(n) \cap \mathrm{Cl}_{n}^{+}
$$ \(\square\)
From now on we will make use of the notation
$$
\operatorname{Pin}(n)^{+}=\operatorname{Spin}(n)=\operatorname{Pin}(n) \cap \mathrm{Cl}_{n}^{+}, \quad \operatorname{Pin}(n)^{-}=\operatorname{Pin}(n) \cap \mathrm{Cl}_{n}^{-} .
$$
We can also characterise $\operatorname{Spin}(n)$ and $\operatorname{Pin}(n)^{-}$using the surjective homomorphism $\rho: \operatorname{Pin}(n) \longrightarrow \mathrm{O}(n)$. If $u_{1}, \ldots, u_{k} \in \mathbb{S}^{n-1}$, then
$$
\operatorname{det} \rho_{u_{1} \cdots u_{k}}=\operatorname{det} \rho_{u_{1}} \cdots \operatorname{det} \rho_{u_{k}}=(-1)^{k},
$$
since each $\rho_{u_{r}}$ is a hyperplane reflection for which $\operatorname{det} \rho_{u_{r}}=-1$. Since
$$
\begin{aligned}
\mathrm{SO}(n)=\mathrm{O}(n)^{+} & =\{A \in \mathrm{O}(n): \operatorname{det} A=1\} \\
\mathrm{O}(n)^{-} & =\{A \in \mathrm{O}(n): \operatorname{det} A=-1\}
\end{aligned}
$$
this means that for $u \in \operatorname{Pin}(n)$,
$$
\begin{aligned}
u \in \operatorname{Spin}(n) & \Longleftrightarrow \rho_{u} \in \mathrm{SO}(n), \\
u \in \operatorname{Pin}(n)^{-} & \Longleftrightarrow \rho_{u} \in \mathrm{O}(n)^{-} .
\end{aligned}
$$
\end{PROOF}




Theorem 5.27

The continuous homomorphism $\rho: \operatorname{Pin}(n) \longrightarrow \mathrm{O}(n)$ is surjective and
$$
\rho^{-1} \mathrm{SO}(n)=\operatorname{Spin}(n), \quad \rho^{-1} \mathrm{O}(n)^{-}=\operatorname{Pin}(n)^{-} .
$$
Hence the restriction of $\rho$ to $\rho^{+}: \operatorname{Spin}(n) \longrightarrow \operatorname{SO}(n)$ is also surjective and the kernels of these homomorphisms are $\operatorname{ker} \rho=\rho^{+}=\{+1,-1\}$.


\begin{THEO}{SG-B02-03-28}{Surjektivität und Zusammenhangskomponenten von $\rho: \Pin\to \Ogroup(n)$}
Der stetige Gruppenhomomorphismus
\[
\rho:\Pin(n)\longrightarrow O(n)
\]
ist surjektiv und es gilt
\[
\rho^{-1}\!\bigl(SO(n)\bigr)=\Spin(n),\qquad
\rho^{-1}\!\bigl(O(n)^{-}\bigr)=\Pin(n)^{-}.
\]
Insbesondere ist die Einschränkung
\[
\rho^{+}:\Spin(n)\longrightarrow SO(n)
\]
surjektiv und die Kerne sind in beiden Fällen
\[
\ker(\rho)=\ker(\rho^{+})=\{+1,-1\}.
\]
\end{THEO}



\begin{PROOF}{SG-B02-03-29}{P: Surjektivität und Zusammenhangskomponenten von $\rho: \Pin\to \Ogroup(n)$}
\textbf{(1) Surjektivität von \(\rho\).}
Nach Prop.\,5.24 ist für jedes \(u\in\mathbb S^{n-1}\subset\Cl_n\) die Abbildung
\(\rho_u(x)=uxu^{-1}=uxu\) die Spiegelung an der Hyperebene \(u^\perp\).
Nach Prop.\,5.25 erzeugen solche Spiegelungen \(O(n)\), und da
\(\Pin(n)=\langle\mathbb S^{n-1}\rangle\), folgt \(\rho(\Pin(n))=O(n)\).

\smallskip
\textbf{(2) Parität = Determinantenzeichen.}
Jedes \(g\in\Pin(n)\) besitzt eine Darstellung als Produkt endlich vieler
Einheitsvektoren \(g=u_1\cdots u_k\) mit \(u_i\in\mathbb S^{n-1}\).
Dann ist
\[
\rho(g)=\rho_{u_1}\circ\cdots\circ\rho_{u_k},
\]
also ein Produkt von \(k\) Spiegelungen. Damit
\(\det \rho(g)=(-1)^k\). Insbesondere:
\[
g\in\Cl_n^{+}\iff k\ \text{gerade}\iff \det\rho(g)=+1,
\]
und
\[
g\in\Cl_n^{-}\iff k\ \text{ungerade}\iff \det\rho(g)=-1.
\]

\smallskip
\textbf{(3) Urbilder der Zusammenhangskomponenten.}
Wegen (2) ist
\[
\rho^{-1}\!\bigl(SO(n)\bigr)=\{g\in\Pin(n):\det\rho(g)=+1\}
=\Pin(n)\cap\Cl_n^{+}=\Spin(n),
\]
und
\[
\rho^{-1}\!\bigl(O(n)^{-}\bigr)=\{g\in\Pin(n):\det\rho(g)=-1\}
=\Pin(n)\cap\Cl_n^{-}=\Pin(n)^{-}.
\]
Damit ist \(\rho^{+}=\rho|_{\Spin(n)}:\Spin(n)\twoheadrightarrow SO(n)\) surjektiv.

\smallskip
\textbf{(4) Kerne.}
Aus Prop.\,5.25 folgt \(\ker\rho=\{\pm1\}\).
Da \(\{\pm1\}\subset\Cl_n^{+}\), liegt der Kern auch in \(\Spin(n)\), also
\(\ker(\rho^{+})=\{\pm1\}\).

\smallskip
\textbf{(5) Stetigkeit.}
Die Abbildung \(\rho(g):x\mapsto g x g^{-1}\) ist in \(g\) polynomial/stetig in der
Clifford-Norm; damit ist \(\rho\) stetig.

Damit sind alle Behauptungen gezeigt.
\end{PROOF}



We can also say more about the topology of $\operatorname{Spin}(n)$ and $\operatorname{Pin}(n)$.



Theorem 5.28

$\operatorname{Spin}(n)$ is a compact, path connected, closed normal subgroup of $\operatorname{Pin}(n)$. Furthermore, if $n \geqslant 3$ the fundamental group of $\operatorname{Spin}(n)$ is trivial, $\pi_{1} \operatorname{Spin}(n)=\{1\}$.



Proof


We only discuss connectivity. Recall that the sphere $\mathbb{S}^{n-1} \subseteq \mathbb{R}^{n} \subseteq \mathrm{Cl}_{n}$ is path connected. Choose a base point $u_{0} \in \mathbb{S}^{n-1}$. Now for an element $u=u_{1} \cdots u_{k} \in \operatorname{Spin}(n)$ with $u_{1}, \ldots, u_{k} \in \mathbb{S}^{n-1}$, as noted in the proof of Proposition 5.25, we must have $k$ even, say $k=2 m$. In fact, we might as well take $m$ to be even since $u=u(-w) w$ for any $w \in \mathbb{S}^{n-1}$. Then there are continuous paths
$$
p_{r}:[0,1] \longrightarrow \mathbb{S}^{n-1} \quad(r=1, \ldots, 2 m),
$$
for which $p_{r}(0)=u_{0}$ and $p_{r}(1)=u_{r}$. Then
$$
p:[0,1] \longrightarrow \mathbb{S}^{n-1} ; \quad p(t)=p_{1}(t) \cdots p_{2 m}(t)
$$
is a continuous path in $\operatorname{Pin}(n)$ with
$$
p(0)=u_{0}^{2 m}=(-1)^{m}=1, \quad p(1)=u
$$
But $t \mapsto \rho(p(t))$ is a continuous path in $\mathrm{O}(n)$ with $\rho(p(0)) \in \mathrm{SO}(n)$, hence $\rho(p(t)) \in \mathrm{SO}(n)$ for all $t$. This shows that $p$ is a path in $\operatorname{Spin}(n)$. So every element $u \in \operatorname{Spin}(n)$ can be connected to 1 and therefore $\operatorname{Spin}(n)$ is path connected.

The final statement involves homotopy theory and is not proved here. It should be compared with the fact that for $n \geqslant 3, \pi_{1} \operatorname{SO}(n) \cong\{1,-1\}$ and in fact the map is an example of a universal covering.



\begin{THEO}{SG-B02-03-30}{Topologische Eigenschaften der Spin-Gruppe}
Die Gruppe \(\Spin(n)\) ist eine kompakte, wegzusammenhängende, abgeschlossene und normale Untergruppe von \(\Pin(n)\). 
Für \(n\ge 3\) ist die Fundamentalgruppe trivial, also
\[
\pi_1(\Spin(n))=\{1\}.
\]
\end{THEO}



\begin{PROOF}{SG-B02-03-31}{P: Topologische Eigenschaften der Spin-Gruppe}
\textbf{(1) Kompaktheit und Abgeschlossenheit.}
Nach der Definition ist \(\Spin(n)=\Pin(n)\cap\Cl_n^+\).
Da \(\Pin(n)\) als Untergruppe der normierten Algebra \(\Cl_n\) abgeschlossen und beschränkt ist, 
ist sie kompakt; das gleiche gilt für \(\Spin(n)\), das als Durchschnitt zweier abgeschlossener Mengen geschlossen ist.

\smallskip
\textbf{(2) Normalität.}
Aus Theorem~5.26(ii) wissen wir, dass \(\Spin(n)\lhd\Pin(n)\) gilt, da \(\rho^{-1}(SO(n))=\Spin(n)\)
und \(SO(n)\lhd O(n)\) normal ist.

\smallskip
\textbf{(3) Wegzusammenhang.}
Die Sphäre \(\mathbb S^{n-1}\subseteq\R^n\subseteq\Cl_n\) ist wegzusammenhängend. 
Sei \(u_0\in\mathbb S^{n-1}\) ein Basispunkt.
Für ein Element \(u=u_1\cdots u_k\in\Spin(n)\) mit \(u_i\in\mathbb S^{n-1}\)
ist \(k\) gerade, sagen wir \(k=2m\). 
Ohne Einschränkung können wir \(m\) gerade wählen, denn
für jedes \(w\in\mathbb S^{n-1}\) gilt \(u=u(-w)w\), was die Parität nicht verändert.
Für jedes \(r=1,\dots,2m\) existiert ein stetiger Weg 
\[
p_r:[0,1]\to\mathbb S^{n-1},\qquad p_r(0)=u_0,\quad p_r(1)=u_r.
\]
Dann definiert
\[
p:[0,1]\to\Pin(n),\qquad p(t)=p_1(t)\cdots p_{2m}(t)
\]
einen stetigen Pfad in \(\Pin(n)\) mit
\[
p(0)=u_0^{2m}=(-1)^m=1,\qquad p(1)=u.
\]
Die Abbildung \(t\mapsto\rho(p(t))\) ist ein stetiger Pfad in \(O(n)\)
mit \(\rho(p(0))\in SO(n)\).
Da \(SO(n)\) eine Zusammenhangskomponente von \(O(n)\) ist, folgt
\(\rho(p(t))\in SO(n)\) für alle \(t\).
Somit ist \(p(t)\in\Spin(n)\) für alle \(t\), also \(u\) durch einen Pfad mit \(1\) verbunden.
Damit ist \(\Spin(n)\) wegzusammenhängend.

\smallskip
\textbf{(4) Zur Fundamentalgruppe.}
Für \(n\ge3\) ist bekannt, dass \(\pi_1(SO(n))\cong\{\pm1\}\)
und dass \(\rho^{+}:\Spin(n)\to SO(n)\) die universelle Überlagerung ist.
Daher ist \(\pi_1(\Spin(n))=\{1\}\).

\smallskip
Damit sind alle Aussagen bewiesen.
\end{PROOF}




The double covering maps $\rho: \operatorname{Spin}(n) \longrightarrow \operatorname{SO}(n)$ generalise the case of $\mathrm{SU}(2) \longrightarrow \mathrm{SO}(3)$ discussed in Section 3.5. In fact, around each element $u \in$ there is an open neighbourhood $N_{u} \subseteq \operatorname{Spin}(n)$ for which $\rho: N_{u} \longrightarrow \rho N_{u}$ is a homeomorphism, and actually a diffeomorphism. 




\begin{THEO}{SG-B02-03-32}{Doppelte Überlagerung von \(\Spin(n)\) über \(\SO(n)\)}
Die Abbildung
\[
\rho:\Spin(n)\longrightarrow SO(n)
\]
ist eine doppelte Überlagerung und verallgemeinert den klassischen Fall
\(\SU(2)\longrightarrow SO(3)\).

Genauer gilt: Für jedes \(u\in\Spin(n)\) existiert eine offene Umgebung 
\(N_u\subseteq\Spin(n)\), sodass
\[
\rho:N_u\longrightarrow \rho(N_u)
\]
ein Homöomorphismus (sogar ein Diffeomorphismus) ist.
\end{THEO}



\begin{PROOF}{SG-B02-03-33}{P: Doppelte Überlagerung von \(\Spin(n)\) über \(\SO(n)\)}
Da \(\ker\rho=\{\pm1\}\) (vgl.\ Theorem~5.27) und \(\Spin(n)\) kompakt sowie
wegzusammenhängend ist, folgt, dass \(\rho\) eine zweiblättrige Überlagerung 
von \(SO(n)\) ist. Die lokale Trivialität ergibt sich daraus, dass 
\(\rho\) als stetiger Gruppenhomomorphismus zwischen Liegruppen 
mit diskretem Kern lokal diffeomorph ist:  
die Differentialabbildung \(d\rho_1:\mathfrak{spin}(n)\to\mathfrak{so}(n)\)
ist ein Isomorphismus der Liealgebren.
Damit existieren für jedes \(u\in\Spin(n)\) offene Umgebungen \(N_u\)
und \(\rho(N_u)\), auf denen \(\rho\) ein Diffeomorphismus ist.

Insbesondere gilt:
\[
\rho:\Spin(n)\twoheadrightarrow SO(n)
\]
ist eine \emph{zweiblättrige Überlagerung von Liegruppen}.
Der Spezialfall \(n=3\) liefert die klassische Abbildung 
\(\SU(2)\to SO(3)\).
\end{PROOF}



This implies the following.



Proposition 5.29

The derivative $\mathrm{d} \rho: \boldsymbol{\operatorname { s p i n }}(n) \longrightarrow \boldsymbol{\mathfrak { s }}(n)$ is an isomorphism of $\mathbb{R}$-Lie algebras and
$$
\operatorname{dim} \operatorname{Spin}(n)=\operatorname{dim} \operatorname{SO}(n)=\binom{n}{2}
$$
The Lie algebra of $\operatorname{Spin}(n)$ can be described as a Lie subalgebra of $\mathrm{Cl}_{n}$.




\begin{THEO}{SG-B02-03-35}{Lie-Algebren-Isomorphismus und Dimension für Spin Gruppe}
Die Ableitung
\[
\mathrm{d}\rho:\mathfrak{spin}(n)\longrightarrow\mathfrak{so}(n)
\]
ist ein Isomorphismus von reellen Lie-Algebren. Insbesondere gilt
\[
\dim\Spin(n)=\dim SO(n)=\binom{n}{2}.
\]
Darüber hinaus lässt sich \(\mathfrak{spin}(n)\) als Lie-Unteralgebra der Clifford-Algebra
\(\Cl_n\) beschreiben.
\end{THEO}




\begin{PROOF}{SG-B02-03-36}{P: Lie-Algebren-Isomorphismus und Dimension für Spin Gruppe}
\textbf{(1) Lie-Algebren-Isomorphismus.}
Da \(\rho:\Spin(n)\to SO(n)\) eine surjektive Liegruppenhomomorphismus mit
diskretem Kern \(\{\pm1\}\) ist, folgt aus der allgemeinen Theorie der
Liegruppen, dass die induzierte Abbildung der Lie-Algebren
\[
\mathrm{d}\rho_1:\mathrm{Lie}(\Spin(n))\to\mathrm{Lie}(SO(n))
\]
ein Isomorphismus ist. 
Hierbei gilt \(\mathrm{Lie}(\Spin(n))=\mathfrak{spin}(n)\) und 
\(\mathrm{Lie}(SO(n))=\mathfrak{so}(n)\).

\smallskip
\textbf{(2) Dimensionsformel.}
Da die Dimension der Lie-Algebra einer Liegruppe gleich der Dimension der Gruppe
als differenzierbare Mannigfaltigkeit ist, ergibt sich
\[
\dim\Spin(n)=\dim\mathfrak{so}(n)=\frac{n(n-1)}{2}=\binom{n}{2}.
\]

\smallskip
\textbf{(3) Realisierung in der Clifford-Algebra.}
Die Lie-Algebra \(\mathfrak{spin}(n)\) kann in \(\Cl_n\) identifiziert werden als
\[
\mathfrak{spin}(n)=\mathrm{span}_{\R}\{\,e_i e_j : 1\le i<j\le n\,\}
\subset \Cl_n^{+},
\]
versehen mit der Lie-Klammer gegeben durch den Kommutator
\([x,y]=xy-yx\).
Diese liegt in der geraden Komponente \(\Cl_n^{+}\) und ist abgeschlossen
unter der Kommutatorbildung, da Produkte von zwei geraden Elementen wieder gerade sind.

Damit ist \(\mathfrak{spin}(n)\subset\Cl_n^{+}\) eine reelle Lie-Unteralgebra,
und \(\mathrm{d}\rho:\mathfrak{spin}(n)\xrightarrow{\cong}\mathfrak{so}(n)\) 
ein Isomorphismus.
\end{PROOF}




Proposition 5.30


For $n \geq 2$, the Lie algebra of $\operatorname{Spin}(n)$ is
$$
\operatorname{spin}(n)=\left\{\sum_{1 \leqslant i<j \leqslant n} t_{i j} e_{i} e_{j}: t_{i j} \in \mathbb{R}\right\} \subseteq \mathrm{Cl}_{n} .
$$



\begin{PROP}{SG-B02-03-37}{Darstellung der Lie-Algebra \(\mathfrak{spin}(n)\) in \(\Cl_n\)}
Für \(n\ge2\) ist die Lie-Algebra der Spin-Gruppe gegeben durch
\[
\mathfrak{spin}(n)
=\left\{
\sum_{1\le i<j\le n} t_{ij}\, e_i e_j
\;\middle|\;
t_{ij}\in\R
\right\}
\subseteq \Cl_n.
\]
\end{PROP}



Proof


\begin{PROOF}{SG-B02-03-38}{P: Darstellung der Lie-Algebra \(\mathfrak{spin}(n)\) in \(\Cl_n\)}
For $1 \leqslant i<j \leqslant n$ and $t \in \mathbb{R}$, consider
$$
\exp \left(t e_{i} e_{j}\right)=\sum_{r=0}^{\infty} \frac{t^{r}}{r!}\left(e_{i} e_{j}\right)^{r}
$$
which converges to an element of $\mathrm{Cl}_{n}^{\times}$. Notice that
$$
\left(e_{i} e_{j}\right)^{r}=(-1)^{\binom{r}{2}} e_{i}^{r} e_{j}^{r}=\left\{\begin{array}{cl}
1 & \text { if } r \equiv 0 \bmod 4, \\
-1 & \text { if } r \equiv 2 \bmod 4, \\
e_{i} e_{j} & \text { if } r \equiv 1 \bmod 4, \\
-e_{i} e_{j} & \text { if } r \equiv 3 \bmod 4 .
\end{array}\right.
$$
Hence we have
$$
\exp \left(t e_{i} e_{j}\right)=\cos t+\sin t e_{i} e_{j}
$$
We also obtain
$$
\begin{aligned}
e_{i}\left(e_{i} e_{j}\right)^{r} & =\left(-e_{i} e_{j}\right)^{r} e_{i}, \\
e_{j}\left(e_{i} e_{j}\right)^{r} & =\left(-e_{i} e_{j}\right)^{r} e_{j},
\end{aligned}
$$
and if $k \neq i, j$,
$$
e_{k}\left(e_{i} e_{j}\right)^{r}=\left(e_{i} e_{j}\right)^{r} e_{k}
$$
Then since $\left(e_{i} e_{j}\right)^{2}=-1$,
$$
\begin{aligned}
\alpha\left(\exp \left(t e_{i} e_{j}\right)\right) e_{i} \exp \left(t e_{i} e_{j}\right)^{-1} & =\exp \left(t e_{i} e_{j}\right) e_{i} \exp \left(-t e_{i} e_{j}\right) \\
& =\left(\cos t+\sin t e_{i} e_{j}\right) e_{i}\left(\cos t-\sin t e_{i} e_{j}\right) \\
& =\left(\cos t+\sin t e_{i} e_{j}\right)^{2} e_{i} \\
& =\left(\cos 2 t+\sin 2 t e_{i} e_{j}\right) e_{i} \\
& =\cos 2 t e_{i}+\sin 2 t e_{j} \\
\alpha\left(\exp \left(t e_{i} e_{j}\right)\right) e_{j} \exp \left(t e_{i} e_{j}\right)^{-1} & =\left(\cos t+\sin t e_{i} e_{j}\right) e_{j}\left(\cos t-\sin t e_{i} e_{j}\right) \\
& =\left(\cos t+\sin t e_{i} e_{j}\right)^{2} e_{j} \\
& =\left(\cos 2 t+\sin 2 t e_{i} e_{j}\right) e_{j} \\
& =\cos 2 t e_{j}-\sin 2 t e_{i}
\end{aligned}
$$
and if $k \neq i, j$,
$$
\begin{aligned}
\alpha\left(\exp \left(t e_{i} e_{j}\right)\right) e_{k} \exp \left(t e_{i} e_{j}\right)^{-1} & =\exp \left(t e_{i} e_{j}\right) \exp \left(t e_{i} e_{j}\right)^{-1} e_{k} \\
& =e_{k}
\end{aligned}
$$
Since the $e_{\ell}$ form a basis of $\mathbb{R}^{n}$, this shows that $\exp \left(t e_{i} e_{j}\right) \in \Gamma_{n} \cap \mathrm{Cl}_{n}^{+}$. But also
$$
\begin{aligned}
\exp \left(t e_{i} e_{j}\right) \overline{\exp \left(t e_{i} e_{j}\right)} & =\exp \left(t e_{i} e_{j}\right) \exp \left((-1)^{2} t e_{j} e_{i}\right) \\
& =\exp \left(t e_{i} e_{j}\right) \exp \left(-t e_{i} e_{j}\right) \\
& =1
\end{aligned}
$$
hence $\exp \left(t e_{i} e_{j}\right) \in \operatorname{Spin}(n)$.

Taking the derivative at 0 of the curve
$$
\alpha_{i j}: \mathbb{R} \longrightarrow \operatorname{Spin}(n) ; \quad \alpha_{i j}(t)=\exp \left(t e_{i} e_{j}\right),
$$
we obtain $e_{i} e_{j}$, so $e_{i} e_{j} \in \mathfrak{s p i n}(n)$. As these $\binom{n}{2}$ elements $e_{i} e_{j}$ are clearly linearly independent, by Proposition 5.29 they must form a basis.
\end{PROOF}




With the aid of the calculations in this proof, the action of the derivative $\mathrm{d} \rho: \operatorname{spin}(n) \longrightarrow \operatorname{so}(n)$ on the basis elements $e_{i} e_{j}$ can be determined, namely
$$
\mathrm{d} \rho\left(e_{i} e_{j}\right)=2 E^{j i}-2 E^{i j},
$$
where we use the basis for $\mathrm{M}_{n}(\mathbb{R})$ provided by Equation (3.24). Similarly, the effect of the homomorphism $\rho: \operatorname{Spin}(n) \longrightarrow \operatorname{SO}(n)$ on $\exp \left(t e_{i} e_{j}\right)$ is
$$
\rho\left(\exp \left(t e_{i} e_{j}\right)\right)=\cos 2 t I_{n}+\sin 2 t\left(E^{j i}-E^{i j}\right) .
$$



\begin{PROP}{SG-B02-03-39}{Ableitung \(\mathrm d\rho\) auf Basisvektoren und Exponentialbilder}
Für \(1\le i<j\le n\) gilt
\[
\mathrm d\rho\!\left(e_i e_j\right)=2E^{ji}-2E^{ij}\in\mathfrak{so}(n),
\]
wobei \((E^{ab})_{pq}=\delta_{ap}\delta_{bq}\) die kanonische Matrixbasis aus \eqref{eq:3.24} bezeichne.
Ferner ist für \(t\in\R\)
\[
\rho\!\left(\exp(te_i e_j)\right)
=\cos(2t)\,I_n+\sin(2t)\,(E^{ji}-E^{ij})\in\SO(n).
\]
\end{PROP}



\begin{PROOF}{SG-B02-03-40}{P: Ableitung \(\mathrm d\rho\) auf Basisvektoren und Exponentialbilder}
Sei \(x\in\R^n\subset\Cl_n\) und \(g(t)=\exp(te_ie_j)\in\Spin(n)\). Dann
\[
\rho(g(t))(x)=g(t)\,x\,g(t)^{-1}.
\]
Ableiten in \(t=0\) liefert
\[
\left.\frac{\mathrm d}{\mathrm dt}\right|_{t=0}\rho(g(t))(x)
=\left.[\,e_i e_j,\,x\,]\right.
= e_i e_j x- x e_i e_j.
\]
Benutze \(e_k e_\ell+e_\ell e_k=-2\delta_{k\ell}\) und schreibe \(x=\sum_{k}x_k e_k\).
Für \(x=e_k\) erhält man die Fälle
\[
[e_i e_j,e_k]=
\begin{cases}
2e_j, & k=i,\\[2pt]
-2e_i,& k=j,\\[2pt]
0,& k\notin\{i,j\}.
\end{cases}
\]
Somit wirkt \(\mathrm d\rho(e_i e_j)\) auf der kanonischen Basis \(\{e_1,\dots,e_n\}\) durch
\[
e_i\mapsto 2e_j,\qquad e_j\mapsto -2e_i,\qquad e_k\mapsto 0\ (k\ne i,j),
\]
also als die schiefsymmetrische Matrix \(2E^{ji}-2E^{ij}\in\mathfrak{so}(n)\).

Für das Exponentialbild genügt es, die ODE
\(\frac{\mathrm d}{\mathrm dt}A(t)=\mathrm d\rho(e_i e_j)\,A(t)\), \(A(0)=I_n\),
in der \((i,j)\)-Ebene zu lösen. Da \(\mathrm d\rho(e_i e_j)\) dort der Standardgenerator
\(\begin{pmatrix}0&-2\\ 2&0\end{pmatrix}\) ist, ergibt sich die Rotation um Winkel \(2t\):
\[
A(t)=\exp\!\big(t\,\mathrm d\rho(e_i e_j)\big)
=\cos(2t)\,I_n+\sin(2t)\,(E^{ji}-E^{ij}).
\]
Außerhalb der \((i,j)\)-Ebene wirkt \(A(t)\) trivial. Damit sind beide Behauptungen gezeigt.
\end{PROOF}



\pagebreak


\subsection{The Centres of Spinor Groups}


Recall that for a group $G$ the centre of $G$ is
$$
\mathrm{Z}(G)=\{c \in G: \forall g \in G, g c=c g\} .
$$
Then $\mathrm{Z}(G) \triangleleft G$. The centres of the special orthogonal groups are easily found and described.


\begin{DEF}{SG-B02-03-41}{Zentrum einer Gruppe}
Sei \(G\) eine Gruppe. Das \emph{Zentrum} von \(G\) ist definiert durch
\[
Z(G)=\{\,c\in G:\forall g\in G,\; gc=cg\,\}.
\]
Das Zentrum \(Z(G)\) ist stets eine normale Untergruppe von \(G\),
d.\,h. \(Z(G)\trianglelefteq G\).
\end{DEF}




\begin{PROP}{SG-B02-03-42}{Zentren von $\SO(n)$}
Sei $n\geq 3$. Für die speziellen orthogonalen Gruppen gilt:
\[
Z(\SO(n))=
\begin{cases}
\{I_n,-I_n\}, & n\ \text{gerade},\\[4pt]
\{I_n\}, & n\ \text{ungerade}.
\end{cases}
\]
\end{PROP}




\begin{PROOF}{SG-B02-03-43}{P: Zentren von $\SO(n)$}
Jedes Element \(A\in\SO(n)\) erfüllt \(A A^T=I_n\) und \(\det A=1\).
Ein \(C\in Z(\SO(n))\) muss mit allen \(A\in\SO(n)\) kommutieren, also \(CA=AC\).

\smallskip
\textbf{(1) Schritt.}  
Aus der Kommutativitätsbedingung folgt, dass \(C\) ein Vielfaches der Identität sein muss:
\(C=\lambda I_n\), da \(\SO(n)\) transitiv auf orthonormalen Basen wirkt
und nur skalare Matrizen mit allen Orthogonalmatrizen kommutieren.

\smallskip
\textbf{(2) Schritt.}  
Da \(C\in\SO(n)\) gilt \(C^T C=I_n\), also \(|\lambda|=1\), und \(\det C=\lambda^n=1\).

\smallskip
\textbf{(3) Schritt.}  
Für \(\lambda=\pm1\) folgt:
\[
Z(\SO(n))=
\begin{cases}
\{I_n,-I_n\}, & \text{falls } n\ \text{gerade},\\[4pt]
\{I_n\}, & \text{falls } n\ \text{ungerade},
\end{cases}
\]
da bei ungeradem \(n\) \(\det(-I_n)=-1\neq1\) gilt.
\end{PROOF}




Proposition 5.31


For $n \geqslant 3$,
$$
\mathrm{Z}(\mathrm{SO}(n))=\left\{t I_{n}: t= \pm 1, t^{n}=1\right\}= \begin{cases}\left\{I_{n}\right\} & \text { if } n \text { is odd } \\ \left\{ \pm I_{n}\right\} & \text { if } n \text { is even }\end{cases}
$$


Before stating our main result on the centres of spinor groups, we note that Spin(1) and Spin(2) are both abelian.


\begin{LEM}{SG-B02-03-44}{Abelianität niedriger Spin-Gruppen}
Die Gruppen \(\Spin(1)\) und \(\Spin(2)\) sind abelsch.
\end{LEM}



\begin{PROOF}{SG-B02-03-45}{P: Abelianität niedriger Spin-Gruppen}
Für \(n=1\) ist \(\Cl_1\cong\C\) und
\(\Spin(1)=\{\pm1\}\subset\R^\times\), also trivialerweise abelsch.

Für \(n=2\) gilt \(\Cl_2\cong\QU\) (den Hamiltonschen Quaternionen) und
\[
\Spin(2)=\{\,\cos t + \sin t\, e_1 e_2 : t\in\R\,\}.
\]
Da \(e_1 e_2\) quadratisch \(-1\) ist, bilden diese Elemente einen isomorphen
Kreis \(\U(1)\subset\QU\), der kommutativ ist. Somit ist auch \(\Spin(2)\) abelsch.
\end{PROOF}


Proposition 5.32


For $n \geqslant 3$,
$$
\begin{aligned}
Z(\operatorname{Spin}(n)) & = \begin{cases}\{ \pm 1\} & \text { if } n \text { is odd } \\
\left\{ \pm 1, \pm e_{1} \cdots e_{n}\right\} & \text { if } n \equiv 2 \bmod 4, \\
\left\{ \pm 1, \pm e_{1} \cdots e_{n}\right\} & \text { if } n \equiv 0 \bmod 4 .\end{cases} \\
& \cong \begin{cases}\mathbb{Z} / 2 & \text { if } n \text { is odd } \\
\mathbb{Z} / 4 & \text { if } n \equiv 2 \bmod 4 \\
\mathbb{Z} / 2 \times \mathbb{Z} / 2 & \text { if } n \equiv 0 \bmod 4 .\end{cases}
\end{aligned}
$$




\begin{THEO}{SG-B02-03-46}{Zentren von \(\Spin(n)\)}
Für \(n\ge 3\) gilt
\[
Z(\Spin(n))=
\begin{cases}
\{\pm 1\}, & n \text{ ungerade},\\[2mm]
\{\pm 1,\ \pm \omega\}, & n \equiv 0 \pmod 4,\\[2mm]
\{\pm 1,\ \pm \omega\}, & n \equiv 2 \pmod 4,
\end{cases}
\qquad
\omega:=e_1 e_2\cdots e_n ,
\]
und damit (als abstrakte Gruppen)
\[
Z(\Spin(n))\cong
\begin{cases}
\Bbb Z/2, & n \text{ ungerade},\\
\Bbb Z/2\times \Bbb Z/2, & n \equiv 0 \pmod 4,\\
\Bbb Z/4, & n \equiv 2 \pmod 4.
\end{cases}
\]
\end{THEO}


Proof


\begin{PROOF}{SG-B02-03-47}{P: Zentren von \(\Spin(n)\)}
If $g \in \mathrm{Z}(\operatorname{Spin}(n))$, then since $\rho: \operatorname{Spin}(n) \longrightarrow \mathrm{SO}(n), \rho(g) \in \mathrm{Z}(\mathrm{SO}(n))$. As $\pm 1 \in \mathrm{Z}(\operatorname{Spin}(n))$, this gives $|\mathrm{Z}(\operatorname{Spin}(n))|=2|\mathrm{Z}(\mathrm{SO}(n))|$ and indeed
$$
\mathrm{Z}(\operatorname{Spin}(n))=\rho^{-1} \mathrm{Z}(\mathrm{SO}(n)) .
$$
For $n$ even,
$$
\begin{aligned}
\left( \pm e_{1} \cdots e_{n}\right)^{2} & =e_{1} \cdots e_{n} e_{1} \cdots e_{n} \\
& =(-1)^{\binom{n}{2}} e_{1}^{2} \cdots e_{n}^{2} \\
& =(-1)^{\binom{n}{2}+n}=(-1)^{\binom{n+1}{2} .}
\end{aligned}
$$
Since
$$
\binom{n+1}{2}=\frac{(n+1) n}{2} \equiv \begin{cases}1 \bmod 2 & \text { if } n \equiv 2 \bmod 4, \\ 0 \bmod 2 & \text { if } n \equiv 0 \bmod 4,\end{cases}
$$
this implies that
$$
\left( \pm e_{1} \cdots e_{n}\right)^{2}=\left\{\begin{aligned}
-1 & \text { if } n \equiv 2 \bmod 4, \\
1 & \text { if } n \equiv 0 \bmod 4 .
\end{aligned}\right.
$$
So for $n$ even, the multiplicative order of $\pm e_{1} \cdots e_{n}$ is 2 or 4 depending on the congruence class of $n$ modulo 4 . This gives the stated groups.
\end{PROOF}



\pagebreak


\subsection{Finite Subgroups of Spinor Groups}

Each orthogonal group $\mathrm{O}(n)$ and $\mathrm{SO}(n)$ contains finite subgroups. For example, when $n=2,3$, these correspond to symmetry groups of compact plane figures\\
and solids. The case of $n=3$ was explored in the exercises of Chapter 4. Here we make some remarks about the symmetric and alternating groups.

Recall that for each $n \geqslant 1$ the symmetric group $\mathrm{S}_{n}$ is the group of all permutations of the set $\mathbf{n}=\{1, \ldots, n\}$. The corresponding alternating group $\mathrm{A}_{n} \leqslant \mathrm{~S}_{n}$ is the subgroup consisting of all even permutations, i.e., the elements $\sigma \in \mathrm{S}_{n}$ for which $\operatorname{sgn}(\sigma)=1$ where $\operatorname{sgn}: \mathrm{S}_{n} \longrightarrow\{ \pm 1\}$ is the sign homomorphism.


\begin{DEF}{SG-B02-03-48}{Symmetrische und Alternierende Gruppen}
Für jedes \(n \ge 1\) bezeichnet man mit
\[
S_n := \{\text{Bijektionen } \sigma : \{1,\ldots,n\}\to\{1,\ldots,n\}\}
\]
die \emph{symmetrische Gruppe} aller Permutationen der Menge
\(\mathbf{n} = \{1,\ldots,n\}\).

\begin{itemize}
\item Die Gruppenoperation ist die Komposition von Abbildungen:
      \((\sigma\circ\tau)(i) = \sigma(\tau(i))\).
\item Die \emph{Ordnung} von \(S_n\) ist \(n!\).
\end{itemize}

Die \emph{Alternierende Gruppe} \(A_n\subseteq S_n\) ist die Untergruppe der
\emph{geraden Permutationen}, d.\,h.
\[
A_n := \{\sigma\in S_n : \sgn(\sigma)=1\},
\]
wobei \(\sgn:S_n\to\{\pm1\}\) der \emph{Signum-Homomorphismus} ist, der jeder
Permutationsdarstellung ihre Parität (gerade oder ungerade Anzahl von
Transpositionen) zuordnet.
\end{DEF}


For a field $\mathbf{k}$, we can make $S_{n}$ act on $\mathbf{k}^{n}$ by linear transformations:
$$
\sigma \cdot\left[\begin{array}{c}
x_{1} \\
x_{2} \\
\vdots \\
x_{n}
\end{array}\right]=\left[\begin{array}{c}
x_{\sigma^{-1}(1)} \\
x_{\sigma^{-1}(2)} \\
\vdots \\
x_{\sigma^{-1}(n)}
\end{array}\right] .
$$
Notice that $\sigma\left(\mathbf{e}_{r}\right)=\mathbf{e}_{\sigma(r)}$. The matrix $[\sigma]$ of the linear transformation induced by $\sigma$ with respect to the basis of $\mathbf{e}_{r}$ 's has all its entries 0 or 1 , with exactly one 1 in each row and column. For example, when $n=3$,
$$
\left[\begin{array}{lll}
1 & 2 & 3
\end{array}\right]=\left[\begin{array}{lll}
0 & 0 & 1 \\
1 & 0 & 0 \\
0 & 1 & 0
\end{array}\right], \quad\left[\left(\begin{array}{ll}
1 & 3
\end{array}\right)\right]=\left[\begin{array}{lll}
0 & 0 & 1 \\
0 & 1 & 0 \\
1 & 0 & 0
\end{array}\right] .
$$
When $\mathbf{k}=\mathbb{R}$ each of these matrices is orthogonal, while when $\mathbf{k}=\mathbb{C}$ it is unitary. For a given $n$ we can view $\mathrm{S}_{n}$ as the subgroup of $\mathrm{O}(n)$ or $\mathrm{U}(n)$ consisting of all such matrices which are usually called permutation matrices.


\begin{DEF}{SG-B02-03-49}{Darstellung der symmetrischen Gruppe als Permutationsmatrizen}
Sei \(\mathbf{k}\) ein Körper. Dann operiert die symmetrische Gruppe
\(S_n\) auf dem \(\mathbf{k}\)-Vektorraum \(\mathbf{k}^n\) durch lineare
Transformationen, definiert durch
\[
\sigma\cdot
\begin{bmatrix}
x_1\\x_2\\\vdots\\x_n
\end{bmatrix}
=
\begin{bmatrix}
x_{\sigma^{-1}(1)}\\x_{\sigma^{-1}(2)}\\\vdots\\x_{\sigma^{-1}(n)}
\end{bmatrix},
\qquad \sigma\in S_n.
\]
Damit gilt insbesondere für die Standardbasisvektoren
\(\mathbf{e}_r\):
\[
\sigma(\mathbf{e}_r)=\mathbf{e}_{\sigma(r)}.
\]

\textbf{Matrixdarstellung.}
Die Matrix \([\sigma]\) der linearen Abbildung, die durch \(\sigma\)
induziert wird, besitzt folgende Eigenschaften:
\begin{itemize}
\item Alle Einträge sind \(0\) oder \(1\),
\item in jeder Zeile und jeder Spalte steht genau eine \(1\),
\item \([\sigma\circ\tau]=[\sigma]\,[\tau]\) für alle \(\sigma,\tau\in S_n\).
\end{itemize}

\textbf{Beispiel:} Für \(n=3\) gilt
\[
\left[(1\,2\,3)\right]=
\begin{bmatrix}
0 & 0 & 1\\
1 & 0 & 0\\
0 & 1 & 0
\end{bmatrix},
\qquad
\left[(1\,3)\right]=
\begin{bmatrix}
0 & 0 & 1\\
0 & 1 & 0\\
1 & 0 & 0
\end{bmatrix}.
\]

\textbf{Interpretation:}
\begin{itemize}
\item Für \(\mathbf{k}=\mathbb{R}\) ist jede solche Matrix orthogonal, also
\([\sigma]\in O(n)\).
\item Für \(\mathbf{k}=\mathbb{C}\) ist sie unitär, also
\([\sigma]\in U(n)\).
\end{itemize}

Somit kann man \(S_n\) als Untergruppe von \(O(n)\) bzw.\ \(U(n)\)
auffassen, bestehend aus allen \emph{Permutationsmatrizen}.
\end{DEF}



Proposition 5.33


For each $\sigma \in \mathrm{S}_{n}$ we have $\operatorname{sgn}(\sigma)=\operatorname{det}([\sigma])$. Hence
$$
A_{n}= \begin{cases}\operatorname{SO}(n) \cap S_{n} & \text { if } \mathbf{k}=\mathbb{R}, \\ \operatorname{SU}(n) \cap S_{n} & \text { if } \mathbf{k}=\mathbb{C} .\end{cases}
$$
Recall that if $n \geqslant 5, \mathrm{~A}_{n}$ is a simple group.


As $\rho: \operatorname{Pin}(n) \longrightarrow \mathrm{O}(n)$ is onto, there are finite subgroups $\widetilde{\mathrm{S}}_{n}=\rho^{-1} \mathrm{~S}_{n} \leqslant \operatorname{Pin}(n)$ and $\widetilde{\mathrm{A}}_{n}=\rho^{-1} \mathrm{~A}_{n} \leqslant \operatorname{Spin}(n)$ for which there are surjective homomorphisms $\rho: \widetilde{\mathrm{S}}_{n} \longrightarrow \mathrm{~S}_{n}$ and $\rho: \widetilde{\mathrm{A}}_{n} \longrightarrow \mathrm{~A}_{n}$ whose kernels contain the elements $\pm 1$. Note that $\left|\widetilde{\mathrm{S}}_{n}\right|=2 \cdot n!$, while $\left|\widetilde{\mathrm{A}}_{n}\right|=n!$. However, for $n \geqslant 4$, there are no homomorphisms $\tau$ : $\mathrm{S}_{n} \longrightarrow \widetilde{\mathrm{~S}}_{n}, \tau: \mathrm{A}_{n} \longrightarrow \widetilde{\mathrm{~A}}_{n}$ for which $\rho \circ \tau=\mathrm{Id}$.
\[
\begin{tikzcd}[row sep=large, column sep=large]
S_n \arrow[r, "\tau"] \arrow[dr, "\mathrm{Id}_{S_n}"'] 
& \widetilde{S}_n \arrow[d, "\rho"] \\
& S_n
\end{tikzcd}
\]
\[
\begin{tikzcd}[row sep=large, column sep=large]
A_n \arrow[r, "\tau"] \arrow[dr, "\mathrm{Id}_{A_n}"'] 
& \widetilde{A}_n \arrow[d, "\rho"] \\
& A_n
\end{tikzcd}
\]
Similar considerations apply to other finite subgroups of $\mathrm{O}(n)$.

In $\mathrm{Cl}_{n}^{\times}$ we have a subgroup $E_{n}$ consisting of all the elements

$$
\pm e_{i_{1}} \cdots e_{i_{r}} \quad\left(1 \leqslant i_{1}<\cdots<i_{r} \leqslant n, 0 \leqslant r\right) .
$$

The order of this group is $\left|E_{n}\right|=2^{n+1}$ and as it contains $\pm 1$, its image under $\rho: \operatorname{Pin}(n) \longrightarrow \mathrm{O}(n)$ is $\bar{E}_{n}=\rho E_{n}$ of order $\left|\bar{E}_{n}\right|=2^{n}$. In fact, $\{ \pm 1\}=\mathrm{Z}\left(E_{n}\right)$ is also the commutator subgroup since $e_{i} e_{j} e_{i}^{-1} e_{j}^{-1}=-1$ and so $\bar{E}_{n}$ is abelian. Every non-trivial element in $\bar{E}_{n}$ has order 2 since $e_{i}^{2}=-1$, hence $\bar{E}_{n} \leqslant \mathrm{O}(n)$ is an elementary 2 -group, i.e., it is isomorphic to $(\mathbb{Z} / 2)^{n}$. Each element $\rho\left(e_{r}\right) \in \mathrm{O}(n)$ is a generalised permutation matrix with all its non-zero entries on the main diagonal. There is also a subgroup $\bar{E}_{n}^{0}=\rho E_{n}^{0} \leqslant \operatorname{SO}(n)$ of order $2^{n-1}$, where

$$
E_{n}^{0}=E_{n} \cap \operatorname{Spin}(n) .
$$

In fact $\bar{E}_{n}^{0}$ is isomorphic to $(\mathbb{Z} / 2)^{n-1}$. These groups $E_{n}$ and $E_{n}^{0}$ are non-abelian and fit into exact sequences of the form

$$
1 \rightarrow \mathbb{Z} / 2 \longrightarrow E_{n} \longrightarrow(\mathbb{Z} / 2)^{n} \rightarrow 1, \quad 1 \rightarrow \mathbb{Z} / 2 \longrightarrow E_{n}^{0} \longrightarrow(\mathbb{Z} / 2)^{n-1} \rightarrow 1,
$$

in which each kernel $\mathbb{Z} / 2$ is equal to the centre of the corresponding group $E_{n}$ or $E_{n}^{0}$. This means they are examples of extraspecial 2 -groups.


\begin{PROP}{SG-B02-03-50}{Permutationsgruppen und ihre Hebungen in $\Pin(n)$ und $\Spin(n)$}
Für jedes $\sigma\in S_n$ gilt
\[
\sgn(\sigma)=\det([\sigma]),
\]
wobei $[\sigma]$ die Permutationsmatrix zu $\sigma$ bezeichnet. Somit folgt
\[
A_n=
\begin{cases}
SO(n)\cap S_n, & \text{falls }\mathbf{k}=\mathbb{R},\\[2mm]
SU(n)\cap S_n, & \text{falls }\mathbf{k}=\mathbb{C}.
\end{cases}
\]
Für $n\ge 5$ ist $A_n$ eine einfache Gruppe.

Da $\rho:\Pin(n)\to O(n)$ surjektiv ist, existieren die endlichen Untergruppen
\[
\widetilde S_n := \rho^{-1}(S_n)\le \Pin(n),
\qquad
\widetilde A_n := \rho^{-1}(A_n)\le \Spin(n),
\]
für die gilt:
\[
\rho:\widetilde S_n\to S_n,\qquad \rho:\widetilde A_n\to A_n
\]
sind surjektive Homomorphismen mit $\ker\rho\supseteq\{\pm1\}$.
Daraus folgt
\[
|\widetilde S_n|=2\cdot n!,\qquad |\widetilde A_n|=n!.
\]
Für $n\ge 4$ existiert jedoch kein Homomorphismus
\[
\tau:S_n\to \widetilde S_n,\qquad
\tau:A_n\to \widetilde A_n,
\]
so dass $\rho\circ\tau=\mathrm{Id}$ gilt.
\[
\begin{tikzcd}[row sep=large, column sep=large]
S_n \arrow[r,"\tau"] \arrow[dr, "\mathrm{Id}_{S_n}"'] & \widetilde S_n \arrow[d,"\rho"]\\
& S_n
\end{tikzcd}
\qquad
\begin{tikzcd}[row sep=large, column sep=large]
A_n \arrow[r,"\tau"] \arrow[dr, "\mathrm{Id}_{A_n}"'] & \widetilde A_n \arrow[d,"\rho"]\\
& A_n
\end{tikzcd}
\]
Analoge Überlegungen gelten auch für andere endliche Untergruppen von $O(n)$.
\end{PROP}




\begin{PROP}{SG-B02-03-51}{Gruppen $E_n$ und $E_n^0$ in der Clifford-Algebra}
In $\Cl_n^{\times}$ sei
\[
E_n := \{\pm e_{i_1}\cdots e_{i_r}\mid 1\le i_1<\cdots<i_r\le n,\;0\le r\le n\}
\]
die Untergruppe, bestehend aus allen Produkten der Basisvektoren $e_i$ und ihren Negativen.

\textbf{Eigenschaften:}
\begin{enumerate}[label=(\roman*)]
\item Der Umfang von $E_n$ ist $|E_n|=2^{n+1}$.
\item Da $\{\pm1\}\subseteq Z(E_n)$ und $e_i e_j e_i^{-1} e_j^{-1}=-1$ gilt, ist
      \[
      \bar E_n := \rho(E_n) \subseteq O(n)
      \]
      abelsch und isomorph zu $(\mathbb Z/2)^n$.
\item Die Elemente $\rho(e_r)$ sind diagonale Permutationsmatrizen mit Einträgen $\pm1$.
\item Die Untergruppe
      \[
      E_n^0 := E_n \cap \Spin(n),\qquad
      \bar E_n^0 := \rho(E_n^0)\subseteq SO(n)
      \]
      hat Ordnung $|E_n^0|=2^n$, $|\bar E_n^0|=2^{n-1}$, und
      \(\bar E_n^0\cong(\mathbb Z/2)^{n-1}\).
\item Beide Gruppen $E_n$ und $E_n^0$ sind nichtabelsch und passen in die exakten Sequenzen
      \[
      1\to \mathbb Z/2\to E_n\to (\mathbb Z/2)^n\to 1,\qquad
      1\to \mathbb Z/2\to E_n^0\to (\mathbb Z/2)^{n-1}\to 1,
      \]
      wobei der Kern $\mathbb Z/2$ jeweils das Zentrum bildet.
\end{enumerate}

Somit sind $E_n$ und $E_n^0$ \emph{extraspeciale $2$-Gruppen}, d.\,h. nichtabelsche Gruppen, deren Zentrum, Kommutator- und Frattinigruppen alle identisch mit $\mathbb Z/2$ sind.
\end{PROP}



\pagebreak


\subsection{EXERCISES}


5.1. In the Clifford algebra $\mathrm{Cl}_{n}$, let $u, v \in \mathbb{R}^{n} \subseteq \mathrm{Cl}_{n}$.\\
a) If $|u|=1$, by expressing $v$ as a sum $v_{1}+v_{2}$ with $v_{1}=t u$ and $u \cdot v_{2}=0$, find a general formula for $u v u$.\\
b) Let $\left\{u_{1}, \ldots, u_{n}\right\}$ be an orthonormal basis for $\mathbb{R}^{n}$. Show that

$$
u_{j} u_{i}= \begin{cases}-1 & \text { if } j=i, \\ -u_{i} u_{j} & \text { if } j \neq i .\end{cases}
$$

Deduce that if $i, j, k, \ell$ are distinct numbers in the range 1 to $n$ then

$$
u_{i} u_{j} u_{k} u_{\ell}=u_{k} u_{\ell} u_{i} u_{j} .
$$

c) Show that each element $A \in \mathrm{O}(n)$ induces an automorphism $A_{*}: \mathrm{Cl}_{n} \longrightarrow \mathrm{Cl}_{n}$ for which $A_{*} x=A x$ if $x \in \mathbb{R}^{n}$. Determine $A_{*}\left(e_{1} \cdots e_{n}\right)$.\\
5.2. Use the Universal Property of Theorem 5.4 in the following.\\
a) Show that the natural embedding

$$
i_{n}: \mathbb{R}^{n} \longrightarrow \mathbb{R}^{n+1} ; \quad\left[\begin{array}{c}
x_{1} \\
\vdots \\
x_{n}
\end{array}\right] \longmapsto\left[\begin{array}{c}
x_{1} \\
\vdots \\
x_{n} \\
0
\end{array}\right]
$$

induces an $\mathbb{R}$-algebra homomorphism $i_{n}^{\prime}: \mathrm{Cl}_{n} \longrightarrow \mathrm{Cl}_{n+1}$ for which $i_{n}^{\prime}(x)=i_{n}(x)$ whenever $x \in \mathbb{R}^{n}$. Show that $i_{n}^{\prime}$ is injective and determine its image $i_{n}^{\prime} \mathrm{Cl}_{n} \subseteq \mathrm{Cl}_{n+1}$.\\
b) Show that the $\mathbb{R}$-linear transformation

$$
k_{n}: \mathbb{R}^{n} \longrightarrow \mathrm{Cl}_{n+1} ; \quad k_{n}(x)=x e_{n+1},
$$

induces an $\mathbb{R}$-algebra homomorphism $k_{n}^{\prime}: \mathrm{Cl}_{n} \longrightarrow \mathrm{Cl}_{n+1}$ for which $k_{n}^{\prime}(x)=k_{n}(x)$ whenever $x \in \mathbb{R}^{n}$. Show that $k_{n}^{\prime}$ is injective with image $k_{n}^{\prime} \mathrm{Cl}_{n}=\mathrm{Cl}_{n+1}^{+}$.\\
5.3. Let $n \geqslant 1$ and $n \equiv 3 \bmod 4$. In the Clifford algebra $\mathrm{Cl}_{n}$, consider the element $\omega_{n}=e_{1} \cdots e_{n}$.\\
a) Show that $\omega_{n}$ is central in $\mathrm{Cl}_{n}$.\\
b) Show that in $\mathrm{Cl}_{n}$, the elements

$$
\Omega_{+}=\frac{1}{2}\left(1+\omega_{n}\right), \quad \Omega_{-}=\frac{1}{2}\left(1-\omega_{n}\right),
$$

are central orthogonal idempotents for which $\Omega_{+}+\Omega_{-}=1$.\\
c) Decompose $\mathrm{Cl}_{n}$ as a product $\mathrm{Cl}_{n}=A_{+} \times A_{-}$.\\
d) Use the decomposition of (c) to show that $\mathrm{Cl}_{3} \cong \mathbb{H} \times \mathbb{H}$ and $\mathrm{Cl}_{7} \cong \mathrm{M}_{8}(\mathbb{R}) \times \mathrm{M}_{8}(\mathbb{R})$.\\
5.4. For $n \geqslant 1$, define the finite Clifford group $\mathrm{ClGp}_{n}$ to be generated by $(n+1)$ elements $\varepsilon_{1}, \ldots, \varepsilon_{n}, \gamma$ with $\gamma$ central and subject to the relations

$$
\varepsilon_{1}^{2}=\cdots=\varepsilon_{n}^{2}=\gamma, \quad \varepsilon_{i} \varepsilon_{j}=\gamma \varepsilon_{j} \varepsilon_{i} \text { if } i \neq j, \quad \gamma^{2}=1 .
$$

a) Show that the order of $\mathrm{ClGp}_{n}$ is $2^{n+1}$, hence the real group algebra $\mathbb{R}\left[\mathrm{ClGp}_{n}\right]$ has dimension $2^{n+1}$.\\
b) Show that $\mathrm{ClGp}_{n}$ is isomorphic to a subgroup of $\mathrm{Cl}_{n}^{\times}$.\\
c) Find an $\mathbb{R}$-algebra homomorphism $\varphi: \mathbb{R}\left[\mathrm{ClGp}_{n}\right] \longrightarrow \mathrm{Cl}_{n}$ which is surjective.\\
d) Show that the $\mathbb{R}$-algebra $\mathbb{R}\left[\mathrm{ClGp}_{n}\right]$ is the product algebra

$$
\mathbb{R}\left[\mathrm{ClGp}_{n}\right]=A \times B,
$$

where $B=\operatorname{ker} \varphi$ and $A \cong \mathrm{Cl}_{n}$.\\
5.5. $\Delta$ Let $S=\left[s_{i j}\right]$ be a $2 m \times 2 m$ real skew symmetric matrix. Working in the Clifford algebra $\mathrm{Cl}_{2 m}$, consider the element

$$
s=\sum_{1 \leqslant i<j \leqslant 2 m} s_{i j} e_{i} e_{j}=\frac{1}{2} \sum_{\substack{1 \leqslant i, j \leqslant m \\ i \neq j}} s_{i j} e_{i} e_{j} .
$$

a) Expressing the element $s^{m}$ in terms of the basis of monomials $e_{i_{1}} \cdots e_{i_{r}}$, show that the coefficient of $e_{1} e_{2} \cdots e_{2 m}$ is the real number

$$
c_{S}=m!\sum_{\sigma \in \Pi(2 m)} \operatorname{sgn} \sigma s_{\sigma(1) \sigma(2)} \cdots s_{\sigma(2 m-1) \sigma(2 m)}
$$

where $\Pi(2 m)$ consists of all permutations $\sigma$ of $\{1,2, \ldots, 2 m\}$ satisfying the conditions

\begin{itemize}
  \item $\sigma(2 r-1)<\sigma(2 s-1)$ if $1 \leqslant r<s \leqslant m$;
  \item $\sigma(2 t-1)<\sigma(2 t)$ if $1 \leqslant t \leqslant m$.
\end{itemize}

Determine $c_{S}$ for a few small values of $m$.\\
b) Determine $c_{S}$ when $S=\left[\begin{array}{cc}O_{m} & I_{m} \\ -I_{m} & O_{m}\end{array}\right]$.\\
c) If $\Sigma_{2 m} \subseteq \mathrm{M}_{2 m}(\mathbb{R})$ is the subspace of all non-singular skew symmetric matrices, deduce that the function

$$
\mathrm{pf}: \Sigma_{2 m} \longrightarrow \mathbb{R}^{\times} ; \quad \mathrm{pf}\left(\left[s_{i j}\right]\right)=\frac{c_{S}}{m!}
$$

is continuous. This function pf is known as a Pfaffian.\\
d) Show that pf takes both positive and negative values and in fact is surjective.\\
e) Let $P=\left[p_{i j}\right] \in \mathrm{GL}_{2 m}(\mathbb{R})$. Using the well-known formula

$$
\operatorname{det} P=\sum_{\sigma \in \mathrm{S}_{2 m}} \operatorname{sgn} \sigma \prod_{k=1}^{2 m} p_{k \sigma(k)}
$$

show that

$$
\operatorname{pf}\left(P^{T} S P\right)=\operatorname{pf}(S) \operatorname{det} P
$$

5.6. a) Verify Proposition 5.31.\\
b) For $n \geqslant 2$, determine $\mathrm{Z}(\mathrm{O}(n))$.



\pagebreak


\section{Lorentz Groups}


We met the Lorentz group Lor in Section 1.5. In this chapter, we will consider this kind of group in greater generality. We merely sketch some of the details, leaving the reader to fill in the more obvious gaps. The most important example is that for which $n=3$ as this provides the geometric setting for Special Relativity. However, many of the main features can be seen in the cases $n=1,2$.

\subsection{Lorentz Groups}




\begin{DEF}{LIE-B02-09-01}{Lorentz Gruppe}
Sei \((\mathbb{R}^{n,1},\langle\cdot,\cdot\rangle)\) der reelle Vektorraum der Dimension
\(n+1\) mit dem Lorentz-Innenprodukt
\[
\langle \mathbf{X},\mathbf{Y}\rangle
= \mathbf{X}^{T} Q_{n,1} \mathbf{Y}
= x_1y_1 + \cdots + x_ny_n - x_{n+1}y_{n+1},
\]
wobei
\[
Q_{n,1} := \operatorname{diag}(\underbrace{1,\ldots,1}_{n},-1)
\in M_{n+1}(\mathbb{R})
\]
die Signatur \((n,1)\) festlegt.

\textbf{Definition.}
Die \emph{Lorentz-Gruppe} ist die Gruppe aller linearen Isometrien des
Lorentzraums \(\mathbb{R}^{n,1}\), d.\,h.
\[
O(n,1)
:= \{A\in GL_{n+1}(\mathbb{R}) \mid A^{T} Q_{n,1} A = Q_{n,1}\}.
\]
Sie ist eine reelle Matrixgruppe und eine geschlossene Untergruppe von
\(GL_{n+1}(\mathbb{R})\). Das zugehörige Lie-Algebra-Objekt ist
\[
\mathfrak{o}(n,1)
:= \{X\in M_{n+1}(\mathbb{R}) \mid X^{T} Q_{n,1} + Q_{n,1} X = 0\}.
\]

\textbf{Bemerkung (Spezielle Lorentz-Gruppe).}
Die Untergruppe
\[
SO(n,1)
:= O(n,1)\cap SL_{n+1}(\mathbb{R})
= \{A\in O(n,1)\mid \det A = 1\}
\]
heißt \emph{spezielle Lorentz-Gruppe}.  
Sie besteht aus allen orientierungserhaltenden Lorentz-Isometrien
und bildet die Identitätskomponente von \(O(n,1)\).

\textbf{Physikalische Konvention.}
In der speziellen Relativitätstheorie wird oft die
\emph{verbundene, zeit-orientierungserhaltende Komponente} der
spezialen Lorentz-Gruppe verwendet, bezeichnet durch
\[
SO^{+}(n,1)
:= \{A\in SO(n,1)\mid A_{n+1,n+1}>0\},
\]
welche den sog.\ \emph{echten Lorentz-Transformationsgruppen}
entspricht.
\end{DEF}





For $n \geqslant 1$, consider the non-singular symmetric matrix

$$
Q_{n, 1}=\operatorname{diag}(\underbrace{1, \ldots, 1}_{n},-1) \in \mathrm{M}_{n+1}(\mathbb{R}),
$$

which defines an inner product on $\mathbb{R}^{n+1}$ given by

$$
\langle\mathbf{x}, \mathbf{y}\rangle=\mathbf{x}^{T} Q_{n, 1} \mathbf{y}=x_{1} y_{1}+\cdots+x_{n} y_{n}-x_{n+1} y_{n+1} .
$$

This is the Lorentz inner product on $\mathbb{R}^{n+1}$ and it is standard to denote this inner product space by $\mathbb{R}^{n, 1}$. Actually, physicists often adopt different sign conventions and use one of the following formulæ in place of Equation (6.1):

$$
\langle\mathbf{x}, \mathbf{y}\rangle=\left\{\begin{array}{c}
-x_{1} y_{1}+\cdots+x_{n} y_{n}+x_{n+1} y_{n+1} \\
-x_{1} y_{1}+\cdots-x_{n} y_{n}+x_{n+1} y_{n+1} \\
x_{1} y_{1}+\cdots-x_{n} y_{n}-x_{n+1} y_{n+1}
\end{array}\right.
$$

This affects the signs in the following definitions. We will also use $t$ in place of $x_{n+1}$ since in Relativity, this coordinate is related to the time measurements while the others are related to spatial ones.

We will often write elements of $\mathbb{R}^{n, 1}$ in the form

$$
\mathbf{X}=\left[\begin{array}{c}
\mathbf{x} \\
t
\end{array}\right] \quad\left(\mathbf{x} \in \mathbb{R}^{n}, t \in \mathbb{R}\right) .
$$

Then

$$
\left\langle\mathbf{X}_{1}, \mathbf{X}_{2}\right\rangle=\mathbf{x}_{1} \cdot \mathbf{x}_{2}-t_{1} t_{2} .
$$




\begin{DEF}{LIE-B02-09-02}{Lorentz-Innenprodukt und Raum $\mathbb{R}^{n,1}$}
Für \(n\ge1\) sei
\[
Q_{n,1} := \operatorname{diag}(\underbrace{1,\ldots,1}_{n},-1)
  \in M_{n+1}(\mathbb{R})
\]
die nicht-singuläre symmetrische Matrix, welche ein \emph{bilineares
Formular} auf \(\mathbb{R}^{n+1}\) definiert durch
\[
\langle \mathbf{x},\mathbf{y}\rangle
= \mathbf{x}^{T} Q_{n,1} \mathbf{y}
= x_1y_1+\cdots+x_ny_n - x_{n+1}y_{n+1}.
\]
Dieses Produkt heißt das \textbf{Lorentz-Innenprodukt} auf
\(\mathbb{R}^{n+1}\). Den so ausgestatteten Raum bezeichnet man mit
\[
\mathbb{R}^{n,1} := (\mathbb{R}^{n+1},\langle\cdot,\cdot\rangle),
\]
also einen reellen Vektorraum der Signatur \((n,1)\).

\textbf{Bemerkung (physikalische Konventionen).}
In der Literatur, insbesondere der Relativitätstheorie, werden häufig
abweichende Vorzeichenkonventionen verwendet, z.\,B.
\[
\langle\mathbf{x},\mathbf{y}\rangle =
\begin{cases}
 -x_1y_1+\cdots+x_ny_n+x_{n+1}y_{n+1},\\[1mm]
 -x_1y_1+\cdots-x_ny_n+x_{n+1}y_{n+1},\\[1mm]
  x_1y_1+\cdots-x_ny_n-x_{n+1}y_{n+1}.
\end{cases}
\]
Diese Konvention beeinflusst nur die Vorzeichen der folgenden Definitionen.

\textbf{Darstellung.}
Man schreibt häufig
\[
\mathbf{X}=
\begin{bmatrix}
\mathbf{x}\\ t
\end{bmatrix}
\qquad
(\mathbf{x}\in\mathbb{R}^n,\; t\in\mathbb{R}),
\]
wobei \(t\) der Zeitkoordinate entspricht, während \(\mathbf{x}\)
den räumlichen Anteil beschreibt. Dann gilt für
\(\mathbf{X}_1,\mathbf{X}_2\in\mathbb{R}^{n,1}\):
\[
\langle \mathbf{X}_1,\mathbf{X}_2\rangle
= \mathbf{x}_1\cdot \mathbf{x}_2 - t_1 t_2,
\]
wobei \(\mathbf{x}_1\cdot\mathbf{x}_2\) das euklidische Skalarprodukt in
\(\mathbb{R}^n\) bezeichnet.
\end{DEF}



Definition 6.1


A non-zero vector $\mathbf{X} \in \mathbb{R}^{n, 1}$ is called
\begin{itemize}
  \item spacelike if $\langle\mathbf{X}, \mathbf{X}\rangle>0$ or equivalently $|\mathbf{x}|^{2}>t^{2}$;
  \item timelike if $\langle\mathbf{X}, \mathbf{X}\rangle<0$ or equivalently $|\mathbf{x}|^{2}<t^{2}$;
  \item a null vector if $\langle\mathbf{X}, \mathbf{X}\rangle=0$ or equivalently $|\mathbf{x}|^{2}=t^{2}$.
\end{itemize}


We can view $\mathbb{R}^{n}$ as sitting inside $\mathbb{R}^{n, 1}$ as the subset of all spacelike vectors of the form
$$
\mathbf{X}=\left[\begin{array}{c}
\mathbf{x} \\
0
\end{array}\right]=\left[\begin{array}{c}
x_{1} \\
\vdots \\
x_{n} \\
0
\end{array}\right]
$$
and we usually identify $\mathbf{x} \in \mathbb{R}^{n}$ with the element $\mathbf{X}=\left[\begin{array}{c}\mathbf{x} \\ 0\end{array}\right] \in \mathbb{R}^{n, 1}$. Of course we then have $\langle\mathbf{X}, \mathbf{X}\rangle=\mathbf{x} \cdot \mathbf{x}$, so this embedding of $\mathbb{R}^{n}$ into $\mathbb{R}^{n, 1}$ is an isometry.




\begin{DEF}{LIE-B02-09-03}{Klassifikation von Vektoren im Lorentzraum}
Ein nichtverschwindender Vektor
\(\mathbf{X}\in\mathbb{R}^{n,1}\) wird nach dem Vorzeichen seines
Lorentz-Skalarprodukts klassifiziert:
\begin{itemize}
  \item \textbf{Spacelike (raumartig)}, falls
  \(\langle \mathbf{X},\mathbf{X}\rangle > 0\),
  äquivalent zu \(|\mathbf{x}|^2 > t^2\).
  \item \textbf{Timelike (zeitartig)}, falls
  \(\langle \mathbf{X},\mathbf{X}\rangle < 0\),
  äquivalent zu \(|\mathbf{x}|^2 < t^2\).
  \item \textbf{Null- oder lichtartig}, falls
  \(\langle \mathbf{X},\mathbf{X}\rangle = 0\),
  äquivalent zu \(|\mathbf{x}|^2 = t^2\).
\end{itemize}

\textbf{Einbettung von \(\mathbb{R}^n\) in \(\mathbb{R}^{n,1}\).}
Wir können den euklidischen Raum \(\mathbb{R}^n\) als Untermenge von
\(\mathbb{R}^{n,1}\) auffassen, bestehend aus allen
raumartigen Vektoren der Form
\[
\mathbf{X}=
\begin{bmatrix}
\mathbf{x}\\0
\end{bmatrix}
=
\begin{bmatrix}
x_1\\
\vdots\\
x_n\\
0
\end{bmatrix}.
\]
Wir identifizieren daher jedes \(\mathbf{x}\in\mathbb{R}^n\)
mit dem Element
\(\mathbf{X}=\begin{bmatrix}\mathbf{x}\\0\end{bmatrix}\in\mathbb{R}^{n,1}\).
Dabei gilt
\[
\langle \mathbf{X},\mathbf{X}\rangle
= \mathbf{x}\cdot \mathbf{x},
\]
sodass diese Einbettung eine \emph{Isometrie}
zwischen \((\mathbb{R}^n,\cdot)\) und der Menge der spacelike Vektoren in
\((\mathbb{R}^{n,1},\langle\cdot,\cdot\rangle)\) ist.
\end{DEF}



If we represent $\mathbb{R}^{n, 1}$ by the ( $x, t$ )-plane with $\mathbb{R}^{n}$ as the $x$-axis, then the $t$-axis is timelike, while the null vectors correspond to the points lying on the line $x=t$. In fact, the null vectors in $\mathbb{R}^{n, 1}$ lie in the set
$$
\mathcal{H}_{n, 1}(0)=\left\{\mathbf{X} \in \mathbb{R}^{n, 1}: \mathbf{X} \neq \mathbf{0},\langle\mathbf{X}, \mathbf{X}\rangle=\mathbf{x} \cdot \mathbf{x}-t^{2}=0\right\} .
$$

For $n>1, \mathcal{H}_{n, 1}(0)$ has two path connected components,
$$
\mathcal{H}_{n, 1}^{+}(0)=\left\{\mathbf{X} \in \mathcal{H}_{n, 1}(0): t>0\right\}, \quad \mathcal{H}_{n, 1}^{-}(0)=\left\{\mathbf{X} \in \mathcal{H}_{n, 1}(0): t<0\right\} .
$$
These are often referred to as the positive or future pointing light cone and the negative or past pointing light cone respectively since in Relativity these represent points moving at the speed of light.

For each positive real number $r$, the hyperboloid
$$
\mathcal{H}_{n, 1}(r)=\left\{\mathbf{X} \in \mathbb{R}^{n, 1}:\langle\mathbf{X}, \mathbf{X}\rangle=\mathbf{x} \cdot \mathbf{x}-t^{2}=-r\right\}
$$
has two path connected components
$$
\mathcal{H}_{n, 1}^{+}(r)=\left\{\mathbf{X} \in \mathcal{H}_{n, 1}(r): t>0\right\}, \quad \mathcal{H}_{n, 1}^{-}(r)=\left\{\mathbf{X} \in \mathcal{H}_{n, 1}(r): t<0\right\} .
$$



\begin{DEF}{LIE-B02-09-04}{Lichtkegel und Hyperboloide im Lorentzraum}
Betrachte den Lorentzraum \(\mathbb{R}^{n,1}\) mit Koordinaten
\(\mathbf{X}=(\mathbf{x},t)\), wobei \(\mathbf{x}\in\mathbb{R}^{n}\) und \(t\in\mathbb{R}\).
Das Lorentz-Innenprodukt ist gegeben durch
\[
\langle \mathbf{X},\mathbf{X}\rangle = \mathbf{x}\cdot\mathbf{x} - t^2.
\]

\textbf{(i) Lichtkegel.}
Die Menge aller \emph{nullartigen} (lichtartigen) Vektoren wird durch den
\emph{Lichtkegel mit Scheitel im Ursprung} beschrieben:
\[
\mathcal{H}_{n,1}(0)
= \{\mathbf{X}\in\mathbb{R}^{n,1}\setminus\{0\} \mid
\langle \mathbf{X},\mathbf{X}\rangle = 0\}
= \{\mathbf{X}=(\mathbf{x},t) \mid |\mathbf{x}|^2 = t^2\}.
\]
Für \(n>1\) zerfällt \(\mathcal{H}_{n,1}(0)\) in zwei wegzusammenhängende Komponenten:
\[
\mathcal{H}_{n,1}^{+}(0)
= \{\mathbf{X}\in\mathcal{H}_{n,1}(0)\mid t>0\}, \qquad
\mathcal{H}_{n,1}^{-}(0)
= \{\mathbf{X}\in\mathcal{H}_{n,1}(0)\mid t<0\}.
\]
\(\mathcal{H}_{n,1}^{+}(0)\) heißt der \textbf{zukünftige Lichtkegel},
\(\mathcal{H}_{n,1}^{-}(0)\) der \textbf{vergangene Lichtkegel}.
In der Relativitätstheorie repräsentieren diese Mengen Punkte,
die sich mit Lichtgeschwindigkeit in Zukunft bzw.\ Vergangenheit bewegen.

\textbf{(ii) Hyperboloide.}
Für ein festes \(r>0\) definiert man die \emph{Hyperboloide konstanter Lorentz-Norm} durch
\[
\mathcal{H}_{n,1}(r)
= \{\mathbf{X}\in\mathbb{R}^{n,1}\mid
\langle \mathbf{X},\mathbf{X}\rangle = -r\}
= \{\mathbf{X}=(\mathbf{x},t)\mid t^2 - |\mathbf{x}|^2 = r\}.
\]
Auch diese besitzen zwei wegzusammenhängende Komponenten:
\[
\mathcal{H}_{n,1}^{+}(r)
= \{\mathbf{X}\in\mathcal{H}_{n,1}(r)\mid t>0\}, \qquad
\mathcal{H}_{n,1}^{-}(r)
= \{\mathbf{X}\in\mathcal{H}_{n,1}(r)\mid t<0\}.
\]
Die obere Komponente \(\mathcal{H}_{n,1}^{+}(r)\) wird als
\textbf{Einblatt-Hyperboloid der Zukunft} bezeichnet und spielt eine
zentrale Rolle in der hyperbolischen Geometrie und der speziellen Relativitätstheorie.
\end{DEF}


For later use we extend the Lorentz inner product to an $\mathbb{R}$-bilinear inner product on $\mathbb{C}^{n+1}$ by setting
$$
\langle\mathbf{U}, \mathbf{V}\rangle=\mathbf{U}^{*} Q_{n, 1} \mathbf{V}=\mathbf{u}^{*} \mathbf{v}-\bar{u} v=\mathbf{u} \cdot \mathbf{v}-\bar{u} v
$$
for each pair of vectors
$$
\mathbf{U}=\left[\begin{array}{l}
\mathbf{u} \\
u
\end{array}\right], \quad \mathbf{V}=\left[\begin{array}{l}
\mathbf{v} \\
v
\end{array}\right] \in \mathbb{C}^{n, 1} .
$$
We will denote the $\mathbb{C}$-vector space $\mathbb{C}^{n+1}$ with this inner product by $\mathbb{C}^{n, 1}$. Then $\langle$,$\rangle has all the properties of a hermitian inner product except that the real$ number $\langle\mathbf{U}, \mathbf{U}\rangle$ is not necessarily non-negative. In particular, for $z, w \in \mathbb{C}$,
$$
\langle z \mathbf{U}, w \mathbf{V}\rangle=\bar{z} w\langle\mathbf{U}, \mathbf{V}\rangle .
$$
We will refer to a non-zero $\mathbf{U} \in \mathbb{C}^{n, 1}$ as spacelike, timelike or null according to whether $\langle\mathbf{U}, \mathbf{U}\rangle>0,\langle\mathbf{U}, \mathbf{U}\rangle<0$ or $\langle\mathbf{U}, \mathbf{U}\rangle=0$ as we do for vectors in $\mathbb{R}^{n, 1}$.



\begin{DEF}{LIE-B02-09-05}{Komplexifizierter Lorentz-Raum und hermitische Fortsetzung}
Zur späteren Verwendung erweitern wir das Lorentz-Skalarprodukt auf einen
\(\mathbb{R}\)-bilinearen Ausdruck auf dem komplexen Vektorraum
\(\mathbb{C}^{n+1}\).
Für Vektoren
\[
\mathbf{U}=
\begin{bmatrix}
\mathbf{u}\\u
\end{bmatrix},
\quad
\mathbf{V}=
\begin{bmatrix}
\mathbf{v}\\v
\end{bmatrix}
\in\mathbb{C}^{n+1},
\]
definieren wir
\[
\langle \mathbf{U},\mathbf{V}\rangle
=\mathbf{U}^{*} Q_{n,1}\mathbf{V}
=\mathbf{u}^{*}\mathbf{v}-\bar{u}v
=\mathbf{u}\cdot\mathbf{v}-\bar{u}v,
\]
wobei \(Q_{n,1}=\operatorname{diag}(1,\dots,1,-1)\in M_{n+1}(\mathbb{R})\)
die Lorentz-Signatur \((n,1)\) festlegt.

\textbf{Bemerkungen.}
\begin{itemize}
  \item Der Raum \((\mathbb{C}^{n+1},\langle\cdot,\cdot\rangle)\) wird als
  \emph{komplexifizierter Lorentzraum} bezeichnet und mit
  \(\mathbb{C}^{n,1}\) bezeichnet.
  \item Das Produkt \(\langle\cdot,\cdot\rangle\) erfüllt alle Eigenschaften eines
  hermiteschen Skalarprodukts, mit Ausnahme der Positivität:
  \(\langle \mathbf{U},\mathbf{U}\rangle\) kann positiv, negativ oder null sein.
  \item Für \(z,w\in\mathbb{C}\) gilt die Hermitizitätsrelation
  \[
  \langle z\mathbf{U}, w\mathbf{V}\rangle
  =\bar{z}w\,\langle\mathbf{U},\mathbf{V}\rangle.
  \]
\end{itemize}

\textbf{Klassifikation der Vektoren.}
Ein nichtverschwindender Vektor
\(\mathbf{U}\in\mathbb{C}^{n,1}\) heißt
\begin{itemize}
  \item \emph{raumartig (spacelike)}, falls \(\langle \mathbf{U},\mathbf{U}\rangle>0\),
  \item \emph{zeitartig (timelike)}, falls \(\langle \mathbf{U},\mathbf{U}\rangle<0\),
  \item \emph{nullartig (lightlike)}, falls \(\langle \mathbf{U},\mathbf{U}\rangle=0\).
\end{itemize}

\textbf{Interpretation.}
Die Raum-Zeit-Geometrie über \(\mathbb{R}^{n,1}\) lässt sich so auf
den komplexen Raum \(\mathbb{C}^{n,1}\) ausdehnen, wobei die Lorentz-Struktur
als \(\mathbb{C}\)-bilineare oder hermitesche Form fortgesetzt wird.
Dies ist insbesondere die Grundlage für die Darstellung komplexer
Spin- und Lorentzgruppen (z.\,B. \(\mathrm{SO}(n,1;\mathbb{C})\),
\(\mathrm{SU}(n,1)\)).
\end{DEF}



We will make repeated use of the following result.




Proposition 6.2


Let $\mathbf{U}=\left[\begin{array}{l}\mathbf{u} \\ u\end{array}\right]$ and $\mathbf{X}=\left[\begin{array}{l}\mathbf{x} \\ t\end{array}\right] \in \mathbb{C}^{n, 1}$ be non-zero vectors with $\langle\mathbf{U}, \mathbf{X}\rangle=0$.\\
i) If $\mathbf{U}$ is timelike then $\langle\mathbf{X}, \mathbf{X}\rangle>0$.\\
ii) If $\mathbf{U}$ is null and $\mathbf{X}$ is not a complex multiple of $\mathbf{U}$, then $\langle\mathbf{X}, \mathbf{X}\rangle>0$.



\begin{PROP}{LIE-B02-09-06}{Orthogonalität zu zeitartigen bzw.\ nullartigen Vektoren}
Seien
\[
\mathbf{U}=
\begin{bmatrix}\mathbf{u}\\ u\end{bmatrix},
\qquad
\mathbf{X}=
\begin{bmatrix}\mathbf{x}\\ t\end{bmatrix}
\in \mathbb{C}^{n,1}\setminus\{0\}
\]
mit $\langle \mathbf{U},\mathbf{X}\rangle=0$, d.\,h.
\[
\mathbf{u}^*\mathbf{x}-\bar u\,t=0\quad\Longleftrightarrow\quad
\mathbf{u}^*\mathbf{x}=\bar u\,t.
\]
Dann gilt:
\begin{enumerate}\itemsep0.3em
\item[(i)] Ist $\mathbf{U}$ \emph{zeitartig}, also $\langle \mathbf{U},\mathbf{U}\rangle
= \|\mathbf{u}\|^2-|u|^2<0$, so ist $\langle \mathbf{X},\mathbf{X}\rangle
=\|\mathbf{x}\|^2-|t|^2>0$.
\item[(ii)] Ist $\mathbf{U}$ \emph{nullartig}, also $\langle \mathbf{U},\mathbf{U}\rangle
=0$, und ist $\mathbf{X}$ \emph{kein} komplexes Vielfaches von $\mathbf{U}$, so ist
$\langle \mathbf{X},\mathbf{X}\rangle>0$.
\end{enumerate}
\end{PROP}


Proof


\begin{PROOF}{LIE-B02-09-07}{P: Orthogonalität zu zeitartigen bzw.\ nullartigen Vektoren}
When (i) holds, then $|\mathbf{u}|^{2}-|u|^{2}<0$, so $u \neq 0$; when (ii) holds then $\mathbf{U} \neq \mathbf{0}$ and $|u|^{2}=|u|^{2}$, also giving $u \neq 0$. Since
$$
\mathbf{u} \cdot \mathbf{x}-\bar{u} t=0
$$
in either case we have
$$
\langle\mathbf{X}, \mathbf{X}\rangle=|\mathbf{x}|^{2}-|t|^{2}=|\mathbf{x}|^{2}-\frac{|\mathbf{u} \cdot \mathbf{x}|^{2}}{|u|^{2}} .
$$
In case (i), this gives
$$
\langle\mathbf{X}, \mathbf{X}\rangle>|\mathbf{x}|^{2}-\frac{|\mathbf{u} \cdot \mathbf{x}|^{2}}{|\mathbf{u}|^{2}} \geqslant 0,
$$
since $|\mathbf{u} \cdot \mathbf{x}| \leqslant|\mathbf{u}||\mathbf{x}|$. In case (ii),
$$
\langle\mathbf{X}, \mathbf{X}\rangle=|\mathbf{x}|^{2}-\frac{|\mathbf{u} \cdot \mathbf{x}|^{2}}{|\mathbf{u}|^{2}} \geqslant 0
$$
with equality if and only if $|\mathbf{u} \cdot \mathbf{x}|=|\mathbf{u}||\mathbf{x}|$, i.e., $\mathbf{x}=z \mathbf{u}$ for some $z \in \mathbb{C}$ with $|z|=1$. In the latter case, from Equation (6.3) we would have
$$
t \bar{u}=z|u|^{2}=z u \bar{u}
$$
so $t=z u$, in turn giving $\mathbf{X}=z \mathbf{U}$ and so $\langle\mathbf{X}, \mathbf{X}\rangle=0$.
\end{PROOF}




We can consider the closed subgroup
$$
\mathrm{O}(n, 1)=\left\{A \in \mathrm{GL}_{n+1}(\mathbb{R}): A^{T} Q_{n, 1} A=Q_{n, 1}\right\} \leqslant \mathrm{GL}_{n+1}(\mathbb{R})
$$


Note that $A \in \mathrm{O}(n, 1)$ if and only if for all $\mathbf{x}, \mathbf{y} \in \mathbb{R}^{n, 1},\langle A \mathbf{x}, A \mathbf{y}\rangle=\langle\mathbf{x}, \mathbf{y}\rangle$, so $\mathrm{O}(n, 1)$ consists of all the Lorentzian isometries of $\mathbb{R}^{n, 1}$. Notice that for $A \in \mathrm{O}(n, 1)$ we have $(\operatorname{det} A)^{2}=1$, so $\operatorname{det} A= \pm 1$. Also $A \in \mathrm{O}(n, 1)$ preserves each of the sets of spacelike, timelike and null vectors; it also either preserves or interchanges the positive and negative light cones. It is standard to set

$$
\mathrm{SO}(n, 1)=\{A \in \mathrm{O}(n, 1): \operatorname{det} A=1\}=\mathrm{O}(n, 1) \cap \mathrm{SL}_{n+1}(\mathbb{R}) \leqslant \mathrm{O}(n, 1)
$$

However, $\mathrm{SO}(n, 1)$ is not connected since its elements can interchange the sets $\mathcal{H}_{n, 1}^{ \pm}(r)$ for $r \geqslant 0$. We define the Lorentz group of $\mathbb{R}^{n, 1}$ to be the closed subgroup of $\operatorname{SO}(n, 1)$ preserving each of the connected sets $\mathcal{H}_{n, 1}^{ \pm}(1)$,

$$
\operatorname{Lor}(n, 1)=\left\{A \in \mathrm{SO}(n, 1): A \mathcal{H}_{n, 1}^{ \pm}(1)=\mathcal{H}_{n, 1}^{ \pm}(1)\right\} \leqslant \mathrm{SO}(n, 1)
$$



\begin{DEF}{LIE-B02-09-08}{Lorentz-Gruppe und spezielle orthogonale Gruppe \(\mathrm{SO}(n,1)\)}
Sei
\[
Q_{n,1}=\operatorname{diag}(1,\dots,1,-1)\in M_{n+1}(\mathbb{R})
\]
die Lorentz-Matrix der Signatur \((n,1)\). Wir definieren die Gruppe aller
reell-linearen Transformationen, die das Lorentz-Skalarprodukt erhalten, durch
\[
\mathrm{O}(n,1)
:=\{A\in \mathrm{GL}_{n+1}(\mathbb{R}) \mid A^{T}Q_{n,1}A = Q_{n,1}\}
\leqslant \mathrm{GL}_{n+1}(\mathbb{R}).
\]

\textbf{Eigenschaften.}
\begin{itemize}
  \item Für alle \(\mathbf{x},\mathbf{y}\in\mathbb{R}^{n,1}\) gilt
  \[
  A\in\mathrm{O}(n,1)\quad\Longleftrightarrow\quad
  \langle A\mathbf{x},A\mathbf{y}\rangle = \langle \mathbf{x},\mathbf{y}\rangle.
  \]
  Damit besteht \(\mathrm{O}(n,1)\) genau aus den Lorentz-Isometrien
  des Minkowski-Raumes \(\mathbb{R}^{n,1}\).
  \item Für jedes \(A\in\mathrm{O}(n,1)\) gilt \((\det A)^2=1\), also
  \(\det A=\pm 1\).
  \item Elemente von \(\mathrm{O}(n,1)\) erhalten die Klassifikation von
  Vektoren (raumartig, zeitartig, nullartig) und können gegebenenfalls
  die Vorzeichenorientierung der Lichtkegel vertauschen.
\end{itemize}

\textbf{Definition der speziellen Lorentz-Gruppe.}
Die \emph{spezielle Lorentz-Gruppe} wird definiert als
\[
\mathrm{SO}(n,1)
:=\{A\in \mathrm{O}(n,1)\mid \det A=1\}
=\mathrm{O}(n,1)\cap \mathrm{SL}_{n+1}(\mathbb{R}).
\]
Diese Gruppe besteht aus allen orientierungserhaltenden Lorentz-Isometrien.

\textbf{Bemerkung zur Zusammenhangskomponente.}
Die Gruppe \(\mathrm{SO}(n,1)\) ist nicht zusammenhängend:
ihre Elemente können die beiden Zusammenhangskomponenten
\(\mathcal{H}_{n,1}^{\pm}(r)\) (für \(r\ge 0\))
vertauschen. Daher definiert man den „zeitorientierten“ Unterraum:

\[
\operatorname{Lor}(n,1)
:=\{A\in \mathrm{SO}(n,1)\mid
A\,\mathcal{H}_{n,1}^{\pm}(1)
=\mathcal{H}_{n,1}^{\pm}(1)\}
\leqslant \mathrm{SO}(n,1).
\]

\textbf{Interpretation.}
\(\operatorname{Lor}(n,1)\) ist die \emph{(zusammenhängende) Lorentz-Gruppe}
des Minkowski-Raumes \(\mathbb{R}^{n,1}\), die sowohl Orientierung als auch
Zeitorientierung erhält. Sie entspricht der Identitätskomponente der
Gruppe \(\mathrm{O}(n,1)\).
\end{DEF}




Proposition 6.3

For $r \geqslant 0$,
$$
\operatorname{Lor}(n, 1)=\left\{A \in \mathrm{SO}(n, 1): A \mathcal{H}_{n, 1}^{ \pm}(r)=\mathcal{H}_{n, 1}^{ \pm}(r)\right\} \leqslant \mathrm{SO}(n, 1)
$$



\begin{PROP}{LIE-B02-09-09}{Unabhängigkeit von \(r\) in der Definition von \(\operatorname{Lor}(n,1)\)}
Für jedes \(r\ge 0\) gilt
\[
\operatorname{Lor}(n,1)
=\{A\in \mathrm{SO}(n,1):A\,\mathcal H_{n,1}^{\pm}(r)=\mathcal H_{n,1}^{\pm}(r)\}
\le \mathrm{SO}(n,1).
\]
\end{PROP}


Proof


\begin{PROOF}{LIE-B02-09-10}{P: Unabhängigkeit von \(r\) in der Definition von \(\operatorname{Lor}(n,1)\)}
It is easy to verify this for $r>0$. For $r=0$, notice that every null vector $\mathbf{U}=\left[\begin{array}{l}\mathbf{u} \\ u\end{array}\right]$ can be expressed as a limit
$$
\mathbf{U}=\lim _{k \rightarrow \infty} \mathbf{V}_{k}
$$
with each $\mathbf{V}_{k}$ timelike. We can even arrange that for every $k, \mathbf{V}_{k}=\left[\begin{array}{l}\mathbf{v}_{k} \\ v_{k}\end{array}\right]$ with $u v_{k}>0$ (or $u v_{k}<0$ ). Since the action of each $A \in \operatorname{Lor}(n, 1)$ on $\mathbb{R}^{n, 1}$ is continuous,
$$
A \mathbf{U}=\lim _{k \rightarrow \infty} A \mathbf{V}_{k},
$$
from which we see that $A$ preserves each of the sets $\mathcal{H}_{n, 1}^{ \pm}(0)$.

On the other hand, suppose that $A \in \mathrm{SO}(n, 1)$ preserves each of the connected sets $\mathcal{H}_{n, 1}^{ \pm}(0)$. If $A$ fails to preserve one (and hence all) of the sets $\mathcal{H}_{n, 1}^{+}(s)$ for $s>0$, then as any $\mathbf{U} \in \mathcal{K}_{n, 1}^{+}(0)$ is a limit
$$
\mathbf{U}=\lim _{k \rightarrow \infty} \mathbf{V}_{k}
$$
with $\mathbf{V}_{k} \in \mathcal{H}_{n, 1}^{+}\left(s_{k}\right)$ for $s_{k}>0$, we have
$$
A \mathbf{U}=\lim _{k \rightarrow \infty} A \mathbf{V}_{k}
$$
where $A \mathbf{V}_{k} \in \mathcal{K}_{n, 1}^{-}\left(s_{k}^{\prime}\right)$ for $s_{k}^{\prime}>0$. This implies that $A \mathbf{U} \in \mathcal{H}_{n, 1}^{-}(0)$, contradicting the original assumption on $A$.
\end{PROOF}




Notice that $A^{*}=A^{T}$ if $A \in \operatorname{Lor}(n, 1)$, so if $\mathbf{U}, \mathbf{V} \in \mathbb{C}^{n, 1}$, we have
$$
\langle A \mathbf{U}, A \mathbf{V}\rangle=\mathbf{U}^{*} A^{*} Q_{n, 1} A \mathbf{V}=\mathbf{U}^{*} Q_{n, 1} \mathbf{V}=\langle\mathbf{U}, \mathbf{V}\rangle
$$


\begin{CONC}{LIE-B02-09-11}{Erhaltung des Lorentz-Produkts durch Elemente der Lorentz-Gruppe}
Für jedes \(A\in \operatorname{Lor}(n,1)\) gilt \(A^{*}=A^{T}\), da \(A\) reell ist.
Damit folgt für alle \(\mathbf{U},\mathbf{V}\in \mathbb{C}^{n,1}\):
\[
\langle A\mathbf{U},A\mathbf{V}\rangle
=\mathbf{U}^{*}A^{*}Q_{n,1}A\mathbf{V}
=\mathbf{U}^{*}Q_{n,1}\mathbf{V}
=\langle \mathbf{U},\mathbf{V}\rangle.
\]
Somit wirken die Elemente von \(\operatorname{Lor}(n,1)\) als
\emph{isometrische lineare Transformationen} auf dem komplexifizierten
Minkowski-Raum \(\mathbb{C}^{n,1}\).
\end{CONC}



When $n=3, \operatorname{Lor}(3,1)=$ Lor in the notation of Section 1.5.



Example 6.4


We have
$$
\begin{aligned}
\operatorname{Lor}(1,1) & =\left\{\left[\begin{array}{cc}
\cosh s & \sinh s \\
\sinh s & \cosh s
\end{array}\right]: s \in \mathbb{R}\right\} \\
\operatorname{SO}(1,1) & =\operatorname{Lor}(1,1) \cup\left[\begin{array}{rr}
-1 & 0 \\
0 & -1
\end{array}\right] \operatorname{Lor}(1,1) \\
\mathrm{O}(1,1) & =\operatorname{SO}(1,1) \cup\left[\begin{array}{rr}
1 & 0 \\
0 & -1
\end{array}\right] \operatorname{SO}(1,1)
\end{aligned}
$$


Proof


If $A \in \mathrm{O}(1,1)$ has $\operatorname{det} A=-1$ then

$$
\left[\begin{array}{rr}
1 & 0 \\
0 & -1
\end{array}\right] A \in \mathrm{SO}(1,1)
$$
so it suffices to determine $\mathrm{SO}(1,1)$. Similarly, if $B \in \mathrm{SO}(1,1)$ with $B \mathcal{H}_{n, 1}^{+}(0)= \mathcal{H}_{n, 1}^{-}(0)$ then
$$
\left[\begin{array}{rr}
-1 & 0 \\
0 & -1
\end{array}\right] B \in \operatorname{Lor}(1,1)
$$
so it suffices to determine $\operatorname{Lor}(1,1)$.


Let $\left[\begin{array}{ll}a & b \\ c & d\end{array}\right] \in \operatorname{Lor}(1,1)$. Since
$$
\left[\begin{array}{ll}
a & b \\
c & d
\end{array}\right]\left[\begin{array}{l}
0 \\
1
\end{array}\right]=\left[\begin{array}{l}
b \\
d
\end{array}\right]
$$
we must have $d>0$. Also, the simultaneous equations
$$
\left\{\begin{array}{l}
a^{2}-c^{2}=1 \\
d^{2}-b^{2}=1 \\
c d-a b=0 \\
a d-b c=1
\end{array}\right.
$$
are satisfied. Putting $d=\cosh s$ and $b=\sinh s$ for suitable $s \in \mathbb{R}$, we easily obtain $a=\cosh s, c=\sinh s$. Hence every element of $\operatorname{Lor}(1,1)$ has the form $\left[\begin{array}{ll}\cosh s & \sinh s \\ \sinh s & \cosh s\end{array}\right]$ for some $s \in \mathbb{R}$.



\begin{EXA}{LIE-B02-09-12}{Lorentz-Gruppe im zweidimensionalen Fall}
Im zweidimensionalen Minkowski-Raum \(\mathbb{R}^{1,1}\) gilt:
\[
\begin{aligned}
\operatorname{Lor}(1,1)
&=\left\{
\begin{bmatrix}
\cosh s & \sinh s\\
\sinh s & \cosh s
\end{bmatrix}
: s\in\mathbb{R}
\right\},\\[1em]
\operatorname{SO}(1,1)
&=\operatorname{Lor}(1,1)\,
\cup\,
\begin{bmatrix}
-1 & 0\\
0 & -1
\end{bmatrix}
\operatorname{Lor}(1,1),\\[1em]
\mathrm{O}(1,1)
&=\operatorname{SO}(1,1)\,
\cup\,
\begin{bmatrix}
1 & 0\\
0 & -1
\end{bmatrix}
\operatorname{SO}(1,1).
\end{aligned}
\]

\begin{proof}
Sei \(A\in \mathrm{O}(1,1)\) mit \(\det A=-1\). Dann gilt
\[
\begin{bmatrix}
1 & 0\\
0 & -1
\end{bmatrix}A\in\mathrm{SO}(1,1),
\]
also genügt es, \(\mathrm{SO}(1,1)\) zu bestimmen.
Ist \(B\in\mathrm{SO}(1,1)\) mit
\(B\,\mathcal H_{1,1}^{+}(0)=\mathcal H_{1,1}^{-}(0)\),
so folgt
\[
\begin{bmatrix}
-1 & 0\\
0 & -1
\end{bmatrix}B\in \operatorname{Lor}(1,1),
\]
also genügt es, \(\operatorname{Lor}(1,1)\) zu bestimmen.

Sei
\[
A=
\begin{bmatrix}
a & b\\
c & d
\end{bmatrix}
\in \operatorname{Lor}(1,1).
\]
Dann folgt aus \(A^{T}Q_{1,1}A=Q_{1,1}\) das Gleichungssystem
\[
\begin{cases}
a^{2}-c^{2}=1,\\
d^{2}-b^{2}=1,\\
cd-ab=0,\\
ad-bc=1.
\end{cases}
\]
Da außerdem
\(
A
\begin{bmatrix}
0\\1
\end{bmatrix}
=
\begin{bmatrix}
b\\d
\end{bmatrix}
\)
und \(d>0\) gelten muss, setzen wir \(d=\cosh s,\ b=\sinh s\)
für ein eindeutiges \(s\in\mathbb{R}\). Daraus folgen \(a=\cosh s,\ c=\sinh s\),
und somit
\[
A=
\begin{bmatrix}
\cosh s & \sinh s\\
\sinh s & \cosh s
\end{bmatrix},\qquad s\in\mathbb{R}.
\]
Dies beschreibt alle Elemente von \(\operatorname{Lor}(1,1)\).
\end{proof}
\end{EXA}




Remark 6.5


The Lie algebra of $\operatorname{Lor}(1,1)$ is easily determined to be
$$
\operatorname{lor}(1,1)=\left\{\left[\begin{array}{ll}
0 & t \\
t & 0
\end{array}\right]: t \in \mathbb{R}\right\}
$$
with trivial Lie bracket. The exponential map $\exp : \mathfrak{l o r}(1,1) \longrightarrow \operatorname{Lor}(1,1)$ is surjective since for $s \in \mathbb{R}$,
$$
\exp \left(\left[\begin{array}{ll}
0 & s \\
s & 0
\end{array}\right]\right)=\left[\begin{array}{ll}
\cosh s & \sinh s \\
\sinh s & \cosh s
\end{array}\right]
$$
As we will eventually see, the exponential map for $\operatorname{Lor}(n, 1)$ is actually surjective for every $n \geqslant 1$.



\begin{EXA}{LIE-B02-09-13}{Lie-Algebra der Lorentz-Gruppe in Dimension \(1+1\)}
Die Lie-Algebra der Lorentz-Gruppe \(\operatorname{Lor}(1,1)\) ist gegeben durch
\[
\operatorname{lor}(1,1)
=\left\{
\begin{bmatrix}
0 & t\\[2pt]
t & 0
\end{bmatrix}
: t\in\mathbb{R}
\right\},
\]
und der Lie-Klammeroperator ist trivial, das heißt
\([X,Y]=0\) für alle \(X,Y\in\operatorname{lor}(1,1)\).

\begin{proof}
Sei
\[
A(s)=
\begin{bmatrix}
\cosh s & \sinh s\\[2pt]
\sinh s & \cosh s
\end{bmatrix},
\qquad s\in\mathbb{R},
\]
eine glatte Kurve in \(\operatorname{Lor}(1,1)\) mit \(A(0)=I_2\).
Dann gilt
\[
\dot A(0)=
\begin{bmatrix}
0 & 1\\[2pt]
1 & 0
\end{bmatrix},
\]
was zeigt, dass \(\operatorname{lor}(1,1)\) eindimensional ist und durch diese Matrix erzeugt wird.

Die Bedingung
\[
X^{T}Q_{1,1}+Q_{1,1}X=0, 
\qquad Q_{1,1}=\operatorname{diag}(1,-1),
\]
liefert genau die angegebene Form der Elemente von \(\operatorname{lor}(1,1)\).

Da die Algebra eindimensional ist, ist die Lie-Klammer notwendigerweise trivial.

Für die Exponentialabbildung folgt unmittelbar
\[
\exp
\begin{bmatrix}
0 & s\\[2pt]
s & 0
\end{bmatrix}
=
\begin{bmatrix}
\cosh s & \sinh s\\[2pt]
\sinh s & \cosh s
\end{bmatrix},
\qquad s\in\mathbb{R},
\]
also ist
\[
\exp:\operatorname{lor}(1,1)\longrightarrow\operatorname{Lor}(1,1)
\]
surjektiv.
\end{proof}
\end{EXA}


\begin{KORO}{LG-06-06}{Surjektivität der Exponentialabbildung in niederen Dimensionen}
Die Exponentialabbildung
\[
\exp:\operatorname{lor}(n,1)\longrightarrow\operatorname{Lor}(n,1)
\]
ist für \(n=1\) surjektiv — und, wie später gezeigt wird, auch für alle \(n\geq 1\).
\end{KORO}



For future use we record the following explicit form for elements of $\operatorname{Lor}(n, 1)$, whose proof is a direct consequence of the definition of the Lorentz group.

Proposition 6.6


A matrix
$$
P=\left[\begin{array}{lll}
\mathbf{U}_{1} & \cdots & \mathbf{U}_{n+1}
\end{array}\right] \in \mathrm{GL}_{n+1}(\mathbb{R})
$$
is in $\operatorname{Lor}(n, 1)$ if and only if its column vectors $\mathbf{U}_{j}=\left[\begin{array}{l}\mathbf{u}_{j} \\ u_{j}\end{array}\right]$ satisfy the orthonormality equations
$$
\left\langle\mathbf{U}_{j}, \mathbf{U}_{k}\right\rangle=\left\{\begin{aligned}
0 & \text { if } j \neq k, \\
1 & \text { if } j=k \leqslant n, \\
-1 & \text { if } j=k=n+1 .
\end{aligned}\right.
$$
In particular $\mathbf{U}_{1}, \ldots, \mathbf{U}_{n}$ are spacelike and $\mathbf{U}_{n+1}$ is timelike with $u_{n+1}>0$.


\begin{PROP}{LIE-B02-09-14}{Charakterisierung der Lorentz-Gruppe durch Spaltenvektoren}
Eine Matrix
\[
P=
\begin{bmatrix}
\mathbf{U}_{1} & \cdots & \mathbf{U}_{n+1}
\end{bmatrix}
\in \mathrm{GL}_{n+1}(\mathbb{R})
\]
gehört genau dann zur Lorentz-Gruppe
\(\operatorname{Lor}(n,1)\),
wenn ihre Spaltenvektoren
\(\mathbf{U}_{j}=
\begin{bmatrix}
\mathbf{u}_{j}\\[2pt]
u_{j}
\end{bmatrix}\)
die Orthogonalitätsrelationen
\[
\left\langle \mathbf{U}_{j}, \mathbf{U}_{k} \right\rangle
=
\begin{cases}
0, & j\neq k,\\[4pt]
1, & j=k\leq n,\\[4pt]
-1, & j=k=n+1
\end{cases}
\]
erfüllen.
\end{PROP}


\begin{PROOF}{LIE-B02-09-15}{P: Charakterisierung der Lorentz-Gruppe durch Spaltenvektoren}
Nach Definition gilt \(P\in \operatorname{Lor}(n,1)\) genau dann, wenn
\[
P^{T} Q_{n,1} P = Q_{n,1}, 
\qquad Q_{n,1}=\operatorname{diag}(1,\ldots,1,-1).
\]
Die linke Seite hat die Einträge
\((P^{T} Q_{n,1} P)_{jk} = \langle \mathbf{U}_{j}, \mathbf{U}_{k} \rangle\),
wobei \(\langle\cdot,\cdot\rangle\) das Lorentz-Skalarprodukt ist.
Damit ist die Bedingung \(P^{T} Q_{n,1} P = Q_{n,1}\) äquivalent zu den angegebenen
Orthonormalitätsgleichungen.

Insbesondere sind die Vektoren \(\mathbf{U}_{1},\ldots,\mathbf{U}_{n}\)
raumartig (spacelike) und \(\mathbf{U}_{n+1}\)
zeitartig (timelike) mit \(u_{n+1}>0\).
\end{PROOF}




Theorem 6.7


\begin{THEO}{LIE-B02-09-16}{$\mathrm{Lor}(n,1)$ ist wegzusammenhängend}
For $n \geqslant 1, \operatorname{Lor}(n, 1)$ is path connected.
\end{THEO}



Proof


\begin{PROOF}{LIE-B02-09-17}{P: $\mathrm{Lor}(n,1)$ ist wegzusammenhängend}
We will show that $\operatorname{Lor}(n, 1)$ has a connected subgroup $H$ for which the homogeneous space $\operatorname{Lor}(n, 1) / H$ is connected. Together with Proposition 9.10, this implies that $\operatorname{Lor}(n, 1)$ is connected.

Consider the continuous action of $\operatorname{Lor}(n, 1)$ on $\mathbb{R}^{n, 1}$ by left multiplication. The stabiliser of the timelike vector $\mathbf{e}_{n+1}$ is
$$
\operatorname{Stab}_{\operatorname{Lor}(n, 1)}\left(\mathbf{e}_{n+1}\right)=\left\{\left[\begin{array}{cc}
A & O_{n-1,1} \\
O_{1, n-1} & 1
\end{array}\right]: A \in \operatorname{SO}(n)\right\} \leqslant \operatorname{Lor}(n, 1)
$$
which is a closed subgroup obviously isomorphic to $\operatorname{SO}(n)$, hence connected by Example 9.14. Next we need to identify the orbit of $\mathbf{e}_{n+1}$.

Suppose that $\mathbf{U}=\left[\begin{array}{l}\mathbf{u} \\ u\end{array}\right] \in \mathcal{H}_{n, 1}^{+}(1)$, so $|\mathbf{u}|^{2}-u^{2}=-1$ and $u>0$. We will show that
$$
\left[\begin{array}{l}
\mathbf{u} \\
u
\end{array}\right] \in \operatorname{Orb}_{\text {Lor }(n, 1)}\left(\mathbf{e}_{n+1}\right) .
$$
Consider the equation
$$
\mathbf{u} \cdot \mathbf{x}-u t=0
$$
for the pair $\mathbf{x} \in \mathbb{R}^{n}$ and $t \in \mathbb{R}$. The solution set is an $n$-dimensional vector subspace of $\mathbb{R}^{n, 1}$; furthermore, by Proposition 6.2 its non-zero vectors are all spacelike. Choosing a basis of this spacelike subspace, then applying the GramSchmidt process, we can produce a basis consisting of vectors
$$
\left[\begin{array}{c}
\mathbf{w}_{1} \\
s_{1}
\end{array}\right], \ldots,\left[\begin{array}{c}
\mathbf{w}_{n} \\
s_{n}
\end{array}\right]
$$
for which
$$
\begin{cases}\mathbf{w}_{i} \cdot \mathbf{w}_{j}-s_{i} s_{j}=\delta_{i j} & (i, j=1, \ldots, n) \\ \mathbf{u} \cdot \mathbf{w}_{k}-u s_{k}=0 & (k=1, \ldots, n)\end{cases}
$$
The resulting matrix
$$
P=\left[\begin{array}{cccc}
\mathbf{w}_{1} & \cdots & \mathbf{w}_{n} & \mathbf{u} \\
s_{1} & \cdots & s_{n} & u
\end{array}\right]
$$
is in $\operatorname{Lor}(n, 1)$ and satisfies
$$
P \mathbf{e}_{n+1}=\left[\begin{array}{l}
\mathbf{u} \\
u
\end{array}\right]
$$
So
$$
\operatorname{Orb}_{\operatorname{Lor}(n, 1)}\left(\mathbf{e}_{n+1}\right)=\mathcal{H}_{n, 1}^{+}(1)
$$
This also shows that $\operatorname{Orb}_{\operatorname{Lor}(n, 1)}\left(\mathbf{e}_{n+1}\right)$ is path connected since we already know that $\mathcal{H}_{n, 1}^{+}(1)$ is. Thus taking $H=\operatorname{Stab}_{\text {Lor }(n, 1)}\left(\mathbf{e}_{n+1}\right) \cong \operatorname{SO}(n)$ we see that $\operatorname{Lor}(n, 1)$ is path connected.

This identification of the homogeneous space $\operatorname{Lor}(n, 1) / \operatorname{SO}(n)$ can also be used to calculate the dimension of $\operatorname{Lor}(n, 1)$.
\end{PROOF}




Proposition 6.8


\begin{PROP}{LIE-B02-09-18}{$\operatorname{dim} \operatorname{Lor}(n, 1)=\binom{n+1}{2}$}
The dimension of the Lorentz group $\operatorname{Lor}(n, 1)$ is
$$
\operatorname{dim} \operatorname{Lor}(n, 1)=\binom{n+1}{2}
$$
\end{PROP}


Proof


\begin{PROOF}{LIE-B02-09-19}{P: $\operatorname{dim} \operatorname{Lor}(n, 1)=\binom{n+1}{2}$}
Using the Implicit Function Theorem 7.11, it is easy to see that
$$
\operatorname{dim} \operatorname{Lor}(n, 1) / \mathrm{SO}(n)=\operatorname{dim} \mathcal{H}_{n, 1}^{+}(1)=n .
$$
Also from Section 3.3 we know that
$$
\operatorname{dim} S O(n)=\binom{n}{2}
$$
Combining these we obtain the result.
\end{PROOF}


\pagebreak


\subsection{A Principal Axis Theorem for Lorentz Groups}


We will make use of the Principal Axis Theorem for special orthogonal groups which will be discussed later as part of Theorem 10.13. This says that every matrix $A \in \mathrm{SO}(m)$ has the form
$$
A= \begin{cases}P R_{2 n}\left(\theta_{1}, \ldots, \theta_{n}\right) P^{T} & \text { if } m=2 n \text { is even }, \\ P R_{2 n+1}\left(\theta_{1}, \ldots, \theta_{n}\right) P^{T} & \text { if } m=2 n+1 \text { is odd },\end{cases}
$$
where for $\theta_{1}, \ldots, \theta_{n}, \theta \in \mathbb{R}$,
$$
\begin{aligned}
R_{2 n}\left(\theta_{1}, \ldots, \theta_{n}\right) & =\left[\begin{array}{cccccc}
R\left(\theta_{1}\right) & O & \cdots & \cdots & \cdots & O \\
O & R\left(\theta_{2}\right) & O & \ddots & \ddots & \vdots \\
\vdots & \ddots & \ddots & \ddots & \ddots & \vdots \\
O & \cdots & \cdots & \cdots & O & R\left(\theta_{n}\right)
\end{array}\right], \\
R_{2 n+1}\left(\theta_{1}, \ldots, \theta_{n}\right) & =\left[\begin{array}{ccccccc}
R\left(\theta_{1}\right) & O & \cdots & \cdots & \cdots & \cdots & O \\
O & R\left(\theta_{2}\right) & O & \ddots & \ddots & \ddots & \ddots \\
\vdots & \ddots & \ddots & \ddots & \ddots & \ddots & \vdots \\
\vdots & \ddots & \ddots & \ddots & O & R\left(\theta_{n}\right) & O \\
O & \cdots & \cdots & \cdots & \cdots & O & 1
\end{array}\right], \\
R(\theta) & =\left[\begin{array}{cc}
\cos \theta & -\sin \theta \\
\sin \theta & \cos \theta
\end{array}\right] .
\end{aligned}
$$



\begin{THEO}{LIE-B02-09-20}{Hauptachsensatz für spezielle orthogonale Gruppen}
Jede Matrix \(A \in \mathrm{SO}(m)\) lässt sich in der Form
\[
A =
\begin{cases}
P\, R_{2n}(\theta_{1},\ldots,\theta_{n})\, P^{T}, & \text{falls } m = 2n \text{ gerade},\\[6pt]
P\, R_{2n+1}(\theta_{1},\ldots,\theta_{n})\, P^{T}, & \text{falls } m = 2n+1 \text{ ungerade},
\end{cases}
\]
darstellen, wobei \(P \in \mathrm{SO}(m)\) orthogonal ist und die Winkel \(\theta_{1},\ldots,\theta_{n} \in \mathbb{R}\) reell sind.

Die Blockmatrizen \(R_{2n}\) und \(R_{2n+1}\) sind definiert durch
\[
\begin{aligned}
R_{2n}(\theta_{1},\ldots,\theta_{n})
&=
\begin{bmatrix}
R(\theta_{1}) & O & \cdots & O\\[2pt]
O & R(\theta_{2}) & \ddots & \vdots\\[2pt]
\vdots & \ddots & \ddots & O\\[2pt]
O & \cdots & O & R(\theta_{n})
\end{bmatrix},
\\[6pt]
R_{2n+1}(\theta_{1},\ldots,\theta_{n})
&=
\begin{bmatrix}
R(\theta_{1}) & O & \cdots & O & O\\[2pt]
O & R(\theta_{2}) & \ddots & \vdots & \vdots\\[2pt]
\vdots & \ddots & \ddots & O & \vdots\\[2pt]
O & \cdots & O & R(\theta_{n}) & O\\[2pt]
O & \cdots & \cdots & O & 1
\end{bmatrix},
\end{aligned}
\]
wobei
\[
R(\theta) =
\begin{bmatrix}
\cos\theta & -\sin\theta\\[2pt]
\sin\theta & \cos\theta
\end{bmatrix}
\]
eine planare Drehmatrix bezeichnet.
\end{THEO}





The following result leads to the analogue for Lorentz groups.


Theorem 6.9

For $n \geqslant 1$, every Lorentz matrix $A \in \operatorname{Lor}(n, 1)$ has one of the forms

$$
\begin{aligned}
& A=P\left[\begin{array}{cc}
B & O_{n-2,2} \\
O_{2, n-2} & \cosh t \\
\sinh t & \sinh t \\
\cosh t
\end{array}\right] P^{-1}, \\
& A=P\left[\begin{array}{cc}
C & O_{n, 1} \\
O_{1, n} & 1
\end{array}\right] P^{-1},
\end{aligned}
$$

where $B \in \mathrm{SO}(n-1), C \in \mathrm{SO}(n), t \in \mathbb{R}$ and $P \in \operatorname{Lor}(n, 1)$.


\begin{THEO}{LIE-B02-09-21}{Kanonische Normalformen von Lorentz-Matrizen}
Für \(n \geq 1\) besitzt jede Lorentz-Matrix
\(A \in \operatorname{Lor}(n,1)\)
eine der folgenden äquivalenten Darstellungen:
\[
\begin{aligned}
A
&=
P
\begin{bmatrix}
B & O_{n-2,2}\\[2pt]
O_{2,n-2} &
\begin{array}{cc}
\cosh t & \sinh t\\[2pt]
\sinh t & \cosh t
\end{array}
\end{bmatrix}
P^{-1},
\\[8pt]
\text{oder} \qquad
A
&=
P
\begin{bmatrix}
C & O_{n,1}\\[2pt]
O_{1,n} & 1
\end{bmatrix}
P^{-1},
\end{aligned}
\]
wobei
\[
B \in \mathrm{SO}(n-1), \qquad
C \in \mathrm{SO}(n), \qquad
t \in \mathbb{R}, \qquad
P \in \operatorname{Lor}(n,1).
\]

Die erste Form entspricht einer \emph{hyperbolischen Boost-Transformation} in einer
zweidimensionalen Unterebene des Minkowski-Raums, kombiniert mit einer Drehung \(B\)
in den verbleibenden \((n-1)\)-raumartigen Richtungen.

Die zweite Form beschreibt eine \emph{reine Drehung} in den raumartigen Dimensionen,
d.\,h. eine Einbettung von \(\mathrm{SO}(n)\) in \(\operatorname{Lor}(n,1)\).
\end{THEO}



Proof


\begin{PROOF}{LIE-B02-09-22}{P: Kanonische Normalformen von Lorentz-Matrizen}
We will prove this by induction on $n \geqslant 1$. The initial case $n=1$ is dealt with in Example 6.4. So we suppose that the result holds for $\operatorname{Lor}(k, 1)$ whenever $1 \leqslant k \leqslant n-1$.

We begin by analysing the eigenvalues and eigenvectors of $A$. Suppose that $\lambda \in \mathbb{C}$ is an eigenvalue for $A \in \operatorname{Lor}(n, 1)$ with eigenvector $\mathbf{U} \in \mathbb{C}^{n, 1}$. If $\mathbf{U}$ is not a null vector, we have
$$
\begin{aligned}
|\lambda|^{2}\langle\mathbf{U}, \mathbf{U}\rangle & =\langle\lambda \mathbf{U}, \lambda \mathbf{U}\rangle \\
& =\langle A \mathbf{U}, A \mathbf{U}\rangle \\
& =\langle\mathbf{U}, \mathbf{U}\rangle,
\end{aligned}
$$
hence $|\lambda|=1$. On the other hand, if $\mathbf{U}$ is a null vector then $|\lambda|$ is potentially unrestricted.

Next suppose that $\lambda$ and $\mu$ are distinct eigenvalues for $A$ with eigenvectors $\mathbf{U}$ and $\mathbf{V}$. Then
$$
\begin{aligned}
\bar{\lambda} \mu\langle\mathbf{U}, \mathbf{V}\rangle & =\langle\lambda \mathbf{U}, \mu \mathbf{V}\rangle \\
& =\langle A \mathbf{U}, A \mathbf{V}\rangle \\
& =\langle\mathbf{U}, \mathbf{V}\rangle
\end{aligned}
$$
hence either $\langle\mathbf{U}, \mathbf{V}\rangle=0$ or $\mu=\bar{\lambda}^{-1}$. In particular, if $|\lambda|=1$ we must also have $\bar{\lambda}^{-1}=\lambda \neq \mu$, yielding $\langle\mathbf{U}, \mathbf{V}\rangle=0$.

Next we will investigate in more detail the situation when $A$ has a null eigenvector $\mathbf{U}=\left[\begin{array}{l}\mathbf{u} \\ u\end{array}\right]$ for an eigenvalue $\lambda$. Since $|\mathbf{u}|^{2}=|u|^{2} \neq 0$, we can multiply by a suitable complex scalar to ensure that $u=1$, so we might as well assume that $\mathbf{U}=\left[\begin{array}{l}\mathbf{u} \\ 1\end{array}\right]$. Notice that
$$
A\left[\begin{array}{c}
\overline{\mathbf{u}} \\
1
\end{array}\right]=\overline{\left(A\left[\begin{array}{l}
\mathbf{u} \\
1
\end{array}\right]\right)}=\bar{\lambda}\left[\begin{array}{l}
\overline{\mathbf{u}} \\
1
\end{array}\right],
$$
showing that $\bar{\lambda}$ is also an eigenvalue with eigenvector $\overline{\mathbf{U}}=\left[\begin{array}{l}\overline{\mathbf{u}} \\ 1\end{array}\right]$. Also,
$$
\begin{aligned}
\langle\overline{\mathbf{U}}, \mathbf{U}\rangle & =\langle A \overline{\mathbf{U}}, A \mathbf{U}\rangle \\
& =\langle\overline{\lambda \mathbf{U}}, \lambda \mathbf{U}\rangle \\
& =\lambda^{2}\langle\overline{\mathbf{U}}, \mathbf{U}\rangle .
\end{aligned}
$$
As $\overline{\mathbf{U}}$ is a null vector, Proposition 6.2 implies that either $\overline{\mathbf{U}}=\mathbf{U}$ or $\lambda^{2}=1$. Each of these possibilities means that $\lambda$ is real and in fact positive, since $A$ preserves the positive light cone. So for a null eigenvector, the corresponding eigenvalue has to be a positive real number.

Continuing with this situation, consider the solution set of the linear equation
$$
\langle\mathbf{U}, \mathbf{X}\rangle=0 .
$$
This forms a vector subspace of $\mathbb{R}^{n, 1}$ of dimension $n$, containing the subspace of real multiples of $\mathbf{U}$ as well as the ( $n-1$ )-dimensional subspace of spacelike elements together with zero. In fact, the spacelike elements are all the vectors of the form
$$
\left[\begin{array}{l}
\mathbf{x} \\
0
\end{array}\right] \quad\left(\mathbf{x} \in \mathbb{R}^{n}, \mathbf{u} \cdot \mathbf{x}=0\right) .
$$
Thus the general solution of Equation (6.6) is
$$
\mathbf{X}=\left[\begin{array}{c}
\mathbf{x}+z \mathbf{u} \\
z
\end{array}\right] \quad\left(\mathbf{x} \in \mathbb{R}^{n}, \mathbf{u} \cdot \mathbf{x}=0, z \in \mathbb{R}\right) .
$$
Then there is a unique null vector $\widetilde{\mathbf{U}} \in \mathbb{R}^{n, 1}$ possessing the two properties\\
$\bullet\langle\tilde{\mathbf{U}}, \mathbf{X}\rangle=0$ for all spacelike solutions of Equation (6.6),\\
$\bullet\langle\mathbf{U}, \tilde{\mathbf{U}}\rangle=-2$.


This vector is easily seen to be $\tilde{\mathbf{U}}=\left[\begin{array}{c}-\mathbf{u} \\ 1\end{array}\right]$. Clearly the spacelike solutions of Equation (6.6) are preserved by $A$, and indeed $A$ acts as an isometry on this space so relative to a basis it is given by an orthogonal matrix. Also
$$
\begin{aligned}
\langle\mathbf{U}, A(\lambda \tilde{\mathbf{U}})\rangle & =\lambda\left\langle A^{-1} \mathbf{U}, \tilde{\mathbf{U}}\right\rangle \\
& =\lambda\left\langle\lambda^{-1} \mathbf{U}, \tilde{\mathbf{U}}\right\rangle \\
& =\lambda \lambda^{-1}\langle\mathbf{U}, \tilde{\mathbf{U}}\rangle \\
& =-2 .
\end{aligned}
$$
Using the above characterisation of $\tilde{\mathbf{U}}$, we see that
$$
A \tilde{\mathbf{U}}=\lambda^{-1} \tilde{\mathbf{U}},
$$
i.e., $\tilde{\mathbf{U}}$ is an eigenvector of $A$ for the eigenvalue $\lambda^{-1}$.


The vectors
$$
\mathbf{U}^{\prime}=\frac{1}{2}(\mathbf{U}-\tilde{\mathbf{U}}), \quad \mathbf{U}^{\prime \prime}=\frac{1}{2}(\mathbf{U}+\tilde{\mathbf{U}}),
$$
satisfy
$$
\begin{aligned}
& A \mathbf{U}^{\prime}=\frac{\left(\lambda+\lambda^{-1}\right)}{2} \mathbf{U}^{\prime}+\frac{\left(\lambda-\lambda^{-1}\right)}{2} \mathbf{U}^{\prime \prime}, \\
& A \mathbf{U}^{\prime \prime}=\frac{\left(\lambda-\lambda^{-1}\right)}{2} \mathbf{U}^{\prime}+\frac{\left(\lambda+\lambda^{-1}\right)}{2} \mathbf{U}^{\prime \prime},
\end{aligned}
$$
and setting $s=\ln \lambda$, we have
$$
\begin{aligned}
& A \mathbf{U}^{\prime}=\cosh s \mathbf{U}^{\prime}+\sinh s \mathbf{U}^{\prime \prime} \\
& A \mathbf{U}^{\prime \prime}=\sinh s \mathbf{U}^{\prime}+\cosh s \mathbf{U}^{\prime \prime}
\end{aligned}
$$
Also
$$
\left\langle\mathbf{U}^{\prime}, \mathbf{U}^{\prime \prime}\right\rangle=0, \quad\left\langle\mathbf{U}^{\prime}, \mathbf{U}^{\prime}\right\rangle=1, \quad\left\langle\mathbf{U}^{\prime \prime}, \mathbf{U}^{\prime \prime}\right\rangle=-1,
$$
so $\mathbf{U}^{\prime}$ is a unit spacelike vector, while $\mathbf{U}^{\prime \prime}$ is timelike and these are orthogonal vectors with respect to the Lorentzian inner product.

Together with an orthonormal basis $\left\{\mathbf{U}_{1}, \ldots, \mathbf{U}_{n-1}\right\}$ of the spacelike solutions of Equation (6.6), we see that $\mathbb{R}^{n, 1}$ has a basis
$$
\left\{\mathbf{U}_{1}, \ldots, \mathbf{U}_{n-1}, \mathbf{U}^{\prime}, \mathbf{U}^{\prime \prime}\right\}
$$
which provides the columns of a Lorentz matrix
$$
P=\left[\begin{array}{lllll}
\mathbf{U}_{1} & \cdots & \mathbf{U}_{n-1} & \mathbf{U}^{\prime} & \mathbf{U}^{\prime \prime}
\end{array}\right]
$$
for which
$$
A=P\left[\begin{array}{ccc}
B & O_{n-2,2} \\
& \cosh t & \sinh t \\
O_{2, n-2} & \sinh t & \cosh t
\end{array}\right] P^{-1}
$$
with $B \in \mathrm{O}(n-1)$; in fact $\operatorname{det} B=1$ since $\operatorname{det} A=1$ and also
$$
\operatorname{det}\left[\begin{array}{cc}
\cosh t & \sinh t \\
\sinh t & \cosh t
\end{array}\right]=\cosh ^{2} t-\sinh ^{2} t=1
$$
Hence $A$ has the form (I) whenever it has a null vector for an eigenvector.


Now we consider the situation where $A$ has a timelike eigenvector U for the eigenvalue $\lambda$. Using similar arguments to those for a null eigenvector, we find that $\lambda=1$ and $\mathbf{U} \in \mathbb{R}^{n, 1}$. Moreover, $A$ acts as an isometry on the solution set of
$$
\langle\mathbf{U}, \mathbf{X}\rangle=\mathbf{0},
$$
which is an $n$-dimensional subspace of $\mathbb{R}^{n, 1}$ with timelike non-zero vectors. Hence $A$ has the form (II).

The remaining case to deal with is that where there are no null or timelike eigenvectors. So suppose that $\mathbf{U}=\left[\begin{array}{l}\mathbf{u} \\ u\end{array}\right]$ is a spacelike eigenvector of unit\\
length with eigenvalue $\lambda$ of unit modulus. By multiplying by a suitable complex number we can assume that $u \in \mathbb{R}$ and $u \geqslant 0$.


When $\mathbf{U} \in \mathbb{R}^{n, 1}$, we have $\lambda= \pm 1$ and $\operatorname{det} A=1$.\\
If $\lambda=1$ then the solutions of the equation
$$
\langle\mathbf{U}, \mathbf{X}\rangle=0
$$
form a subspace of $\mathbb{R}^{n, 1}$ closed under the action of $A$ and indeed, on choosing a basis for this subspace we find that it is isomorphic to $\mathbb{R}^{n-1,1}$ and the action of $A$ is given by an element of $\operatorname{Lor}(n-1,1)$. Hence there is a $Q \in \operatorname{Lor}(n, 1)$ and $B \in \operatorname{Lor}(n-1,1)$ for which
$$
A=Q\left[\begin{array}{cc}
1 & O_{1, n} \\
O_{n, 1} & B
\end{array}\right] Q^{-1}
$$
By the induction hypothesis, $B$ has one of the forms (I) or (II), from which we deduce that $A$ does too.

If $\lambda=-1$, this is an eigenvalue of multiplicity at least 2 since non-real complex eigenvalues come in complex conjugate pairs. With the aid of the GramSchmidt process we may produce a second eigenvector $\mathbf{V}$ for this eigenvalue of unit length and orthogonal to $\mathbf{U}$. The solutions of the pair of simultaneous equations
$$
\langle\mathbf{U}, \mathbf{X}\rangle=0=\langle\mathbf{V}, \mathbf{X}\rangle
$$
form a subspace of $\mathbb{R}^{n, 1}$ closed under the action of $A$. On choosing a basis for this subspace we find that it is isomorphic to $\mathbb{R}^{n-2,1}$ and the action of $A$ is given by an element of $\operatorname{Lor}(n-2,1)$. Hence there is a $Q \in \operatorname{Lor}(n, 1)$ and $B^{\prime} \in \operatorname{Lor}(n-2,1)$ for which
$$
A=Q\left[\begin{array}{cc}
-I_{2} & O_{2, n-1} \\
O_{n-1,2} & B^{\prime}
\end{array}\right] Q^{-1}
$$
Again the induction hypothesis implies that $B^{\prime}$ has one of the forms (I) or (II), and therefore $A$ does too.

When $\mathbf{U} \notin \mathbb{R}^{n, 1}, \overline{\mathbf{U}}=\left[\begin{array}{l}\overline{\mathbf{u}} \\ u\end{array}\right]$ is also a spacelike eigenvector of unit length for the eigenvalue $\bar{\lambda}$. The vectors
$$
\mathbf{U}^{\prime}=\frac{1}{2}(\mathbf{U}+\overline{\mathbf{U}}), \quad \mathbf{U}^{\prime \prime}=\frac{1}{2 i}(\mathbf{U}-\overline{\mathbf{U}}),
$$
are orthonormal and satisfy the equations
$$
A \mathbf{U}^{\prime}=\cos \theta \mathbf{U}^{\prime}+\sin \theta \mathbf{U}^{\prime \prime}, \quad A \mathbf{U}^{\prime \prime}=-\sin \theta \mathbf{U}^{\prime}+\cos \theta \mathbf{U}^{\prime \prime}
$$
where $\lambda=e^{\theta i}$. As in the previous case, the set of solutions of the pair of simultaneous equations
$$
\left\langle\mathbf{U}^{\prime}, \mathbf{X}\right\rangle=0=\left\langle\mathbf{U}^{\prime \prime}, \mathbf{X}\right\rangle
$$
forms a Lorentzian subspace which can be identified with $\operatorname{Lor}(n-2,1)$ after choosing a basis, and $A$ acts on it as an element of $\operatorname{Lor}(n-2,1)$. Again there is a $Q \in \operatorname{Lor}(n, 1)$ and $B^{\prime \prime} \in \operatorname{Lor}(n-2,1)$ for which
$$
A=Q\left[\begin{array}{cc}
R(\theta) & O_{2, n-1} \\
O_{n-1,2} & B^{\prime \prime}
\end{array}\right] Q^{-1} .
$$
By the induction hypothesis, $B^{\prime \prime}$ has one of the forms (I) or (II), from which we deduce that $A$ does too.
\end{PROOF}




Combining this result with the Principal Axis Theorem 10.13 for $\operatorname{SO}(k)$, we obtain the following.

Theorem 6.10 (Principal Axis Theorem for Lorentz groups)
For $n \geqslant 1$, each matrix $A \in \operatorname{Lor}(n, 1)$ has one of the forms

$$
\begin{aligned}
& A=Q\left[\begin{array}{cc}
R_{n-1}\left(\theta_{1}, \ldots, \theta_{\ell}\right) & O_{n-2,2} \\
O_{2, n-2} & \cosh t \\
\sinh t & \sinh t \\
\cosh t
\end{array}\right] Q^{-1} \\
& A=Q\left[\begin{array}{cc}
R_{n}\left(\theta_{1}, \ldots, \theta_{m}\right) & O_{n, 1} \\
O_{1, n} & 1
\end{array}\right] Q^{-1}
\end{aligned}
$$

where $t \in \mathbb{R}, Q \in \operatorname{Lor}(n, 1)$ and in case ( $\mathrm{I}^{\prime}$ ), $n=2 \ell+1$ or $n=2 \ell+2$, while in case (III'), $n=2 m$ or $n=2 m+1$.



\begin{THEO}{LIE-B02-09-23}{Hauptachsensatz für Lorentz-Gruppen}
Für \(n \ge 1\) besitzt jede Lorentz-Matrix
\(A \in \operatorname{Lor}(n,1)\)
eine der beiden kanonischen Darstellungsformen:
\[
\begin{aligned}
\text{(I$'$)}\quad
A
&=
Q
\begin{bmatrix}
R_{n-1}(\theta_{1},\ldots,\theta_{\ell}) & O_{n-2,2}\\[4pt]
O_{2,n-2} &
\begin{array}{cc}
\cosh t & \sinh t\\[2pt]
\sinh t & \cosh t
\end{array}
\end{bmatrix}
Q^{-1},\\[10pt]
\text{(III$'$)}\quad
A
&=
Q
\begin{bmatrix}
R_{n}(\theta_{1},\ldots,\theta_{m}) & O_{n,1}\\[4pt]
O_{1,n} & 1
\end{bmatrix}
Q^{-1},
\end{aligned}
\]
wobei
\[
t \in \mathbb{R}, \qquad
Q \in \operatorname{Lor}(n,1),
\]
und gilt:
\[
\begin{cases}
n = 2\ell + 1 \text{ oder } n = 2\ell + 2, & \text{im Fall (I$'$)},\\[4pt]
n = 2m \text{ oder } n = 2m + 1, & \text{im Fall (III$'$)}.
\end{cases}
\]

Hierbei bezeichnet
\[
R_{n}(\theta_{1},\ldots,\theta_{m})
=
\operatorname{diag}\big(R(\theta_{1}),\ldots,R(\theta_{m}),I_{r}\big),
\quad
R(\theta)=
\begin{bmatrix}
\cos\theta & -\sin\theta\\[2pt]
\sin\theta & \cos\theta
\end{bmatrix},
\]
die Blockdiagonalmatrix aus planaren Rotationen.


Fall (I$'$) beschreibt eine \emph{Boost-Transformation} in einer zweidimensionalen Lorentz-Unterebene
kombiniert mit einer Rotation \(R_{n-1}\) im raumartigen Unterraum.  
Fall (III$'$) beschreibt eine \emph{reine Rotation}, also eine Einbettung der
speziellen orthogonalen Gruppe \(\mathrm{SO}(n)\) in \(\operatorname{Lor}(n,1)\).

Der Satz ist das exakte Analogon des Hauptachsensatzes für
\(\mathrm{SO}(n)\) (vgl.~Proposition~\ref{LG-06-08}), erweitert auf den
pseudo-euklidischen Minkowski-Raum.
\end{THEO}






\begin{EXA}{LIE-B02-09-24}{Beispiele zu den Normalformen der Lorentz-Gruppe}
\noindent\textbf{(a) Boost im \((1+1)\)-dimensionalen Minkowski-Raum.}\\[2pt]
Für \(n=1\) ist \(\operatorname{Lor}(1,1)\) eindimensional.
Jede Matrix der Form
\[
A(t)=
\begin{bmatrix}
\cosh t & \sinh t\\[2pt]
\sinh t & \cosh t
\end{bmatrix},
\qquad t\in\mathbb{R},
\]
gehört zu \(\operatorname{Lor}(1,1)\) und beschreibt eine \emph{hyperbolische Rotation}
(„Boost“) zwischen Raum- und Zeitkoordinate.
Der Parameter \(t\) entspricht der \emph{Rapidität}, wobei
\(\tanh t = v/c\) der dimensionslose Geschwindigkeitsparameter ist.

\vspace{0.5em}
\noindent\textbf{(b) Rotation in \(\operatorname{Lor}(2,1)\).}\\[2pt]
Für \(n=2\) kann jede Lorentz-Transformation durch eine Drehung in der raumartigen
Ebene und/oder einen Boost in der Zeitebene dargestellt werden.
Eine reine Raumrotation ergibt sich durch
\[
A(\theta)=
\begin{bmatrix}
\cos\theta & -\sin\theta & 0\\[2pt]
\sin\theta & \cos\theta & 0\\[2pt]
0 & 0 & 1
\end{bmatrix},
\qquad \theta\in\mathbb{R},
\]
was dem Fall (III$'$) in Theorem~\ref{LG-06-10} entspricht.

Eine reine Zeit-Raum-Transformation (Boost) entspricht dagegen dem Fall (I$'$):
\[
B(t)=
\begin{bmatrix}
1 & 0 & 0\\[2pt]
0 & \cosh t & \sinh t\\[2pt]
0 & \sinh t & \cosh t
\end{bmatrix}.
\]
Beide Transformationen erfüllen
\(B^{T}Q_{2,1}B = A^{T}Q_{2,1}A = Q_{2,1}\)
mit \(Q_{2,1}=\mathrm{diag}(1,1,-1)\).

\vspace{0.5em}
\noindent\textbf{(c) Allgemeiner Fall in \(\operatorname{Lor}(3,1)\).}\\[2pt]
Für \(n=3\) kombiniert eine allgemeine Lorentz-Transformation
eine Drehung in der raumartigen \(3\)-Ebene und einen Boost
in einer \(2\)-dimensionalen Lorentz-Unterebene:
\[
A = P
\begin{bmatrix}
R(\theta) & 0 & 0\\[2pt]
0 & \cosh t & \sinh t\\[2pt]
0 & \sinh t & \cosh t
\end{bmatrix}
P^{-1},
\qquad P\in\operatorname{Lor}(3,1),
\]
wobei \(R(\theta)\in\mathrm{SO}(2)\) eine planare Rotation und
\(\begin{bmatrix}\cosh t & \sinh t\\[2pt]\sinh t & \cosh t\end{bmatrix}\)
ein Boost in einer Zeit-Raum-Ebene darstellt.
Dies illustriert geometrisch, dass jede Lorentz-Transformation als Kombination
von Raumrotation und Boost interpretiert werden kann.
\end{EXA}





\begin{DEF}{LIE-B02-09-25}{Rotation und Boost}
Sei $\mathbb{R}^{n,1}$ der Minkowski-Raum mit Skalarprodukt 
$\langle X,Y\rangle = x_1y_1+\cdots+x_ny_n - t_X t_Y$ und 
$\operatorname{Lor}(n,1)=\{A\in \mathrm{SO}(n,1): A\,\mathcal H^{\pm}_{n,1}(1)=\mathcal H^{\pm}_{n,1}(1)\}$.
\begin{enumerate}
\item Eine \emph{Rotation} ist ein $A\in \operatorname{Lor}(n,1)$, das nur raumartige Koordinaten mischt und $t$ fixiert. In geeigneter Basis hat $A$ Blockform
\[
A=\begin{bmatrix}
R & 0\\[2pt]
0 & 1
\end{bmatrix}, 
\qquad R\in \mathrm{SO}(n),
\]
und in einer zweidimensionalen raumartigen Ebene $(x_i,x_j)$ konkret
\(
R(\theta)=\begin{bmatrix}\cos\theta & -\sin\theta\\ \sin\theta & \cos\theta\end{bmatrix}.
\)
\item Ein \emph{Boost} ist ein $A\in \operatorname{Lor}(n,1)$, das genau eine raumartige Richtung mit der Zeitrichtung mischt. In geeigneter Basis (Boost entlang $x_1$)
\[
B(\eta)=
\begin{bmatrix}
\cosh\eta & 0 & \cdots & 0 & \sinh\eta\\
0 & 1 &  &  & 0\\
\vdots &  & \ddots &  & \vdots\\
0 &  &  & 1 & 0\\
\sinh\eta & 0 & \cdots & 0 & \cosh\eta
\end{bmatrix},
\]
wobei die \emph{Rapidität} $\eta\in\mathbb{R}$ durch $\tanh\eta = v/c$ (Relativgeschwindigkeit $v$) gegeben ist.
\end{enumerate}
\end{DEF}



\begin{CONC}{LIE-B02-09-26}{Geometrischer Charakter von Rotation und Boost}
Rotationen sind \emph{euklidische Drehungen} in raumartigen Ebenen (Kreise als Bahnen); Boosts sind \emph{hyperbolische Drehungen} in zeit-raumartigen Ebenen (Hyperbeln als Bahnen). Beides erhält die Minkowski-Norm, d.\,h.\ $A^T Q_{n,1} A=Q_{n,1}$.
\end{CONC}



\begin{PROP}{LIE-B02-09-27}{Invarianz und Wirkung auf Vierergrößen}
Sei $X=(\mathbf x,t)\in\mathbb{R}^{n,1}$ ein Vierervektor.
\begin{enumerate}
\item Für jede Rotation $R\in\mathrm{SO}(n)$ gilt 
$\langle (R\mathbf x,t),(R\mathbf x,t)\rangle = \langle (\mathbf x,t),(\mathbf x,t)\rangle$ und $t$ bleibt unverändert.
\item Für jeden Boost $B(\eta)$ entlang $x_1$ gilt
\[
\begin{bmatrix}
x_1'\\ t'
\end{bmatrix}
=
\begin{bmatrix}
\cosh\eta & \sinh\eta\\
\sinh\eta & \cosh\eta
\end{bmatrix}
\begin{bmatrix}
x_1\\ t
\end{bmatrix},
\qquad 
x_j'=x_j\ (j\ge 2),
\]
und $\langle X',X'\rangle=\langle X,X\rangle$.
\end{enumerate}
\end{PROP}



\begin{PROP}{LIE-B02-09-28}{Zeitdilatation und Längenkontraktion (1+1)}
Für ein zeitartiges Intervall mit Eigenzeit $\Delta\tau$ und Relativgeschwindigkeit $v$ (mit $\gamma=\cosh\eta=1/\sqrt{1-v^2/c^2}$) folgen unmittelbar:
\[
\Delta t = \gamma\,\Delta\tau,
\qquad
\Delta x_{\text{proper}} = \frac{1}{\gamma}\,\Delta x_{\text{gemessen}}.
\]
\end{PROP}



\begin{CONC}{LIE-B02-09-29}{„System vs.\ Beschreibung“}
Lorentz-Transformationen (Rotationen und Boosts) ändern nicht das physikalische \emph{System}, sondern die \emph{Beschreibung} desselben in verschiedenen Inertialsystemen:
\begin{enumerate}
\item Ein Ereignis $p\in\mathbb{R}^{n,1}$ ist beobachterunabhängig; seine Koordinaten $(\mathbf x,t)$ hängen vom Bezugssystem ab und transformieren mit $A\in\operatorname{Lor}(n,1)$.
\item Physikalische Gesetze sind kovariant: Tensor-/Spinorgrößen transformieren repräsentationsgetreu (z.\,B.\ Vektoren mit der Standarddarstellung, elektromagnetischer Tensor mit der adjungierten Darstellung), ihre \emph{Form} bleibt invariant.
\end{enumerate}
\end{CONC}



\begin{CONC}{LIE-B02-09-30}{Kanonnische Blöcke aus dem Hauptachsensatz}
Jedes $A\in\operatorname{Lor}(n,1)$ zerfällt (nach geeigneter Wahl einer Lorentz-Basis) in einen direkten Block aus:
\begin{enumerate}
\item Rotationsblöcken $R(\theta_k)$ in paarweisen orthogonalen raumartigen Ebenen;
\item einem einzelnen Boost-Block $B(\eta)$ in einer $(x_i,t)$-Ebene;
\item bzw.\ einem reinen Raum-Block $R\in \mathrm{SO}(n)$ (falls kein Boost auftritt).
\end{enumerate}
\emph{Interpretation:} Jede Lorentz-Transformation ist (bis zur Konjugation in $\operatorname{Lor}(n,1)$) eine Kombination aus Rotationen und höchstens einem Boost entlang einer Raumrichtung.
\end{CONC}



\begin{CONC}{LIE-B02-09-31}{Rapidität als additiver Parameter}
Für kollineare Boosts entlang derselben Raumachse addieren sich Rapiditäten:
\(
B(\eta_1)B(\eta_2)=B(\eta_1+\eta_2).
\)
Dies ist der Grund, warum $\eta$ rechnerisch „bequemer“ als $v$ ist.
\end{CONC}




We can also deduce information on the exponential for Lorentz groups.

Theorem 6.11


For $n \geqslant 1$, the exponential map $\exp : \operatorname{lor}(n, 1) \longrightarrow \operatorname{Lor}(n, 1)$ is surjective.


\begin{THEO}{LIE-B02-09-32}{Surjektivität der Exponentialabbildung der Lorentz-Gruppe}
Für alle $n \geqslant 1$ ist die Exponentialabbildung
\[
\exp : \operatorname{lor}(n,1) \longrightarrow \operatorname{Lor}(n,1)
\]
\emph{surjektiv}.
\end{THEO}



Proof


\begin{PROOF}{LIE-B02-09-33}{P: Surjektivität der Exponentialabbildung der Lorentz-Gruppe}
When $n=1$, we know from Remark 6.5 that every element of $\operatorname{Lor}(1,1)$ has the form given in (I). Also, by Theorem 10.16, $\exp : s o(k) \longrightarrow S O(k)$ is surjective for $k \geqslant 1$. Let $A \in \operatorname{Lor}(n, 1)$ for $n \geqslant 2$. When $A$ has the form (I),
$$
A=P\left[\begin{array}{cc}
\exp (S) & O_{n-2,2} \\
O_{2, n-2} & \exp \left(\left[\begin{array}{cc}
0 & t \\
t & 0
\end{array}\right]\right)
\end{array}\right] P^{-1}=\exp \left(P\left[\begin{array}{ccc}
S & O_{n-2,2} \\
O_{2, n-2} & 0 & t \\
& t & 0
\end{array}\right] P^{-1}\right)
$$
for some $S \in \mathfrak{s o}(n-1)$. On the other hand, when $A$ has the form (II),
$$
A=P\left[\begin{array}{cc}
\exp (S) & O_{n, 1} \\
O_{1, n} & 1
\end{array}\right] P^{-1}=\exp \left(P\left[\begin{array}{cc}
S & O_{n, 1} \\
O_{1, n} & 0
\end{array}\right] P^{-1}\right)
$$
for some $S \in \mathfrak{s} 0(n)$. In either case, $A \in \exp \operatorname{lor}(n, 1)$.
\end{PROOF}





\pagebreak





\subsection{$\mathbf{S L}_{\mathbf{2}}(\mathbb{C})$ and the Lorentz $\operatorname{Group} \operatorname{Lor}(\mathbf{3}, \mathbf{1})$}

Let us now consider the Lie algebra of $\mathrm{SL}_{2}(\mathbb{C}), \mathfrak{s} \mathfrak{l}_{2}(\mathbb{C})$. By Equation (3.8),
$$
\mathfrak{s l}_{2}(\mathbb{C})=\operatorname{ker} \operatorname{tr} \subseteq \mathbf{M}_{2}(\mathbb{C})
$$
and so $\operatorname{dim}_{\mathbb{C}} \mathfrak{s l}_{2}(\mathbb{C})=3$. The following matrices form a $\mathbb{C}$-basis for $\mathfrak{s l}_{2}(\mathbb{C})$:
$$
H^{\prime}=\left[\begin{array}{cc}
1 & 0 \\
0 & -1
\end{array}\right], \quad E^{\prime}=\left[\begin{array}{ll}
0 & 1 \\
0 & 0
\end{array}\right], \quad F^{\prime}=\left[\begin{array}{ll}
0 & 0 \\
1 & 0
\end{array}\right] .
$$
The elements $H^{\prime}, i H^{\prime}, E^{\prime}, i E^{\prime}, F^{\prime}, i F^{\prime}$ form an $\mathbb{R}$-basis and $\operatorname{dim} \mathfrak{s l}_{2}(\mathbb{C})=6$. 

Notice also that $\mathfrak{s u}(2) \subseteq \mathfrak{s l}_{2}(\mathbb{C})$ and the elements $H, E, F \in \mathfrak{s u}(2)$ form a $\mathbb{C}$-basis of $\mathfrak{s l}_{2}(\mathbb{C})$, so $H, i H, E, i E, F, i F$ form an $\mathbb{R}$-basis. 

The Lie brackets of $H^{\prime}, E^{\prime}, F^{\prime}$ are determined by
$$
\left[H^{\prime}, E^{\prime}\right]=2 E^{\prime}, \quad\left[H^{\prime}, F^{\prime}\right]=-2 F^{\prime}, \quad\left[E^{\prime}, F^{\prime}\right]=H^{\prime} .
$$


\begin{DEF}{LIE-B02-09-34}{Lie-Algebra \(\mathfrak{sl}_2(\mathbb{C})\)}
Die Lie-Algebra der komplexen speziellen linearen Gruppe
\[
\mathrm{SL}_2(\mathbb{C})
=\left\{A\in \mathrm{M}_2(\mathbb{C}) : \det A = 1\right\}
\]
ist
\[
\mathfrak{sl}_2(\mathbb{C})
=\left\{
X\in \mathrm{M}_2(\mathbb{C}) : \operatorname{tr}(X)=0
\right\}.
\]
Damit gilt \(\dim_{\mathbb{C}}\mathfrak{sl}_2(\mathbb{C})=3\) und
\(\dim_{\mathbb{R}}\mathfrak{sl}_2(\mathbb{C})=6\).

\begin{itemize}
\item Eine \(\mathbb{C}\)-Basis ist gegeben durch:
\[
H'=\begin{bmatrix}1 & 0 \\ 0 & -1\end{bmatrix},
\quad
E'=\begin{bmatrix}0 & 1 \\ 0 & 0\end{bmatrix},
\quad
F'=\begin{bmatrix}0 & 0 \\ 1 & 0\end{bmatrix}.
\]
\item Eine \(\mathbb{R}\)-Basis ist entsprechend
\[
H',\, iH',\, E',\, iE',\, F',\, iF'.
\]
\end{itemize}

Die Lie-Klammern der Basisvektoren erfüllen die Relationen
\[
[H',E']=2E',\qquad
[H',F']=-2F',\qquad
[E',F']=H'.
\]
Damit ist \(\mathfrak{sl}_2(\mathbb{C})\) die komplexifizierte Version der
reellen Lie-Algebra \(\mathfrak{su}(2)\), deren Elemente
\(H,E,F\) eine \(\mathbb{C}\)-Basis von \(\mathfrak{sl}_2(\mathbb{C})\) bilden.
\end{DEF}



Notice that the subspaces spanned by each of the pairs $H^{\prime}, E^{\prime}$ and $H^{\prime}, F^{\prime}$ are $\mathbb{C}$-Lie subalgebras. In fact, $H^{\prime}, E^{\prime}$ span the Lie algebra $\mathrm{ut}_{2}^{0}(\mathbb{C})$ of the group of upper triangular complex matrices of determinant 1 , while $H^{\prime}, F^{\prime}$ span the Lie algebra of the group of lower triangular complex matrices of determinant 1.


\begin{DEF}{LIE-B02-09-35}{Dreiecksunteralgebren von \(\mathfrak{sl}_2(\mathbb{C})\)}
Die von den Paaren \( (H',E') \) bzw.\ \( (H',F') \) erzeugten Unterräume sind
\(\mathbb{C}\)-Lie-Unteralgebren von \(\mathfrak{sl}_2(\mathbb{C})\).

\begin{itemize}
\item
Die von \(H'\) und \(E'\) erzeugte Unteralgebra
\[
\mathfrak{b}_+ = \operatorname{span}_{\mathbb{C}}\{H',E'\}
\]
ist die Lie-Algebra
\[
\mathfrak{ut}_2^{0}(\mathbb{C})
=\left\{
\begin{bmatrix}
a & b\\
0 & -a
\end{bmatrix}
: a,b\in\mathbb{C}
\right\},
\]
die zur Gruppe der oberen Dreiecksmatrizen mit Determinante \(1\) gehört.

\item
Analog ist
\[
\mathfrak{b}_- = \operatorname{span}_{\mathbb{C}}\{H',F'\}
=
\left\{
\begin{bmatrix}
a & 0\\
b & -a
\end{bmatrix}
: a,b\in\mathbb{C}
\right\}
\]
die Lie-Algebra der unteren Dreiecksmatrizen von Determinante \(1\).
\end{itemize}

Somit enthält \(\mathfrak{sl}_2(\mathbb{C})\) zwei natürliche, maximal
lösbare Unteralgebren – die sogenannten \emph{Borel-Unteralgebren} –,
die den oberen bzw.\ unteren Dreiecksteil der Gruppe
\(\mathrm{SL}_2(\mathbb{C})\) widerspiegeln.
\end{DEF}




Given the existence of the double covering homomorphism $\overline{\mathrm{Ad}}: \mathrm{SU}(2) \longrightarrow$ SO(3) of Section 3.5, it seems reasonable to ask if a similar homomorphism exists for $\mathrm{SL}_{2}(\mathbb{C})$. It does, but we need to use the Lorentz group $\operatorname{Lor}(3,1)$ and obtain an important double covering homomorphism $\operatorname{SL}_{2}(\mathbb{C}) \longrightarrow \operatorname{Lor}(3,1)$ which appears in Physics in connection with spinors and twistors.


\begin{THEO}{LIE-B02-09-36}{Doppelüberlagerungs-Homomorphismus \(\mathrm{SL}_2(\mathbb{C})\to \operatorname{Lor}(3,1)\)}
Es gibt einen kanonischen surjektiven Gruppenhomomorphismus
\[
\rho:\ \mathrm{SL}_2(\mathbb{C})\longrightarrow \operatorname{SO}^+(3,1)\subset \operatorname{Lor}(3,1),
\qquad \ker\rho=\{\pm I_2\},
\]
so dass \(\mathrm{SL}_2(\mathbb{C})\) die zusammenhängende, orientierungs- und
zeitorientierungserhaltende Lorentzgruppe \(\operatorname{SO}^+(3,1)\) zweifach überdeckt.
Insbesondere gilt (als Lie-Gruppen)
\[
\operatorname{Spin}(3,1)\ \cong\ \mathrm{SL}_2(\mathbb{C})\quad\text{und}\quad
\operatorname{Spin}(3,1)\xrightarrow{\ \rho\ }\operatorname{SO}^+(3,1)
\]
ist die kanonische Doppelüberlagerung.
\end{THEO}


\begin{PROOF}{LIE-B02-09-37}{P: Doppelüberlagerungs-Homomorphismus \(\mathrm{SL}_2(\mathbb{C})\to \operatorname{Lor}(3,1)\)}
[Konstruktion und Begründung]
Identifiziere Minkowski-Vektoren \(X=(\mathbf{x},t)\in\mathbb{R}^{3,1}\) mit hermiteschen
\(2\times2\)-Matrizen
\[
\widehat X\;=\;t\,\sigma_0+\sum_{j=1}^3 x_j\,\sigma_j
\;=\;
\begin{bmatrix}
t+x_3 & x_1-ix_2\\
x_1+ix_2 & t-x_3
\end{bmatrix},
\]
wobei \(\sigma_0=I_2\) und \(\sigma_1,\sigma_2,\sigma_3\) die Pauli-Matrizen sind.
Dann gilt
\[
\det(\widehat X)=t^2-\|\mathbf{x}\|^2=\langle X,X\rangle_{3,1}.
\]
Für \(g\in \mathrm{SL}_2(\mathbb{C})\) definiere eine Wirkung auf hermiteschen Matrizen durch
\[
\widehat X\ \longmapsto\ \widehat X' \;=\; g\,\widehat X\, g^{*},
\]
wobei \(g^{*}=\overline g^{\,T}\).
Diese Abbildung erhält Hermitizität und Determinante:
\(\det(\widehat X')=\det(g)\det(\widehat X)\det(g^{*})=\det(\widehat X)\).
Über der obigen Identifikation induziert das eine lineare Transformation
\(\Lambda(g)\) auf \(\mathbb{R}^{3,1}\) mit
\(\langle \Lambda(g)X,\Lambda(g)X\rangle=\langle X,X\rangle\), also
\(\Lambda(g)\in \operatorname{O}(3,1)\).
Zusätzlich ist \(\det \Lambda(g)=1\) und \(\Lambda(g)\) ist orthochron (erhält die Zukunfts-Lichtkegel),
weshalb \(\Lambda(g)\in \operatorname{SO}^+(3,1)\).
Setze \(\rho(g)=\Lambda(g)\).

\emph{Homomorphie:} \((g_1g_2)\widehat X(g_1g_2)^*=g_1(g_2\widehat X g_2^*)g_1^*\) zeigt
\(\rho(g_1g_2)=\rho(g_1)\rho(g_2)\).

\emph{Kern:} Aus \(g\widehat X g^*=\widehat X\) für alle \(\widehat X\) folgt,
dass \(g\) die adjungierte Wirkung auf dem \(4\)-dimensionalen reellen Raum
hermitescher Matrizen trivial macht; das geschieht genau für \(g=\pm I_2\).
Also \(\ker\rho=\{\pm I_2\}\).

\emph{Surjektivität:} \(\mathrm{SL}_2(\mathbb{C})\) ist einfach zusammenhängend und
\(\rho\) ist ein lokaler Diffeomorphismus auf Identitätskomponenten.
Da \(\dim \mathrm{SL}_2(\mathbb{C})=3\cdot 2=6=\dim \operatorname{SO}^+(3,1)\)
und \(\rho\) auf Lie-Algebra-Ebene ein Isomorphismus
\(\mathfrak{sl}_2(\mathbb{C})\cong \mathfrak{so}(3,1)\) ist, folgt die Surjektivität
auf der Identitätskomponente.
\end{PROOF}



\begin{CONC}{LIE-B02-09-38}{Physikalische Bedeutung Doppelüberlagerungs-Homomorphismus \(\mathrm{SL}_2(\mathbb{C})\to \operatorname{Lor}(3,1)\)}
Die Identifikation \(X\leftrightarrow \widehat X\) ist die \emph{Pauli-Vektor}-Kodierung.
Die Wirkung \(g:\widehat X\mapsto g\widehat X g^*\) beschreibt genau, wie sich
Vierervektoren (Energie-Impuls, Raumzeit-Koordinaten, elektromagnetischer Feldtensor in geeigneter Darstellung)
unter einem Wechsel des Inertialsystems transformieren.
Die \emph{Spinoren} sind die (fundamentale) Darstellung von \(\mathrm{SL}_2(\mathbb{C})\);
über \(\rho\) induziert sie die Vektor-Darstellung der Lorentzgruppe.
Dies erklärt, warum Fermionfelder (Spin-\(\tfrac12\)) unter \(\mathrm{SL}_2(\mathbb{C})\)
transformieren, während Beobachter-Größen (Vierervektoren) unter \(\operatorname{SO}^+(3,1)\)
transformieren.
\end{CONC}




Next we will describe the $\mathbb{R}$-Lie algebra $\operatorname{lor}(3,1)$ of $\operatorname{Lor}(3,1) \leqslant \operatorname{SL}_{4}(\mathbb{R})$. Let $\alpha:(-\varepsilon, \varepsilon) \longrightarrow \operatorname{Lor}(3,1)$ be a differentiable curve with $\alpha(0)=I$. By definition, for $t \in(-\varepsilon, \varepsilon)$ we have

$$
\alpha(t) Q \alpha(t)^{T}=Q,
$$

where

$$
Q=\left[\begin{array}{rrrr}
1 & 0 & 0 & 0 \\
0 & 1 & 0 & 0 \\
0 & 0 & 1 & 0 \\
0 & 0 & 0 & -1
\end{array}\right]
$$

Differentiating and setting $t=0$ we obtain

$$
\alpha^{\prime}(0) Q+Q \alpha^{\prime}(0)^{T}=O
$$

giving

$$
\begin{gathered}
{\left[\begin{array}{llll}
\alpha^{\prime}(0)_{11} & \alpha^{\prime}(0)_{12} & \alpha^{\prime}(0)_{13} & -\alpha^{\prime}(0)_{14} \\
\alpha^{\prime}(0)_{21} & \alpha^{\prime}(0)_{22} & \alpha^{\prime}(0)_{23} & -\alpha^{\prime}(0)_{24} \\
\alpha^{\prime}(0)_{31} & \alpha^{\prime}(0)_{32} & \alpha^{\prime}(0)_{33} & -\alpha^{\prime}(0)_{34} \\
\alpha^{\prime}(0)_{41} & \alpha^{\prime}(0)_{42} & \alpha^{\prime}(0)_{43} & -\alpha^{\prime}(0)_{44}
\end{array}\right]+} \\
{\left[\begin{array}{cccc}
\alpha^{\prime}(0)_{11} & \alpha^{\prime}(0)_{21} & \alpha^{\prime}(0)_{31} & \alpha^{\prime}(0)_{41} \\
\alpha^{\prime}(0)_{12} & \alpha^{\prime}(0)_{22} & \alpha^{\prime}(0)_{32} & \alpha^{\prime}(0)_{42} \\
\alpha^{\prime}(0)_{13} & \alpha^{\prime}(0)_{23} & \alpha^{\prime}(0)_{33} & \alpha^{\prime}(0)_{43} \\
-\alpha^{\prime}(0)_{14} & -\alpha^{\prime}(0)_{42} & -\alpha^{\prime}(0)_{34} & -\alpha^{\prime}(0)_{44}
\end{array}\right]=O .}
\end{gathered}
$$

So we have

$$
\alpha^{\prime}(0)=\left[\begin{array}{cccc}
0 & \alpha^{\prime}(0)_{12} & \alpha^{\prime}(0)_{13} & \alpha^{\prime}(0)_{14} \\
-\alpha^{\prime}(0)_{12} & 0 & \alpha^{\prime}(0)_{23} & \alpha^{\prime}(0)_{24} \\
-\alpha^{\prime}(0)_{13} & -\alpha^{\prime}(0)_{23} & 0 & \alpha^{\prime}(0)_{34} \\
\alpha^{\prime}(0)_{14} & \alpha^{\prime}(0)_{24} & \alpha^{\prime}(0)_{34} & 0
\end{array}\right]
$$



\begin{DEF}{LIE-B02-09-39}{Explicit form of \(\operatorname{lor}(3,1)\)}
Let \(Q=\mathrm{diag}(1,1,1,-1)\). The Lie algebra of the proper,
orthochronous Lorentz group \(\operatorname{Lor}(3,1)\subset \mathrm{SL}_4(\mathbb{R})\) is
\[
\operatorname{lor}(3,1)
=\bigl\{\,X\in \mathrm{M}_4(\mathbb{R})\;:\; X^T Q+QX=0\,\bigr\}.
\]
Equivalently, every \(X\in \operatorname{lor}(3,1)\) has the block form
\[
X=
\begin{bmatrix}
A & \mathbf{b}\\[2pt]
\mathbf{b}^{\!T} & 0
\end{bmatrix},
\qquad
A\in\mathfrak{so}(3)\ (\text{i.e. }A^T=-A),\ \ \mathbf{b}\in\mathbb{R}^{3}.
\]
Writing
\[
A=
\begin{bmatrix}
0 & a_{12} & a_{13}\\
-a_{12} & 0 & a_{23}\\
-a_{13} & -a_{23} & 0
\end{bmatrix},
\qquad
\mathbf{b}=
\begin{bmatrix} b_1\\ b_2\\ b_3\end{bmatrix},
\]
we obtain the coordinate description
\[
X=
\begin{bmatrix}
0 & a_{12} & a_{13} & b_1\\
-a_{12} & 0 & a_{23} & b_2\\
-a_{13} & -a_{23} & 0 & b_3\\
b_1 & b_2 & b_3 & 0
\end{bmatrix}.
\]
Thus \(\dim\operatorname{lor}(3,1)=3+3=6\).
\end{DEF}



\begin{DEF}{LIE-B02-09-40}{Rotation and boost generators for $\mathfrak{so}(3)$ and $\mathfrak{lor}(3,1)$}
Define the rotation generators \(J_1,J_2,J_3\) (coming from \(A\in\mathfrak{so}(3)\))
and the boost generators \(K_1,K_2,K_3\) (coming from \(\mathbf{b}\)) by
\[
J_1=
\begin{bmatrix}
0&0&0&0\\
0&0&1&0\\
0&-1&0&0\\
0&0&0&0
\end{bmatrix},\quad
J_2=
\begin{bmatrix}
0&0&-1&0\\
0&0&0&0\\
1&0&0&0\\
0&0&0&0
\end{bmatrix},\quad
J_3=
\begin{bmatrix}
0&1&0&0\\
-1&0&0&0\\
0&0&0&0\\
0&0&0&0
\end{bmatrix},
\]
\[
K_1=
\begin{bmatrix}
0&0&0&1\\
0&0&0&0\\
0&0&0&0\\
1&0&0&0
\end{bmatrix},\quad
K_2=
\begin{bmatrix}
0&0&0&0\\
0&0&0&1\\
0&0&0&0\\
0&1&0&0
\end{bmatrix},\quad
K_3=
\begin{bmatrix}
0&0&0&0\\
0&0&0&0\\
0&0&0&1\\
0&0&1&0
\end{bmatrix}.
\]
Every \(X\in\operatorname{lor}(3,1)\) decomposes uniquely as
\[
X=\omega_1 J_1+\omega_2 J_2+\omega_3 J_3+\beta_1 K_1+\beta_2 K_2+\beta_3 K_3,
\]
with \(\omega_i=a_{jk}\) (cyclic) and \(\beta_i=b_i\).
\end{DEF}



\begin{DEF}{LIE-B02-09-41}{Commutation relations for rotation and boost generators for $\mathfrak{so}(3)$}
With \(\varepsilon_{ijk}\) the Levi-Civita symbol, the nonzero Lie brackets are
\[
[J_i,J_j]=\varepsilon_{ijk}\,J_k,\qquad
[J_i,K_j]=\varepsilon_{ijk}\,K_k,\qquad
[K_i,K_j]=-\varepsilon_{ijk}\,J_k.
\]
Thus \(\operatorname{lor}(3,1)\cong \mathfrak{so}(3,1)\) is generated by spatial rotations \(J_i\)
and boosts \(K_i\).
\end{DEF}



\begin{DEF}{LIE-B02-09-42}{Exponentials = rotations and boosts for $\SO(3), \mathrm{Lor}(3,1)$}
For \(\boldsymbol{\omega}\in\mathbb{R}^3\) and \(\boldsymbol{\beta}\in\mathbb{R}^3\),
\[
\exp\!\Big(\sum_i \omega_i J_i\Big)\in \mathrm{SO}(3)\subset \operatorname{Lor}(3,1)
\quad\text{(a spatial rotation),}
\]
\[
\exp\!\Big(\sum_i \beta_i K_i\Big)\in \operatorname{Lor}(3,1)
\quad\text{(a pure boost along direction \(\boldsymbol{\beta}\)).}
\]
In particular, \(\exp(tK_1)=
\begin{bmatrix}
\cosh t&0&0&\sinh t\\
0&1&0&0\\
0&0&1&0\\
\sinh t&0&0&\cosh t
\end{bmatrix}\),
and analogously for \(K_2,K_3\).
\end{DEF}




Notice that the trace of such a matrix is zero. In fact, every matrix of the form

$$
A=\left[\begin{array}{cccc}
0 & a_{12} & a_{13} & a_{14} \\
-a_{12} & 0 & a_{23} & a_{24} \\
-a_{13} & -a_{23} & 0 & a_{34} \\
a_{14} & a_{24} & a_{34} & 0
\end{array}\right]
$$

or equivalently satisfying

$$
A Q+Q A^{T}=O \quad(\text { and hence } \operatorname{tr} A=0)
$$

is in $\operatorname{lor}(3,1)$. To see this, consider the curve

$$
\alpha: \mathbb{R} \longrightarrow \mathrm{GL}_{4}(\mathbb{R}) ; \quad \alpha(t)=\exp (t A),
$$

for which $\alpha^{\prime}(0)=A$. Since $Q A^{T}=-A Q$, we have

$$
\exp (t A) Q \exp (t A)^{T}=\exp (t A) Q \exp \left(t A^{T}\right)=\exp (t A) \exp (-t A) Q=Q
$$

By Lemma 3.23,

$$
\operatorname{det} \exp (t A)=e^{\operatorname{tr}(t A)}=1 .
$$

Finally, for each $t \in \mathbb{R}, \alpha(t)$ preserves each of the components $\mathcal{H}_{3,1}^{ \pm}(1)$ of the hyperboloid

$$
x_{1}^{2}+x_{2}^{2}+x_{3}^{2}-x_{4}^{2}=-1 .
$$

All of this shows that $\exp (t A) \in \operatorname{Lor}(3,1)$ so we might as well redefine

$$
\alpha: \mathbb{R} \longrightarrow \operatorname{Lor}(3,1) ; \quad \alpha(t)=\exp (t A) .
$$



\begin{PROP}{LIE-B02-09-43}{Exponentials of \(Q\)-skew matrices lie in \(\operatorname{Lor}(3,1)\)}
Let \(Q=\mathrm{diag}(1,1,1,-1)\) and let
\[
A=
\begin{bmatrix}
0 & a_{12} & a_{13} & a_{14}\\
-a_{12} & 0 & a_{23} & a_{24}\\
-a_{13} & -a_{23} & 0 & a_{34}\\
a_{14} & a_{24} & a_{34} & 0
\end{bmatrix}
\quad\Longleftrightarrow\quad
AQ+QA^{T}=0 .
\]
Then \(\operatorname{tr}A=0\) and for all \(t\in\mathbb{R}\),
\[
\alpha(t):= e^{tA}\in \operatorname{Lor}(3,1),
\qquad\text{i.e.}\quad
\alpha(t)^{T}Q\,\alpha(t)=Q,\ \ \det\alpha(t)=1,
\]
and \(\alpha(t)\) preserves each connected sheet \(\mathcal{H}_{3,1}^{\pm}(1)\) of the unit
hyperboloid \(x_1^2+x_2^2+x_3^2-x_4^2=-1\).
\end{PROP}



\begin{PROOF}{LIE-B02-09-44}{P: Exponentials of \(Q\)-skew matrices lie in \(\operatorname{Lor}(3,1)\)}
The displayed form of \(A\) is equivalent to \(AQ+QA^{T}=0\), so \(A\in\mathfrak{so}(3,1)\).
Its diagonal entries are zero, hence \(\operatorname{tr}A=0\).

Set \(\alpha(t)=e^{tA}\). Consider
\(
f(t):=\alpha(t)Q\alpha(t)^{T}.
\)
Differentiating and using \(AQ+QA^{T}=0\),
\[
f'(t)=\alpha(t)\bigl(AQ+QA^{T}\bigr)\alpha(t)^{T}=0,
\]
so \(f(t)\equiv f(0)=Q\). Thus \(\alpha(t)^{T}Q\,\alpha(t)=Q\) for all \(t\), i.e. \(\alpha(t)\in O(3,1)\).

By Lemma on determinants of exponentials,
\(
\det\alpha(t)=\det e^{tA}=e^{\operatorname{tr}(tA)}=1,
\)
hence \(\alpha(t)\in SO(3,1)\) for all \(t\).

Finally, \(\alpha(0)=I\) preserves each connected component \(\mathcal{H}_{3,1}^{\pm}(1)\),
and \(t\mapsto \alpha(t)\) is continuous in \(SO(3,1)\). Since \(\alpha(t)\) preserves the
Lorentz norm, it maps the unit timelike vector \(e_4=(0,0,0,1)\) to another unit timelike
vector \(\alpha(t)e_4\). The time component of \(\alpha(t)e_4\) depends continuously on \(t\)
and is \(+1\) at \(t=0\); thus it stays positive for all \(t\). Therefore \(\alpha(t)\) preserves
the future sheet \(\mathcal{H}_{3,1}^{+}(1)\), and by determinant \(=1\), it also preserves
\(\mathcal{H}_{3,1}^{-}(1)\). Hence \(\alpha(t)\in \operatorname{Lor}(3,1)\).
\end{PROOF}






We also have $A \in \mathfrak{l o r}(3,1)$. So we have
$$
\begin{aligned}
\mathfrak{l o r}(3,1) & =\left\{A \in \mathrm{M}_{4}(\mathbb{R}): A Q+Q A^{T}=O\right\} \\
& =\left\{A \in \mathrm{M}_{4}(\mathbb{R}): A=\left[\begin{array}{cccc}
0 & a_{12} & a_{13} & a_{14} \\
-a_{12} & 0 & a_{23} & a_{24} \\
-a_{13} & -a_{23} & 0 & a_{34} \\
a_{14} & a_{24} & a_{34} & 0
\end{array}\right]\right\}
\end{aligned}
$$
and therefore
$$
\operatorname{dim} \operatorname{Lor}(3,1)=\operatorname{dim} \operatorname{lor}(3,1)=6 .
$$


\begin{DEF}{LIE-B02-09-45}{Struktur und Dimension der Lie-Algebra \(\mathfrak{lor}(3,1)\)}
Die Lie-Algebra der Lorentz-Gruppe \(\operatorname{Lor}(3,1)\) ist
\[
\mathfrak{lor}(3,1)
=\bigl\{\,A\in \mathrm{M}_4(\mathbb{R})\;:\; A Q+Q A^{T}=0\,\bigr\},
\qquad
Q=\mathrm{diag}(1,1,1,-1).
\]
Explizit besitzt jedes \(A\in\mathfrak{lor}(3,1)\) die Form
\[
A=
\begin{bmatrix}
0 & a_{12} & a_{13} & a_{14}\\
-a_{12} & 0 & a_{23} & a_{24}\\
-a_{13} & -a_{23} & 0 & a_{34}\\
a_{14} & a_{24} & a_{34} & 0
\end{bmatrix},
\qquad a_{ij}\in\mathbb{R}.
\]
Jedes Element erfüllt \(A^{T}Q=-QA\) und \(\operatorname{tr}A=0\).
Die freien Parameter \(a_{12},a_{13},a_{23}\) beschreiben die
\emph{Rotationsgeneratoren}, während \(a_{14},a_{24},a_{34}\)
den \emph{Boost-Generatoren} entsprechen.

Somit gilt:
\[
\dim_{\mathbb{R}}\mathfrak{lor}(3,1)=3+3=6
\qquad\text{und}\qquad
\dim_{\mathbb{R}}\operatorname{Lor}(3,1)=6.
\]
\end{DEF}



\begin{CONC}{LIE-B02-09-46}{Interpretation der sechs Parameter der Lorentz Algebra}
\begin{itemize}
  \item Die drei Parameter \(a_{12},a_{13},a_{23}\) entsprechen infinitesimalen
  Rotationen im räumlichen Teil \(\mathbb{R}^3\).
  \item Die drei Parameter \(a_{14},a_{24},a_{34}\) entsprechen infinitesimalen
  \emph{Boosts} (Lorentz-Transformationen zwischen Inertialsystemen, die sich mit
  konstanter Geschwindigkeit relativ zueinander bewegen).
\end{itemize}
Damit zerfällt die Lie-Algebra \(\mathfrak{lor}(3,1)\) als
\[
\mathfrak{lor}(3,1)
=\operatorname{span}\{J_1,J_2,J_3\}\oplus\operatorname{span}\{K_1,K_2,K_3\},
\]
wobei \(J_i\) die Generatoren der Raumrotationen und \(K_i\) die der Boosts sind.
\end{CONC}




An $\mathbb{R}$-basis for $\operatorname{lor}(3,1)$ consists of the elements
$$
\begin{array}{lll}
P_{12}=\left[\begin{array}{cccc}
0 & -1 & 0 & 0 \\
1 & 0 & 0 & 0 \\
0 & 0 & 0 & 0 \\
0 & 0 & 0 & 0
\end{array}\right], & P_{13}=\left[\begin{array}{cccc}
0 & 0 & -1 & 0 \\
0 & 0 & 0 & 0 \\
1 & 0 & 0 & 0 \\
0 & 0 & 0 & 0
\end{array}\right], & P_{14}=\left[\begin{array}{llll}
0 & 0 & 0 & 1 \\
0 & 0 & 0 & 0 \\
0 & 0 & 0 & 0 \\
1 & 0 & 0 & 0
\end{array}\right], \\
P_{23}=\left[\begin{array}{cccc}
0 & 0 & 0 & 0 \\
0 & 0 & -1 & 0 \\
0 & 1 & 0 & 0 \\
0 & 0 & 0 & 0
\end{array}\right], & P_{24}=\left[\begin{array}{cccc}
0 & 0 & 0 & 0 \\
0 & 0 & 0 & 1 \\
0 & 0 & 0 & 0 \\
0 & 1 & 0 & 0
\end{array}\right], & P_{34}=\left[\begin{array}{llll}
0 & 0 & 0 & 0 \\
0 & 0 & 0 & 0 \\
0 & 0 & 0 & 1 \\
0 & 0 & 1 & 0
\end{array}\right] .
\end{array}
$$
The non-trivial brackets for these are
$$
\begin{array}{lllll}
{\left[P_{12}, P_{13}\right]=P_{23},} & {\left[P_{12}, P_{14}\right]=P_{24},} & {\left[P_{12}, P_{23}\right]=-P_{13},} & {\left[P_{12}, P_{24}\right]=-P_{14},} \\
{\left[P_{12}, P_{34}\right]=0,} & {\left[P_{13}, P_{14}\right]=P_{34},} & {\left[P_{13}, P_{23}\right]=P_{12},} & {\left[P_{13}, P_{24}\right]=0,} \\
{\left[P_{13}, P_{34}\right]=0,} & {\left[P_{14}, P_{23}\right]=0,} & {\left[P_{14}, P_{24}\right]=-P_{12},} & {\left[P_{14}, P_{34}\right]=-P_{13},} \\
{\left[P_{23}, P_{24}\right]=P_{34},} & {\left[P_{23}, P_{34}\right]=-P_{24},} & {\left[P_{24}, P_{34}\right]=-P_{23} .} &
\end{array}
$$



\begin{DEF}{LIE-B02-09-47}{Basis und Kommutatoren von \(\mathfrak{lor}(3,1)\)}
Eine \(\mathbb{R}\)-Basis von \(\mathfrak{lor}(3,1)\) ist gegeben durch
\[
\{\,P_{12},P_{13},P_{14},P_{23},P_{24},P_{34}\,\},
\]
wobei
\[
\begin{aligned}
P_{12}&=\begin{bmatrix}0&-1&0&0\\ 1&0&0&0\\ 0&0&0&0\\ 0&0&0&0\end{bmatrix},&
P_{13}&=\begin{bmatrix}0&0&-1&0\\ 0&0&0&0\\ 1&0&0&0\\ 0&0&0&0\end{bmatrix},&
P_{14}&=\begin{bmatrix}0&0&0&1\\ 0&0&0&0\\ 0&0&0&0\\ 1&0&0&0\end{bmatrix},\\[2mm]
P_{23}&=\begin{bmatrix}0&0&0&0\\ 0&0&-1&0\\ 0&1&0&0\\ 0&0&0&0\end{bmatrix},&
P_{24}&=\begin{bmatrix}0&0&0&0\\ 0&0&0&1\\ 0&0&0&0\\ 0&1&0&0\end{bmatrix},&
P_{34}&=\begin{bmatrix}0&0&0&0\\ 0&0&0&0\\ 0&0&0&1\\ 0&0&1&0\end{bmatrix}.
\end{aligned}
\]
Die nichttrivialen Lie-Klammern sind:
\[
\begin{array}{llll}
[P_{12},P_{13}]=P_{23}, & [P_{12},P_{14}]=P_{24}, & [P_{12},P_{23}]=-P_{13}, & [P_{12},P_{24}]=-P_{14},\\[1mm]
[P_{12},P_{34}]=0, & [P_{13},P_{14}]=P_{34}, & [P_{13},P_{23}]=P_{12}, & [P_{13},P_{24}]=0,\\[1mm]
[P_{13},P_{34}]=0, & [P_{14},P_{23}]=0, & [P_{14},P_{24}]=-P_{12}, & [P_{14},P_{34}]=-P_{13},\\[1mm]
[P_{23},P_{24}]=P_{34}, & [P_{23},P_{34}]=-P_{24}, & [P_{24},P_{34}]=-P_{23}. &
\end{array}
\]

Die drei Elemente \(P_{12},P_{13},P_{23}\) erzeugen die räumlichen Rotationen,
die drei Elemente \(P_{14},P_{24},P_{34}\) die Boosts. Mit der Signatur
\(\eta=\mathrm{diag}(1,1,1,-1)\) und den Standardgeneratoren
\(M_{\mu\nu}=-M_{\nu\mu}\) (identifiziert durch \(M_{ij}=P_{ij}\)) gilt die
kompakte Lorentz-Algebra-Formel
\[
[M_{\mu\nu},M_{\rho\sigma}]
=\eta_{\nu\rho}M_{\mu\sigma}-\eta_{\mu\rho}M_{\nu\sigma}
-\eta_{\nu\sigma}M_{\mu\rho}+\eta_{\mu\sigma}M_{\nu\rho}.
\]
Setzt man
\[
J_1:=P_{23},\quad J_2:=P_{31}=-P_{13},\quad J_3:=P_{12},\qquad
K_1:=P_{14},\ K_2:=P_{24},\ K_3:=P_{34},
\]
ergeben sich die vertrauten Relationen
\[
[J_i,J_j]=\varepsilon_{ijk}\,J_k,\qquad
[J_i,K_j]=\varepsilon_{ijk}\,K_k,\qquad
[K_i,K_j]=-\varepsilon_{ijk}\,J_k,
\]
wobei \(\varepsilon_{ijk}\) das Levi-Civita-Symbol ist.
\end{DEF}




We will now define the homomorphism $\mathrm{SL}_{2}(\mathbb{C}) \longrightarrow \operatorname{Lor}(3,1)$. To do this we will identify the $2 \times 2$ skew hermitian matrices Sk - $\operatorname{Herm}_{2}(\mathbb{C})$ with $\mathbb{R}^{4}$ using the correspondence
$$
\left[\begin{array}{cc}
(t+x) i & y+z i \\
-y+z i & (t-x) i
\end{array}\right] \longleftrightarrow x \mathbf{e}_{1}+y \mathbf{e}_{2}+z \mathbf{e}_{3}+t \mathbf{e}_{4} .
$$

Define an $\mathbb{R}$-bilinear inner product on Sk - $\operatorname{Herm}_{2}(\mathbb{C})$ by the formula
$$
\left\langle S_{1} \mid S_{2}\right\rangle=\frac{1}{4}\left(\operatorname{det}\left(S_{1}+S_{2}\right)-\operatorname{det}\left(S_{1}-S_{2}\right)\right)
$$

When $S_{1}=S_{2}=S$ we have
$$
\langle S \mid S\rangle=\frac{1}{4}(\operatorname{det} 2 S-\operatorname{det} O)=\operatorname{det} S
$$

It is easy to check that
$$
\begin{aligned}
\left\langle\left[\begin{array}{cc}
\left(t_{1}+x_{1}\right) i & y_{1}+z_{1} i \\
-y_{1}+z_{1} i & \left(t_{1}-x_{1}\right) i
\end{array}\right]\right|\left[\begin{array}{cc}
\left(t_{2}+x_{2}\right) i & y_{2}+z_{2} i \\
-y_{2}+z_{2} i & \left(t_{2}-x_{2}\right) i
\end{array}\right] & =x_{1} x_{2}+y_{1} y_{2}+z_{1} z_{2}-t_{1} t_{2}
\end{aligned}
$$
which is the Lorentzian inner product on $\mathbb{R}^{3,1}=\mathbb{R}^{4}$ given by
$$
x_{1} x_{2}+y_{1} y_{2}+z_{1} z_{2}-t_{1} t_{2}=\left[\begin{array}{llll}
x_{1} & y_{1} & z_{1} & t_{1}
\end{array}\right]\left[\begin{array}{rrrr}
1 & 0 & 0 & 0 \\
0 & 1 & 0 & 0 \\
0 & 0 & 1 & 0 \\
0 & 0 & 0 & -1
\end{array}\right]\left[\begin{array}{l}
x_{2} \\
y_{2} \\
z_{2} \\
t_{2}
\end{array}\right]
$$

Now observe that for $A \in \mathrm{SL}_{2}(\mathbb{C})$ and $S \in \mathrm{Sk}-\operatorname{Herm}_{2}(\mathbb{C})$,
$$
\left(A S A^{*}\right)^{*}=A S^{*} A^{*}=-A S A^{*}
$$
so $A S A^{*} \in \operatorname{Sk}-\operatorname{Herm}_{2}(\mathbb{C})$. By Equation (6.9a), for $S_{1}, S_{2} \in \operatorname{Sk}-\operatorname{Herm}_{2}(\mathbb{C})$, and the fact that $\operatorname{det} A=1=\operatorname{det} A^{*}$,
$$
\begin{aligned}
\left\langle A S_{1} A^{*} \mid A S_{2} A^{*}\right\rangle & =\frac{1}{4}\left(\operatorname{det} A\left(S_{1}+S_{2}\right) A^{*}-\operatorname{det} A\left(S_{1}-S_{2}\right) A^{*}\right) \\
& =\frac{1}{4}\left(\operatorname{det} A \operatorname{det}\left(S_{1}+S_{2}\right) \operatorname{det} A^{*}-\operatorname{det} A \operatorname{det}\left(S_{1}-S_{2}\right) \operatorname{det} A^{*}\right) \\
& =\frac{1}{4}\left(\operatorname{det}\left(S_{1}+S_{2}\right)-\operatorname{det}\left(S_{1}-S_{2}\right)\right) \\
& =\left\langle S_{1} \mid S_{2}\right\rangle
\end{aligned}
$$
Hence the function
$$
\mathrm{Sk}-\mathrm{Herm}_{2}(\mathbb{C}) \longrightarrow \mathrm{Sk}^{-} \mathrm{Herm}_{2}(\mathbb{C}) ; \quad S \mapsto A S A^{*},
$$
is an $\mathbb{R}$-linear transformation preserving the inner product $\langle\mid\rangle$. We can identify this with an $\mathbb{R}$-linear transformation $\widetilde{\operatorname{Ad}}_{A}: \mathbb{R}^{3,1} \rightarrow \mathbb{R}^{3,1}$ which preserves the Lorentzian inner product. In fact, $\operatorname{det} \widetilde{\mathrm{Ad}}_{A}=1$ and $\widetilde{\mathrm{Ad}}_{A}$ preserves the components of the hyperboloid $x^{2}+y^{2}+z^{2}-t^{2}=-1$. Let
$$
\widetilde{\mathrm{Ad}}: \mathrm{SL}_{2}(\mathbb{C}) \longrightarrow \operatorname{Lor}(3,1) ; \quad \widetilde{\mathrm{Ad}}(A)=\widetilde{\mathrm{Ad}}_{A}
$$
Then $\widetilde{\text { Ad }}$ is a homomorphism since
$$
\widetilde{\operatorname{Ad}}_{A B}(S)=A B(S)(A B)^{*}=A B(S) B^{*} A^{*}=\widetilde{\operatorname{Ad}}_{A}\left(\widetilde{\operatorname{Ad}}_{B}(S)\right)=\widetilde{\operatorname{Ad}}_{A} \widetilde{\operatorname{Ad}}_{B}(S)
$$
It is clearly continuous. Also, $A \in \operatorname{ker} \widetilde{\operatorname{Ad}}$ if and only if $A S A^{*}=S$ for all $S \in \mathrm{Sk}-\mathrm{Herm}_{2}(\mathbb{C})$, and it is easy to see that this occurs exactly when $A= \pm I$. Thus we have ker $\widetilde{\mathrm{Ad}}=\{ \pm I\}$.



\begin{DEF}{LIE-B02-09-48}{Kanonische Homomorphismus \(\widetilde{\operatorname{Ad}}:\mathrm{SL}_2(\mathbb{C})\to \operatorname{Lor}(3,1)\)}
Dieser Abschnitt konstruiert systematisch die kanonische Doppelüberlagerung
\[
\widetilde{\operatorname{Ad}}:\mathrm{SL}_2(\mathbb{C})\longrightarrow \operatorname{Lor}(3,1),
\]
welche in der Physik die fundamentale Verbindung zwischen komplexen Spinoren und Lorentz-Transformationen darstellt.

\paragraph{(1) Identifikation von \(\mathbb{R}^{3,1}\) mit schiefhermiteschen Matrizen.}
Sei
\[
\operatorname{SkHerm}_2(\mathbb{C})
=\{S\in \mathrm{M}_2(\mathbb{C}) : S^*=-S\}
\]
der reelle Vektorraum der \(2\times2\)-schiefhermiteschen Matrizen.
Wir identifizieren \(\operatorname{SkHerm}_2(\mathbb{C})\) mit \(\mathbb{R}^{4}\) durch die Zuordnung
\[
\begin{bmatrix}
(t+x)i & y+zi\\
-y+zi & (t-x)i
\end{bmatrix}
\longleftrightarrow
x\,\mathbf{e}_1+y\,\mathbf{e}_2+z\,\mathbf{e}_3+t\,\mathbf{e}_4.
\]

\paragraph{(2) Das bilineare Skalarprodukt.}
Auf \(\operatorname{SkHerm}_2(\mathbb{C})\) wird ein reell-bilineares Skalarprodukt definiert durch
\[
\langle S_1\mid S_2\rangle
=\tfrac{1}{4}\bigl(\det(S_1+S_2)-\det(S_1-S_2)\bigr),
\]
woraus für \(S_1=S_2=S\) unmittelbar folgt:
\[
\langle S\mid S\rangle=\det S.
\]
Setzt man
\[
S_i=
\begin{bmatrix}
(t_i+x_i)i & y_i+z_i i\\
-y_i+z_i i & (t_i-x_i)i
\end{bmatrix},
\]
so erhält man
\[
\langle S_1\mid S_2\rangle
=x_1x_2+y_1y_2+z_1z_2-t_1t_2,
\]
das heißt, das durch diese Zuordnung induzierte Skalarprodukt ist genau das Lorentz-Produkt auf \(\mathbb{R}^{3,1}\):
\[
\langle\mathbf{X}_1,\mathbf{X}_2\rangle
=\mathbf{X}_1^T
\begin{bmatrix}
1 & 0 & 0 & 0\\
0 & 1 & 0 & 0\\
0 & 0 & 1 & 0\\
0 & 0 & 0 & -1
\end{bmatrix}
\mathbf{X}_2.
\]

\paragraph{(3) Wirkung von \(\mathrm{SL}_2(\mathbb{C})\) auf \(\operatorname{SkHerm}_2(\mathbb{C})\).}
Für \(A\in \mathrm{SL}_2(\mathbb{C})\) und \(S\in \operatorname{SkHerm}_2(\mathbb{C})\) definieren wir
\[
\Phi_A(S)=A S A^*.
\]
Da
\[
(ASA^*)^*=A S^* A^*=-A S A^*,
\]
liegt \(\Phi_A(S)\) wieder in \(\operatorname{SkHerm}_2(\mathbb{C})\).
Das Skalarprodukt wird durch diese Abbildung erhalten:
\[
\begin{aligned}
\langle A S_1 A^*\mid A S_2 A^*\rangle
&=\tfrac{1}{4}\bigl(\det(A(S_1+S_2)A^*)-\det(A(S_1-S_2)A^*)\bigr)\\
&=\tfrac{1}{4}\bigl(\det(S_1+S_2)-\det(S_1-S_2)\bigr)
=\langle S_1\mid S_2\rangle,
\end{aligned}
\]
da \(\det A=\det A^*=1\).
Somit ist \(\Phi_A\) eine reell-lineare Isometrie von
\(\operatorname{SkHerm}_2(\mathbb{C})\) bezüglich des Lorentz-Produkts.

\paragraph{(4) Induzierter Lorentz-Homomorphismus.}
Über die Identifikation \(\operatorname{SkHerm}_2(\mathbb{C})\cong \mathbb{R}^{3,1}\)
definieren wir für jedes \(A\in\mathrm{SL}_2(\mathbb{C})\) die lineare Abbildung
\[
\widetilde{\operatorname{Ad}}_A:\mathbb{R}^{3,1}\to\mathbb{R}^{3,1},
\qquad
\widetilde{\operatorname{Ad}}_A(S)=A S A^*.
\]
Diese Abbildung erfüllt:
\[
\det\widetilde{\operatorname{Ad}}_A=1,
\qquad
\widetilde{\operatorname{Ad}}_A\ \text{erhält die Hyperboloide}\ 
x^2+y^2+z^2-t^2=-1.
\]
Daher gilt \(\widetilde{\operatorname{Ad}}_A\in \operatorname{Lor}(3,1)\).

Dies definiert einen Gruppenhomomorphismus
\[
\boxed{
\widetilde{\operatorname{Ad}}:\mathrm{SL}_2(\mathbb{C})
\longrightarrow \operatorname{Lor}(3,1),
\qquad
A\mapsto \widetilde{\operatorname{Ad}}_A.
}
\]
Die Homomorphie folgt aus:
\[
\widetilde{\operatorname{Ad}}_{AB}(S)
=AB\,S(AB)^*
=A(BSB^*)A^*
=\widetilde{\operatorname{Ad}}_A\bigl(\widetilde{\operatorname{Ad}}_B(S)\bigr).
\]
Kontinuität ist offensichtlich.

\paragraph{(5) Der Kern.}
Ein \(A\in\mathrm{SL}_2(\mathbb{C})\) liegt genau dann im Kern von \(\widetilde{\operatorname{Ad}}\),
wenn
\[
A S A^*=S
\quad\forall\, S\in \operatorname{SkHerm}_2(\mathbb{C}).
\]
Dies gilt nur für \(A=\pm I\), also
\[
\ker\widetilde{\operatorname{Ad}}=\{\pm I\}.
\]

\paragraph{(6) Schlussfolgerung.}
Damit ergibt sich die exakte Sequenz
\[
1\longrightarrow\{\pm I\}\longrightarrow
\mathrm{SL}_2(\mathbb{C})
\overset{\widetilde{\operatorname{Ad}}}{\longrightarrow}
\operatorname{Lor}(3,1)
\longrightarrow 1,
\]
d.\,h. \(\mathrm{SL}_2(\mathbb{C})\) ist eine zweifache Überlagerung von \(\operatorname{Lor}(3,1)\).
\end{DEF}


\begin{CONC}{LIE-B02-09-49}{Physikalische Bedeutung der Kanonische Homomorphismus \(\widetilde{\operatorname{Ad}}:\mathrm{SL}_2(\mathbb{C})\to \operatorname{Lor}(3,1)\)}
Diese Abbildung realisiert die fundamentale Beziehung zwischen Spinoren und Lorentz-Transformationen:
\(\mathrm{SL}_2(\mathbb{C})\) beschreibt die \emph{Spin(3,1)}-Struktur,
deren Darstellung auf \(\mathbb{R}^{3,1}\) genau die gewöhnlichen Lorentz-Transformationen liefert.
Daher treten Spinoren physikalisch als „Quadratwurzeln“ von Lorentz-Vektoren auf.
\end{CONC}





Theorem 6.12


i) The derivative $\mathrm{d} \widetilde{\mathrm{Ad}}: \mathfrak{s l}_{2}(\mathbb{C}) \longrightarrow \mathfrak{l o r}(3,1)$ is an isomorphism of $\mathbb{R}$-Lie algebras. ii) $\widetilde{\mathrm{Ad}}: \mathrm{SL}_{2}(\mathbb{C}) \longrightarrow \operatorname{Lor}(3,1)$ is a continuous surjective $\operatorname{Lie}$ homomorphism for which $\operatorname{ker} \widehat{\mathrm{Ad}}=\{ \pm I\}$, hence

$$
\mathrm{SL}_{2}(\mathbb{C}) /\{ \pm I\} \cong \operatorname{Lor}(3,1)
$$





\begin{THEO}{LIE-B02-09-50}{Differential und Doppelüberlagerung \(\mathrm{SL}_2(\mathbb{C})\to \operatorname{Lor}(3,1)\)}
\begin{enumerate}[label=\textnormal{(\roman*)}]
\item Das Differential der konstruierten kanonischen Abbildung
\[
\widetilde{\Ad}:\SL_2(\C)\longrightarrow \Lor(3,1),
\qquad
S\mapsto ASA^*,
\]
ist ein Isomorphismus reeller Lie-Algebren
\[
\mathrm{d}\widetilde{\Ad}:\mathfrak{sl}_2(\C)\to\mathfrak{lor}(3,1)
\]
\item \(\widetilde{\Ad}:\SL_2(\C)\to \Lor(3,1)\) ist ein stetiger surjektiver Lie-Gruppenhomomorphismus
mit
\(\ker\widetilde{\Ad}=\{\pm I\}\). Insbesondere gilt die kanonische Isomorphie
\[
\SL_2(\C)/\{\pm I\}\;\cong\;\Lor(3,1).
\]
\end{enumerate}
\end{THEO}



Proof


\begin{PROOF}{LIE-B02-09-51}{P: Differential und Doppelüberlagerung \(\mathrm{SL}_2(\mathbb{C})\to \operatorname{Lor}(3,1)\)}
(i) As in the case of $\mathrm{SU}(2)$ and $\mathrm{SO}(3)$, we can determine the derivative $\mathrm{d} \widetilde{\mathrm{Ad}}$ by considering for each $C \in \mathfrak{s l}_{2}(\mathbb{C})$, the curve
$$
\gamma: \mathbb{R} \longrightarrow \mathrm{SL}_{2}(\mathbb{C}) ; \quad \gamma(t)=\exp (t C),
$$
which gives rise to the curve
$$
\bar{\gamma}: \mathbb{R} \longrightarrow \operatorname{Lor}(3,1) ; \quad \bar{\gamma}(t)=\widetilde{\operatorname{Ad}}_{\gamma(t)}
$$
Using as an $\mathbb{R}$-basis for $\mathrm{Sk}-\mathrm{Herm}_{2}(\mathbb{C})$ the vectors
$$
V_{1}=\left[\begin{array}{rr}
i & 0 \\
0 & -i
\end{array}\right], \quad V_{2}=\left[\begin{array}{rr}
0 & 1 \\
-1 & 0
\end{array}\right], \quad V_{3}=\left[\begin{array}{cc}
0 & i \\
i & 0
\end{array}\right], \quad V_{4}=\left[\begin{array}{cc}
i & 0 \\
0 & i
\end{array}\right]
$$
we can determine the action of $\widetilde{\mathrm{Ad}}_{\gamma(t)}$ on Sk-Herm ${ }_{2}(\mathbb{C})$ and interpret it as an element of $\operatorname{Lor}(3,1)$. Differentiating at $t=0$, we obtain the action of $C$ as an element of $\mathfrak{l o r}(3,1)$ and so $\mathrm{d} \widetilde{\mathrm{Ad}}(C)$. For $X \in \mathrm{Sk}-\mathrm{Herm}_{2}(\mathbb{C})$ we have
$$
\widetilde{\operatorname{Ad}}_{\gamma(t)}(X)=\exp (t C) X \exp (t C)^{*}=\exp (t C) X \exp \left(t C^{*}\right)
$$
hence
$$
\left.\frac{\mathrm{d}}{\mathrm{~d} t} \widetilde{\mathrm{Ad}}_{\gamma(t)}(X) \right\rvert\, t=0=C X+X C^{*}
$$
So for the $\mathbb{R}$-basis $H, i H, E, i E, F, i F$ of $\mathfrak{s l}_{2}(\mathbb{C})$, we have
$$
H\left(x_{1} V_{1}+x_{2} V_{2}+x_{3} V_{3}+x_{4} V_{4}\right)+\left(x_{1} V_{1}+x_{2} V_{2}+x_{3} V_{3}+x_{4} V_{4}\right) H^{*}=x_{2} V_{3}-x_{3} V_{2},
$$
so
$$
\mathrm{d} \widetilde{\operatorname{Ad}}(H)=\left[\begin{array}{rrrr}
0 & 0 & 0 & 0 \\
0 & 0 & -1 & 0 \\
0 & 1 & 0 & 0 \\
0 & 0 & 0 & 0
\end{array}\right] .
$$
Here is the complete list written in terms of the matrices $P_{r s}$ which we know form an $\mathbb{R}$-basis of $\operatorname{Lor}(3,1)$ :
$$\mathrm{d} \widetilde{\mathrm{Ad}}(H)=\left[\begin{array}{rrrr}0 & 0 & 0 & 0 \\ 0 & 0 & -1 & 0 \\ 0 & 1 & 0 & 0 \\ 0 & 0 & 0 & 0\end{array}\right]=P_{23}, \quad \mathrm{~d} \widetilde{\mathrm{Ad}}(i H)=\left[\begin{array}{rrrr}0 & 0 & 0 & -1 \\ 0 & 0 & 0 & 0 \\ 0 & 0 & 0 & 0 \\ -1 & 0 & 0 & 0\end{array}\right]=-P_{14}$$
$$\mathrm{d} \widetilde{\mathrm{Ad}}(E)=\left[\begin{array}{rrrr}0 & 0 & 1 & 0 \\ 0 & 0 & 0 & 0 \\ -1 & 0 & 0 & 0 \\ 0 & 0 & 0 & 0\end{array}\right]=-P_{13}, \quad \mathrm{~d} \widetilde{\mathrm{Ad}}(i E)=\left[\begin{array}{rrrr}0 & 0 & 0 & 0 \\ 0 & 0 & 0 & -1 \\ 0 & 0 & 0 & 0 \\ 0 & -1 & 0 & 0\end{array}\right]=-P_{24}$$
$$\mathrm{d} \widetilde{\mathrm{Ad}}(F)=\left[\begin{array}{rrrr}0 & -1 & 0 & 0 \\ 1 & 0 & 0 & 0 \\ 0 & 0 & 0 & 0 \\ 0 & 0 & 0 & 0\end{array}\right]=P_{12}, \quad \mathrm{~d} \widetilde{\mathrm{Ad}}(i F)=\left[\begin{array}{rrrr}0 & 0 & 0 & 0 \\ 0 & 0 & 0 & 0 \\ 0 & 0 & 0 & -1 \\ 0 & 0 & -1 & 0\end{array}\right]=-P_{34}$$.\\
This shows that $\mathrm{d} \widetilde{\mathrm{Ad}}(C)$ maps a basis for $\mathrm{Sk}-\operatorname{Herm}_{2}(\mathbb{C})$ to one for $\mathfrak{l o r}(3,1)$, thus it is an isomorphism of Lie algebras.


(ii) The surjectivity of $\widetilde{\mathrm{Ad}}$ follows immediately from the surjectivity of the derivative in (i) together with Theorem 6.11.
\end{PROOF}





\begin{DEF}{LIE-B02-09-52}{Lorentzsche Clifford-Algebren und Spin-Gruppen}
\textbf{(a) Definition der Lorentz-Clifford-Algebra.}
Für \(n\ge 1\) sei \(\R^{n-1,1}\) der Vektorraum \(\R^n\) mit dem Lorentz-Skalarprodukt
\(\langle\cdot,\cdot\rangle\).
Die \emph{Lorentz-Clifford-Algebra} \(\Cl_{n-1,1}\) ist die \(\R\)-Algebra,
erzeugt von Elementen \(e_1,\dots,e_n\) mit Relationen
\[
\begin{cases}
e_s e_r = - e_r e_s, & s\ne r,\\[4pt]
e_t^2 = -1, & 1\le t\le n-1,\\[4pt]
e_n^2 = +1.
\end{cases}
\]
Die kanonische lineare Einbettung \(j:\R^{n-1,1}\to\Cl_{n-1,1}\) wird durch
\(x_1 e_1+\dots+x_n e_n\mapsto \sum x_i e_i\) gegeben.

\textbf{(b) Universelle Eigenschaft.}
Es gilt ein Analogon der universellen Eigenschaft aus Theorem~5.4:
Ist \(A\) eine \(\R\)-Algebra und \(f:\R^{n-1,1}\to A\) eine lineare Abbildung mit
\(f(x)^2=-\langle x,x\rangle\,1\), so existiert genau ein
\(\R\)-Algebrenhomomorphismus
\[
F:\Cl_{n-1,1}\longrightarrow A,\qquad F\circ j=f.
\]

\textbf{(c) Rekursive Darstellung.}
Für \(n\ge 2\) definiere
\[
f_{n-1,1}:\R^{n-1,1}\to M_2(\Cl_{n-2}),\qquad
f_{n-1,1}\!\left(\sum_{r=1}^n a_r e_r\right)
=\begin{bmatrix}
\sum_{r=1}^{n-2} a_r e_r + a_n & -a_{n-1}\\[4pt]
a_{n-1} & -\sum_{r=1}^{n-2} a_r e_r - a_n
\end{bmatrix}.
\]
Dann induziert \(f_{n-1,1}\) mittels der universellen Eigenschaft
einen Isomorphismus von \(\R\)-Algebren
\[
F_{n-1,1}:\Cl_{n-1,1}\xrightarrow{\ \cong\ } M_2(\Cl_{n-2}),
\qquad
F_{n-1,1}\circ j_{n-1,1}=f_{n-1,1}.
\]

\textbf{(d) Die Lorentzsche Spin-Gruppe.}
Analog zur euklidischen Spin-Gruppe sei
\[
\Spin(n-1,1)=\{u_1\cdots u_{2k}:\ u_i\in \R^{n-1,1},\ |u_i|=1\}
\subset \Cl_{n-1,1}^\times.
\]
Dann existiert eine zusammenhängende Untergruppe
\[
\Spin^0(n-1,1)\subseteq \Spin(n-1,1)
\]
sowie eine zweifache Überlagerung der Lorentz-Gruppe:
\[
\Spin^0(n-1,1)\longrightarrow \Lor(n-1,1).
\]
\end{DEF}



\begin{THEO}{LIE-B02-09-53}{Dirac-Faktorisierung und Darstellung von \(\Cl_{3,1}\)}
\textbf{(a) Isomorphie der Algebren.}
Im Fall \(n=4\) gilt
\[
\Cl_{3,1}\ \cong\ M_2(\mathbb{H}).
\]
Über die reguläre Darstellung \(M_2(\mathbb{H})\hookrightarrow M_4(\C)\)
erhält man eine \(\C\)-lineare Darstellung von \(\Cl_{3,1}\)
auf dem \(\C\)-Vektorraum \(\C^4\).
Unter der Identifikation
\[
\SL_2(\C)\ \cong\ \Spin^0(3,1)
\]
entsteht eine stetige \(\C\)-lineare Wirkung von \(\SL_2(\C)\)
auf \(\C^4\).

\textbf{(b) Faktorisiertes Wellengleichungs-Operator.}
Für eine glatte Abbildung
\(\varphi:\R^{3,1}\to\C^4,\ \varphi=(\varphi_1,\varphi_2,\varphi_3,\varphi_4)\),
sei der D’Alembert-Operator definiert durch
\[
\square^2\varphi
=\frac{\partial^2\varphi}{\partial x_1^2}
+\frac{\partial^2\varphi}{\partial x_2^2}
+\frac{\partial^2\varphi}{\partial x_3^2}
-\frac{\partial^2\varphi}{\partial t^2}.
\]
Dann lässt sich dieser schreiben als
\[
\square^2
=\left(e_1\frac{\partial}{\partial x_1}
+e_2\frac{\partial}{\partial x_2}
+e_3\frac{\partial}{\partial x_3}
+e_4\frac{\partial}{\partial t}\right)^2,
\]
wobei die \(e_i\) die Generatoren der Clifford-Algebra \(\Cl_{3,1}\) sind,
die über die in (a) definierte Matrixdarstellung auf \(\C^4\) wirken.

\textbf{Bemerkung.}
Diese Faktorisierung stammt von \textsc{Dirac} (1928), der damit die
relativistische Schrödingergleichung für Elektronen formulierte.
Der Zusammenhang zwischen dieser Darstellung und Clifford-Algebren
wurde später von \textsc{Atiyah–Bott–Shapiro} im Rahmen der Index-Theorie
hergestellt.
\end{THEO}




\pagebreak


\subsection{EXERCISES}

6.1. a) For $t \in \mathbb{R}$, show that the matrix

$$
A_{t}=\left[\begin{array}{cccc}
\cosh t & 0 & 0 & -\sinh t \\
0 & \cos t & -\sin t & 0 \\
0 & \sin t & \cos t & 0 \\
-\sinh t & 0 & 0 & \cosh t
\end{array}\right]
$$

is in $\operatorname{Lor}(3,1)$.\\
b) Express $A_{t}$ in one of the forms predicted by Theorem 6.10.\\
c) Express $A_{t}$ in the form $\exp (t U)$ for some $U \in \mathfrak{l o r}(3,1)$.\\
d) Find a differential equation satisfied by the function

$$
\alpha: \mathbb{R} \longrightarrow \operatorname{Lor}(3,1) ; \quad \alpha(t)=A_{t} .
$$

6.2. In this question, we freely use the notation of Section 6.3.\\
a) Consider the differential equation

$$
\alpha^{\prime}(t)=\left(P_{24}+P_{34}\right) \alpha(t),
$$

where $\alpha: \mathbb{R} \longrightarrow \mathrm{GL}_{4}(\mathbb{R})$. Solve this to obtain a one-parameter subgroup $\alpha$ in $\mathrm{GL}_{4}(\mathbb{R})$ and show that $\alpha(t) \in \operatorname{Lor}(3,1)$ for all $t \in \mathbb{R}$.\\
b) Consider the differential equation

$$
\widetilde{\alpha}^{\prime}(t)=-i(E+F) \widetilde{\alpha}(t),
$$

where $\widetilde{\alpha}: \mathbb{R} \longrightarrow \mathrm{GL}_{2}(\mathbb{C})$. Solve this to obtain a one-parameter subgroup in $\mathrm{GL}_{2}(\mathbb{C})$ and show that $\widetilde{\alpha}(t) \in \mathrm{SL}_{2}(\mathbb{C})$ for all $t \in \mathbb{R}$.\\
c) Explain the relationship between these two curves.\\
6.3. $\boldsymbol{\Delta} \boldsymbol{\Delta}$ For $n \geqslant 1$, let $\mathbb{R}^{n-1,1}$ be $\mathbb{R}^{n}$ equipped with the Lorentz inner product $\langle$,$\rangle . Define a Lorentz Clifford algebra \mathrm{Cl}_{n-1,1}$ in a similar fashion to the way $\mathrm{Cl}_{n}$ was defined in Section 5.1, except that the algebra generators $e_{1}, \ldots, e_{n-1}, e_{n}$ satisfy the relations

$$
\begin{cases}e_{s} e_{r}=-e_{s} e_{r} & \text { if } s \neq r, \\ e_{t}^{2}=-1 & \text { if } 1 \leqslant t \leqslant n-1, \\ e_{n}^{2}=1 . & \end{cases}
$$

Let $j: \mathbb{R}^{n-1,1} \longrightarrow \mathrm{Cl}_{n-1,1}$ be the $\mathbb{R}$-linear transformation embed$\operatorname{ding} \mathbb{R}^{n-1,1}$ into $\mathrm{Cl}_{n-1,1}$.\\
a) Show that there is an analogue for $\mathrm{Cl}_{n-1,1}$ of the Universal Property of Theorem 5.4.\\
b) For $n \geqslant 2$, define $f_{n-1,1}: \mathbb{R}^{n-1,1} \longrightarrow \mathrm{M}_{2}\left(\mathrm{Cl}_{n-2}\right)$ by

$$
f_{n-1,1}\left(\sum_{1 \leqslant r \leqslant n} a_{r} e_{r}\right)=\left[\begin{array}{cc}
\sum_{1 \leqslant r \leqslant n-2} a_{r} e_{r}+a_{n} & -a_{n-1} \\
a_{n-1} & -\sum_{1 \leqslant r \leqslant n-2} a_{r} e_{r}-a_{n}
\end{array}\right] .
$$

Using the Universal Property of (a), show that $f_{n-1,1}$ induces an isomorphism of $\mathbb{R}$-algebras

$$
F_{n-1,1}: \mathrm{Cl}_{n-1,1} \longrightarrow \mathrm{M}_{2}\left(\mathrm{Cl}_{n-2}\right)
$$

for which $F_{n-1,1} \circ j_{n-1,1}=f_{n-1,1}$.\\
c) Discuss the analogue $\operatorname{Spin}(n-1,1) \leqslant \mathrm{Cl}_{n-1,1}^{\times}$of the spinor group $\operatorname{Spin}(n) \leqslant \mathrm{Cl}_{n}^{\times}$. In particular, deduce that there is a connected subgroup $\operatorname{Spin}^{0}(n-1,1) \leqslant \operatorname{Spin}(n-1,1)$ and a double covering homomorphism

$$
\operatorname{Spin}^{0}(n-1,1) \longrightarrow \operatorname{Lor}(n-1,1) .
$$

6.4. $\Delta \boldsymbol{\Delta}$ a) Using the previous question in the case where $n=4$, deduce that there is an isomorphism of $\mathbb{R}$-algebras

$$
\mathrm{Cl}_{3,1} \longrightarrow \mathrm{M}_{2}(\mathbb{H})
$$

Using the regular representation $\mathrm{M}_{2}(\mathbb{H}) \longrightarrow \mathrm{M}_{4}(\mathbb{C})$, this provides an action of $\mathrm{Cl}_{3,1}$ by $\mathbb{C}$-linear transformations on the $\mathbb{C}$-vector space $\mathbb{C}^{4}$. Identifying $\operatorname{SL}_{2}(\mathbb{C})$ with $\operatorname{Spin}^{0}(3,1)$, this gives a continuous $\mathbb{C}$-linear action of $\mathrm{SL}_{2}(\mathbb{C})$ on $\mathbb{C}^{4}$.\\
b) Let $\varphi: \mathbb{R}^{3,1} \longrightarrow \mathbb{C}^{4}$ be a smooth function with

$$
\varphi=\left(\varphi_{1}, \varphi_{2}, \varphi_{3}, \varphi_{4}\right)
$$

Show that the $D^{\prime}$ Alembertian operator $\square^{2}$ given by

$$
\square^{2} \varphi=\frac{\partial^{2} \varphi_{1}}{\partial x_{1}{ }^{2}}+\frac{\partial^{2} \varphi_{2}}{\partial x_{2}{ }^{2}}+\frac{\partial^{2} \varphi_{3}}{\partial x_{3}{ }^{2}}-\frac{\partial^{2} \varphi_{4}}{\partial t^{2}}
$$

can be expressed as

$$
\square^{2}=\left(e_{1} \frac{\partial}{\partial x_{1}}+e_{2} \frac{\partial}{\partial x_{2}}+e_{3} \frac{\partial}{\partial x_{3}}+e_{4} \frac{\partial}{\partial t}\right)^{2}
$$

where the terms inside the brackets are to be thought of as differential operators with coefficients in $\mathrm{Cl}_{3,1}$ which act on $\mathbb{C}^{4}$ via the matrix representation defined in (a).\\[0pt]
[This was originally observed by the physicist Dirac, who used it to factorise a relativistic version of the Schrödinger equation in order to study the electron. The connection with Clifford algebras seems to have been observed by Atiyah, Bott \& Shapiro in their work on Index Theory [3].]







\pagebreak
\printbibliography
\end{document}